\documentclass[11pt]{article}

\usepackage[margin=1.1in]{geometry}

\usepackage{amssymb,mathtools,natbib}

% Anchored formal environments for causalsmith papers.
% First mandatory argument = obj_id (machine anchor joining the paper to the
% Lean crosswalk; invisible in the rendered PDF). Optional argument = name.
% Each environment also defines \label{obj:<obj_id>} so prose can use
% \cref{obj:T-1}; cleveref derives the printed kind and number from the target.
\usepackage{amsmath}
\usepackage{mathrsfs}
\usepackage{amsthm}
\usepackage{booktabs}
% Long displays may break across pages; let TeX stretch a little harder before an
% overfull \hbox so wide formulas/tables are less likely to run off the page.
\allowdisplaybreaks
\setlength{\emergencystretch}{3em}
\newtheorem{theorem}{Theorem}
\newtheorem{lemma}{Lemma}
\newtheorem{assumption}{Assumption}
\newtheorem{definition}{Definition}
\newtheorem{proposition}{Proposition}
\newtheorem{remark}{Remark}
\newtheorem{citedresult}{Cited result}
% The amsthm counter is `algorithmbox` (not `algorithm`) so a later `\usepackage{algorithm}` cannot clash.
\newcounter{algorithmbox}
\renewcommand{\thealgorithmbox}{\arabic{algorithmbox}}
\NewDocumentEnvironment{theoremv}{m o s}{\begin{theorem}\label{obj:#1}{\normalfont[\texttt{\detokenize{#1}}]\IfValueT{#2}{\ (#2)}\IfBooleanT{#3}{\verificationfootnotemark}\ }}{\end{theorem}}
\NewDocumentEnvironment{lemmav}{m o s}{\begin{lemma}\label{obj:#1}{\normalfont[\texttt{\detokenize{#1}}]\IfValueT{#2}{\ (#2)}\IfBooleanT{#3}{\verificationfootnotemark}\ }}{\end{lemma}}
\NewDocumentEnvironment{assumptionv}{m o s}{\begin{assumption}\label{obj:#1}{\normalfont[\texttt{\detokenize{#1}}]\IfValueT{#2}{\ (#2)}\IfBooleanT{#3}{\verificationfootnotemark}\ }}{\end{assumption}}
\NewDocumentEnvironment{definitionv}{m o s}{\begin{definition}\label{obj:#1}{\normalfont[\texttt{\detokenize{#1}}]\IfValueT{#2}{\ (#2)}\IfBooleanT{#3}{\verificationfootnotemark}\ }}{\end{definition}}
% A `propositionv` is a proved result tiered as a Proposition (theorem-grade, machine-verified).
\NewDocumentEnvironment{propositionv}{m o s}{\begin{proposition}\label{obj:#1}{\normalfont[\texttt{\detokenize{#1}}]\IfValueT{#2}{\ (#2)}\IfBooleanT{#3}{\verificationfootnotemark}\ }}{\end{proposition}}
% A `remarkv` holds NON-load-bearing interpretive / open-question content (e.g. a causalsmith `oeq:`
% open question). It is neither proved nor assumed; no theorem may depend on it.
\NewDocumentEnvironment{remarkv}{m o s}{\begin{remark}\label{obj:#1}{\normalfont[\texttt{\detokenize{#1}}]\IfValueT{#2}{\ (#2)}\IfBooleanT{#3}{\verificationfootnotemark}\ }}{\end{remark}}
% An `algorithmv` is a DEFINITION-kind object presented as a boxed, numbered-step procedure (an
% estimator / selection rule built in stages). Same anchor contract as `definitionv`; its body is
% ordinary LaTeX whose steps are `\begin{enumerate}` items, so tex2html needs no extra machinery.
% Presented in the CONVENTIONAL algorithm style readers expect from `algorithm`+`algorithmic`:
% a rule above, an "Algorithm N (Title)" caption line, a rule under the caption, the numbered
% steps, and a closing rule. It is NOT a float — the anchor contract requires the object to stay
% where the outline places it, and floats drift. The obj-id tag is suppressed here (it is
% bookkeeping, and a caption line carrying it reads as debug output in a finished paper).
\NewDocumentEnvironment{algorithmv}{m o s}%
  {\par\addvspace{\medskipamount}\noindent\rule{\linewidth}{0.8pt}\par\nobreak\vspace{2pt}%
   \refstepcounter{algorithmbox}\label{obj:#1}%
   \noindent\textbf{Algorithm \thealgorithmbox}\IfValueT{#2}{ \textnormal{(#2)}}%
   \IfBooleanT{#3}{\verificationfootnotemark}\par\nobreak\vspace{2pt}%
   \noindent\rule{\linewidth}{0.4pt}\par\nobreak\vspace{2pt}\normalfont}%
  {\par\nobreak\vspace{2pt}\noindent\rule{\linewidth}{0.8pt}\par\addvspace{\medskipamount}}
% A `citedv` is an IMPORTED external result the paper relies on but does NOT prove
% (a causalsmith `gate_class:"cited"` node). It is machine-checked against the cited
% source at F2.5, never proved here; the fixed provenance note makes that explicit
% so the environment can never be read as a contribution of this paper. The \citep
% attribution is supplied by the surrounding prose (P2), keyed to references.bib.
\NewDocumentEnvironment{citedv}{m o s}{\begin{citedresult}\label{obj:#1}{\normalfont[\texttt{\detokenize{#1}}]\IfValueT{#2}{\ (#2)}\IfBooleanT{#3}{\verificationfootnotemark}\ \textit{(imported result; machine-checked against the cited source, not proved here.)}\ }}{\end{citedresult}}
% \leanref{obj_id}{display text}: an inline first-mention link to a from-note object's
% verified Lean code. In the PDF it renders the display text plainly (the Lean link is a
% web-only affordance); tex2html turns it into a drawer-opening span.
\newcommand{\leanref}[2]{#2}
% For a theorem/lemma whose Lean declaration accepts a source-matched published proposition as an
% input, the starred anchored environment puts the mark in its heading; the generated text remains
% outside the frozen mathematical environment. Keep the combined macro for older paper bundles.
\newcommand{\verificationfootnotemark}{\footnotemark}
\newcommand{\verificationfootnotetext}[1]{\footnotetext{#1}}
\newcommand{\verificationfootnote}[1]{\verificationfootnotemark\verificationfootnotetext{#1}}
% xcolor before hyperref (hyperref would otherwise pull in plain `color` for colorlinks, and a
% later xcolor load clashes). The page-1 identity stamp at the end of this file needs a grey.
\usepackage{xcolor}
\usepackage{graphicx}
\usepackage{eso-pic}
% hyperref follows natbib and the environment macros; cleveref must follow hyperref.
\usepackage[colorlinks=true,linkcolor=blue,citecolor=blue,urlcolor=blue]{hyperref}
\usepackage[capitalise,nameinlink,noabbrev]{cleveref}
% Pin the names explicitly: labels are placed inside the underlying amsthm environments, and these
% declarations make the PDF contract independent of cleveref's theorem-name inference. House style:
% reference names always print with a capital first letter ("Theorem 1", "Definition 2"), in any
% sentence position — hence `capitalise` above and capitalized \crefname forms below.
\crefname{theorem}{Theorem}{Theorems}
\Crefname{theorem}{Theorem}{Theorems}
\crefname{lemma}{Lemma}{Lemmas}
\Crefname{lemma}{Lemma}{Lemmas}
\crefname{assumption}{Assumption}{Assumptions}
\Crefname{assumption}{Assumption}{Assumptions}
\crefname{definition}{Definition}{Definitions}
\Crefname{definition}{Definition}{Definitions}
\crefname{proposition}{Proposition}{Propositions}
\Crefname{proposition}{Proposition}{Propositions}
\crefname{remark}{Remark}{Remarks}
\Crefname{remark}{Remark}{Remarks}
\crefname{citedresult}{Cited result}{Cited results}
\Crefname{citedresult}{Cited result}{Cited results}
\crefname{algorithmbox}{Algorithm}{Algorithms}
\Crefname{algorithmbox}{Algorithm}{Algorithms}
% ---------------------------------------------------------------------------
% Page-1 identity stamp (arXiv style) and the pinned title-block date.
%
% Split of responsibility: THIS file owns the FORMAT (placement, font, colour, punctuation),
% the generated `paper_stamp.tex` owns only the VALUES (number+version, area, revised date,
% DOI, link target). P4 refreshes this template on every emit, so a layout fix reaches every
% bundle from a recompile; and because `paper_stamp.tex` is not one of the P2 assembly
% digest's authored inputs, a DOI that arrives after publication needs only that one small
% file rewritten plus a recompile — never a redraft, never a reassembly.
%
% Defaults first, so a bundle WITHOUT a stamp file compiles exactly as it always did:
% `\cspaperdate` falls back to `\today` and no stamp is drawn.
\providecommand*{\cspaperdate}{\today}  % the \date{} target
\providecommand*{\csstampid}{}          % e.g. CSWP-2026-008v3 (empty ⇒ no stamp at all)
\providecommand*{\csstamparea}{}        % e.g. Stat
\providecommand*{\csstampdate}{}        % e.g. 27 Aug 2026
\providecommand*{\csstampdoi}{}         % e.g. 10.5281/zenodo.123 (empty ⇒ DOI part omitted)
\providecommand*{\csstamplink}{}        % href target (DOI url, else the paper's page; may be empty)
\IfFileExists{paper_stamp.tex}{% GENERATED by CausalSmith P4 — paper identity stamp values. Do not hand-edit.
%
% Field order on the stamp (owned by paper_macros.tex):
%   <number>v<version> · doi:<doi> · [<Area>] <date>   — the DOI is never the last token, so a
%   text extractor cannot run it into whatever follows. Without a DOI that part is omitted.
% Format lives in paper_macros.tex; only the values are here, and this file is NOT one of
% the P2 assembly digest's authored inputs. A DOI arriving after publication therefore needs
% this file rewritten plus a `latexmk` recompile — never a redraft or a reassembly.
\renewcommand*{\csstampid}{CSWP-2026-041v3}
\renewcommand*{\csstamparea}{Stat}
\renewcommand*{\csstampdate}{4 Oct 2026}
\renewcommand*{\csstampdoi}{}
\renewcommand*{\csstamplink}{https://causalsmith.org/papers/stat_finitep_homogeneity_densegamma_v1}
\renewcommand*{\cspaperdate}{4 October 2026}
}{}
\makeatletter
% Field order is deliberate: the DOI is NEVER the last token on the line. A trailing identifier
% runs into whatever a text extractor emits next, which makes it unsafe to match mechanically
% (the Zenodo stamp matcher reads exactly this line); the date is a safe terminator.
%   <number>v<version> · doi:<doi> · [<Area>] <date>      — with a DOI
%   <number>v<version> · [<Area>] <date>                  — without
\newcommand{\cs@stampsep}{\,\textperiodcentered\,}
\newcommand{\cs@stamptext}{%
  \csstampid
  \ifx\csstampdoi\@empty\else\cs@stampsep doi:\csstampdoi\fi
  \ifx\csstamparea\@empty\else\cs@stampsep[\csstamparea]\fi
  \ifx\csstampdate\@empty\else\quad\csstampdate\fi
}
\ifx\csstampid\@empty\else
  \definecolor{csstampgrey}{gray}{0.45}
  % Background shipout material: it occupies no space, so the text block and the \thanks
  % footnote are byte-identical to an unstamped compile. The starred form applies to the
  % NEXT page shipped only — i.e. page 1. Placed in the left margin (the text block starts
  % at 1.1in), reading bottom-to-top, in a Times-like face at 8pt.
  \AddToShipoutPictureBG*{%
    \AtPageLowerLeft{%
      \put(\LenToUnit{0.42in},\LenToUnit{1.1in}){%
        \rotatebox{90}{%
          \fontfamily{ptm}\fontsize{8}{10}\selectfont
          \ifx\csstamplink\@empty
            \textcolor{csstampgrey}{\cs@stamptext}%
          \else
            \expandafter\href\expandafter{\csstamplink}{\textcolor{csstampgrey}{\cs@stamptext}}%
          \fi
        }%
      }%
    }%
  }
\fi
\makeatother


\title{Testing treatment effect homogeneity under finite conditional moments}

\author{CausalSmith\thanks{Machine-generated by the CausalSmith research pipeline.}}

\date{\cspaperdate}

\begin{document}

\maketitle

\begin{abstract}
We characterize the minimax separation rate for testing whether a conditional mean treatment effect is constant under a known uniform covariate distribution on the unit interval, fixed overlap, independent Hölder restrictions on the propensity, baseline mean, and effect, fixed mean envelopes, and an armwise conditional outcome-moment bound of order between one and two. Separation is measured by the population \(L_2\) distance of the effect from its average. Matching upper and lower bounds identify two rate mechanisms: effect approximation under outcome tails and the interaction of rough propensity and baseline functions. An explicit test combines outcome clipping, histogram projection corrections, independent evaluation blocks, and profiling over the unknown common effect. It controls size over the entire composite null at every stipulated sample size and attains uniform power at the stated separation whenever that separation lies below the model maximum, with public thresholds ensuring this condition. At matched smoothness, every moment order below two yields a polynomially larger critical radius than the bounded-outcome benchmark, whose minimax risk equals that of a signed-binary submodel. Supplying the entire true continuous propensity yields a matching frontier over a broader continuous-nuisance class retaining uniform design, overlap, effect smoothness, mean caps, and the conditional moment envelope. Comparing the two experiments on the original class identifies the exact region where the unknown-propensity experiment attains the supplied-propensity separation order.
\end{abstract}

\section{Introduction}

This paper establishes matching minimax separation bounds for conditional mean treatment-effect homogeneity under finite conditional outcome moments and independently specified nuisance smoothness. The statistical target is constancy of the difference between the treated and untreated conditional means, with an unknown common effect. Under a potential-outcome representation satisfying consistency and conditional independence of treatment assignment, this contrast is the conditional average treatment effect. We measure heterogeneity by its population \(L_2\) distance from its average, which equals the standard deviation of the conditional effect under the stipulated covariate distribution. The resulting frontier quantifies how outcome tails and assignment information jointly determine the smallest uniformly detectable departure from homogeneity.

The model has a known uniform covariate distribution on \([0,1]\), a propensity between \(1/4\) and \(3/4\), and separate Hölder restrictions on the propensity, baseline conditional mean, and mean effect. Write \(p\), the conditional moment order, for an exponent satisfying \(1<p\le2\); write \(\alpha\) and \(\beta\), the propensity and baseline smoothness exponents, for numbers in \((0,1]\); and write \(\gamma\), the effect smoothness exponent, for a number in \([1/4,1]\). These exponents are public. The Hölder radii are fixed at \(20\), the baseline mean and effect have supremum norms at most \(1/2\), and each treatment arm has conditional raw \(p\)-moment at most \(10\), uniformly in the covariate up to null sets. These restrictions define the class in \cref{obj:def:model}. Testing errors are assessed uniformly over that class: size covers every admissible constant effect and nuisance function, while power covers every law separated by the stipulated population distance, as specified in \cref{obj:def:testing-risk}.

Our main result gives the critical separation radius up to positive constants depending on the public tuple. Define \(q=(p-1)/p\), the moment-derived exponent, and \(S=\alpha+\beta\), the sum of nuisance smoothness exponents. The two separation exponents are
\[
E_0(p)=\frac{2\gamma q}{2\gamma+q},
\qquad
E_4(p)=
\frac{2S}{1+2\alpha+\beta/q+S/(2\gamma)}.
\]
For sample size \(n\ge2\), the capped critical radius has order
\[
n^{-\min\{E_0(p),E_4(p)\}},
\]
as established in \cref{obj:thm:heavy-tail-comparison}. The first exponent balances effect approximation with outcome-moment control. The second incorporates the interaction of propensity and baseline regularity. Their exact equality boundary is determined by
\[
F(p)=
\frac{2\alpha}{p-1}
+\frac{\beta}{q}
+\frac{S}{2\gamma},
\]
the phase-comparison quantity: effect smoothness determines the rate when \(F(p)\ge1\), and rough confounding determines it when \(F(p)<1\). The bounds are pure powers of sample size, including at equality, with constants locally uniform on compact subsets of the admissible exponent domain.

The attaining procedure operates directly on the original records. It clips outcomes, forms histogram scores with ordered-pair projection corrections, and takes the inner product of scores from two independent evaluation blocks. Profiling this quadratic statistic over the admissible common-effect interval accommodates the composite null. The public construction in \cref{obj:def:explicit-score-ledger} specifies all ranks, clipping levels, bias and covariance bounds, and rejection cutoffs. Its rejection probability is at most \(0.0075\) throughout the null for every \(n\ge2\). Its type-II error is at most \(0.0075\) at the stated upper separation when that separation lies below the model's maximal heterogeneity distance; fully evaluable sample-size thresholds ensure a nonempty power domain. These guarantees appear in \cref{obj:thm:whole-null-calibration-resolved,obj:thm:original-record-attainable-rate}. Matching converses apply to all admissible randomized tests observing the complete records.

The matched outcome comparison retains the same design, overlap, smoothness, and mean envelopes while imposing bounded outcomes. Its separation exponent is the original exponent evaluated at \(p=2\). Every admissible \(p<2\) strictly reduces that exponent, so the finite-moment critical radius exceeds its bounded counterpart by a polynomial factor. At \(p=2\), the ratio remains bounded across sample sizes. Moreover, mean-preserving sign randomization gives exact equality of bounded and signed-binary minimax risks and capped radii, as shown in \cref{obj:prop:bounded-submodel-comparison}. The full-record witnesses in \cref{obj:thm:tail-essential-full-record-converse} connect the slower finite-moment frontier to growing outcome marks: for each fixed tuple with \(p<2\), close finite-moment null and alternative mixtures coexist with eventual total-variation separation of every normalized finite prior pair supported on the bounded null and bounded alternatives at the same lower separation.

Supplying the entire true continuous propensity gives a complementary frontier. Over the continuous-carrier class in \cref{obj:def:oracle-broad-model}, the propensity and conditional arm means may have unrestricted moduli of continuity, while uniform design, overlap, effect Hölder smoothness, mean caps, and the conditional moment envelope are retained. A clipped inverse-propensity test attains the matching order \(n^{-E_0(p)}\) on this broader class, with whole-null calibration and a public nonempty-power threshold, by \cref{obj:thm:oracle-frontier-broad-class}. Comparing supplied and unknown propensity on the original shared class then yields an exact preservation boundary: their critical-radius ratio remains bounded over \(n\ge2\) precisely when \(F(p)\ge1\). For \(F(p)<1\), supplying the propensity improves separation by a factor of order \(n^{E_0(p)-E_4(p)}\). This comparison, including its equality case, follows from the matched frontiers in \cref{obj:thm:oracle-frontier-resolved,obj:thm:exact-propensity-preservation-boundary}.

These results connect econometric homogeneity testing with quadratic nonparametric testing and higher-order causal methods. \citet[Section III]{CrumpHotzImbensMitnik2008Homogeneity} develop sieve Wald tests of conditional mean-effect homogeneity. Recent orthogonal-score procedures provide asymptotic calibration and power under their respective outcome and nuisance-learning conditions \citep[Assumptions 1--5, Theorems 3.1--3.3 and 4.1]{LuSong2026OrthogonalICM}; \citep[Assumptions 3.3--3.4, Proposition 3.1]{LapentaStrittmatterVergara2026Omnibus}. Directional testing also explicitly pursues uniform asymptotic calibration, with power tied to alignment of a learned direction and the population signal \citep[Theorem 1, Corollaries 1--2]{DhawanGuoShah2026Directional}. Our contribution concerns finite-sample whole-null calibration and uniform population separation over the primitive class specified above. The approximation--variability balance has antecedents in quadratic goodness-of-fit theory \citep[Sections 2--4, Theorems 1--2]{IngsterSapatinas2009GoodnessOfFit}, while projection corrections connect to higher-order conditional-effect estimation \citep[Sections 3--4]{KennedyBalakrishnanRobinsWasserman2024MinimaxCATE}.



The paper proceeds through related work, setup and assumptions, the matched separation rates and outcome-tail comparison, the calibrated original-record construction, and the supplied-propensity frontiers and assignment-information comparison. A discussion considers interpretation and extensions. The appendices develop causal identification and benchmark equivalence, calibration and attainment, full-record lower bounds, supplied-propensity arguments and comparison algebra, followed by proofs of the main results.


\section{Related work}

The closest literature tests whether conditional average treatment effects are constant across covariate-defined subpopulations. \citet[Section III]{CrumpHotzImbensMitnik2008Homogeneity} develop a sieve Wald approach that compares series estimates of the two conditional outcome means and removes the constant component when testing homogeneity. Their analysis combines smoothness and sieve-growth restrictions with conditions on conditional variances and third moments of outcome disturbances to obtain asymptotic calibration and consistency against specified smooth alternatives. This establishes an important precedent for treating homogeneity as a global restriction on conditional means. The present paper studies that restriction through finite-sample control over the entire composite null and a population separation criterion that is uniform over the stipulated model class; the corresponding guarantees are stated in \cref{obj:thm:whole-null-calibration-resolved,obj:thm:original-record-attainable-rate}.

Orthogonal conditional moments provide a second route to the same substantive question. \citet[Assumptions 1--5, Proposition 2, and Theorems 3.1--3.3 and 4.1]{LuSong2026OrthogonalICM} construct a marked empirical process from a doubly robust score and centered covariate indicators. Their theory specifies series estimators for the propensity score and outcome regressions, with primitive smoothness and series-growth conditions ensuring a feasible-to-oracle approximation uniform over the process index. The outcome conditions include finite second moments and uniformly bounded conditional variances. The resulting tests have asymptotic Gaussian-process calibration, multiplier-bootstrap validity, consistency against fixed alternatives, and nontrivial power along specified local alternatives approaching the null at the inverse square root of sample size. Uniformity over the process index supports these distributional approximations. Our separation criterion instead takes a supremum over all admissible population laws whose effect function is sufficiently far from its average, as specified in \cref{obj:def:testing-risk}.

\citet[Assumptions 3.3--3.4, Proposition 3.1, and Section 3.3]{LapentaStrittmatterVergara2026Omnibus} integrate orthogonal unconditional moments to obtain an omnibus test that accommodates flexible nuisance estimation. Their formulation separates the controls used for identification from the covariates used to assess heterogeneity. Under unconfoundedness, the stated theory assumes bounded outcomes and nuisance estimators converging faster than the inverse fourth root of sample size in population \(L_2\), together with boundedness and overlap conditions on the fitted functions. It establishes an asymptotic null distribution, consistency against fixed alternatives, and a bootstrap implementation that retains the fitted nuisance functions across replications. These conditions make explicit the nuisance-learning and outcome envelopes supporting their calibration. Here, the primitive restrictions in \cref{obj:def:model} determine the public tuning and finite-sample bounds directly.

A closely related scalar target is the variance of the conditional average treatment effect. \citet[Assumptions 3--4, Algorithm 1, and Theorem 1]{Yu2026VarianceHTE} address its zero boundary using intra-fold sample splitting: distinct evaluation subsamples share the same out-of-fold nuisance estimates, producing cancellation of nuisance-bias terms. The stated asymptotic normality result uses finite fourth outcome moments, a nondegeneracy condition, and convergence faster than the inverse fourth root of sample size for the propensity, outcome-regression, and conditional-effect estimators. Under the uniform covariate marginal in \cref{obj:ass:uniform}, the squared population \(L_2\) distance of the effect from its average also equals its variance. Our testing risk therefore measures separation using a closely related population quantity, with calibration and power assessed uniformly over the model's conditional-moment envelope.

Earlier variance-target work studies estimation and confidence intervals rather than minimax separation. \citet[Sections 2.1--2.3]{Levy2021} define the variance of the full conditional treatment-effect function and develop TMLE and cross-validated TMLE estimators; their efficiency conditions include control of an outcome-regression remainder that is not removed merely by knowing the treatment mechanism. \citet[Sections 3--4, Theorems 6--7]{SanchezBecerra2023} instead work in experiments with known assignment probabilities and a predictive first stage. They establish efficiency for the full variance under additional learning conditions and uniform asymptotic coverage for a fold-specific proxy variance under their stated fourth-moment and regularity conditions. Our target is the same full population variance under the uniform design, but the criterion is finite-sample testing risk over Hölder and conditional-moment classes; it does not replace the variance by a fold-specific proxy or assume a learned first stage.

Directional testing offers a complementary way to concentrate power. \citet[Section 2.2.1, Assumptions 1--5, Theorem 1, and Corollaries 1--2]{DhawanGuoShah2026Directional} learn a promising score direction on one sample and evaluate a debiased score on an independent sample. Their general theory incorporates conditional score-moment bounds, variance conditions, uniform nuisance-estimation requirements, and the additional moment conditions stated in their asymptotic theorem. It delivers uniform asymptotic null calibration and a power guarantee requiring sufficient correlation between the learned direction and the population signal. Thus, uniform asymptotic inference is already an explicit objective in this literature. The present analysis adds finite-sample whole-null calibration and uniform separation guarantees indexed by primitive moment and smoothness restrictions. The distinction concerns both the error criterion and the alternatives over which power is guaranteed.

The quadratic structure connects these causal tests to nonparametric goodness-of-fit theory. \citet[Sections 2--4 and Theorems 1--2]{IngsterSapatinas2009GoodnessOfFit} analyze minimax separation in multivariate regression using increasing-dimensional U-statistics and ellipsoidal smoothness classes. \citet[Sections 2.3--2.4 and Theorems 1--2]{CommingesDalalyan2013QuadraticNull} study composite regression nulls defined by quadratic functionals, obtaining sharp asymptotic testing results for nonnegative diagonal functionals under their stated conditions. These analyses explain how approximation error and the stochastic cost of a quadratic statistic jointly determine detectable signal size. In the present setting, that balance also incorporates unknown treatment assignment and conditional outcome tails. The comparisons in \cref{obj:thm:heavy-tail-comparison,obj:thm:exact-propensity-preservation-boundary} characterize their consequences for critical separation radii within the stipulated classes.

The projection corrections also have a close methodological antecedent in higher-order causal estimation. \citet[Sections 3--4]{KennedyBalakrishnanRobinsWasserman2024MinimaxCATE} derive minimax bounds for pointwise conditional-effect estimation in Hölder models and construct a higher-order local polynomial R-learner, with attainment under the stated additional conditions. Its second-order corrections use projections to control nuisance-product bias, linking the accuracy of causal estimation to separate smoothness restrictions on treatment assignment and outcome means. Our corrected scores in \cref{obj:def:score-family,obj:def:explicit-score-ledger} use related projection ideas for a global homogeneity test. The population targets and loss criteria distinguish the two analyses: pointwise estimation measures recovery of the effect function at a covariate value, whereas \cref{obj:def:testing-risk,obj:def:critical-radius} assess rejection errors for constancy against population-separated alternatives.

Outcome-tail handling connects the construction to robust U-statistics. \citet[Section 2 and Theorem 3]{JolyLugosi2016RobustU} use median-of-means aggregation to estimate U-statistic expectations; their canonical-statistic result permits moment orders strictly above one and at most two. This supplies a precise comparator for finite-moment control of nonlinear statistics. Our construction combines clipping, projection corrections, and independent evaluation blocks, with explicit bias and covariance bounds in \cref{obj:def:explicit-score-ledger}. Together with the full-record converse in \cref{obj:thm:tail-essential-full-record-converse}, the analysis relates the attainable separation scale to the conditional outcome-moment restriction.

Other heterogeneity procedures address distinct scientific targets. \citet[Sections 2--4]{DukesStensrudBrioschiHudson2026Homogeneity} study quantitative and qualitative effect heterogeneity through contrasts tied to the population value of treatment rules, with asymptotic testing guarantees under their regularity conditions. \citet[Sections 2--3 and Theorem 3.3]{JainLuedtke2026Distributional} target equality of conditional potential-outcome distributions and establish asymptotic bootstrap calibration and consistency against fixed alternatives under their stated conditions. \citet[Sections 2--3]{AshouriHenderson2026Predictive} assess the predictive advantage of fitted models that allow heterogeneous effects over models imposing a constant effect. Policy-value contrasts, conditional distributional equality, and predictive performance each give heterogeneity a specific operational meaning. The contribution developed here concerns constancy of the conditional mean contrast under fixed overlap, a known uniform covariate marginal, stipulated smoothness restrictions, and a conditional raw-moment envelope. Within that scope, it combines whole-null finite-sample calibration with separation-rate comparisons that quantify the roles of outcome tails and supplied propensity information.


\section{Setup and assumptions}

Throughout, generic positive constants may depend on the fixed public exponents and numerical envelopes, while remaining independent of the sample size \leanref{sym:n}{\ensuremath{n}}. We write \(\|f\|_\infty=\sup_{x\in[0,1]}|f(x)|\) and \(\|f\|_2=(\int f^2\,d\lambda)^{1/2}\), where \leanref{sym:lambda}{\ensuremath{\lambda}} is Lebesgue probability measure on \([0,1]\). The relations \(\lesssim\) and \(\asymp\) denote inequalities up to such constants and their two-sided counterpart, respectively. Sample indices range over \(i=1,\ldots,n\), treatment-arm indices over \(a\in\{0,1\}\), and all sample sizes satisfy \(n\in\mathbb N\) with \(n\ge2\).

\subsection{Records, public exponents, and conditional means}

An \leanref{S-1}{original record} is \leanref{sym:O}{\ensuremath{O=(X,A,Y)}}, with covariate \leanref{sym:X}{\ensuremath{X\in[0,1]}}, treatment indicator \leanref{sym:A}{\ensuremath{A\in\{0,1\}}}, and observed outcome \leanref{sym:Y}{\ensuremath{Y\in\mathbb R}}. An original-record law \leanref{sym:P}{\ensuremath{P}} is a Borel probability measure on this product space; the restrictions defining the statistical model are collected in \cref{obj:def:model}. Its covariate marginal is denoted by \leanref{sym:PX}{\ensuremath{P_X}}.

The public exponent tuple \leanref{sym:v}{\ensuremath{v=(p,\alpha,\beta,\gamma)}} consists of the conditional moment order \leanref{sym:p}{\ensuremath{p}}, the propensity smoothness exponent \leanref{sym:alpha}{\ensuremath{\alpha}}, the baseline smoothness exponent \leanref{sym:beta}{\ensuremath{\beta}}, and the effect smoothness exponent \leanref{sym:gamma}{\ensuremath{\gamma}}. The admissible public domain \leanref{sym:V}{\ensuremath{\mathcal V}} is
\[
\mathcal V
=
\left\{
(p,\alpha,\beta,\gamma)\in\mathbb R^4:
1<p\le2,\quad
0<\alpha\le1,\quad
0<\beta\le1,\quad
1/4\le\gamma\le1
\right\}.
\]
For benchmark comparisons, the matched smoothness tuple \leanref{sym:w}{\ensuremath{w=(\alpha,\beta,\gamma)}} belongs to the domain \leanref{sym:Wdomain}{\ensuremath{\mathcal W}}, where
\[
\mathcal W=(0,1]\times(0,1]\times[1/4,1].
\]
These exponents are public parameters indexing the classes over which testing errors are controlled.

For \(s\in(0,1]\), the Hölder ball \leanref{sym:Holder}{\ensuremath{\mathcal H^s(20)}} uses the convention
\[
[f]_s
=
\sup_{\substack{x,z\in[0,1]\\x\ne z}}
\frac{|f(x)-f(z)|}{|x-z|^s},
\qquad
\mathcal H^s(20)
=
\left\{
f\in C([0,1]):
\|f\|_\infty\le20,\ [f]_s\le20
\right\},
\]
where \(C([0,1])\) denotes continuous real functions with the uniform topology. Thus the convention at \(s=1\) is Lipschitz continuity. We write \leanref{sym:EqClass}{\ensuremath{[f]_\lambda}} for the equivalence class of a function under Lebesgue almost-everywhere equality; the same convention applies to sums of functions and to each conditional mean.

Let \leanref{sym:Q}{\ensuremath{Q_{a,P}(dy\mid x)}} be a conditional outcome probability kernel given \(X=x\) and \(A=a\). The propensity equivalence class \leanref{sym:eClass}{\ensuremath{\mathcal E_P}} is the class of \(x\mapsto P(A=1\mid X=x)\), and the conditional arm-mean classes \leanref{sym:mzeroClass}{\ensuremath{\mathcal G_{0,P}}} and \leanref{sym:moneClass}{\ensuremath{\mathcal G_{1,P}}} are the classes of
\[
x\longmapsto \int y\,Q_{0,P}(dy\mid x)
\quad\text{and}\quad
x\longmapsto \int y\,Q_{1,P}(dy\mid x),
\]
respectively. Under the restrictions below, these classes have continuous representatives: the propensity \leanref{sym:e}{\ensuremath{e_P}}, the baseline mean \leanref{sym:mzero}{\ensuremath{m_{0,P}}}, and the treated mean \leanref{sym:mone}{\ensuremath{m_{1,P}}}, as specified in \cref{obj:def:model}. Their canonical mean contrast is \leanref{sym:tau}{\ensuremath{\tau_P=m_{1,P}-m_{0,P}}}. Continuity and the full support of Lebesgue measure make each representative unique, so pointwise restrictions on these functions have an unambiguous meaning.

\subsection{Primitive restrictions}

The design conditions fix the distribution over which heterogeneity is measured and ensure that both treatment arms are represented throughout the covariate interval.

\begin{assumptionv}{ass:uniform}[Uniform covariate distribution]
Let \(\lambda\) denote Lebesgue probability measure on \([0,1]\). The covariate marginal \(P_X\) satisfies
\[
P_X=\lambda.
\]
\end{assumptionv}

The known uniform random design condition supplies a fixed integration measure for the population target and for approximation errors, following the random-design goodness-of-fit formulation \citep{IngsterSapatinas2009GoodnessOfFit}. Treatment assignment is further restricted by a fixed overlap envelope.

\begin{assumptionv}{ass:overlap}[Fixed propensity overlap]
The continuous propensity representative \(e_P\) satisfies
\[
e_P([0,1])\subseteq[1/4,3/4].
\]
\end{assumptionv}

Uniform overlap keeps both conditional arm distributions observable across the entire interval. The overlap form follows the conditional-effect model in \citet[Section 2]{KennedyBalakrishnanRobinsWasserman2024MinimaxCATE}; the numerical bounds \(1/4\) and \(3/4\) are fixed normalizations stipulated for the present class.

Smoothness is indexed separately for treatment assignment, the baseline mean, and the mean contrast. We begin with regularity of the propensity.

\begin{assumptionv}{ass:propensity-smoothness}[Propensity Hölder smoothness]
Let \(\mathcal H^\alpha(20)\) denote the continuous functions on \([0,1]\) with supremum norm and \(\alpha\)-Hölder constant at most \(20\), and let \([f]_\lambda\) denote the equivalence class of \(f\) under almost-everywhere equality for Lebesgue probability measure \(\lambda\). The propensity equivalence class \(\mathcal E_P\) satisfies
\[
\mathcal E_P\in\{[f]_\lambda:f\in\mathcal H^\alpha(20)\}.
\]
\end{assumptionv}

In our model, independent propensity Hölder smoothness controls variation in assignment probabilities through its own exponent, with radius \(20\) fixing the scale of that variation. Baseline outcome variation receives a separate regularity restriction.

\begin{assumptionv}{ass:baseline-smoothness}[Baseline mean smoothness]
Let \(\mathcal H^\beta(20)\) denote the continuous functions on \([0,1]\) with supremum norm and \(\beta\)-Hölder constant at most \(20\), and let \([g]_\lambda\) denote the equivalence class of \(g\) under almost-everywhere equality for Lebesgue probability measure \(\lambda\). The baseline conditional arm-mean equivalence class \(\mathcal G_{0,P}\) satisfies
\[
\mathcal G_{0,P}\in\{[g]_\lambda:g\in\mathcal H^\beta(20)\}.
\]
\end{assumptionv}

In our model, the independent baseline Hölder condition controls the conditional mean in the untreated arm, allowing its regularity to differ from propensity regularity. The third exponent governs the conditional mean contrast directly.

\begin{assumptionv}{ass:effect-smoothness}[Mean effect smoothness]
Let \(\mathcal H^\gamma(20)\) denote the continuous functions on \([0,1]\) with supremum norm and \(\gamma\)-Hölder constant at most \(20\), and let \([f]_\lambda\) denote the equivalence class of \(f\) under almost-everywhere equality for Lebesgue probability measure \(\lambda\). With \(m_{0,P}\) the continuous baseline arm-mean representative, the treated conditional arm-mean equivalence class \(\mathcal G_{1,P}\) satisfies
\[
\mathcal G_{1,P}\in\{[m_{0,P}+t]_\lambda:t\in\mathcal H^\gamma(20)\}.
\]
\end{assumptionv}

Independent effect Hölder smoothness describes the treated mean as the baseline mean plus a regular contrast. This separates the target regularity from the two nuisance regularities. By comparison, the pointwise conditional-effect model of \citet[Sections 2--3]{KennedyBalakrishnanRobinsWasserman2024MinimaxCATE} summarizes propensity and baseline smoothness through an average exponent. The Hölder radius \(20\) used here is a fixed normalization of the present model.

We also fix the scales of the baseline mean and the mean effect. The first envelope bounds the baseline level.

\begin{assumptionv}{ass:baseline-cap}[Baseline mean bound]
The continuous baseline arm-mean representative \(m_{0,P}\) satisfies
\[
\|m_{0,P}\|_\infty\le1/2.
\]
\end{assumptionv}

The bounded conditional baseline mean condition uses the fixed normalization \(1/2\), placing baseline levels on a common scale across the indexed models. This numerical envelope is stipulated here. A corresponding envelope fixes the scale of the contrast.

\begin{assumptionv}{ass:effect-cap}[Uniform mean effect bound]
The mean effect \(\tau_P\), the canonical difference between the treated and baseline conditional means, satisfies
\[
\|\tau_P\|_\infty\le \frac12.
\]
\end{assumptionv}

The bounded conditional effect condition uses a second stipulated normalization \(1/2\), keeping the target on the same scale across smoothness and moment classes. Together, the two mean envelopes imply \(\|m_{1,P}\|_\infty\le1\).

Outcome tails are controlled separately from conditional mean levels. The following restriction bounds a raw conditional moment in each arm.

\begin{assumptionv}{ass:raw-moment}[Uniform conditional moment bound]
Let \(\lambda\) denote Lebesgue probability measure on \([0,1]\), and let \(Q_{a,P}(dy\mid x)\) denote the conditional outcome probability kernel in treatment arm \(a\). The conditional raw moment of order \(p\) satisfies
\[
\max_{a\in\{0,1\}}
\operatorname*{ess\,sup}_{x\sim\lambda}
\int |y|^p\,Q_{a,P}(dy\mid x)\le 10.
\]
\end{assumptionv}

This conditional raw-moment envelope is specific to the present analysis. It controls outcome magnitudes uniformly across covariate values, up to null sets, and across treatment arms. The range \(1<p\le2\) ensures integrability of the arm means while allowing the moment order to index the strength of tail control.

The original-record model \leanref{sym:M}{\ensuremath{\mathcal M_v}} in \cref{obj:def:model} collects these design, smoothness, mean-envelope, and moment restrictions.

\begin{definitionv}{def:model}[Original-record model \(\mathcal M_v\)]
The model \(\mathcal M_v\) is the set of original-record probability laws \(P\) satisfying the following restrictions, indexed by \(v=(p,\alpha,\beta,\gamma)\in\mathcal V\), where \(\mathcal V\) denotes the admissible domain of public moment and primitive smoothness tuples:
\begin{itemize}
\item \textbf{(Covariate law.)} \cref{obj:ass:uniform}.
\item \textbf{(Overlap.)} \cref{obj:ass:overlap}.
\item \textbf{(Propensity smoothness.)} \cref{obj:ass:propensity-smoothness}.
\item \textbf{(Baseline smoothness.)} \cref{obj:ass:baseline-smoothness}.
\item \textbf{(Effect smoothness.)} \cref{obj:ass:effect-smoothness}.
\item \textbf{(Baseline bound.)} \cref{obj:ass:baseline-cap}.
\item \textbf{(Effect bound.)} \cref{obj:ass:effect-cap}.
\item \textbf{(Conditional moment bound.)} \cref{obj:ass:raw-moment}.
\end{itemize}
Here \(P_X\) denotes the covariate marginal of \(P\). For a function \(f\), write \([f]_\lambda\) for its equivalence class under Lebesgue almost-everywhere equality. Let \(\mathcal E_P\) denote the propensity equivalence class, and let \(\mathcal G_{0,P}\) and \(\mathcal G_{1,P}\) denote the baseline and treated conditional arm-mean equivalence classes. Their continuous representatives are defined by
\[
\mathcal E_P=[e_P]_\lambda,\qquad
\mathcal G_{0,P}=[m_{0,P}]_\lambda,\qquad
\mathcal G_{1,P}=[m_{1,P}]_\lambda,
\qquad
e_P,m_{0,P},m_{1,P}\in C([0,1]),
\]
where \(C([0,1])\) is the space of continuous real functions on \([0,1]\) with the uniform topology. The canonical mean effect is
\[
\tau_P=m_{1,P}-m_{0,P}.
\]
\end{definitionv}

\subsection{Causal interpretation and the homogeneity target}

A potential-outcome representation connects the canonical mean contrast to a conditional average treatment effect. Let \leanref{sym:Yzero}{\ensuremath{Y(0)}} and \leanref{sym:Yone}{\ensuremath{Y(1)}} denote the baseline and treated potential-outcome coordinates. The completion \leanref{sym:Causal}{\ensuremath{\mathbb Q_P}} in \cref{obj:def:causal-completion} combines the observed arm distributions with assignment that is conditionally independent of these coordinates given the covariate.

\begin{assumptionv}{ass:null-constancy}[Constant mean effect restriction]
There exists \(c\in[-1/2,1/2]\) such that
\[
\tau_P(x)=c
\qquad\text{for every }x\in[0,1].
\]
\end{assumptionv}

Under this completion, consistency and conditional independence of assignment give
\[
\mathbb E_{\mathbb Q_P}[Y(1)-Y(0)\mid X=x]=\tau_P(x)
\quad\text{for }\lambda\text{-almost every }x.
\]
The observed-record marginal is \(P\), so the same conditional mean contrast admits the potential-outcome interpretation associated with consistency and unconfounded assignment \citep[Section 3]{Rubin1974}; \citep[Section 2]{Rosenbaum1983}. The auxiliary justification for this completion is given in \cref{obj:lem:causal-nonempty}.

To measure heterogeneity, define the average contrast \leanref{sym:meanTau}{\ensuremath{\bar\tau_P}} and its population distance from constancy \leanref{sym:d}{\ensuremath{d(P)}} by
\[
\bar\tau_P=\int_{[0,1]}\tau_P(x)\,d\lambda(x),
\qquad
d(P)
=
\left(
\int_{[0,1]}\{\tau_P(x)-\bar\tau_P\}^{2}\,d\lambda(x)
\right)^{1/2}.
\]
Equivalently, \(d(P)=\inf_{b\in\mathbb R}\|\tau_P-b\|_2\). The uniform design fixes the population weighting, and continuity makes zero distance equivalent to constancy on the entire interval. Homogeneity therefore permits an unknown common effect \leanref{sym:c}{\ensuremath{c}} in the public mean envelope.

\begin{definitionv}{def:null}[Constant-effect null $H_0(v)$]
\[
H_0(v)
=
\left\{
P\in\mathcal M_v:
\exists c\in[-1/2,1/2]\ 
\forall x\in[0,1],\ \tau_P(x)=c
\right\}.
\]
The model restrictions and null constancy condition are given in \cref{obj:ass:uniform,obj:ass:overlap,obj:ass:propensity-smoothness,obj:ass:baseline-smoothness,obj:ass:effect-smoothness,obj:ass:baseline-cap,obj:ass:effect-cap,obj:ass:raw-moment,obj:ass:null-constancy}.
\end{definitionv}

The unknown-constant conditional-average-effect homogeneity condition treats the common effect level as a nuisance parameter, matching the mean-homogeneity target of \citep{CrumpHotzImbensMitnik2008Homogeneity}. The constant-effect null \leanref{sym:Hnull}{\ensuremath{H_0(v)}} in \cref{obj:def:null} imposes this condition within the original-record model.

\begin{definitionv}{def:causal-completion}[Potential-outcome completion $\mathbb Q_P$]
The potential-outcome completion \(\mathbb Q_P\) is defined using Lebesgue probability measure \(\lambda\) on \([0,1]\) and the conditional outcome probability kernels \(Q_{a,P}(dy\mid x)\) in arms \(a\in\{0,1\}\) by
\[
\mathbb Q_P(dx,dy_0,dy_1,da)
=
\lambda(dx)\,Q_{0,P}(dy_0\mid x)\,Q_{1,P}(dy_1\mid x)\,
e_P(x)^a(1-e_P(x))^{1-a}.
\]
Writing \(Y(0)\) and \(Y(1)\) for the baseline and treated potential-outcome coordinates, respectively, the observed outcome satisfies
\[
Y=(1-A)Y(0)+AY(1).
\]
Changes to the conditional outcome kernels on \(\lambda\)-null sets leave \(\mathbb Q_P\) unchanged.
\end{definitionv}

Thus the null consists precisely of model laws with \(d(P)=0\). Under the causal completion, its target is constancy of the conditional average treatment effect.

\subsection{The original-record testing experiment}

Testing uses an independent sample \leanref{sym:Data}{\ensuremath{\mathcal D_n=((X_i,A_i,Y_i))_{i=1}^n}} from \(P\), together with an independent public uniform variable \leanref{sym:U}{\ensuremath{U}} on \([0,1]\). An \leanref{S-2}{original-record randomized test} is a jointly Borel rejection-probability map \leanref{sym:phi}{\ensuremath{\phi}} with domain and range
\[
\phi:
\bigl([0,1]\times\{0,1\}\times\mathbb R\bigr)^n\times[0,1]
\longrightarrow[0,1].
\]
Its expected rejection probability is denoted by \leanref{sym:Risk}{\ensuremath{\mathsf R_n(P,\phi)}}, the expectation over both the sample and the public randomization. The experiment and this expectation are specified together in \cref{obj:def:original-record-experiment}.

\begin{definitionv}{def:original-record-experiment}[Original-record experiment \(\mathsf S_{n,P}\)]
The original-record experiment is the probability law
\[
\mathsf S_{n,P}=P^{\otimes n}\otimes\lambda
\]
on
\[
\bigl([0,1]\times\{0,1\}\times\mathbb R\bigr)^n\times[0,1],
\]
where \(P\) is the original-record law and \(\lambda\) is Lebesgue probability measure on \([0,1]\). Each original record has coordinates \(X\in[0,1]\), \(A\in\{0,1\}\), and \(Y\in\mathbb R\), representing the covariate, treatment, and observed outcome. Write
\[
\mathcal D_n=\bigl((X_i,A_i,Y_i)\bigr)_{i=1}^n
\]
for the independent sample and \(U\) for the independent public uniform randomization variable. For a randomized rejection-probability map \(\phi\), its rejection expectation is
\[
\mathsf R_n(P,\phi)
=
\int\phi(\mathcal D_n,u)\,
\mathsf S_{n,P}(d\mathcal D_n,du).
\]
\end{definitionv}

Tests may depend on the public tuple and sample size. Their size is controlled uniformly over all laws in the constant-effect null, including every admissible common effect level and nuisance function.

Let \leanref{sym:D}{\ensuremath{D_v}} denote the original-model distance supremum,
\[
D_v=\sup_{P\in\mathcal M_v}d(P),
\]
and let \leanref{sym:r}{\ensuremath{r}} be a separation satisfying \(0<r<D_v\). The minimax type-II error \leanref{sym:B}{\ensuremath{B_n(v,r)}} in \cref{obj:def:testing-risk} compares tests of size at most \(1/10\) against every model law whose population heterogeneity distance is at least \(r\).

\begin{definitionv}{def:testing-risk}[Minimax testing risk $B_n(v,r)$]
\[
B_n(v,r)=
\inf_{\phi:\ \sup_{P\in H_0(v)}\mathsf R_n(P,\phi)\le 1/10}
\ \sup_{P\in\mathcal M_v:\ d(P)\ge r}
\bigl\{1-\mathsf R_n(P,\phi)\bigr\}.
\]
The parameters satisfy:
\begin{itemize}
\item \textbf{(Sample size.)} \(n\ge 2\).
\item \textbf{(Public tuple.)} \(v\in\mathcal V\), the admissible domain of moment and primitive smoothness tuples.
\item \textbf{(Separation.)} \(0<r<D_v\), where \(D_v\) is the supremum of \(d(P)\) over \(\mathcal M_v\).
\end{itemize}
Here \(\mathcal M_v\) is the original-record model, \(H_0(v)\) consists of its laws with an unknown constant mean effect, and \(d(P)\) is the population Lebesgue \(L_2\) distance of the effect from its average. The infimum ranges over randomized rejection-probability maps \(\phi\), and \(\mathsf R_n(P,\phi)\) is their expected rejection probability in the original-record experiment.
\end{definitionv}

The size constraint and the alternative supremum give the risk a uniform interpretation over the stipulated model. To summarize the separation at which this risk reaches \(1/10\), we use the capped critical radius \leanref{sym:rstar}{\ensuremath{r_n^*(v)}} in \cref{obj:def:critical-radius}.

\begin{definitionv}{def:critical-radius}[Capped critical radius $r_n^*(v)$]
\[
r_n^*(v)=
\inf\Bigl(
\bigl\{r\in(0,D_v):B_n(v,r)\le 1/10\bigr\}
\cup\{D_v\}
\Bigr).
\]
Here \(B_n(v,r)\) is the original-model minimax type-II error under the size constraint \(1/10\), and \(D_v\) is the supremum of the population effect distance over the original model.
\end{definitionv}

Including the model supremum makes the defining set nonempty for every admissible sample size and public tuple. When the risk threshold is achieved at an interior separation, the radius records the infimum of such separations; otherwise, the cap supplies the value \(D_v\).

\subsection{Matched bounded and signed-binary benchmarks}

The benchmark family retains the same design, overlap, smoothness exponents, and mean envelopes while imposing public outcome support restrictions. Its bounded member uses the following unit envelope.

\begin{assumptionv}{ass:bounded-outcome}[Almost sure outcome bound]
The observed outcome \(Y\) satisfies
\[
P(|Y|\le 1)=1.
\]
\end{assumptionv}

The bounded-outcome condition places every observed outcome within a public unit envelope, a fixed-envelope version of bounded-outcome regularity used in homogeneity testing \citep{LapentaStrittmatterVergara2026Omnibus}. The bounded model \leanref{sym:Mbound}{\ensuremath{\mathcal M^{\mathrm b}_w}} in \cref{obj:def:bounded-model} combines this support restriction with the primitive restrictions at moment order two.

\begin{definitionv}{def:bounded-model}[Bounded-outcome model $\mathcal M^{\mathrm b}_w$]
The bounded-outcome model \(\mathcal M^{\mathrm b}_w\) is defined for the matched smoothness tuple \(w=(\alpha,\beta,\gamma)\) in \(\mathcal W\), the admissible domain of matched smoothness tuples, by
\[
\mathcal M^{\mathrm b}_w
=
\left\{
P\in\mathcal M_{(2,\alpha,\beta,\gamma)}:
P(|Y|\le 1)=1
\right\},
\qquad
w=(\alpha,\beta,\gamma)\in\mathcal W.
\]
Its restrictions are given in \cref{obj:ass:uniform,obj:ass:overlap,obj:ass:propensity-smoothness,obj:ass:baseline-smoothness,obj:ass:effect-smoothness,obj:ass:baseline-cap,obj:ass:effect-cap,obj:ass:raw-moment,obj:ass:bounded-outcome}.
\end{definitionv}

The signed-binary member specializes the support further.

\begin{assumptionv}{ass:binary-outcome}[Signed binary outcome support]
The observed outcome \(Y\) satisfies
\[
P\bigl(Y\in\{-1,1\}\bigr)=1.
\]
\end{assumptionv}

Signed-binary support is the affine recoding \(Y=2B-1\) of a conventional binary outcome \(B\in\{0,1\}\), the outcome convention used in the binary conditional-effect model of \citet[Section 2]{KennedyBalakrishnanRobinsWasserman2024MinimaxCATE}. Either support restriction implies a conditional second raw moment bounded by one almost everywhere, and hence satisfies the moment envelope at \(p=2\). The signed-binary model \leanref{sym:Mbin}{\ensuremath{\mathcal M^{\mathrm{bin}}_w}} in \cref{obj:def:binary-model} is consequently indexed by the matched smoothness tuple.

\begin{definitionv}{def:binary-model}[Signed-binary model $\mathcal M^{\mathrm{bin}}_w$]
The signed-binary model \(\mathcal M^{\mathrm{bin}}_w\) is defined for the matched smoothness tuple \(w=(\alpha,\beta,\gamma)\) in \(\mathcal W\), the admissible domain of matched smoothness tuples, by
\[
\mathcal M^{\mathrm{bin}}_w
=
\left\{
P\in\mathcal M_{(2,\alpha,\beta,\gamma)}:
P(Y\in\{-1,1\})=1
\right\},
\qquad
w=(\alpha,\beta,\gamma)\in\mathcal W.
\]
Its restrictions are given in \cref{obj:ass:uniform,obj:ass:overlap,obj:ass:propensity-smoothness,obj:ass:baseline-smoothness,obj:ass:effect-smoothness,obj:ass:baseline-cap,obj:ass:effect-cap,obj:ass:raw-moment,obj:ass:binary-outcome}.
\end{definitionv}

These definitions give
\[
\mathcal M^{\mathrm{bin}}_w
\subseteq
\mathcal M^{\mathrm b}_w
\subseteq
\mathcal M_{(2,\alpha,\beta,\gamma)}.
\]
Both benchmarks preserve the unknown common effect as a nuisance parameter. Their nulls, denoted by \leanref{sym:Hnullbound}{\ensuremath{H_0^{\mathrm b}(w)}} and \leanref{sym:Hnullbin}{\ensuremath{H_0^{\mathrm{bin}}(w)}}, are defined in \cref{obj:def:bounded-null,obj:def:binary-null}.

\begin{definitionv}{def:bounded-null}[Bounded constant-effect null $H_0^{\mathrm b}(w)$]
\[
H_0^{\mathrm b}(w)
=
\left\{
P\in\mathcal M^{\mathrm b}_w:
\exists c\in[-1/2,1/2]\ 
\forall x\in[0,1],\ \tau_P(x)=c
\right\}.
\]
The model restrictions and null constancy condition are given in \cref{obj:ass:uniform,obj:ass:overlap,obj:ass:propensity-smoothness,obj:ass:baseline-smoothness,obj:ass:effect-smoothness,obj:ass:baseline-cap,obj:ass:effect-cap,obj:ass:raw-moment,obj:ass:bounded-outcome,obj:ass:null-constancy}.
\end{definitionv}

\begin{definitionv}{def:binary-null}[Signed-binary constant-effect null $H_0^{\mathrm{bin}}(w)$]
\[
H_0^{\mathrm{bin}}(w)
=
\left\{
P\in\mathcal M^{\mathrm{bin}}_w:
\exists c\in[-1/2,1/2]\ 
\forall x\in[0,1],\ \tau_P(x)=c
\right\}.
\]
The model restrictions and null constancy condition are given in \cref{obj:ass:uniform,obj:ass:overlap,obj:ass:propensity-smoothness,obj:ass:baseline-smoothness,obj:ass:effect-smoothness,obj:ass:baseline-cap,obj:ass:effect-cap,obj:ass:raw-moment,obj:ass:binary-outcome,obj:ass:null-constancy}.
\end{definitionv}

For their separation domains, define the bounded and signed-binary distance suprema \leanref{sym:Dbound}{\ensuremath{D_w^{\mathrm b}}} and \leanref{sym:Dbin}{\ensuremath{D_w^{\mathrm{bin}}}} by
\[
D_w^{\mathrm b}
=
\sup_{P\in\mathcal M^{\mathrm b}_w}d(P),
\qquad
D_w^{\mathrm{bin}}
=
\sup_{P\in\mathcal M^{\mathrm{bin}}_w}d(P).
\]
The corresponding minimax type-II errors \leanref{sym:Bbound}{\ensuremath{B_n^{\mathrm b}(w,r)}} and \leanref{sym:Bbin}{\ensuremath{B_n^{\mathrm{bin}}(w,r)}}, defined in \cref{obj:def:bounded-risk,obj:def:binary-risk}, use the same original-record experiment, population distance, and size threshold.

\begin{definitionv}{def:bounded-risk}[Bounded-model testing risk]
The bounded-model minimax type-II error \(B_n^{\mathrm b}(w,r)\) is defined by
\[
B_n^{\mathrm b}(w,r)
=
\inf_{\phi:\,\sup_{P\in H_0^{\mathrm b}(w)}\mathsf R_n(P,\phi)\le 1/10}
\sup_{P\in\mathcal M^{\mathrm b}_w:\,d(P)\ge r}
\{1-\mathsf R_n(P,\phi)\},
\]
where the tests \(\phi\), rejection probabilities \(\mathsf R_n(P,\phi)\), and heterogeneity distances \(d(P)\) are as in \cref{obj:def:testing-risk}, and the bounded model and null are those of \cref{obj:def:bounded-model,obj:def:bounded-null}. The parameter domain is:
\begin{itemize}
\item \textbf{(Sample size.)} \(n\in\mathbb N\) with \(n\ge2\).
\item \textbf{(Smoothness.)} \(w=(\alpha,\beta,\gamma)\in\mathcal W\), where the admissible smoothness domain is
\[
\mathcal W
=
\{(\alpha,\beta,\gamma)\in\mathbb R^3:
0<\alpha\le1,\quad 0<\beta\le1,\quad 1/4\le\gamma\le1\}.
\]
\item \textbf{(Separation.)} \(0<r<D_w^{\mathrm b}\), where the bounded-model distance supremum is
\[
D_w^{\mathrm b}
=
\sup_{P\in\mathcal M^{\mathrm b}_w}d(P).
\]
\end{itemize}
\end{definitionv}

\begin{definitionv}{def:binary-risk}[Signed-binary minimax testing risk]
The signed-binary minimax type-II error \(B_n^{\mathrm{bin}}(w,r)\) is defined by
\[
B_n^{\mathrm{bin}}(w,r)
=
\inf_{\phi:\,\sup_{P\in H_0^{\mathrm{bin}}(w)}
\mathsf R_n(P,\phi)\le 1/10}
\sup_{P\in\mathcal M^{\mathrm{bin}}_w:\,d(P)\ge r}
\{1-\mathsf R_n(P,\phi)\},
\]
where the tests and their expected rejection probabilities \(\mathsf R_n(P,\phi)\) follow \cref{obj:def:testing-risk}, and the signed-binary model and null are specified in \cref{obj:def:binary-model,obj:def:binary-null}. Here \(\lambda\) denotes Lebesgue probability measure on \([0,1]\), and the population heterogeneity distance is
\[
d(P)
=
\left(
\int_{[0,1]}
\left\{\tau_P(x)-\int_{[0,1]}\tau_P(z)\,d\lambda(z)\right\}^{2}
\,d\lambda(x)
\right)^{1/2}.
\]
The public parameter domain is:
\begin{itemize}
\item \textbf{(Sample size.)} \(n\in\mathbb N\) with \(n\ge2\).
\item \textbf{(Smoothness.)} \(w=(\alpha,\beta,\gamma)\in\mathcal W\), where
\[
\mathcal W=(0,1]\times(0,1]\times[1/4,1].
\]
\item \textbf{(Separation.)} \(0<r<D_w^{\mathrm{bin}}\), where the supremum of legal signed-binary heterogeneity distances is
\[
D_w^{\mathrm{bin}}
=
\sup_{P\in\mathcal M^{\mathrm{bin}}_w}d(P).
\]
\end{itemize}
\end{definitionv}

Finally, the benchmark critical radii \leanref{sym:rstarbound}{\ensuremath{r_n^{*,\mathrm b}(w)}} and \leanref{sym:rstarbin}{\ensuremath{r_n^{*,\mathrm{bin}}(w)}} in \cref{obj:def:bounded-radius,obj:def:binary-radius} apply the same capping convention to the respective risks.

\begin{definitionv}{def:bounded-radius}[Bounded-model critical separation radius]
The capped bounded-model critical separation radius \(r_n^{*,\mathrm b}(w)\) is defined for \(n\in\mathbb N\) and \(w=(\alpha,\beta,\gamma)\in\mathbb R^3\). Write \(d(P)\) for the population Lebesgue \(L_2\) distance of the mean effect under \(P\) from its average, and define the bounded-model cap by
\[
D_w^{\mathrm b}
=
\sup_{P\in\mathcal M^{\mathrm b}_w} d(P),
\]
where \(\mathcal M^{\mathrm b}_w\) is the bounded model in \cref{obj:def:bounded-model}. With \(B_n^{\mathrm b}(w,r)\), the bounded-model testing risk defined in \cref{obj:def:bounded-risk}, set
\[
r_n^{*,\mathrm b}(w)
=
\inf\left(
\left\{r\in(0,D_w^{\mathrm b}):B_n^{\mathrm b}(w,r)\le \frac{1}{10}\right\}
\cup\{D_w^{\mathrm b}\}
\right).
\]
\end{definitionv}

\begin{definitionv}{def:binary-radius}[Capped signed-binary critical radius]
The capped signed-binary critical separation radius \(r_n^{*,\mathrm{bin}}(w)\), for sample size \(n\in\mathbb N\) and primitive smoothness tuple \(w=(\alpha,\beta,\gamma)\in\mathbb R^3\), is
\[
r_n^{*,\mathrm{bin}}(w)
=
\inf\left(
\left\{r\in(0,D_w^{\mathrm{bin}}):
B_n^{\mathrm{bin}}(w,r)\le \frac{1}{10}\right\}
\cup\{D_w^{\mathrm{bin}}\}
\right),
\]
where \(D_w^{\mathrm{bin}}\) is the supremum of the population heterogeneity distances over the signed-binary model in \cref{obj:def:binary-model}, and \(B_n^{\mathrm{bin}}(w,r)\) is the signed-binary testing risk in \cref{obj:def:binary-risk}.
\end{definitionv}

Throughout the benchmark comparisons, these quantities are evaluated at \(n\ge2\) and \(w\in\mathcal W\). Holding the primitive smoothness tuple fixed aligns the population target and error thresholds across the original, bounded, and signed-binary families, providing a common basis for assessing the role of outcome tails.


\section{Separation rates and the cost of outcome tails}

Holding the primitive smoothness tuple fixed gives a direct comparison of the precision attainable under the finite conditional moment restriction and the bounded-outcome benchmark. The comparison concerns the population heterogeneity distance and the capped critical radii defined in \cref{obj:def:critical-radius,obj:def:bounded-radius}. Two exponents determine the original-record separation rate: one depends on effect smoothness and moment order, while the other also reflects the smoothness of the propensity and baseline mean.

We collect the scales and their finite-sample multipliers before stating the comparison. Write \(C^{\mathrm{att}}_v\) and \(N^{\mathrm{att}}_v\) for the fully evaluable attainment multiplier and nonempty-power threshold specified in \cref{obj:thm:original-record-attainable-rate}; evaluation at \((2,w)\) gives their matched bounded counterparts. The following definition combines these attainment quantities with the lower-bound multipliers.

\begin{definitionv}{def:sharp-frontier-scales}[Separation scales \(\rho_n(v)\) and \(\rho_n^{\mathrm b}(w)\)]
The separation scales and lower multipliers are
\[
 \rho_n(v)=n^{-E(p)},\qquad
 \rho_n^{\mathrm b}(w)=n^{-E(2)},\qquad
 c_w^{\mathrm b}=c_{(2,w)},
\]
\[
 c_v=
 \begin{cases}
 c_0^{\mathrm e},&F(p)\ge1,\\
 \min\!\left\{k_v,\,
 c_0^{\mathrm e}N_{\mathrm{low}}(v)^{-(E_0(p)-E_4(p))}\right\},
 &F(p)<1.
 \end{cases}
\]
For an admissible public tuple \(v=(p,\alpha,\beta,\gamma)\in\mathcal V\) and its matched smoothness tuple \(w=(\alpha,\beta,\gamma)\), the quantities in these formulas are defined by
\[
 q=\frac{p-1}{p},\qquad S=\alpha+\beta,\qquad
 D_p=1+2\alpha+\frac{\beta}{q}+\frac{S}{2\gamma},
\]
\[
 E_0(p)=\frac{2\gamma q}{2\gamma+q},\qquad
 E_4(p)=\frac{2S}{D_p},\qquad
 F(p)=\frac{2\alpha}{p-1}+\frac{\beta}{q}+\frac{S}{2\gamma},
 \qquad E(p)=\min\{E_0(p),E_4(p)\},
\]
\[
 H_{\mathrm{fine}}=2^{40},\qquad
 N_{\mathrm{low}}(v)=
 \left\lceil\max\left\{2,32^{D_p\gamma/(2S)}\right\}\right\rceil,
\]
\[
 k_v=2^{-17}(2H_{\mathrm{fine}})^{-S},\qquad
 c_0^{\mathrm e}=\frac{4^{-\gamma}}{16\sqrt3}.
\]
The arguments \(p\) in \(E_0,E_4,F,E\) vary with \((\alpha,\beta,\gamma)\) fixed; \(c_{(2,w)}\) is \(c_v\) evaluated at \(v=(2,\alpha,\beta,\gamma)\).

The upper multipliers and sample-size thresholds are the fully evaluable attainment quantities:
\[
 C_v=C^{\mathrm{att}}_v,\qquad N_v=N^{\mathrm{att}}_v,\qquad
 C_w^{\mathrm b}=C^{\mathrm{att}}_{(2,w)},\qquad
 N_w^{\mathrm b}=N^{\mathrm{att}}_{(2,w)}.
\]
The test \(\phi^*_{n,v}\) is the total original-record test specified by the explicit score construction, and its bounded-model counterpart is
\[
 \phi^{*,\mathrm b}_{n,w}=\phi^*_{n,(2,w)}.
\]
\end{definitionv}

The original-model scale \leanref{sym:rho}{\ensuremath{\rho_n(v)}} and the matched bounded scale \leanref{sym:rhobound}{\ensuremath{\rho_n^{\mathrm b}(w)}} summarize the sample-size exponents. Their lower multipliers, \leanref{sym:clower}{\ensuremath{c_v}} and \leanref{sym:clowerbound}{\ensuremath{c_w^{\mathrm b}}}, and upper multipliers, \leanref{sym:Cupper}{\ensuremath{C_v}} and \leanref{sym:Cupperbound}{\ensuremath{C_w^{\mathrm b}}}, retain the finite-sample content of the comparison. The associated public thresholds \leanref{sym:Npublic}{\ensuremath{N_v}} and \leanref{sym:Npublicbound}{\ensuremath{N_w^{\mathrm b}}} place the upper separations below a fixed attainable heterogeneity distance. Thus the power statements below are attached to an explicit nonempty separation domain.

To measure the cost of outcome tails at the same smoothness tuple, we use the ratio of the capped radii.

\begin{definitionv}{def:comparison-handle}[Capped-radius ratio $\Delta_n(v)$]
\[
\Delta_n(v)=\frac{r_n^*(v)}{r_n^{*,\mathrm b}(w)},
\qquad w=(\alpha,\beta,\gamma).
\]
The numerator and denominator are the capped critical separation radii for the original and bounded original-record experiments, respectively, at the same sample size \(n\) and matched smoothness tuple \(w\).
\end{definitionv}

The ratio \leanref{sym:Delta}{\ensuremath{\Delta_n(v)}} in \cref{obj:def:comparison-handle} compares experiments with the same mean-effect target and error thresholds. We distinguish persistent enlargement of the original-model radius from a comparison that remains bounded over all sample sizes.

\begin{definitionv}{def:tail-region}[Tail region \(\mathcal T\)]
The tail region \(\mathcal T\) consists of admissible public tuples \(v=(p,\alpha,\beta,\gamma)\), with moment exponent \(p\) and primitive smoothness exponents \(\alpha,\beta,\gamma\), for which
\[
\mathcal T
=
\left\{
v\in\mathcal V:
p<2,\ 
\lim_{n\to\infty}\Delta_n(v)=\infty
\right\},
\]
where \(\mathcal V\) is the admissible tuple domain and \(\Delta_n(v)\) is the original-to-bounded capped-radius ratio at matched smoothness.
\end{definitionv}

\begin{definitionv}{def:equality-region}[Equality region \(\mathcal E\)]
The equality region \(\mathcal E\) is defined using \(\mathcal V\), the admissible domain of public moment and smoothness tuples, by
\[
\mathcal E
=
\left\{
v\in\mathcal V:
\sup_{n\ge 2}\Delta_n(v)<\infty
\right\},
\]
where \(\Delta_n(v)\) is the original-to-bounded capped-radius ratio at matched smoothness.
\end{definitionv}

The tail region \leanref{sym:TailRegion}{\ensuremath{\mathcal T}} and equality region \leanref{sym:EqualityRegion}{\ensuremath{\mathcal E}} describe the behavior of this matched comparison. The equality region concerns bounded radius ratios; its definition permits different finite-sample constants.

Before comparing rates, the bounded benchmark can be identified exactly with its signed-binary submodel. Let \(\lambda\) denote Lebesgue probability measure on \([0,1]\), and let \leanref{sym:dzero}{\ensuremath{d_0=1/(8\sqrt{2})}} be the fixed positive witness distance supplied by \cref{obj:lem:causal-nonempty}. The next result shows that both benchmarks have the same attainable heterogeneity distances and the same minimax testing risks.

\begin{propositionv}{prop:bounded-submodel-comparison}[Bounded and binary submodel comparison]
Let \(v=(p,\alpha,\beta,\gamma)\) be an admissible parameter tuple as in \cref{obj:def:model}, and let \(w=(\alpha,\beta,\gamma)\). Use the model classes in \cref{obj:def:binary-model,obj:def:bounded-model,obj:def:model}. Let \(\lambda\) denote Lebesgue probability measure on \([0,1]\), and define the heterogeneity distance, its model suprema, and the witness distance by
\[
d(P)=
\left[
\int_{[0,1]}
\left(\tau_P(x)-\int_{[0,1]}\tau_P(z)\,\lambda(dz)\right)^2
\,\lambda(dx)
\right]^{1/2},
\]
\[
D_w^{\mathrm{bin}}=\sup_{P\in\mathcal M^{\mathrm{bin}}_w}d(P),
\qquad
D_w^{\mathrm b}=\sup_{P\in\mathcal M^{\mathrm b}_w}d(P),
\qquad
D_v=\sup_{P\in\mathcal M_v}d(P),
\qquad
d_0=\frac{1}{8\sqrt{2}}.
\]
Then
\[
\mathcal M^{\mathrm{bin}}_w\subseteq\mathcal M^{\mathrm b}_w
\subseteq\mathcal M_v,
\qquad
D_w^{\mathrm{bin}}=D_w^{\mathrm b}=D_v\ge d_0.
\]
For every integer \(n\ge2\) and every real \(r\) satisfying \(0<r<D_v\), the testing risks of \cref{obj:def:binary-risk,obj:def:bounded-risk,obj:def:testing-risk} satisfy
\[
B_n^{\mathrm{bin}}(w,r)=B_n^{\mathrm b}(w,r)\le B_n(v,r).
\]
For every integer \(n\ge2\), the capped critical radii of \cref{obj:def:binary-radius,obj:def:bounded-radius,obj:def:critical-radius} satisfy
\[
0<r_n^{*,\mathrm{bin}}(w)
=r_n^{*,\mathrm b}(w)\le r_n^*(v).
\]

For every \(P\in\mathcal M_{(2,\alpha,\beta,\gamma)}\), its signed-binary conversion belongs to \(\mathcal M^{\mathrm{bin}}_w\). This conversion retains the distribution of \((X,A)\) and uses conditional signed-binary outcomes whose probabilities of \(+1\) in arms \(0\) and \(1\), respectively, are
\[
\frac{1+m_{0,P}(x)}{2},
\qquad
\frac{1+m_{1,P}(x)}{2},
\]
with the complementary probabilities assigned to \(-1\). It preserves the propensity \(e_P\), the baseline mean \(m_{0,P}\), the effect \(\tau_P\), and the heterogeneity distance \(d(P)\). The converted law satisfies the null constancy condition of \cref{obj:ass:null-constancy} if and only if the original law does.

For every \(n\in\mathbb N\), every test \(\phi\) in the sense of \cref{obj:def:testing-risk}, and every \(P\in\mathcal M^{\mathrm b}_w\), independently convert each original record by retaining \((X,A)\) and drawing an outcome sign with conditional probabilities
\[
\Pr(+1\mid X,A,Y)=\frac{1+Y}{2},
\qquad
\Pr(-1\mid X,A,Y)=\frac{1-Y}{2}.
\]
Integrating this conversion randomization into \(\phi\) gives a rejection-probability map on the original experiment. Its expected rejection probability under \(P\) equals the expected rejection probability of \(\phi\) under the signed-binary conversion of \(P\).
\end{propositionv}

The sign randomization in \cref{obj:prop:bounded-submodel-comparison} preserves the conditional outcome means, and hence the effect function, its heterogeneity distance, and null membership. Integrating the randomization into a test transfers its rejection probabilities to the bounded experiment. The resulting equality of minimax risks holds throughout the common nonempty separation domain \(0<r<D_v\), and equality of capped radii holds for every \(n\ge2\). This gives an exact decision-theoretic benchmark for the mean-effect target.

The lower-bound comparison uses finite mixtures of original-record laws. Their public tuning selects a copula-frame construction in the rough-confounding regime once its sample-size threshold is reached, and an opposite paired-tent construction otherwise. The following selection rule fixes the priors used in the main result; their support laws and complete-record comparison bounds are developed in \cref{obj:lem:paired-tent-full-record-lower,obj:def:copula-frame-priors,obj:lem:copula-frame-legality,obj:lem:full-record-copula-component-bound,obj:lem:rough-full-record-sharp-lower}.



The selected null and alternative priors \leanref{sym:piNull}{\ensuremath{\pi_{0,n,v}}} and \leanref{sym:piAlt}{\ensuremath{\pi_{1,n,v}}} in \cref{obj:def:converse-handle} induce the complete independent-record mixtures \leanref{sym:MixNull}{\ensuremath{\mathbb P_{0,n,v}}} and \leanref{sym:MixAlt}{\ensuremath{\mathbb P_{1,n,v}}}. Their rare-mark probability \leanref{sym:rarity}{\ensuremath{\varepsilon_{n,v}}} and outcome magnitude \leanref{sym:magnitude}{\ensuremath{L_{n,v}}} describe how the lower witnesses use the finite moment allowance.

For the next statement, \(\mathsf R_n(P,\phi)\) denotes expected rejection in the original-record experiment, \(Q_{a,P}(dy\mid x)\) denotes the conditional outcome probability kernel in arm \(a\), and \leanref{sym:TV}{\ensuremath{\operatorname{TV}}} denotes the supremum of the absolute difference in event probabilities. The threshold \(N(v)\) identifies the sample sizes at which every bounded-model finite prior pair at the original-model lower separation has mixture total variation at least one half. The theorem combines the matched separation bounds, the power guarantees, and this tail-essential comparison.



The exponent comparison in \cref{obj:thm:heavy-tail-comparison} explains the two statistical regimes. When \(F(p)\ge1\), the effect-smoothness exponent \(E_0(p)\) determines the separation scale: moment order and effect smoothness govern its sample-size dependence. When \(F(p)<1\), the rough-confounding exponent \(E_4(p)\) determines the scale, bringing the propensity and baseline smoothness exponents into the rate. The identity for \(E_4(p)-E_0(p)\) makes \(F(p)=1\) their exact equality boundary. Across both regimes, the strict inequality \(E(p)<E(2)\) for \(p<2\) translates the finite moment allowance into a polynomial enlargement of the matched critical radius.

The two denominators reflect different balances. At histogram rank \(J\), approximating the effect contributes the usual Hölder term of order \(J^{-\gamma}\); clipping at level \(T\) contributes a moment-tail term of order \(T^{1-p}\), while the clipped quadratic variability grows with \(T^{2-p}\). Balancing these terms gives \(E_0(p)\). When the nuisance functions are rough, the ordered-pair correction removes the first-order propensity error, leaving a product of the fine-scale propensity and baseline approximation errors. Its rank cost depends on \(2\alpha+\beta/q\): propensity regularity enters twice through treatment weighting and projection correction, whereas the baseline appears once and is amplified by the tail factor \(1/q\). The additional term \(S/(2\gamma)\) records the coarse resolution needed to retain the effect signal. Correction-specific clipping prevents the fine increments from paying the main cutoff at every resolution, which is what makes the multiresolution covariance sum compatible with \(E_4(p)\). These statements summarize the verified mean and covariance bounds in \cref{obj:lem:original-score-mean-ledger,obj:lem:original-score-covariance-ledger}; their full calculations are in the attainment appendix.

For example, \(v=(3/2,1,1,1)\) has \(F(p)=8\), so the effect channel governs with \(E=E_0=2/7\). In contrast, \(v=(3/2,1/20,1/20,1)\) has \(F(p)=2/5\), so rough confounding governs with \(E=E_4=2/13\). The first tuple uses smooth nuisance functions and the second deliberately makes them much rougher than the effect. These statistical regimes are distinct from the procedure's representation choice in \cref{obj:def:explicit-score-ledger}, which uses the conditions \(S\ge\gamma\) or \(\alpha\ge q\) to choose between a single correction and a multiresolution expansion.

The regimes can be summarized by substituting the formulas in \cref{obj:def:sharp-frontier-scales}. In the following table, \(q=(p-1)/p\) and \(S=\alpha+\beta\), with the primitive smoothness tuple held fixed.
\begin{center}
\begin{tabular}{lll}
\hline
Model and regime & Condition & Separation exponent \\
\hline
Original, effect smoothness
& \(F(p)>1\)
& \(\displaystyle E_0(p)=\frac{2\gamma q}{2\gamma+q}\) \\[6pt]
Original, equality boundary
& \(F(p)=1\)
& \(\displaystyle E_0(p)=E_4(p)=\frac{2\gamma q}{2\gamma+q}\) \\[6pt]
Original, rough confounding
& \(F(p)<1\)
& \(\displaystyle E_4(p)=\frac{2S}{1+2\alpha+\beta/q+S/(2\gamma)}\) \\[6pt]
Bounded, effect smoothness or boundary
& \(F(2)\ge1\)
& \(\displaystyle E_0(2)=\frac{2\gamma}{4\gamma+1}\) \\[6pt]
Bounded, rough confounding
& \(F(2)<1\)
& \(\displaystyle E_4(2)=\frac{2S}{1+2S+S/(2\gamma)}\) \\[6pt]
\hline
\end{tabular}
\end{center}
Here \(F(2)=2S+S/(2\gamma)\). Each row describes the power of \(n\) in the corresponding scale; the finite-sample radius bounds retain the multipliers and capping convention of \cref{obj:thm:heavy-tail-comparison}. The original and bounded exponents are selected separately at their respective moment orders.

For every fixed admissible tuple, the two-sided bounds on \(\Delta_n(v)\) quantify the cost of tails through \(n^{E(2)-E(p)}\). The moment-order-two face has equal original and bounded exponents and a uniformly bounded ratio over \(n\ge2\), together with the ordering \(\Delta_n(v)\ge1\). Every admissible tuple with \(p<2\) instead belongs to the tail region, where the ratio diverges. This distinction between moment orders is separate from the equality boundary \(F(p)=1\), which identifies the meeting of the two rate mechanisms within a given moment class.

The finite-sample cap remains part of the interpretation. The radius inequalities hold for every \(n\ge2\), while the power guarantees for the attaining tests \leanref{sym:phistar}{\ensuremath{\phi^*_{n,v}}} and \leanref{sym:phistarbound}{\ensuremath{\phi^{*,\mathrm b}_{n,w}}} from \cref{obj:def:sharp-frontier-scales} apply when their upper separations lie below the corresponding model distance suprema. The public thresholds ensure this condition through the fixed witness distance \(d_0\). Size is controlled over the whole null at every \(n\ge2\). On compact subsets of the admissible exponent domains, common lower and upper multipliers and a common bound on the public thresholds make these comparisons locally uniform.

Finally, the tail-essential witnesses identify the statistical content of the enlargement. For \(p<2\), their outcome magnitudes grow with sample size, while the selected null and separated alternative mixtures remain within total variation one half in the complete-record experiment. At the same original-model lower separation, every bounded-model finite prior pair has mixture total variation at least one half once the tuple-dependent threshold is reached. Thus the witnesses connect the slower finite-moment rate to outcome tails while preserving the stipulated primitive restrictions and the mean-effect target. The full support-law constructions and the quantified finite-prior comparison appear in the proof appendix through \cref{obj:def:marked-table,obj:def:frame-handle,obj:def:copula-frame-priors,obj:thm:tail-essential-full-record-converse}.


\section{A calibrated test from the original records}

The upper separation guarantee in \cref{obj:thm:heavy-tail-comparison} is attained by a procedure whose inputs are the original records and the public moment and smoothness tuple. The construction combines histogram projection, clipped outcomes, and ordered-pair corrections within two independent evaluation blocks. Minimization over the candidate constant then accommodates the composite null of a homogeneous mean effect. We first specify the projection geometry and score ingredients, and then give the complete calculation.

At a positive integer rank \leanref{sym:J}{\ensuremath{J}}, the histogram projection averages functions over equal-width cells. Its endpoint convention is part of the procedure.

\begin{definitionv}{def:projection}[Histogram projection \(\Pi_J\)]
For a positive integer \(J\), define the histogram cells, feature vector, and projection kernel by
\[
I_{j,J}=
\begin{cases}
[(j-1)/J,j/J),&1\le j<J,\\
[(J-1)/J,1],&j=J,
\end{cases}
\qquad
F_J(x)=\bigl(\sqrt J\,\mathbf 1_{I_{j,J}}(x)\bigr)_{j=1}^J,
\]
\[
\Pi_J(x,z)=\langle F_J(x),F_J(z)\rangle.
\]
The corresponding projection of an integrable function \(f\) is
\[
(\Pi_J f)(x)=\int_0^1\Pi_J(x,z)f(z)\,dz.
\]
\end{definitionv}

The normalized feature vector \leanref{sym:F}{\ensuremath{F_J(x)}} represents cellwise constant functions in Euclidean coordinates, and the projection kernel \leanref{sym:Pi}{\ensuremath{\Pi_J(x,z)}} selects pairs in the same cell. An internal bin endpoint belongs to the cell on its right, while the final cell contains \(1\). These conventions assign a value to every record, including records at cell boundaries.

To obtain independent score vectors, partition the independent sample into two deterministic blocks of equal size.

\begin{definitionv}{def:blocks}[Evaluation blocks \(\mathcal I_1,\mathcal I_2\)]
The two evaluation blocks and their common size are
\[
s_n=\lfloor n/2\rfloor,
\qquad
\mathcal I_1=\{1,\ldots,s_n\},
\qquad
\mathcal I_2=\{s_n+1,\ldots,2s_n\}.
\]
\end{definitionv}

The evaluation blocks \leanref{sym:blocks}{\ensuremath{\mathcal I_1,\mathcal I_2}} each contain \leanref{sym:bsize}{\ensuremath{s_n}} records. When the sample size is odd, the last record is left outside these blocks. Independence between the blocks permits their score inner product to estimate a squared population mean.

Within each block, a single-record score is paired with an ordered-pair correction. Both depend affinely on the candidate constant \(c\in[-1/2,1/2]\). The following ingredients use a coarse histogram rank \leanref{sym:Jcoarse}{\ensuremath{J_{\mathrm c}}}, whose public selection is specified in \cref{obj:def:explicit-score-ledger}.

\begin{definitionv}{def:score-family}[Score families \(H_{b,T}\) and \(U_{b,G,T}\)]
For \(n\ge4\), the single-record and ordered-pair score vectors are
\[
H_{b,T}(c)=
\frac{1}{s_n}\sum_{i\in\mathcal I_b}
F_{J_{\mathrm c}}(X_i)A_i\{\ell_T(Y_i)-c\},
\]
\[
U_{b,G,T}(c)=
\frac{1}{s_n(s_n-1)}
\sum_{\substack{i,j\in\mathcal I_b\\i\ne j}}
F_{J_{\mathrm c}}(X_i)A_iG(X_i,X_j)
\{\ell_T(Y_j)-cA_j\}.
\]
Here \(b\in\{1,2\}\), \(G:[0,1]^2\to\mathbb R\) is the pair-score kernel, and the clipping function is
\[
\ell_T(y)=\max\{-T,\min\{y,T\}\}.
\]
For \(n=2,3\), set
\[
H_{b,T}(c)=U_{b,G,T}(c)=0.
\]
\end{definitionv}

The clipping function \leanref{sym:clip}{\ensuremath{\ell_T(y)}} bounds each outcome contribution before averaging. The single-record score \leanref{sym:Hscore}{\ensuremath{H_{b,T}(c)}} uses treated records, while the ordered-pair score \leanref{sym:Uscore}{\ensuremath{U_{b,G,T}(c)}} combines the treatment indicator of one record with the clipped outcome and treatment indicator of a distinct record. The correction kernel \leanref{sym:G}{\ensuremath{G}} supplies the projection used in this pairing. The explicit dependence on \(c\) makes profiling a finite quadratic calculation.

The construction uses two resolutions for different tasks. The coarse rank \(M\) retains the effect signal tested by the cross-block inner product, while the fine rank \(K\) controls the product of propensity and baseline approximation errors left after the ordered-pair correction. A single correction suffices when \(S\ge\gamma\) or \(\alpha\ge q\). Otherwise the dyadic increments between \(M\) and \(K\) receive separate cutoffs \(T_j\): coarse increments can tolerate more clipping variability, while fine increments use smaller cutoffs to keep their covariance contribution summable. This representation condition determines how the same statistic is assembled; it is separate from the phase condition \(F(p)=1\), which determines which separation exponent is smaller.

The complete procedure now selects these ranks and clipping thresholds, subtracts the pair corrections, and computes the bias and covariance bounds used for calibration. The fine histogram rank \leanref{sym:Jfine}{\ensuremath{J_{\mathrm f}}} governs the correction resolution. Each selection depends on the sample size and public tuple.

\begin{algorithmv}{def:explicit-score-ledger}[Explicit score and ledgers \(Z_b(c),\mathfrak b_{n,v},\Lambda_{n,v}\)]
Inputs are the original records and the public tuple \(v=(p,\alpha,\beta,\gamma)\); \(H_{b,T}(c)\) and \(U_{b,G,T}(c)\) denote the blockwise single-record histogram score and ordered-pair correction score, respectively.
\begin{enumerate}
\item Define the public exponents
\[
q=\frac{p-1}{p},
\qquad
S=\alpha+\beta,
\qquad
t=\frac{2-p}{p-1},
\]
\[
E_0=\frac{2\gamma q}{2\gamma+q},
\qquad
E_4=\frac{2S}{1+2\alpha+\beta/q+S/(2\gamma)},
\qquad
E=\min\{E_0,E_4\}.
\]
For \(n=2,3\), set
\[
J_{\mathrm c}=J_{\mathrm f}=1,
\qquad
T_j=1\quad\text{for every cutoff},
\qquad
Z_b(c)=0,
\]
\[
\mathfrak b_{n,v}=\Lambda_{n,v}=h_{n,v}=0,
\qquad
\phi^*_{n,v}=0.
\]
The remaining steps apply for \(n\ge4\).
\item Let \(s_n\) be the evaluation-block size. Set
\[
s=s_n,
\qquad
a=s^{-E},
\qquad
M=J_{\mathrm c}
=
2^{\left\lceil\log_2(a^{-1/\gamma})\right\rceil},
\]
\[
K=J_{\mathrm f}
=
2^{\left\lceil\log_2\!\left(\max\{M,a^{-1/S}\}\right)\right\rceil}.
\]
Define upper dyadic rounding and the main clipping quantities by
\[
\mathcal R(x)=2^{\lceil\log_2x\rceil},
\quad x>0,
\qquad
T_0=\mathcal R(a^{-1/(p-1)}),
\qquad
V_0=32T_0^{2-p}.
\]
\item If \(S\ge\gamma\) or \(\alpha\ge q\), set
\[
Z_b(c)=H_{b,T_0}(c)-U_{b,\Pi_K,T_0}(c),
\]
\[
B=400K^{-S}+20T_0^{1-p},
\qquad
L_1=16\sqrt{V_0},
\qquad
L_2=4\sqrt{KV_0}.
\]
\item Otherwise, set
\[
\lambda_0=\frac{(p-1)^2}{4p},
\qquad
L=\log_2(K/M),
\qquad
R_j=2^jM\quad(0\le j\le L),
\]
\[
D_j=\Pi_{R_j}-\Pi_{R_{j-1}},
\qquad 1\le j\le L.
\]
For each increment define
\[
T_j=
\mathcal R\!\left(
\max\!\left\{
1,\,
\min\!\left\{
a^{-1/(p-1)},
\left[
\frac{R_j^{-\alpha}}{a(R_j/K)^{\lambda_0}}
\right]^{1/(p-1)}
\right\}
\right\}
\right),
\qquad
V_j=32T_j^{2-p},
\]
\[
a_j=40R_{j-1}^{-\alpha},
\qquad
d_j=
110R_{j-1}^{-\min\{\alpha,\beta,\gamma\}}
+20T_j^{1-p}.
\]
Here \(L\) is an integer, and empty sums equal zero. Set
\[
Z_b(c)=
H_{b,T_0}(c)-U_{b,\Pi_M,T_0}(c)
-\sum_{j=1}^{L}U_{b,D_j,T_j}(c),
\]
\[
B=
400K^{-S}+20T_0^{1-p}
+10\sum_{j=1}^{L}a_jT_j^{1-p},
\]
\[
L_1=
16\sqrt{V_0}
+8\sum_{j=1}^{L}\bigl(d_j+a_j\sqrt{V_j}\bigr),
\]
\[
L_2=
4\left(
\sqrt{MV_0}+\sum_{j=1}^{L}\sqrt{R_jV_j}
\right).
\]
\item In either branch, define the public bias ledger, covariance ledger, and rejection cutoff
\[
\mathfrak b_{n,v}=B,
\qquad
\Lambda_{n,v}
=
\frac{L_1^2}{s}
+\frac{2L_2^2}{s(s-1)},
\]
\[
h_{n,v}
=
\mathcal R\!\left(
2B^2+2^{16}\sqrt M\,\Lambda_{n,v}
\right).
\]
Use the public rank and outcome grids
\[
\left\{2^j:0\le j\le\left\lceil\log_2(2s^2)\right\rceil\right\},
\qquad
\left\{2^j:0\le j\le\left\lceil\log_2(2s)\right\rceil\right\},
\]
respectively. The cutoff is the upper endpoint of the two-point dyadic grid bracketing its unrounded formula.
\item Form the quadratic statistic
\[
W(c)=\langle Z_1(c),Z_2(c)\rangle.
\]
Evaluate its minimum over \([-1/2,1/2]\) at the two endpoints and, when the quadratic coefficient is positive and the vertex is interior, at the vertex. Output
\[
\phi^*_{n,v}
=
\mathbf 1\!\left\{
\min_{|c|\le1/2}W(c)>h_{n,v}
\right\}.
\]
Every sum is finite. Histogram and ordered-pair scores take their stipulated values at bin endpoints and are evaluated directly from the original records, without division by observed cell counts.
\end{enumerate}
\end{algorithmv}

The corrected block score \leanref{sym:Zscore}{\ensuremath{Z_b(c)}} uses either a single fine-rank correction or a sum of dyadic projection increments. In the latter construction, each increment receives its own clipping threshold; the main clipping threshold is \leanref{sym:T}{\ensuremath{T_0}}. This allocation enters both the public bias ledger \leanref{sym:Bias}{\ensuremath{\mathfrak b_{n,v}}} and the public covariance ledger \leanref{sym:Covbound}{\ensuremath{\Lambda_{n,v}}}. The rejection cutoff \leanref{sym:Cutoff}{\ensuremath{h_{n,v}}} combines their contributions and is rounded upward.

The statistical comparison and the computational selection have distinct roles. The comparison of \(E_0\) and \(E_4\), as defined in \cref{obj:def:sharp-frontier-scales,obj:def:explicit-score-ledger}, determines the exponent used for tuning. The condition \(S\ge\gamma\) or \(\alpha\ge q\) selects the single-correction representation; when \(S<\gamma\) and \(\alpha<q\), the procedure uses the multiresolution representation. Both representations use the same minimum exponent \(E\).

All calculations in \cref{obj:def:explicit-score-ledger} are finite. The public dyadic grids specify the rank and outcome selections, and the upper-dyadic rule specifies the rejection cutoff. Scores are normalized by the block size and number of ordered pairs. Consequently, empty histogram cells and observations at bin endpoints have stipulated values throughout the calculation. For \(n=2,3\), the zero-score convention and zero rejection rule complete the procedure.

The cross-block statistic \leanref{sym:Wscore}{\ensuremath{W(c)}} is quadratic because each block score is affine in \(c\). Its minimum is therefore obtained by comparing the endpoints and, when applicable, the interior vertex. Rejection requires the statistic to exceed its cutoff for every admissible constant, which gives the profiling step its role in testing the whole null.

For the calibration statement, write \(\psi_{n,v}:=\phi^*_{n,v}\) for the rule in \cref{obj:def:explicit-score-ledger}, and define the population score mean \leanref{sym:theta}{\ensuremath{\theta_P(c)}} by
\[
\theta_P(c):=\mathbb E_{P^{\otimes n}}[Z_b(c)].
\]
Let \(\lambda\) be Lebesgue probability measure on \([0,1]\), and use
\[
d(P):=
\left\{\int_0^1
\left(\tau_P(x)-\int_0^1\tau_P(t)\,d\lambda(t)\right)^2
\,d\lambda(x)\right\}^{1/2},
\qquad
D_v:=\sup_{P\in\mathcal M_v}d(P),
\]
for the heterogeneity distance and its model supremum. Set \(C:=128\) for the simultaneous profiling multiplier, and define the calibrated separation by
\[
R_{\mathrm{att}}(n,v):=\frac{16}{3}
\left\{
16J_{\mathrm c}^{-\gamma}
+5\mathfrak b_{n,v}
+1024J_{\mathrm c}^{1/4}\sqrt{\Lambda_{n,v}}
\right\}.
\]
The following result connects these public quantities to uniform size control and power.

\begin{theoremv}{thm:whole-null-calibration-resolved}[Whole-null calibration and power]
Let \(v=(p,\alpha,\beta,\gamma)\) be an admissible public tuple as in \cref{obj:def:model}, and let \(n\ge2\) be an integer. Use the ranks, cutoffs, scores, bias ledger \(\mathfrak b_{n,v}\), and covariance ledger \(\Lambda_{n,v}\) of \cref{obj:def:explicit-score-ledger}. Write \(\psi_{n,v}\) for the original-record test constructed there, and define its rejection expectation by
\[
\mathsf R_n(P,\psi_{n,v})
:=\mathbb E\!\left[\psi_{n,v}(\text{sample},\text{seed})\right],
\]
where the sample consists of \(n\) independent records with law \(P\), and the seed is independent and uniform on the unit interval. Define the population score mean by
\[
\theta_P(c):=\mathbb E_{P^{\otimes n}}[Z_b(c)].
\]
Let \(\lambda\) denote Lebesgue probability measure on \([0,1]\), and define the heterogeneity distance and its model supremum by
\[
d(P):=
\left\{\int_0^1
\left(\tau_P(x)-\int_0^1\tau_P(t)\,d\lambda(t)\right)^2
\,d\lambda(x)\right\}^{1/2},
\qquad
D_v:=\sup_{P\in\mathcal M_v}d(P).
\]

For every \(P\in\mathcal M_v\), each affine coefficient of each block score is square-integrable, and its scalar projection in any direction \(u\in\mathbb R^{J_{\mathrm c}}\) has variance at most
\[
\Lambda_{n,v}\|u\|_2^2.
\]
For each block, the absolute covariance between projections of its two affine coefficients in directions \(u,z\in\mathbb R^{J_{\mathrm c}}\) is at most
\[
\Lambda_{n,v}\|u\|_2\|z\|_2.
\]
The simultaneous profiling event for \(W(c)\) specified in \cref{obj:def:explicit-score-ledger}, with profiling multiplier \(C:=128\), is measurable and has probability at least \(0.9925\) under \(P^{\otimes n}\). If \(\Lambda_{n,v}=0\), then, almost surely, simultaneously for every \(|c|\le1/2\),
\[
W(c)=\|\theta_P(c)\|_2^2.
\]

For every \(P\in H_0(v)\), with the null class defined in \cref{obj:def:null},
\[
\inf_{|c|\le1/2}\|\theta_P(c)\|_2\le\mathfrak b_{n,v},
\qquad
\mathsf R_n(P,\psi_{n,v})\le0.0075.
\]
Define the calibrated separation by
\[
R_{\mathrm{att}}(n,v)
:=\frac{16}{3}
\left\{
16J_{\mathrm c}^{-\gamma}
+5\mathfrak b_{n,v}
+1024J_{\mathrm c}^{1/4}\sqrt{\Lambda_{n,v}}
\right\}.
\]
If \(n\ge4\) and \(R_{\mathrm{att}}(n,v)<D_v\), then every \(P\in\mathcal M_v\) satisfying \(d(P)\ge R_{\mathrm{att}}(n,v)\) obeys
\[
1-\mathsf R_n(P,\psi_{n,v})\le0.0075.
\]
For \(n=2\) or \(n=3\), \(\psi_{n,v}\) equals zero at every sample and seed.
\end{theoremv}

Under a homogeneous mean effect, an admissible constant makes the population score norm at most the bias ledger. Profiling retains this comparison while the covariance bounds control fluctuations simultaneously over the candidate constants. Independence between the evaluation blocks gives the population identity \(\mathbb E[W(c)]=\|\theta_P(c)\|_2^2\). The calibrated separation combines coarse-projection approximation, original-mean bias, and stochastic fluctuation. For \(n\ge4\) with \(R_{\mathrm{att}}(n,v)<D_v\), effects at or above that separation are detected with type-II error at most \(0.0075\); the size bound applies throughout the null for every \(n\ge2\).

The appendix develops the three parts of this calibration. The original-mean analysis in \cref{obj:lem:original-score-mean-ledger} controls the bias introduced by projection and clipping. The coefficient covariance bounds are supplied by \cref{obj:lem:original-score-covariance-ledger}, and \cref{obj:lem:uniform-affine-profile-bound} controls the simultaneous profiling operation. Their combination yields the attainable separation result in \cref{obj:thm:original-record-attainable-rate}, connecting the explicit procedure to the power scale \(\rho_n(v)\) in \cref{obj:def:sharp-frontier-scales} and the upper separation bound in \cref{obj:thm:heavy-tail-comparison}.


\section{Supplied propensity scores and assignment information}

The calibrated procedure in \cref{obj:def:explicit-score-ledger} obtains assignment information from the original records. We now consider an experiment in which the entire true continuous propensity is supplied together with those records. This experiment supports two complementary conclusions: an attaining test over a broader class with continuous propensity and conditional mean functions, and an exact comparison of testing sensitivity on the original model.

A supplied-function test \leanref{sym:phior}{\ensuremath{\phi^{\mathrm e}}} is a jointly Borel map
\[
\phi^{\mathrm e}:
C([0,1])\times
\bigl([0,1]\times\{0,1\}\times\mathbb R\bigr)^n
\times[0,1]\longrightarrow[0,1].
\]
The function space carries the uniform topology. The remaining arguments are the independent original-record sample and an independent public uniform randomization variable. Under a law \(P\), the supplied function is \(e_P\); the formal risk definition below specifies the corresponding rejection expectation.

\subsection{The continuous-carrier experiment}

Let \(\mathcal V\) denote the admissible domain of public moment and smoothness tuples in \cref{obj:def:model}, and let \(\lambda\) be Lebesgue probability measure on \([0,1]\). The broader model retains uniform design, overlap, effect smoothness, the mean caps, and the conditional moment envelope. Its propensity and both conditional arm means have continuous representatives. The following definition specifies these restrictions and the maximum heterogeneity distance.

\begin{definitionv}{def:oracle-broad-model}[Continuous-carrier model \(\mathcal M_v^{\mathrm{e,br}}\)]
The class \(\mathcal M_v^{\mathrm{e,br}}\) consists of the original-record laws satisfying the six conditions below and admitting continuous representatives of their propensity and both conditional arm means. Here \(v=(p,\alpha,\beta,\gamma)\in\mathcal V\) is an admissible public tuple, \(\lambda\) is Lebesgue probability measure on \([0,1]\), and \(Q_{a,P}(dy\mid x)\) is the conditional outcome probability kernel in treatment arm \(a\).
\begin{itemize}
\item \textbf{(Uniform covariates.)} The covariate marginal satisfies
\[
 P_X=\lambda,
\]
as stipulated in \cref{obj:ass:uniform}.
\item \textbf{(Overlap.)} The continuous propensity representative satisfies
\[
 e_P([0,1])\subseteq[1/4,3/4],
\]
as stipulated in \cref{obj:ass:overlap}.
\item \textbf{(Effect smoothness.)} The continuous mean effect satisfies
\[
 \tau_P=m_{1,P}-m_{0,P}\in\mathcal H^\gamma(20),
\]
where \(\mathcal H^\gamma(20)\) consists of continuous functions on \([0,1]\) whose sup norm and \(\gamma\)-Hölder constant are at most \(20\), as stipulated in \cref{obj:ass:effect-smoothness}.
\item \textbf{(Baseline cap.)} The baseline mean satisfies
\[
 \|m_{0,P}\|_\infty\le1/2,
\]
as stipulated in \cref{obj:ass:baseline-cap}.
\item \textbf{(Effect cap.)} The mean effect satisfies
\[
 \|\tau_P\|_\infty\le1/2,
\]
as stipulated in \cref{obj:ass:effect-cap}.
\item \textbf{(Raw moment.)} The conditional outcome kernels satisfy
\[
 \max_{a\in\{0,1\}}\operatorname*{ess\,sup}_{x\sim\lambda}
 \int |y|^p Q_{a,P}(dy\mid x)\le10,
\]
as stipulated in \cref{obj:ass:raw-moment}.
\end{itemize}
For a function \(h\), write \([h]_\lambda\) for its equivalence class under Lebesgue almost-everywhere equality. Let \(\mathcal E_P\) denote the propensity equivalence class and let \(\mathcal G_{0,P},\mathcal G_{1,P}\) denote the conditional arm-mean equivalence classes. The continuous-representative requirement is
\[
 \exists f,g_0,g_1\in C([0,1]):\qquad
 \mathcal E_P=[f]_\lambda,\qquad
 \mathcal G_{0,P}=[g_0]_\lambda,\qquad
 \mathcal G_{1,P}=[g_1]_\lambda.
\]
Here \(C([0,1])\) is the space of continuous real functions on the unit interval with the uniform topology. These representatives are the existing functionals \(e_P,m_{0,P},m_{1,P}\). Membership depends on \(p\) and \(\gamma\), with the full tuple \(v\) retained as the public index.

Writing \(d(P)\) for the population Lebesgue \(L_2\) distance of \(\tau_P\) from its average, define
\[
 D_v^{\mathrm{e,br}}=\sup_{P\in\mathcal M_v^{\mathrm{e,br}}}d(P).
\]
\end{definitionv}

The continuous-carrier class \leanref{sym:MoracleBroad}{\ensuremath{\mathcal M_v^{\mathrm{e,br}}}} in \cref{obj:def:oracle-broad-model} permits continuous propensity and baseline mean functions with unrestricted moduli of continuity. Effect smoothness controls the approximation of heterogeneity, while overlap bounds the inverse-propensity weights. The conditional moment envelope controls outcome clipping uniformly across covariate values and treatment arms. The caps bound the scale of the conditional means and effect. Within this class, the null retains the unknown constant effect used in the original experiment.

\begin{definitionv}{def:oracle-broad-null}[Constant-effect null \(H_0^{\mathrm{e,br}}(v)\)]
The constant-effect null is
\[
 H_0^{\mathrm{e,br}}(v)=
 \left\{P\in\mathcal M_v^{\mathrm{e,br}}:
 \exists c\in[-1/2,1/2]\ \forall x\in[0,1],\
 \tau_P(x)=c\right\}.
\]
Membership in the continuous-carrier model imposes the restrictions in
\cref{obj:ass:uniform,obj:ass:overlap,obj:ass:effect-smoothness,obj:ass:baseline-cap,obj:ass:effect-cap,obj:ass:raw-moment}.
The unknown-constant effect condition is that of \cref{obj:ass:null-constancy}.
\end{definitionv}

Size is controlled uniformly over the constant-effect null \leanref{sym:HnulloracleBroad}{\ensuremath{H_0^{\mathrm{e,br}}(v)}} in \cref{obj:def:oracle-broad-null}. For a supplied-function test, write \leanref{sym:Riskor}{\ensuremath{\mathsf R_n^{\mathrm e}(P,\phi^{\mathrm e})}} for its rejection expectation when it receives \(e_P\). The following risk and radius measure the separation required to attain type-II error at most \(1/10\), subject to the same bound on size.

\begin{definitionv}{def:oracle-broad-risk}[Supplied-propensity risk \(B_n^{\mathrm{e,br}}(v,r)\)]
The supplied-propensity minimax type-II error is
\[
 B_n^{\mathrm{e,br}}(v,r)=
 \inf_{\phi^{\mathrm e}:\,
     \sup_{P\in H_0^{\mathrm{e,br}}(v)}
       \mathsf R_n^{\mathrm e}(P,\phi^{\mathrm e})\le1/10}
 \ \sup_{P\in\mathcal M_v^{\mathrm{e,br}}:\,d(P)\ge r}
       \left\{1-\mathsf R_n^{\mathrm e}(P,\phi^{\mathrm e})\right\},
\]
for parameters satisfying
\begin{itemize}
\item \textbf{(Sample size.)} \(n\ge2\).
\item \textbf{(Public tuple.)} \(v\in\mathcal V\), the admissible domain of moment and primitive smoothness tuples.
\item \textbf{(Separation.)} \(0<r<D_v^{\mathrm{e,br}}\), where \(D_v^{\mathrm{e,br}}\) is the supremum of the population effect distance over \(\mathcal M_v^{\mathrm{e,br}}\).
\end{itemize}
Here \(d(P)\) is the population Lebesgue \(L_2\) distance of the mean effect from its average. The admissible tests \(\phi^{\mathrm e}\) are precisely the jointly Borel maps taking a supplied continuous function, an original-record sample, and a public randomization variable to a rejection probability in \([0,1]\). Their rejection expectation is
\[
 \mathsf R_n^{\mathrm e}(P,\phi^{\mathrm e})
 =\int \phi^{\mathrm e}(e_P,\mathcal D_n,u)\,
       P^{\otimes n}(d\mathcal D_n)\,du.
\]
The sample \(\mathcal D_n\) consists of \(n\) independent original records with law \(P\); the integration in \(u\) is over an independent uniform variable on \([0,1]\). The supplied function \(e_P\) is the whole true continuous propensity.
\end{definitionv}

The broader-class risk \leanref{sym:BoracleBroad}{\ensuremath{B_n^{\mathrm{e,br}}(v,r)}} in \cref{obj:def:oracle-broad-risk} uses the original outcome records throughout. Its capped radius includes the model supremum as a saturation value, making the sensitivity measure well defined at every stipulated sample size.

\begin{definitionv}{def:oracle-broad-radius}[Capped radius \(r_n^{*,\mathrm{e,br}}(v)\)]
The capped supplied-propensity critical radius is
\[
 r_n^{*,\mathrm{e,br}}(v)=
 \inf\left(
 \left\{r\in(0,D_v^{\mathrm{e,br}}):
 B_n^{\mathrm{e,br}}(v,r)\le1/10\right\}
 \cup\left\{D_v^{\mathrm{e,br}}\right\}\right).
\]
Here \(D_v^{\mathrm{e,br}}\) is the supremum of the population effect distance over the continuous-carrier model and serves as the saturation value in the capped-radius convention.
\end{definitionv}

The capped radius \leanref{sym:rstaroracleBroad}{\ensuremath{r_n^{*,\mathrm{e,br}}(v)}} and its saturation value \leanref{sym:DoracleBroad}{\ensuremath{D_v^{\mathrm{e,br}}}} are defined in \cref{obj:def:oracle-broad-radius,obj:def:oracle-broad-model}. To attain the separation rate, the test uses clipped inverse-propensity outcomes and the histogram features of \cref{obj:def:projection}. Removing the constant histogram direction targets deviations from an unknown constant effect. The procedure then takes the inner product of two independent block averages.

\begin{algorithmv}{def:oracle-handle}[Supplied-propensity test \(\phi^{*,\mathrm e}_{n,v}\)]
Inputs are \(n\ge2\) original records, the public tuple \(v=(p,\alpha,\beta,\gamma)\), and a supplied continuous propensity \(e\).
\begin{enumerate}
\item Set the block size, moment exponent, separation exponent, and target scale to
\[
s=\lfloor n/2\rfloor,
\qquad
q=\frac{p-1}{p},
\qquad
E_0=\frac{2\gamma q}{2\gamma+q},
\qquad
a=s^{-E_0}.
\]
Use the evaluation blocks
\[
\mathcal I_1=\{1,\ldots,s\},
\qquad
\mathcal I_2=\{s+1,\ldots,2s\}.
\]
\item Select the clipping cutoff and histogram rank
\[
T=2^{\left\lceil\log_2(a^{-1/(p-1)})\right\rceil},
\qquad
J=2^{\left\lceil\log_2(a^{-1/\gamma})\right\rceil}.
\]
Define outcome clipping and the overlap-clipped propensity by
\[
\ell_T(y)=\max\{-T,\min\{y,T\}\},
\qquad
\widetilde e(x)=\max\{1/4,\min\{e(x),3/4\}\}.
\]
For every legal supplied true propensity, \(\widetilde e=e\).
\item Form the original-record inverse-propensity scores
\[
R_i=
\frac{A_i\ell_T(Y_i)}{\widetilde e(X_i)}
-
\frac{(1-A_i)\ell_T(Y_i)}{1-\widetilde e(X_i)}.
\]
Let \(F_J\) be the normalized histogram feature vector. Its constant direction and the centered block vectors are
\[
u=J^{-1/2}(1,\ldots,1)^T,
\qquad
V_b=\frac1s\sum_{i\in\mathcal I_b}
\bigl(F_J(X_i)-u\bigr)R_i,
\quad b\in\{1,2\}.
\]
These vectors represent the block scores obtained by applying the constant-centered histogram projection \(\Pi_J-1\).
\item Set the public clipping-bias and covariance bounds to
\[
b=20T^{1-p},
\qquad
\Lambda=\frac{80T^{2-p}}s.
\]
Output the rejection rule
\[
\phi^{*,\mathrm e}_{n,v}
=
\mathbf 1\!\left\{
\langle V_1,V_2\rangle
>
2b^2+1024\sqrt J\,\Lambda
\right\}.
\]
Its calibration and necessity properties are given by its singleton and canonical quadratic projections and the normalized common-propensity paired-tent converse.
\end{enumerate}
\end{algorithmv}

The supplied-propensity rule \leanref{sym:phistaror}{\ensuremath{\phi^{*,\mathrm e}_{n,v}}} in \cref{obj:def:oracle-handle} uses public clipping-bias and covariance bounds to set its rejection threshold. The tuning depends on the moment order and effect smoothness. The next result gives matching bounds for the broader-class critical radius, uniform calibration over its entire null, and power whenever the stated separation lies below the saturation value.



The conditional-mean identity in \cref{obj:thm:oracle-frontier-broad-class} explains the role of supplied assignment information: inverse weighting recovers the mean effect directly, and constant centering isolates its heterogeneity. Clipping controls the contribution of outcome tails, while effect smoothness controls histogram approximation. Their balance yields the supplied-propensity separation scale \leanref{sym:rhoor}{\ensuremath{\rho_n^{\mathrm e}(v)}} defined in that result. The matching lower bound uses null and alternative laws with the same constant propensity and close complete-record distributions. Thus the stated separation order remains necessary even when assignment information is identical across the comparison laws.

The radius bounds apply to every \(n\ge2\), with the cap accounting for finite-sample saturation. The procedure's power guarantee applies when its upper separation is strictly below the broader-model supremum. The public sample-size threshold in \cref{obj:thm:oracle-frontier-broad-class} ensures this condition by placing that separation at most halfway to the fixed witness distance. These statements distinguish an all-sample radius bound from a power guarantee over a nonempty range of alternatives.

\subsection{Comparison on the original model}

To measure the value of supplied assignment information on a fixed statistical class, we now use the original model of \cref{obj:def:model}. Its quantitative propensity and baseline smoothness restrictions are retained in both experiments being compared. The supplied-function test carrier remains the same, while the null and alternative suprema are taken over the original model. The following definitions make this shared-class comparison precise.

\begin{definitionv}{def:oracle-risk}[Supplied-propensity risk $B_n^{\mathrm e}(v,r)$]
\[
B_n^{\mathrm e}(v,r)=
\inf_{\phi^{\mathrm e}:\ \sup_{P\in H_0(v)}
\mathsf R_n^{\mathrm e}(P,\phi^{\mathrm e})\le 1/10}
\ \sup_{P\in\mathcal M_v:\ d(P)\ge r}
\bigl\{1-\mathsf R_n^{\mathrm e}(P,\phi^{\mathrm e})\bigr\}.
\]
The parameters satisfy:
\begin{itemize}
\item \textbf{(Sample size.)} \(n\ge 2\).
\item \textbf{(Public tuple.)} \(v\in\mathcal V\), the admissible domain of moment and primitive smoothness tuples.
\item \textbf{(Separation.)} \(0<r<D_v\), where \(D_v\) is the supremum of \(d(P)\) over \(\mathcal M_v\).
\end{itemize}
Here \(\mathcal M_v\) is the original-record model, \(H_0(v)\) consists of its laws with an unknown constant mean effect, and \(d(P)\) is the population Lebesgue \(L_2\) distance of the effect from its average. The infimum ranges over randomized rejection-probability maps \(\phi^{\mathrm e}\) receiving the whole true continuous propensity, and \(\mathsf R_n^{\mathrm e}(P,\phi^{\mathrm e})\) is their expected rejection probability with that propensity supplied.
\end{definitionv}

The shared-class risk \leanref{sym:Bor}{\ensuremath{B_n^{\mathrm e}(v,r)}} in \cref{obj:def:oracle-risk} uses the original-model null for its size constraint. Its radius therefore uses the original-model supremum as the saturation value.

\begin{definitionv}{def:oracle-radius}[Supplied-propensity radius $r_n^{*,\mathrm e}(v)$]
\[
r_n^{*,\mathrm e}(v)=
\inf\Bigl(
\bigl\{r\in(0,D_v):B_n^{\mathrm e}(v,r)\le 1/10\bigr\}
\cup\{D_v\}
\Bigr).
\]
Here \(B_n^{\mathrm e}(v,r)\) is the supplied-propensity minimax type-II error on the original model under the size constraint \(1/10\), and \(D_v\) is the supremum of the population effect distance over that model.
\end{definitionv}

The shared-class capped radius \leanref{sym:rstaror}{\ensuremath{r_n^{*,\mathrm e}(v)}} of \cref{obj:def:oracle-radius} is compatible with the broader-class attainment result. The original model is contained in the continuous-carrier class, and the finite lower-bound witnesses in \cref{obj:thm:oracle-frontier-broad-class} belong to both classes. The next result records the matching separation bounds and attainment on the original model, providing the supplied-propensity benchmark for the radius comparison.



The compatibility in \cref{obj:thm:oracle-frontier-resolved} rests on both attainment and necessity. The supplied-propensity procedure controls size and power on the original class, and common-propensity witnesses establish the corresponding lower separation bound within that class. Consequently, the original-model comparison uses matched statistical classes in its numerator and denominator. Its finite-sample power condition is expressed relative to the original-model supremum; the broader-class result uses its own supremum.

We measure preservation of separation order by a uniformly bounded ratio of the unknown- and supplied-propensity capped radii. The following region includes every admissible tuple for which that ratio remains bounded across sample sizes.

\begin{definitionv}{def:preservation-region}[Propensity preservation region $\mathcal R$]
\[
\mathcal R=
\left\{
v\in\mathcal V:
\sup_{n\ge 2}
\frac{r_n^*(v)}{r_n^{*,\mathrm e}(v)}
<\infty
\right\}.
\]
Here \(\mathcal V\) is the admissible domain of public moment and primitive smoothness tuples. The radii \(r_n^*(v)\) and \(r_n^{*,\mathrm e}(v)\) are the capped critical separation radii on the same original-record model with unknown and supplied true continuous propensity, respectively.
\end{definitionv}

The preservation region \leanref{sym:Region}{\ensuremath{\mathcal R}} in \cref{obj:def:preservation-region} follows from comparing the matched rates in \cref{obj:thm:heavy-tail-comparison,obj:thm:oracle-frontier-resolved}. The scalar comparison quantity \(F(p)\) defined in \cref{obj:def:sharp-frontier-scales} determines whether the separation exponents coincide. The next result gives the exact boundary, including its equality case and bounds for every capped finite-sample radius ratio.



The boundary in \cref{obj:thm:exact-propensity-preservation-boundary} expresses the value of assignment information through the difference between the two separation exponents. When \(F(p)\ge1\), the unknown-propensity experiment attains the supplied-propensity separation order on the same original class. At \(F(p)=1\), the ratio remains between positive, sample-size-independent constants. When \(F(p)<1\), supplying the true propensity improves sensitivity by a polynomial factor whose exponent is stated explicitly in the result.

These comparisons concern orders of critical separation, with finite-sample constants and saturation incorporated into the capped radii. They hold for every \(n\ge2\), and their constants admit uniform bounds on compact subsets of the admissible parameter domain. Together with the broader-class attainment result, they distinguish testing with supplied assignment information from the exact region where the original unknown-propensity experiment attains the same separation order. Requirements for a particular estimated-propensity procedure are outside this experiment comparison.


\section{Discussion and extensions}

The separation comparisons provide two benchmarks for the sensitivity of homogeneity testing: one varies the outcome envelope at matched primitive smoothness, and the other varies the assignment information available to the test on a shared model class. The first comparison, in \cref{obj:thm:heavy-tail-comparison}, measures the cost of allowing outcomes subject to a finite conditional moment bound relative to bounded outcomes. The second, in \cref{obj:thm:exact-propensity-preservation-boundary}, measures the value of supplying the true propensity function. Together, they characterize how outcome tails and assignment information affect the smallest detectable departure from a constant conditional mean effect.

Matching primitive smoothness makes the outcome comparison interpretable. The propensity and baseline smoothness exponents describe nuisance complexity, whereas the effect smoothness exponent describes the complexity of the departure being tested. The bounded benchmark in \cref{obj:def:bounded-model} retains these three restrictions. Consequently, the radius ratio in \cref{obj:def:comparison-handle} compares outcome envelopes while keeping the stipulated effect and nuisance complexity fixed. The regions in \cref{obj:def:tail-region,obj:def:equality-region} distinguish a diverging separation penalty from a uniformly bounded radius ratio. These are benchmarks for worst-case sensitivity over the specified classes: they identify when the conditional moment envelope changes the order of the testing problem.

The assignment comparison separates the complexity of the effect from the information needed to handle the nuisance functions. On the shared original class, \cref{obj:thm:exact-propensity-preservation-boundary} identifies the exact region where the unknown-propensity experiment attains the supplied-propensity separation order. The broader-class result in \cref{obj:thm:oracle-frontier-broad-class} gives a complementary interpretation. Its model, defined in \cref{obj:def:oracle-broad-model}, permits continuous propensity and baseline mean functions subject to the retained overlap and mean envelopes, together with the conditional moment and effect smoothness restrictions. Attainment over this class supports a broader nuisance scope when the true propensity is supplied. The resulting separation order is governed by the moment order and effect smoothness, clarifying the role of assignment information in accommodating nuisance variation.

The population target also has a direct variance interpretation. Writing the distance from a constant effect as
\[
d(P)=
\left[
\int_0^1
\left\{\tau_P(x)-\int_0^1\tau_P(z)\,dz\right\}^{2}\,dx
\right]^{1/2},
\]
the uniform covariate design in \cref{obj:ass:uniform} gives
\[
d(P)^2=\operatorname{Var}_{P}\{\tau_P(X)\}.
\]
Thus separation is measured in the standard deviation of the conditional mean treatment effect across covariate values. Under the causal completion in \cref{obj:def:causal-completion}, this contrast equals the conditional expectation of the potential-outcome difference. The target therefore quantifies heterogeneity in conditional mean effects, with the covariate distribution determining the population weighting.

Public finite-sample quantities connect these comparisons to calibrated rejection rules. The construction in \cref{obj:def:explicit-score-ledger} specifies the projection ranks, clipping levels, bias and covariance bounds, and rejection cutoff. Whole-null calibration in \cref{obj:thm:whole-null-calibration-resolved} accounts for the unknown constant effect through profiling. Combined with the attainment guarantees in \cref{obj:thm:heavy-tail-comparison}, the public sample-size thresholds identify when the stated power separation lies inside the model's nonempty alternative domain. The analogous supplied-propensity guarantees appear in \cref{obj:thm:oracle-frontier-broad-class,obj:thm:oracle-frontier-resolved}.

The explicit sufficient constants are deliberately conservative. The following values come directly from the displayed formulas; logarithms make their scale readable. Here \(N_v^{\mathrm{att}}\) is the original-record nonempty-power threshold and \(N_v^{\mathrm e}\) is its supplied-propensity counterpart.
\begin{center}
\begin{tabular}{llllll}
\hline
\(p\) & \((\alpha,\beta,\gamma)\) & regime & \(E\) & \(\log_{10}C_v^{\mathrm{att}}\) & \(\log_{10}N_v^{\mathrm{att}}\; /\;\log_{10}N_v^{\mathrm e}\) \\
\hline
\(2\) & \((1,1,1)\) & effect & \(0.400\) & \(8.785\) & \(25.35\;/\;12.86\) \\
\(3/2\) & \((1,1,1)\) & effect & \(0.286\) & \(8.786\) & \(35.49\;/\;18.00\) \\
\(3/2\) & \((1/20,1/20,1)\) & rough & \(0.154\) & \(9.934\) & \(73.37\;/\;18.00\) \\
\hline
\end{tabular}
\end{center}
For the first row, for example, the supplied-propensity formula gives \(N_v^{\mathrm e}=7{,}183{,}754{,}725{,}838\). These are analytic sufficient thresholds for the stated finite-sample power guarantee, rather than estimates of the sample size needed in practice. The theorem's main quantitative content is the minimax power of \(n\) and the matching lower bound; practical calibration of the constants would require implementation and simulation beyond the present analysis.

\subsection{Limitations and future work}

The analysis treats a known uniform covariate design on the unit interval and public moment and smoothness exponents. Extending the construction to an unknown covariate distribution would require handling the population weighting and the estimation of projection geometry together. Higher-dimensional covariates raise a related question about how approximation complexity and nuisance smoothness interact in the separation comparison.

Adaptive selection of the moment and smoothness exponents is another research direction. A useful extension would quantify the cost of choosing clipping levels and projection ranks from the data while maintaining calibration over the entire constant-effect null and establishing nonempty power domains.

Finally, the supplied-function experiment provides a benchmark for studying estimated assignment functions. Such an extension would quantify how propensity approximation error and the dependence introduced by its estimation enter the bias bounds, rejection threshold, and attainable separation. Comparing those guarantees with the exact supplied-function benchmark would clarify the assignment accuracy needed to retain its sensitivity to effect heterogeneity.


\appendix

\section*{Appendices}

\section{Causal interpretation and benchmark equivalence}

The population target in \cref{obj:def:model,obj:def:null} uses continuous representatives of conditional means. Its interpretation therefore requires both invariance to conditional-distribution versions and a causal representation that reproduces the observed law. The following result establishes these properties and supplies legal constant-effect and heterogeneous-effect witnesses for every admissible parameter tuple.

\begin{lemmav}{lem:causal-nonempty}[Nonemptiness and causal identification]
Fix an admissible parameter tuple \(v=(p,\alpha,\beta,\gamma)\) as in \cref{obj:def:model}. Let \(\lambda\) denote Lebesgue probability measure on \([0,1]\), and define the heterogeneity distance, its model supremum, and the witness distance by
\[
d(P)=
\left\{\int_0^1
\left(\tau_P(x)-\int_0^1\tau_P(z)\,d\lambda(z)\right)^2
\,d\lambda(x)\right\}^{1/2},
\qquad
D_v=\sup_{P\in\mathcal M_v}d(P),
\qquad
d_0=\frac{1}{8\sqrt{2}}.
\]
The set \(\{d(P):P\in\mathcal M_v\}\) is bounded above, and
\[
d_0\le D_v\le\frac12.
\]

For every \(P\in\mathcal M_v\), with model membership understood as in \cref{obj:def:model}, the continuous propensity, baseline mean, and contrast representatives are determined uniquely by the observed law: every other observational representation of the same probability law has exactly the same \(e_P\), \(m_{0,P}\), and \(\tau_P\), pointwise on \([0,1]\). Consequently, the treated mean representative \(m_{1,P}=m_{0,P}+\tau_P\) is also uniquely determined, and
\[
|m_{1,P}(x)-m_{1,P}(z)|
\le 40|x-z|^{\min\{\beta,\gamma\}}
\qquad (x,z\in[0,1]).
\]

In the completion \(\mathbb Q_P\) of \cref{obj:def:causal-completion}, let \(Y(0)\) and \(Y(1)\) denote the baseline and treated potential-outcome coordinates. The distribution under \(\mathbb Q_P\) of
\[
\bigl(X,A,(1-A)Y(0)+AY(1)\bigr)
\]
is \(P\), and
\[
(Y(0),Y(1))\perp A\mid X.
\]
Both potential outcomes are integrable under \(\mathbb Q_P\), and
\[
\mathbb E_{\mathbb Q_P}\!\left[Y(1)-Y(0)\mid X\right]
=\tau_P(X)
\qquad \mathbb Q_P\text{-almost surely}.
\]
Moreover,
\[
d(P)=0\quad\Longleftrightarrow\quad P\in H_0(v),
\]
where \(H_0(v)\) is defined in \cref{obj:def:null}.

The law with uniform covariate distribution, propensity and mean functions
\[
e_P(x)=\frac12+\frac{\sin(2\pi x)}{10},
\qquad
m_{0,P}(x)=\frac{\cos(2\pi x)}8,
\qquad
\tau_P(x)=\frac{\cos(2\pi x)}8,
\]
and deterministic conditional outcome
\[
Y=m_{0,P}(X)+A\tau_P(X)
\]
belongs to \(\mathcal M_v\) and has \(d(P)=d_0\). Replacing its contrast function by \(\tau_P(x)=1/8\) for every \(x\in[0,1]\), while retaining the displayed propensity, baseline mean, and deterministic conditional-outcome construction, yields a law in \(H_0(v)\).
\end{lemmav}

\begin{proof}[Proof of \cref{obj:lem:causal-nonempty}]
Fix an admissible tuple \(v\). Thus
\[
1<p\le2,\qquad
0<\alpha,\beta\le1,\qquad
\frac14\le\gamma\le1.
\]
We establish the witness construction and distance bounds first, followed by the identification and completion properties.

\begin{enumerate}
\item Define the separated witness primitives by
\[
e_{\mathrm{sep}}(x)=\frac12+\frac{\sin(2\pi x)}{10},
\qquad
m_{0,\mathrm{sep}}(x)=\frac{\cos(2\pi x)}8,
\qquad
\tau_{\mathrm{sep}}(x)=\frac{\cos(2\pi x)}8,
\qquad x\in[0,1].
\]
For \(y_\ast\in\mathbb R\), write \(\delta_{y_\ast}\) for the unit point mass, defined by
\[
\delta_{y_\ast}(\mathsf E)=\mathbf 1_{\mathsf E}(y_\ast)
\qquad
\bigl(\mathsf E\in\mathcal B(\mathbb R)\bigr),
\]
and use the same notation for point masses on the treatment space \(\{0,1\}\).
The witness law is
\[
\begin{aligned}
P_{\mathrm{sep}}(dx,da,dy)
=\lambda(dx)\bigl[
 &(1-e_{\mathrm{sep}}(x))\,
 \delta_0(da)\delta_{m_{0,\mathrm{sep}}(x)}(dy)\\
 &+e_{\mathrm{sep}}(x)\,
 \delta_1(da)\delta_{m_{0,\mathrm{sep}}(x)+\tau_{\mathrm{sep}}(x)}(dy)
\bigr].
\end{aligned}
\]
The two weights are nonnegative and sum to one, so this construction is a probability law with covariate marginal \(\lambda\). Its arm kernels and means satisfy
\[
Q_{0,P_{\mathrm{sep}}}(dy\mid x)
=\delta_{m_{0,\mathrm{sep}}(x)}(dy),
\qquad
Q_{1,P_{\mathrm{sep}}}(dy\mid x)
=\delta_{m_{0,\mathrm{sep}}(x)+\tau_{\mathrm{sep}}(x)}(dy),
\]
\[
\int y\,Q_{0,P_{\mathrm{sep}}}(dy\mid x)=m_{0,\mathrm{sep}}(x),
\qquad
\int y\,Q_{1,P_{\mathrm{sep}}}(dy\mid x)
=m_{0,\mathrm{sep}}(x)+\tau_{\mathrm{sep}}(x).
\]
The sine and cosine bounds give
\[
\frac14\le e_{\mathrm{sep}}(x)\le\frac34,
\qquad
|m_{0,\mathrm{sep}}(x)|=|\tau_{\mathrm{sep}}(x)|\le\frac18.
\]
Their increment bounds are
\[
\begin{aligned}
|e_{\mathrm{sep}}(x)-e_{\mathrm{sep}}(z)|
&\le\frac{\pi}{5}|x-z|
\le20|x-z|,\\
|m_{0,\mathrm{sep}}(x)-m_{0,\mathrm{sep}}(z)|
=|\tau_{\mathrm{sep}}(x)-\tau_{\mathrm{sep}}(z)|
&\le\frac{\pi}{4}|x-z|
\le20|x-z|.
\end{aligned}
\]
For every exponent \(\sigma_{\mathrm H}\in(0,1]\),
\[
|x-z|\le |x-z|^{\sigma_{\mathrm H}}
\qquad(x,z\in[0,1]).
\]
Consequently these continuous functions satisfy the supremum and increment bounds in
\cref{obj:ass:propensity-smoothness,obj:ass:baseline-smoothness,obj:ass:effect-smoothness}
at their respective exponents.

For the moment restriction, the baseline and effect caps imply
\[
|m_{0,\mathrm{sep}}(x)|\le1,
\qquad
|m_{0,\mathrm{sep}}(x)+\tau_{\mathrm{sep}}(x)|
\le |m_{0,\mathrm{sep}}(x)|+|\tau_{\mathrm{sep}}(x)|
\le1.
\]
Since \(p>0\), both deterministic arm kernels therefore satisfy
\[
\int |y|^p\,Q_{a_{\mathrm{arm}},P_{\mathrm{sep}}}(dy\mid x)
\le1\le10
\qquad
\bigl(a_{\mathrm{arm}}\in\{0,1\},\ x\in[0,1]\bigr).
\]
This verifies every restriction in \cref{obj:def:model}, and hence
\[
P_{\mathrm{sep}}\in\mathcal M_v.
\]
% lean: CausalSmith.Stat.FinitepHomogeneityDensegamma.separatedWitness_inModel

\item\label{proof:causal-nonempty:distance-bounds} For each \(P\in\mathcal M_v\), define its average contrast by
\[
\bar\tau_P=\int_0^1\tau_P(x)\,d\lambda(x).
\]
By \cref{obj:ass:effect-cap},
\[
|\bar\tau_P|
\le\int_0^1|\tau_P(x)|\,d\lambda(x)
\le\frac12,
\qquad
|\tau_P(x)-\bar\tau_P|\le1.
\]
The bounded functions below are integrable. Expanding the square and using
\(\lambda([0,1])=1\) gives
\[
\begin{aligned}
\int_0^1(\tau_P(x)-\bar\tau_P)^2\,d\lambda(x)
&=\int_0^1\tau_P(x)^2\,d\lambda(x)
 -2\bar\tau_P\int_0^1\tau_P(x)\,d\lambda(x)
 +\bar\tau_P^2\\
&=\int_0^1\tau_P(x)^2\,d\lambda(x)-\bar\tau_P^2.
\end{aligned}
\]
Moreover,
\[
\int_0^1\tau_P(x)^2\,d\lambda(x)
\le\int_0^1\frac14\,d\lambda(x)=\frac14.
\]
Thus
\[
0\le d(P)\le\frac12
\qquad(P\in\mathcal M_v).
\]
The distance set is therefore bounded above. It is nonempty by the witness from the preceding step, so its supremum satisfies
\[
D_v\le\frac12.
\]

For the witness, substitution in the cosine integrals gives
\[
\bar\tau_{P_{\mathrm{sep}}}
=\frac{1}{16\pi}\int_0^{2\pi}\cos\vartheta\,d\vartheta
=\frac{1}{16\pi}[\sin\vartheta]_0^{2\pi}
=0
\]
and
\[
\begin{aligned}
\int_0^1\tau_{\mathrm{sep}}(x)^2\,d\lambda(x)
&=\frac{1}{128\pi}
  \int_0^{2\pi}\cos^2\vartheta\,d\vartheta\\
&=\frac{1}{128\pi}
 \left[\frac{\vartheta}{2}
       +\frac{\sin\vartheta\cos\vartheta}{2}\right]_0^{2\pi}
=\frac1{128}.
\end{aligned}
\]
Since
\[
d_0>0,\qquad d_0^2=\frac1{128},
\]
taking the nonnegative square root yields
\[
d(P_{\mathrm{sep}})=d_0.
\]
Membership of this distance in the bounded distance set then gives
\[
d_0=d(P_{\mathrm{sep}})\le D_v\le\frac12.
\]
% lean: CausalSmith.Stat.FinitepHomogeneityDensegamma.model_distance_upper
% lean: CausalSmith.Stat.FinitepHomogeneityDensegamma.separatedWitness_distance
% lean: CausalSmith.Stat.FinitepHomogeneityDensegamma.model_distance_lower

\item Define the constant witness by
\[
e_{\mathrm{const}}=e_{\mathrm{sep}},
\qquad
m_{0,\mathrm{const}}=m_{0,\mathrm{sep}},
\qquad
\tau_{\mathrm{const}}(x)=\frac18
\quad(x\in[0,1]),
\]
\[
\begin{aligned}
P_{\mathrm{const}}(dx,da,dy)
=\lambda(dx)\bigl[
 &(1-e_{\mathrm{const}}(x))\,
 \delta_0(da)\delta_{m_{0,\mathrm{const}}(x)}(dy)\\
 &+e_{\mathrm{const}}(x)\,
 \delta_1(da)\delta_{m_{0,\mathrm{const}}(x)+1/8}(dy)
\bigr].
\end{aligned}
\]
The constant contrast satisfies
\[
|\tau_{\mathrm{const}}(x)|=\frac18\le\frac12,
\qquad
|\tau_{\mathrm{const}}(x)-\tau_{\mathrm{const}}(z)|
=0\le20|x-z|^\gamma.
\]
The preceding deterministic construction again gives uniform design, the displayed primitives, and the required arm means. The propensity and baseline restrictions have already been verified. The same moment argument applies because
\[
|m_{0,\mathrm{const}}(x)|\le\frac12,
\qquad
|m_{0,\mathrm{const}}(x)+\tau_{\mathrm{const}}(x)|\le1.
\]
Thus \(P_{\mathrm{const}}\in\mathcal M_v\), and its contrast is everywhere the admissible constant \(1/8\). By \cref{obj:def:null},
\[
P_{\mathrm{const}}\in H_0(v).
\]
% lean: CausalSmith.Stat.FinitepHomogeneityDensegamma.constantWitness_inNull
% lean: CausalSmith.Stat.FinitepHomogeneityDensegamma.deterministicLaw_primitives

\item Fix \(P\in\mathcal M_v\), and consider any other observational representation of the same probability law. Denote its continuous primitives and arm kernels by
\[
\widetilde e_P,\widetilde m_{0,P},\widetilde\tau_P\in C([0,1]),
\qquad
\widetilde m_{1,P}
=\widetilde m_{0,P}+\widetilde\tau_P,
\]
\[
\widetilde Q_{a_{\mathrm{arm}},P}(\,\cdot\mid x)
\in\mathcal P(\mathbb R)
\qquad
\bigl(a_{\mathrm{arm}}\in\{0,1\},\ x\in[0,1]\bigr),
\]
where \(\mathcal P(\mathbb R)\) denotes the Borel probability measures on \(\mathbb R\).
Define the two conditional record kernels by
\[
\begin{aligned}
\mathsf K_P(x,da,dy)
&=(1-e_P(x))\delta_0(da)Q_{0,P}(dy\mid x)
  +e_P(x)\delta_1(da)Q_{1,P}(dy\mid x),\\
\widetilde{\mathsf K}_P(x,da,dy)
&=(1-\widetilde e_P(x))\delta_0(da)\widetilde Q_{0,P}(dy\mid x)
  +\widetilde e_P(x)\delta_1(da)\widetilde Q_{1,P}(dy\mid x).
\end{aligned}
\]
The second representation has covariate marginal \(\lambda\), because its probability law is \(P\). Their observational factorizations therefore give
\[
P(dx,da,dy)
=\lambda(dx)\mathsf K_P(x,da,dy)
=\lambda(dx)\widetilde{\mathsf K}_P(x,da,dy).
\]
Uniqueness of probability kernels in this factorization implies
\[
\widetilde{\mathsf K}_P(x,\cdot)
=\mathsf K_P(x,\cdot)
\qquad\text{for \(\lambda\)-almost every }x.
\]
Evaluating the equality on \(\{1\}\times\mathbb R\) yields
\[
\widetilde e_P(x)=e_P(x)
\qquad\text{for \(\lambda\)-almost every }x.
\]
Continuous functions agreeing \(\lambda\)-almost everywhere agree throughout \([0,1]\): a nonzero difference at any point would, by continuity, remain nonzero on a relative neighborhood of positive Lebesgue measure. Hence
\[
\widetilde e_P=e_P
\qquad\text{pointwise on }[0,1].
\]

For every Borel \(\mathsf E\subseteq\mathbb R\), evaluation on the two arm slices now gives, at almost every \(x\),
\[
\begin{aligned}
(1-e_P(x))\widetilde Q_{0,P}(\mathsf E\mid x)
&=(1-e_P(x))Q_{0,P}(\mathsf E\mid x),\\
e_P(x)\widetilde Q_{1,P}(\mathsf E\mid x)
&=e_P(x)Q_{1,P}(\mathsf E\mid x).
\end{aligned}
\]
Both weights are strictly positive by \cref{obj:ass:overlap}. Cancelling them proves
\[
\widetilde Q_{0,P}(\,\cdot\mid x)=Q_{0,P}(\,\cdot\mid x),
\qquad
\widetilde Q_{1,P}(\,\cdot\mid x)=Q_{1,P}(\,\cdot\mid x)
\qquad\text{for \(\lambda\)-almost every }x.
\]
The arm-mean version identities consequently imply
\[
\begin{aligned}
\widetilde m_{0,P}(x)
&=\int y\,\widetilde Q_{0,P}(dy\mid x)
 =\int y\,Q_{0,P}(dy\mid x)
 =m_{0,P}(x),\\
\widetilde m_{1,P}(x)
&=\int y\,\widetilde Q_{1,P}(dy\mid x)
 =\int y\,Q_{1,P}(dy\mid x)
 =m_{1,P}(x)
\end{aligned}
\qquad\text{for \(\lambda\)-almost every }x.
\]
Continuity upgrades both equalities to pointwise equalities. Subtraction then gives
\[
\widetilde\tau_P
=\widetilde m_{1,P}-\widetilde m_{0,P}
=m_{1,P}-m_{0,P}
=\tau_P
\qquad\text{pointwise on }[0,1].
\]
% lean: CausalSmith.Stat.FinitepHomogeneityDensegamma.observed_continuous_versions_unique

\item The smoothness restrictions in
\cref{obj:ass:baseline-smoothness,obj:ass:effect-smoothness}, together with
\(m_{1,P}=m_{0,P}+\tau_P\), give
\[
\begin{aligned}
|m_{1,P}(x)-m_{1,P}(z)|
&=|(m_{0,P}(x)-m_{0,P}(z))
   +(\tau_P(x)-\tau_P(z))|\\
&\le |m_{0,P}(x)-m_{0,P}(z)|
    +|\tau_P(x)-\tau_P(z)|\\
&\le20|x-z|^\beta+20|x-z|^\gamma\\
&\le40|x-z|^{\min\{\beta,\gamma\}}.
\end{aligned}
\]
The last inequality uses \(0\le|x-z|\le1\) and
\(\min\{\beta,\gamma\}>0\). It includes \(x=z\), when all displayed increments and powers vanish.
% lean: CausalSmith.Stat.FinitepHomogeneityDensegamma.treated_mean_holder

\item In the completion of \cref{obj:def:causal-completion}, define
\[
Y_{\mathrm{obs}}=(1-A)Y(0)+AY(1).
\]
For any nonnegative Borel function
\[
f_{\mathrm{obs}}:[0,1]\times\{0,1\}\times\mathbb R
\longrightarrow[0,\infty],
\]
integration over the completion gives
\[
\begin{aligned}
&\int f_{\mathrm{obs}}(X,A,Y_{\mathrm{obs}})\,d\mathbb Q_P\\
&=\int_0^1
 \int_{\mathbb R}\int_{\mathbb R}
 \bigl[
 (1-e_P(x))f_{\mathrm{obs}}(x,0,y_0)
 +e_P(x)f_{\mathrm{obs}}(x,1,y_1)
 \bigr]\,
 Q_{1,P}(dy_1\mid x)Q_{0,P}(dy_0\mid x)\,d\lambda(x)\\
&=\int_0^1\left[
 (1-e_P(x))\int_{\mathbb R}f_{\mathrm{obs}}(x,0,y)\,Q_{0,P}(dy\mid x)
 +e_P(x)\int_{\mathbb R}f_{\mathrm{obs}}(x,1,y)\,Q_{1,P}(dy\mid x)
 \right]d\lambda(x)\\
&=\int f_{\mathrm{obs}}\,dP.
\end{aligned}
\]
Indeed, in each arm the unused potential-outcome kernel integrates to one; the last equality is the observational factorization with uniform covariate marginal. Equality for all such functions proves
\[
\mathcal L_{\mathbb Q_P}(X,A,Y_{\mathrm{obs}})=P.
\]
% lean: CausalSmith.Stat.FinitepHomogeneityDensegamma.completion_observed_marginal

\item\label{proof:causal-nonempty:conditional-law} Integrating out the probability kernels first gives
\[
\mathcal L_{\mathbb Q_P}(X)=\lambda.
\]
The conditional law of the remaining coordinates is the defining product kernel:
\[
\begin{aligned}
&\mathcal L_{\mathbb Q_P}(A,Y(0),Y(1)\mid X=x)\\
&\quad=
\bigl((1-e_P(x))\delta_0+e_P(x)\delta_1\bigr)
\otimes Q_{0,P}(\,\cdot\mid x)
\otimes Q_{1,P}(\,\cdot\mid x)
\qquad\text{for \(\lambda\)-almost every }x.
\end{aligned}
\]
Projecting this identity onto treatment and onto the potential-outcome pair yields
\[
\begin{aligned}
\mathcal L_{\mathbb Q_P}(A\mid X=x)
&=(1-e_P(x))\delta_0+e_P(x)\delta_1,\\
\mathcal L_{\mathbb Q_P}(Y(0),Y(1)\mid X=x)
&=Q_{0,P}(\,\cdot\mid x)\otimes Q_{1,P}(\,\cdot\mid x).
\end{aligned}
\]
Thus the joint conditional law equals the product of these two conditional marginals. The conditional-independence criterion gives
\[
(Y(0),Y(1))\perp A\mid X.
\]
% lean: CausalSmith.Stat.FinitepHomogeneityDensegamma.completion_conditional_independence

\item For \(y\in\mathbb R\), the elementary bound needed for integrability is
\[
|y|\le1+|y|^p.
\]
If \(|y|\le1\), this follows from \(|y|^p\ge0\). If \(|y|>1\), the condition \(p>1\) gives
\[
|y|\le|y|^p\le1+|y|^p.
\]
Hence \cref{obj:ass:raw-moment} implies, for each
\(a_{\mathrm{arm}}\in\{0,1\}\),
\[
\int_{\mathbb R}|y|\,Q_{a_{\mathrm{arm}},P}(dy\mid x)
\le1+\int_{\mathbb R}|y|^p\,Q_{a_{\mathrm{arm}},P}(dy\mid x)
\le11
\qquad\text{for \(\lambda\)-almost every }x.
\]
Integrating the assignment kernel and the unused potential-outcome kernel, both of total mass one, gives
\[
\begin{aligned}
\int |Y(a_{\mathrm{arm}})|\,d\mathbb Q_P
&=\int_0^1\int_{\mathbb R}
 |y|\,Q_{a_{\mathrm{arm}},P}(dy\mid x)\,d\lambda(x)\\
&\le\int_0^1 11\,d\lambda(x)=11<\infty.
\end{aligned}
\]
This proves integrability for \(Y(0)\) and \(Y(1)\).
% lean: CausalSmith.Stat.FinitepHomogeneityDensegamma.completion_potential_integrable

\item For each arm, integration against the joint conditional law from \cref{proof:causal-nonempty:conditional-law} reduces to the corresponding arm mean, since the other two factors are probability measures. The conditional-expectation formula for an integrable function therefore gives
\[
\mathbb E_{\mathbb Q_P}[Y(a_{\mathrm{arm}})\mid X]
=\int_{\mathbb R}y\,Q_{a_{\mathrm{arm}},P}(dy\mid X)
\qquad\mathbb Q_P\text{-almost surely},
\quad a_{\mathrm{arm}}\in\{0,1\}.
\]
Because \(X\) has law \(\lambda\), the observational arm-mean version identities hold also at \(X\), almost surely under \(\mathbb Q_P\):
\[
\int y\,Q_{0,P}(dy\mid X)=m_{0,P}(X),
\qquad
\int y\,Q_{1,P}(dy\mid X)=m_{0,P}(X)+\tau_P(X).
\]
Linearity of conditional expectation, justified by the integrability established above, now yields
\[
\begin{aligned}
\mathbb E_{\mathbb Q_P}[Y(1)-Y(0)\mid X]
&=\mathbb E_{\mathbb Q_P}[Y(1)\mid X]
 -\mathbb E_{\mathbb Q_P}[Y(0)\mid X]\\
&=(m_{0,P}(X)+\tau_P(X))-m_{0,P}(X)
=\tau_P(X)
\end{aligned}
\qquad\mathbb Q_P\text{-almost surely}.
\]
% lean: CausalSmith.Stat.FinitepHomogeneityDensegamma.completion_conditional_effect

\item Finally, suppose \(d(P)=0\). The defining integral is nonnegative, so
\[
\int_0^1(\tau_P(x)-\bar\tau_P)^2\,d\lambda(x)=0.
\]
Its integrand is nonnegative and bounded by one, as established in \cref{proof:causal-nonempty:distance-bounds}. The zero-integral criterion therefore gives
\[
(\tau_P(x)-\bar\tau_P)^2=0
\quad\text{and hence}\quad
\tau_P(x)=\bar\tau_P
\qquad\text{for \(\lambda\)-almost every }x.
\]
Continuity and the full support of \(\lambda\) upgrade this to
\[
\tau_P(x)=\bar\tau_P
\qquad(x\in[0,1]).
\]
The average bound in \cref{proof:causal-nonempty:distance-bounds} gives
\[
-\frac12\le\bar\tau_P\le\frac12.
\]
Together with \(P\in\mathcal M_v\), this is exactly the membership condition in \cref{obj:def:null}.

Conversely, if \(P\in H_0(v)\), that definition supplies a constant
\[
c_{\mathrm{null}}\in[-1/2,1/2],
\qquad
\tau_P(x)=c_{\mathrm{null}}
\quad(x\in[0,1]).
\]
Since \(\lambda\) is a probability measure,
\[
\bar\tau_P=c_{\mathrm{null}},
\qquad
d(P)
=\left\{\int_0^1(c_{\mathrm{null}}-c_{\mathrm{null}})^2
\,d\lambda(x)\right\}^{1/2}
=0.
\]
Thus
\[
d(P)=0\quad\Longleftrightarrow\quad P\in H_0(v),
\]
completing the proof.
% lean: CausalSmith.Stat.FinitepHomogeneityDensegamma.hetDist_zero_iff_null
\end{enumerate}
\end{proof}

The uniqueness assertion makes every pointwise restriction on the propensity and mean functions a property of the observed law. The uniform design and overlap conditions in \cref{obj:ass:uniform,obj:ass:overlap} give the conditional arm distributions a common Lebesgue almost-everywhere interpretation; continuity makes their canonical mean representatives unambiguous across the whole interval. In particular, the treated mean inherits the regularity stated in \cref{obj:lem:causal-nonempty} from the separate baseline and contrast restrictions.

Version invariance also applies to the potential-outcome completion. As specified in \cref{obj:def:causal-completion}, changing either conditional outcome distribution on a Lebesgue-null set leaves \(\mathbb Q_P\) unchanged. The completion combines the observed arm distributions with conditionally independent assignment and reproduces the original-record marginal. Its integrable potential-outcome difference has conditional mean \(\tau_P(X)\), so \(d(P)\) measures variation in conditional average treatment effects under this completion. The equivalence between zero distance and membership in \cref{obj:def:null} connects the population criterion to the pointwise constant-effect restriction.

The two witnesses in \cref{obj:lem:causal-nonempty} give the testing problem a nonempty null and a positive separation domain for every admissible tuple. The heterogeneous witness attains the same fixed distance \(d_0\) throughout the parameter domain, while the constant-contrast witness permits a nonzero common effect. Together with the upper bound on \(D_v\), these properties give the cap in \cref{obj:def:critical-radius} a finite, positive population scale.

\subsection{Mean preservation and the benchmark experiment}

The bounded and signed-binary models in \cref{obj:def:bounded-model,obj:def:binary-model} retain the primitive restrictions on assignment and conditional means. Their comparison in \cref{obj:prop:bounded-submodel-comparison} uses two related operations: a model-level conversion specified by conditional arm means, and a record-level randomization available for bounded observations. At the model level, the converted outcome has probabilities
\[
\Pr(+1\mid X=x,A=a)=\frac{1+m_{a,P}(x)}2,
\qquad
\Pr(-1\mid X=x,A=a)=\frac{1-m_{a,P}(x)}2,
\qquad a\in\{0,1\}.
\]
These probabilities describe the signed-binary conversion in \cref{obj:prop:bounded-submodel-comparison}. It retains the propensity, both conditional arm means, and hence the contrast and its distance from constancy. The same result establishes that the bounded, signed-binary, and original models have equal distance suprema.

For a bounded original record, the conversion can instead use its observed outcome directly. To express this randomization on the full real-outcome test carrier, define the measurable sign map
\[
s(y,t)
=
2\mathbf 1\!\left\{
t\le \frac{1+\max\{-1,\min\{y,1\}\}}2
\right\}-1,
\qquad (y,t)\in\mathbb R\times[0,1].
\]
For \(|y|\le1\), a uniform \(t\) gives the probabilities \((1+y)/2\) and \((1-y)/2\) for the two signs, with conditional mean \(y\). The truncation in this definition specifies the map on the entire original-record carrier. Its restriction to the bounded model is precisely the record-level conversion in \cref{obj:prop:bounded-submodel-comparison}.

Write \(\boldsymbol V=(V_1,\ldots,V_n)\) for independent uniform variables on \([0,1]\), independent of the sample and of the public seed \(U\) in \cref{obj:def:original-record-experiment}. The randomized signed-binary sample is
\[
\mathcal D_n^{\,s}(\boldsymbol V)
=
\bigl((X_i,A_i,s(Y_i,V_i))\bigr)_{i=1}^n.
\]
For a test \(\phi\) on the carrier of \cref{obj:def:testing-risk}, its integrated rejection-probability map is
\[
\widetilde\phi(\mathcal D_n,u)
=
\int_{[0,1]^n}
\phi\!\left(
\bigl((X_i,A_i,s(Y_i,t_i))\bigr)_{i=1}^n,u
\right)
\,\lambda^{\otimes n}(d\boldsymbol t),
\qquad
\boldsymbol t=(t_1,\ldots,t_n).
\]
Thus \(\widetilde\phi\) is a jointly Borel map from
\[
\bigl([0,1]\times\{0,1\}\times\mathbb R\bigr)^n\times[0,1]
\quad\text{to}\quad[0,1].
\]
It uses the original-record experiment and its public seed; the additional sign randomization is averaged into the rejection probability.

For \(P\in\mathcal M_w^{\mathrm b}\), let \(P^{\,s}\) denote its signed-binary conversion from \cref{obj:prop:bounded-submodel-comparison}. The rejection-probability preservation assertion there takes the explicit form
\[
\mathsf R_n(P,\widetilde\phi)
=
\mathsf R_n(P^{\,s},\phi).
\]
This identity expresses the statistical role of mean-preserving randomization: rejection probabilities under each bounded law transfer exactly to those under its signed-binary conversion, including laws in the constant-effect null and laws at a specified heterogeneity distance.

\Cref{obj:prop:bounded-submodel-comparison} consequently gives the exact benchmark identity
\[
B_n^{\mathrm{bin}}(w,r)=B_n^{\mathrm b}(w,r),
\qquad n\ge2,\quad 0<r<D_v,
\]
for the risks in \cref{obj:def:binary-risk,obj:def:bounded-risk}. With the common distance supremum and the capping conventions in \cref{obj:def:binary-radius,obj:def:bounded-radius}, the same comparison gives equal capped critical radii for every \(n\ge2\). The signed-binary benchmark therefore represents the bounded mean-effect testing problem at the level of attainable distances, minimax risks, and critical radii.


\section{Bias, covariance, profiling, and attainment}

The calibration in \cref{obj:thm:whole-null-calibration-resolved} rests on three population quantities: the mean of the original-record score, the covariance of its affine coefficients, and the fluctuation of the cross-block inner product uniformly over candidate constants. This appendix organizes those quantities around the public construction in \cref{obj:def:explicit-score-ledger}. The resulting bounds connect the finite score calculation to the explicit separation guarantee.

\subsection{Calibration and population decomposition}

The deterministic calibration handle collects the score construction, its affine coefficients, the public bounds, and the profiling operation. Keeping these components together makes their dependence on the sample size and public exponent tuple explicit.

\begin{algorithmv}{def:calibration-handle}[Deterministic calibration handle]
The calibration handle, indexed by a sample size \(n\in\mathbb N\) and a public exponent tuple \(v=(p,\alpha,\beta,\gamma)\in\mathbb R^4\), is the ordered collection of the following seven components. Its branch, ranks, cutoffs, ledgers, and public grids are the deterministic selections from \(n\) and \(v\) in \cref{obj:def:explicit-score-ledger}.

\begin{enumerate}
\item The tuning collection
\[
\bigl(M,K,T_0,(T_j)_{j=1}^{L}\bigr),
\qquad M=J_{\mathrm c},\quad K=J_{\mathrm f},
\]
with the number of increments \(L\) specified in \cref{obj:def:explicit-score-ledger}.

\item The block-score family, taking a block index, a real candidate constant \(c\), and a dataset of \(n\) original records to a vector in \(\mathbb R^M\). For \(n\ge4\), it is
\[
Z_b(c)=
\begin{cases}
H_{b,T_0}(c)-U_{b,\Pi_K,T_0}(c),
&\text{single-correction branch},\\[1mm]
H_{b,T_0}(c)-U_{b,\Pi_M,T_0}(c)
-\displaystyle\sum_{j=1}^{L}U_{b,D_j,T_j}(c),
&\text{multiresolution branch},
\end{cases}
\]
where the block and kernel conventions are those of
\cref{obj:def:blocks,obj:def:score-family,obj:def:explicit-score-ledger}.
For \(n<4\), \(Z_b(c)=0\).

\item The two affine-coefficient maps, indexed by the block and a binary coefficient selector, with values
\[
Z_{b,0}=Z_b(0),
\qquad
Z_{b,A}=Z_b(0)-Z_b(1).
\]
Thus, for every real \(c\) and every dataset,
\[
Z_b(c)=Z_{b,0}-cZ_{b,A}.
\]

\item The deterministic ledger triple
\[
\bigl(\mathfrak b_{n,v},\Lambda_{n,v},h_{n,v}\bigr),
\]
whose bias, covariance, and rejection-cutoff selections are exactly those of
\cref{obj:def:explicit-score-ledger}.

\item The dataset-to-real profiling map
\[
\inf_{c\in[-1/2,1/2]}W(c),
\qquad
W(c)=\langle Z_1(c),Z_2(c)\rangle,
\]
using the two-block indexing convention of
\cref{obj:def:explicit-score-ledger}.

\item The original-record rejection rule specified in
\cref{obj:def:explicit-score-ledger}, namely, for \(n\ge4\),
\[
\mathbf 1\!\left\{
\inf_{c\in[-1/2,1/2]}W(c)>h_{n,v}
\right\},
\]
and the zero rejection rule for \(n<4\).

\item The ordered triple of finite public rank, outcome-cutoff, and rejection-cutoff grids selected in \cref{obj:def:explicit-score-ledger}. The first two grids depend on \(n\), and the third depends on \(n\) and \(v\).
\end{enumerate}
\end{algorithmv}

The affine coefficient \(Z_{b,A}\) records the score's response to the candidate constant, while \(Z_{b,0}\) records its outcome contribution. Their joint population behavior is therefore sufficient to describe the entire profiled statistic. The mean handle below retains both coefficients, their singleton and canonical pair projections, and all four combinations of single-record and pair contributions.

\begin{algorithmv}{def:mean-handle}[Original score mean handle]
The original score mean handle is the ordered pair of multiresolution and single-correction handles associated with the two score members of \cref{obj:def:calibration-handle}. Its parameters are an original-record law \(P\), ranks \(J,K,L\in\mathbb N\), a main cutoff \(T_{\mathrm{main}}\in\mathbb R\), and increment cutoffs \(T_j^{\mathrm{inc}}\in\mathbb R\) indexed by \(1\le j\le L\). Each branch retains, in order, a mean map indexed by \(c\in\mathbb R\), coefficient means, singleton projections, canonical pair projections, four population contractions for each coefficient pair, and a sample statistic indexed by \(n\) and \(c\).

The multiresolution branch uses coarse rank \(J_{\mathrm c}=J\) and fine rank \(J_{\mathrm f}=2^LJ\); the single-correction branch uses
\[
J_{\mathrm c}=J,\qquad J_{\mathrm f}=K.
\]
Both branches use \(T_{\mathrm{main}}\); the multiresolution branch also uses the increment cutoffs \(T_j^{\mathrm{inc}}\), \(1\le j\le L\). Define the projected propensity and propensity weight by
\[
e_K=\Pi_{J_{\mathrm f}}e_P,
\qquad
w_K(x)=e_P(x)\{1-e_K(x)\}.
\]
The retained mean map is
\[
\int_0^1 F_{J_{\mathrm c}}(x)
\left[
w_K(x)\{\tau_P(x)-c\}
+\{e_P(x)-e_K(x)\}m_{0,P}(x)
\right]\,dx
\]
plus the exact multiresolution or single-correction clipping-tail remainder, respectively, specified in \cref{obj:lem:original-score-mean-ledger}.

This retained map equals the population expectation of the corresponding corrected block score for every \(c\in\mathbb R\) and either evaluation block under the following conditions:
\begin{itemize}
\item \textbf{(Model.)} The public tuple \(v\) is admissible and \(P\in\mathcal M_v\), as specified in \cref{obj:def:model}.
\item \textbf{(Sample and coarse rank.)} \(n\ge4\) and \(J>0\).
\item \textbf{(Main cutoff.)} \(T_{\mathrm{main}}\ge1\).
\item \textbf{(Multiresolution cutoffs.)} In the multiresolution branch, \(T_j^{\mathrm{inc}}\ge1\) for every integer \(1\le j\le L\).
\item \textbf{(Single-correction refinement.)} In the single-correction branch, \(K>0\) and \(J\mid K\).
\end{itemize}

For each affine coefficient \(r\in\{0,1\}\), let \(h_r\) and \(g_r\) be the selected branch's raw single-record and symmetric pair kernels from \cref{obj:def:calibration-handle}. All expectations below use independent original records with law \(P\). Define the centered pair-kernel singleton projection by
\[
g_{r,1}(o)
=
\mathbb E[g_r(o,O')]
-\mathbb E[g_r(O,O')].
\]
The retained coefficient means, singleton projections, and canonical pair projections are, respectively,
\[
\mathbb E[h_r(O)]+\mathbb E[g_r(O,O')],
\]
\[
h_r(o)-\mathbb E[h_r(O)]+2g_{r,1}(o),
\]
\[
g_r(o,z)-\mathbb E[g_r(O,O')]
-g_{r,1}(o)-g_{r,1}(z).
\]
For every \(r,r'\in\{0,1\}\), the retained ordered quadruple of population contractions is defined by
\[
h_r\!\times h_{r'}
=
\mathbb E\langle h_r(O_1),h_{r'}(O_2)\rangle,
\]
\[
h_r\!\times g_{r'}
=
\mathbb E\langle h_r(O_1),g_{r'}(O_2,O_3)\rangle,
\]
\[
g_r\!\times h_{r'}
=
\mathbb E\langle g_r(O_1,O_2),h_{r'}(O_3)\rangle,
\]
\[
g_r\!\times g_{r'}
=
\mathbb E\langle g_r(O_1,O_2),g_{r'}(O_3,O_4)\rangle,
\]
where \(O_1,O_2,O_3,O_4\) are independent with law \(P\).

For each candidate constant \(c\in\mathbb R\), set
\[
h_c=h_0-ch_1,
\qquad
g_c=g_0-cg_1.
\]
For \(n\ge4\), use the block size \(s_n=\lfloor n/2\rfloor\) and the two deterministic blocks \(\mathcal I_1,\mathcal I_2\) of \cref{obj:def:blocks}. Define
\[
Z_b(c)=
\frac1{s_n}\sum_{i\in\mathcal I_b}h_c(O_i)
+
\frac1{s_n(s_n-1)}
\sum_{\substack{i,j\in\mathcal I_b\\i\ne j}}
g_c(O_i,O_j),
\qquad b\in\{1,2\}.
\]
The retained sample output is the single function \(c\mapsto W(c)\), given by
\[
\begin{aligned}
W(c)
&=\langle Z_1(c),Z_2(c)\rangle\\
&=\frac1{s_n^2}
\sum_{\substack{i\in\mathcal I_1\\j\in\mathcal I_2}}
\langle h_c(O_i),h_c(O_j)\rangle\\
&\quad+\frac1{s_n^2(s_n-1)}
\sum_{\substack{i\in\mathcal I_1\\j,k\in\mathcal I_2,\ j\ne k}}
\langle h_c(O_i),g_c(O_j,O_k)\rangle\\
&\quad+\frac1{s_n^2(s_n-1)}
\sum_{\substack{i,j\in\mathcal I_1,\ i\ne j\\k\in\mathcal I_2}}
\langle g_c(O_i,O_j),h_c(O_k)\rangle\\
&\quad+\frac1{s_n^2(s_n-1)^2}
\sum_{\substack{i,j\in\mathcal I_1,\ i\ne j\\k,l\in\mathcal I_2,\ k\ne l}}
\langle g_c(O_i,O_j),g_c(O_k,O_l)\rangle.
\end{aligned}
\]
For \(n<4\), its defining value is
\[
W(c)=0.
\]
The two mixed affine-coefficient contributions are retained through their sum in \(W(c)\).
\end{algorithmv}

The conventions in \cref{obj:def:mean-handle,obj:def:calibration-handle} use the same one-based blocks \(\mathcal I_1,\mathcal I_2\). In the mean calculation, \(T_{\mathrm{main}}\) denotes the procedure's main cutoff \(T_0\), while \(T_j^{\mathrm{inc}}\) denotes its increment cutoff \(T_j\) for \(1\le j\le L\). Thus the two cutoff families and the two block scores have distinct names without a positional shift. At the public selections, the coarse rank is \(J=M\) and the multiresolution fine rank is \(2^LJ=K\).

For clarity about the symmetric pair convention, write \(o=(x,a,y)\) and \(z=(x',a',y')\). For a correction function \(G\) and cutoff \(T\), define its signed, symmetrized contribution by
\[
k_{G,T,c}(o,z)
:=-\frac12\left[
F_M(x)aG(x,x')\{\ell_T(y')-ca'\}
+
F_M(x')a'G(x',x)\{\ell_T(y)-ca\}
\right].
\]
The single-record contribution is
\[
h_c(o):=F_M(x)a\{\ell_{T_0}(y)-c\}.
\]
In the single-correction branch, \(g_c:=k_{\Pi_K,T_0,c}\). In the multiresolution branch, using the increment indices of \cref{obj:def:explicit-score-ledger},
\[
g_c:=k_{\Pi_M,T_0,c}
+\sum_{j=1}^{L}k_{D_j,T_j,c}.
\]
The affine coefficients in \cref{obj:def:mean-handle} are thus \(h_0=h_c|_{c=0}\), \(h_1=h_c|_{c=0}-h_c|_{c=1}\), and the corresponding two coefficients of \(g_c\). Symmetrization preserves the ordered-pair average because that average contains both orientations of every distinct pair.

The four contractions in \cref{obj:def:mean-handle} retain the single--single, single--pair, pair--single, and pair--pair contributions separately. Their independent-record arguments match the four finite sums defining \(W(c)\). Both orientations of the mixed contribution are retained, including their dependence on the two affine coefficients.

\subsection{Original-mean bias}

The population calculation separates a weighted effect contrast, a baseline approximation term, and clipping tails. To specify the latter explicitly, let \(\lambda\) denote Lebesgue probability measure on \([0,1]\), and define the conditional clipping discrepancy and its unconditional conditional mean by
\[
\delta_{a,T}(x)
:=\mathbb E_P[\ell_T(Y)-Y\mid X=x,A=a],
\qquad
r_T(x):=e_P(x)\delta_{1,T}(x)
+\{1-e_P(x)\}\delta_{0,T}(x).
\]
These quantities are defined up to \(\lambda\)-almost-everywhere equality; their integrated contributions use that convention. Under the mean handle's conditions, the clipping-tail remainders for the two branches are
\[
\begin{aligned}
\mathcal T_{\mathrm{single}}
&:=\int_0^1 F_M(x)
\left[e_P(x)\delta_{1,T_0}(x)
-e_K(x)r_{T_0}(x)\right]\,d\lambda(x),\\
\mathcal T_{\mathrm{multi}}
&:=\int_0^1 F_M(x)
\left[
e_P(x)\delta_{1,T_0}(x)
-(\Pi_Me_P)(x)r_{T_0}(x)
-\sum_{j=1}^{L}(D_je_P)(x)r_{T_j}(x)
\right]\,d\lambda(x).
\end{aligned}
\]
Here \((D_je_P)(x):=\int_0^1D_j(x,z)e_P(z)\,d\lambda(z)\), and the increment indices follow \cref{obj:def:explicit-score-ledger}. Each remainder is independent of the candidate constant.

The retained mean in \cref{obj:def:mean-handle} therefore identifies the baseline approximation contribution through \((e_P-e_K)m_{0,P}\), while the clipping remainder records the outcome-tail contribution at each selected resolution. The public bound \(B\) in \cref{obj:def:explicit-score-ledger} aggregates these terms. The next statement controls that bound uniformly over the candidate constants and relates the score mean to heterogeneity.

\begin{lemmav}{lem:original-score-mean-ledger}[Original score mean bounds]
Let \(v=(p,\alpha,\beta,\gamma)\) be an admissible tuple and \(P\in\mathcal M_v\), as specified in \cref{obj:def:model}, and let \(n\ge4\) be an integer. Use the explicit score ledger's coarse rank \(M\), fine rank \(K\), corrected block scores \(Z_b(c)\), and deterministic bias bound \(B\).

Let \(\lambda\) denote Lebesgue probability measure on \([0,1]\), and define the projected propensity and its associated weight by
\[
e_K(x)=(\Pi_K e_P)(x),
\qquad
w_K(x)=e_P(x)\bigl(1-e_K(x)\bigr).
\]
Define the population heterogeneity distance by
\[
d(P)=
\left(
\int_0^1
\left[\tau_P(x)-\int_0^1\tau_P(z)\,d\lambda(z)\right]^2
\,d\lambda(x)
\right)^{1/2}.
\]

For every block index \(b\) and every \(c\in\mathbb R\), \(Z_b(c)\) is integrable under the original independent-record law \(P^{\otimes n}\). Its population mean is
\[
\theta_P(c)
=
\mathbb E_{P^{\otimes n}}\!\left[Z_b(c)\right].
\]
For every block index \(b\) and every \(c\in\mathbb R\) with \(|c|\le\tfrac12\),
\[
\left\|
\theta_P(c)
-
\int_0^1
w_K(x)\bigl(\tau_P(x)-c\bigr)F_M(x)\,d\lambda(x)
\right\|_2
\le B,
\qquad
\|\theta_P(c)\|_2
\ge
\frac{3}{16}d(P)-16M^{-\gamma}-B.
\]
Furthermore, if \(P\in H_0(v)\) as defined in \cref{obj:def:null}, then for every block index \(b\) and every \(c_0\in\mathbb R\) satisfying \(\tau_P(x)=c_0\) for all \(x\in[0,1]\),
\[
\|\theta_P(c_0)\|_2\le B.
\]
\end{lemmav}

\begin{proof}[Proof of \cref{obj:lem:original-score-mean-ledger}]
\begin{enumerate}
\item
Fix a block index \(b\). The clipped outcomes, treatment indicators, histogram features, and histogram kernels are bounded measurable functions for every fixed choice of the public ranks and cutoffs. In each branch of \cref{obj:def:explicit-score-ledger}, \(Z_b(c)\) is a finite linear combination of products of these functions. Consequently, for every \(c\in\mathbb R\),
\[
\int \|Z_b(c)\|_2\,dP^{\otimes n}<\infty,
\qquad
\theta_P(c)=\int Z_b(c)\,dP^{\otimes n}.
\]
The mean calculations below have the same right-hand side for both blocks.
% lean: CausalSmith.Stat.FinitepHomogeneityDensegamma.integrable_ledgerScore

\item
The ranks in \cref{obj:def:explicit-score-ledger} are positive powers of two, and the definition of the fine rank gives
\[
0<M\le K,\qquad M\mid K.
\]
Indeed, the dyadic exponent of \(K\) is at least that of \(M\), so their quotient is an integer power of two. The cutoffs satisfy
\[
T_0\ge1,\qquad T_j\ge1\quad(1\le j\le L)
\]
in the branch where the increment cutoffs occur.

We record the histogram identities needed for the mean calculation. For every positive integer \(J\) and measurable integrable function \(f:[0,1]\to\mathbb R\), define
\[
\operatorname{coef}_J(f)
=
\int_0^1 f(x)F_J(x)\,d\lambda(x)\in\mathbb R^J.
\]
The cells in \cref{obj:def:projection} partition \([0,1]\), with the stipulated convention at the right endpoint, and each has measure \(1/J\). Thus, for \(x\in I_{j,J}\),
\[
(\Pi_J f)(x)=J\int_{I_{j,J}}f(z)\,d\lambda(z),
\qquad
\bigl(\operatorname{coef}_J(f)\bigr)_j
=
\sqrt J\int_{I_{j,J}}f(z)\,d\lambda(z).
\]
In particular, if \(\mathfrak c\in[0,\infty)\) and \(|f|\le\mathfrak c\) almost everywhere, then
\[
\left|\bigl(\operatorname{coef}_J(f)\bigr)_j\right|
\le\frac{\mathfrak c}{\sqrt J},
\qquad
\|\operatorname{coef}_J(f)\|_2^2
\le\sum_{j=1}^J\frac{\mathfrak c^2}{J}
=\mathfrak c^2.
\]
Hence
\[
\|\operatorname{coef}_J(f)\|_2\le\mathfrak c.
\]
Cell averaging also preserves almost-everywhere bounds and constants.

For a function \(f\) with Hölder exponent \(\mathfrak s\in(0,1]\) and Hölder constant at most \(20\), the diameter of a cell gives, for \(x\in I_{j,J}\),
\[
\begin{aligned}
|f(x)-(\Pi_Jf)(x)|
&\le J\int_{I_{j,J}}|f(x)-f(z)|\,d\lambda(z)\\
&\le20J\int_{I_{j,J}}|x-z|^{\mathfrak s}\,d\lambda(z)
\le20J^{-\mathfrak s}.
\end{aligned}
\]

If \(M\mid J\), the fine cells partition each coarse cell. Averaging their integrals therefore yields
\[
\Pi_M(\Pi_Jf)=\Pi_Mf,
\qquad
\Pi_J(\Pi_Jf)=\Pi_Jf.
\]
Moreover, the support of \(\Pi_J(x,z)\) places \(x\) and \(z\) in the same coarse cell, so
\[
\Pi_J(x,z)F_M(x)=\Pi_J(z,x)F_M(z).
\]
For a bounded measurable function \(\varphi_{\mathrm{bd}}:[0,1]\to\mathbb R\) and a measurable integrable function \(\varphi_{\mathrm{int}}:[0,1]\to\mathbb R\), Fubini's theorem consequently gives
\[
\begin{aligned}
\operatorname{coef}_M\bigl(\varphi_{\mathrm{bd}}\Pi_J\varphi_{\mathrm{int}}\bigr)
&=\int_0^1\int_0^1
\varphi_{\mathrm{bd}}(x)\Pi_J(x,z)\varphi_{\mathrm{int}}(z)F_M(x)\,d\lambda(z)\,d\lambda(x)\\
&=\int_0^1 \varphi_{\mathrm{int}}(z)F_M(z)
\left[\int_0^1\Pi_J(z,x)\varphi_{\mathrm{bd}}(x)\,d\lambda(x)\right]d\lambda(z)\\
&=\operatorname{coef}_M\bigl((\Pi_J\varphi_{\mathrm{bd}})\varphi_{\mathrm{int}}\bigr).
\end{aligned}
\]
The double integral is absolutely integrable because \(\varphi_{\mathrm{bd}}\), the kernel, and the features are bounded.
% lean: CausalSmith.Stat.FinitepHomogeneityDensegamma.ledger_increment_rank_bound
% lean: CausalSmith.Stat.FinitepHomogeneityDensegamma.ledger_coarse_dvd_fine
% lean: CausalSmith.Stat.FinitepHomogeneityDensegamma.ledgerT0_ge_one
% lean: CausalSmith.Stat.FinitepHomogeneityDensegamma.ledgerT_ge_one
% lean: CausalSmith.Stat.FinitepHomogeneityDensegamma.histogram_coefficient_norm_le
% lean: CausalSmith.Stat.FinitepHomogeneityDensegamma.projection_holder_error
% lean: CausalSmith.Stat.FinitepHomogeneityDensegamma.projection_nesting
% lean: CausalSmith.Stat.FinitepHomogeneityDensegamma.histogram_feature_selfAdjoint_integrable

\item
For every positive integer \(J\) divisible by \(M\) and every \(c\in\mathbb R\), define the original mean vector
\[
\theta^{\mathrm{orig}}_{P,J}(c)
=
\operatorname{coef}_M\!\left(
e_P(1-\Pi_Je_P)(\tau_P-c)
+(e_P-\Pi_Je_P)m_{0,P}
\right).
\]
Also define the weighted effect vector
\[
\theta^{\mathrm{wt}}_P(c)
=
\operatorname{coef}_M\bigl(w_K(\tau_P-c)\bigr).
\]
Their difference at the fine rank is
\[
\theta^{\mathrm{orig}}_{P,K}(c)-\theta^{\mathrm{wt}}_P(c)
=
\operatorname{coef}_M\bigl((e_P-e_K)m_{0,P}\bigr).
\]
Idempotence and the self-adjointness identity established above imply
\[
\Pi_K(e_P-e_K)=0,
\qquad
\operatorname{coef}_M\bigl((e_P-e_K)\Pi_Km_{0,P}\bigr)=0.
\]
Subtracting the latter zero vector gives the residual identity
\[
\operatorname{coef}_M\bigl((e_P-e_K)m_{0,P}\bigr)
=
\operatorname{coef}_M\bigl(
(e_P-e_K)(m_{0,P}-\Pi_Km_{0,P})
\bigr).
\]
The smoothness restrictions in
\cref{obj:ass:propensity-smoothness,obj:ass:baseline-smoothness},
together with the Hölder estimate above, give
\[
|e_P-e_K|\le20K^{-\alpha},
\qquad
|m_{0,P}-\Pi_Km_{0,P}|\le20K^{-\beta}.
\]
Since \(S=\alpha+\beta\) in \cref{obj:def:explicit-score-ledger},
\[
\left|(e_P-e_K)(m_{0,P}-\Pi_Km_{0,P})\right|
\le400K^{-S}.
\]
Applying the coefficient norm bound proves, for every real \(c\),
\[
\|\theta^{\mathrm{orig}}_{P,K}(c)-\theta^{\mathrm{wt}}_P(c)\|_2
\le400K^{-S}.
\]
% lean: CausalSmith.Stat.FinitepHomogeneityDensegamma.baseline_projection_residual_identity
% lean: CausalSmith.Stat.FinitepHomogeneityDensegamma.baseline_projection_residual_norm_le
% lean: CausalSmith.Stat.FinitepHomogeneityDensegamma.origMeanVector_weighted_error_le

\item
First consider the single-correction branch of
\cref{obj:def:explicit-score-ledger}. We derive the mean identity for any positive \(J\) divisible by \(M\), so that it can also be used at successive ranks in the other branch.

Use the clipping function \(\ell_T\) and the score formulas defined in \cref{obj:def:score-family}. For \(T\ge1\), \(x\in[0,1]\), and \(\mathfrak a\in\{0,1\}\), define the arm-specific and pooled clipping remainders by
\[
r_T^{(\mathfrak a)}(x)
=
\int_{\mathbb R}\bigl(y-\ell_T(y)\bigr)
\,Q_{\mathfrak a,P}(dy\mid x),
\qquad
r_T(x)
=
e_P(x)r_T^{(1)}(x)+(1-e_P(x))r_T^{(0)}(x).
\]
At a covariate where the integrand in an arm-specific integral is nonintegrable, assign that integral the value zero. The moment restriction makes these integrands integrable for almost every covariate.

Indeed, admissibility gives \(p>1\), and
\[
|y|\le1+|y|^p.
\]
Thus \cref{obj:ass:raw-moment} supplies integrability of the original outcome in each conditional arm almost everywhere. For \(|y|\le T\), the clipping remainder vanishes; for \(|y|>T\),
\[
|y-\ell_T(y)|
\le |y|
=|y|^p|y|^{1-p}
\le |y|^pT^{1-p}.
\]
Integration therefore gives, almost everywhere,
\[
|r_T^{(\mathfrak a)}(x)|
\le
T^{1-p}\int_{\mathbb R}|y|^pQ_{\mathfrak a,P}(dy\mid x)
\le10T^{1-p}.
\]
Taking the propensity-weighted average yields
\[
|r_T(x)|
\le e_P(x)\,10T^{1-p}
+(1-e_P(x))\,10T^{1-p}
=10T^{1-p}.
\]

For \(c\in\mathbb R\), define
\[
f^{\mathrm{mark}}_{P,c}(x)
=
m_{0,P}(x)+e_P(x)(\tau_P(x)-c).
\]
The conditional arm means specified in \cref{obj:def:model} give, almost everywhere,
\[
\mathbb E_P[\ell_T(Y)-cA\mid X=x]
=
f^{\mathrm{mark}}_{P,c}(x)-r_T(x),
\]
\[
\mathbb E_P[A(\ell_T(Y)-c)\mid X=x]
=
e_P(x)\bigl(m_{0,P}(x)+\tau_P(x)-c-r_T^{(1)}(x)\bigr).
\]
By \cref{obj:def:blocks}, \(s_n\ge2\), and each block contains \(s_n\) records and \(s_n(s_n-1)\) ordered pairs of distinct records. Identical distribution cancels the single-record normalization, and independence cancels the ordered-pair normalization. Conditioning and
\cref{obj:ass:uniform} consequently give
\[
\begin{aligned}
\mathbb E_{P^{\otimes n}}[H_{b,T}(c)]
&=
\operatorname{coef}_M\bigl(e_P(m_{0,P}+\tau_P-c)\bigr)
-\operatorname{coef}_M\bigl(e_Pr_T^{(1)}\bigr),\\
\mathbb E_{P^{\otimes n}}[U_{b,\Pi_J,T}(c)]
&=
\operatorname{coef}_M\bigl(e_P\Pi_Jf^{\mathrm{mark}}_{P,c}\bigr)
-\operatorname{coef}_M\bigl(e_P\Pi_Jr_T\bigr).
\end{aligned}
\]
Self-adjointness moves the projection in the first correction term:
\[
\operatorname{coef}_M\bigl(e_P\Pi_Jf^{\mathrm{mark}}_{P,c}\bigr)
=
\operatorname{coef}_M\bigl((\Pi_Je_P)f^{\mathrm{mark}}_{P,c}\bigr).
\]
The resulting original-outcome difference is exactly
\[
\begin{aligned}
&\operatorname{coef}_M\bigl(
e_P(m_{0,P}+\tau_P-c)-(\Pi_Je_P)f^{\mathrm{mark}}_{P,c}
\bigr)\\
&\hspace{2em}
=
\operatorname{coef}_M\bigl(
e_P(1-\Pi_Je_P)(\tau_P-c)+(e_P-\Pi_Je_P)m_{0,P}
\bigr)
=
\theta^{\mathrm{orig}}_{P,J}(c).
\end{aligned}
\]
Define, for these \(J\) and \(T\),
\[
\theta^{\mathrm{tail}}_{P,J}(T)
=
-\operatorname{coef}_M\bigl(e_Pr_T^{(1)}\bigr)
+\operatorname{coef}_M\bigl(e_P\Pi_Jr_T\bigr).
\]
We have proved the exact identity
\[
\mathbb E_{P^{\otimes n}}
[H_{b,T}(c)-U_{b,\Pi_J,T}(c)]
=
\theta^{\mathrm{orig}}_{P,J}(c)
+\theta^{\mathrm{tail}}_{P,J}(T).
\]

Since \(0\le e_P\le1\), and cell averaging preserves the almost-everywhere bound on \(r_T\),
\[
|e_Pr_T^{(1)}|\le10T^{1-p}\quad\text{almost everywhere},
\qquad
|e_P\Pi_Jr_T|\le10T^{1-p}\quad\text{everywhere}.
\]
The coefficient norm bound and the triangle inequality give
\[
\begin{aligned}
\|\theta^{\mathrm{tail}}_{P,J}(T)\|_2
&\le
\|\operatorname{coef}_M(e_Pr_T^{(1)})\|_2
+\|\operatorname{coef}_M(e_P\Pi_Jr_T)\|_2\\
&\le10T^{1-p}+10T^{1-p}
=20T^{1-p}.
\end{aligned}
\]
In the single-correction branch, take \(J=K\) and \(T=T_0\). Combining this estimate with the original-mean error bound gives
\[
\begin{aligned}
\|\theta_P(c)-\theta^{\mathrm{wt}}_P(c)\|_2
&\le
\|\theta^{\mathrm{orig}}_{P,K}(c)-\theta^{\mathrm{wt}}_P(c)\|_2
+\|\theta^{\mathrm{tail}}_{P,K}(T_0)\|_2\\
&\le400K^{-S}+20T_0^{1-p}
=B.
\end{aligned}
\]
This calculation applies to every \(c\in\mathbb R\).
% lean: CausalSmith.Stat.FinitepHomogeneityDensegamma.conditional_truncation_moments
% lean: CausalSmith.Stat.FinitepHomogeneityDensegamma.tailRemainder_abs_le
% lean: CausalSmith.Stat.FinitepHomogeneityDensegamma.singleMean_exact
% lean: CausalSmith.Stat.FinitepHomogeneityDensegamma.singleMeanTail_norm_le
% lean: CausalSmith.Stat.FinitepHomogeneityDensegamma.ledgerScore_weighted_mean_error_le

\item
Now consider the multiresolution branch. Because \(M\) and \(K\) are ordered powers of two, the increment chain in
\cref{obj:def:explicit-score-ledger} satisfies
\[
L=\log_2(K/M)\in\mathbb N_0,
\qquad
R_0=M,\qquad R_j=2R_{j-1},\qquad R_L=2^LM=K.
\]
For a measurable integrable \(f:[0,1]\to\mathbb R\), the action of its increment kernel is
\[
(D_jf)(x)
=
\int_0^1D_j(x,z)f(z)\,d\lambda(z)
=
(\Pi_{R_j}f)(x)-(\Pi_{R_{j-1}}f)(x),
\qquad 1\le j\le L.
\]
Linearity of the ordered-pair score gives
\[
U_{b,D_j,T_j}(c)
=
U_{b,\Pi_{R_j},T_j}(c)
-
U_{b,\Pi_{R_{j-1}},T_j}(c).
\]
The exact single-correction identity established above applies at both ranks. Also, subtracting the definitions of the two tail vectors gives
\[
\theta^{\mathrm{tail}}_{P,R_j}(T_j)
-
\theta^{\mathrm{tail}}_{P,R_{j-1}}(T_j)
=
\operatorname{coef}_M(e_PD_jr_{T_j}).
\]
Consequently,
\[
\mathbb E_{P^{\otimes n}}[U_{b,D_j,T_j}(c)]
=
\theta^{\mathrm{orig}}_{P,R_{j-1}}(c)
-
\theta^{\mathrm{orig}}_{P,R_j}(c)
-
\operatorname{coef}_M(e_PD_jr_{T_j}).
\]
The original means telescope:
\[
\sum_{j=1}^L
\left(
\theta^{\mathrm{orig}}_{P,R_{j-1}}(c)
-
\theta^{\mathrm{orig}}_{P,R_j}(c)
\right)
=
\theta^{\mathrm{orig}}_{P,M}(c)
-
\theta^{\mathrm{orig}}_{P,K}(c).
\]
Taking the expectation of the multiresolution score therefore yields
\[
\theta_P(c)
=
\theta^{\mathrm{orig}}_{P,K}(c)
+
\theta^{\mathrm{tail}}_{P,M}(T_0)
+
\sum_{j=1}^L\operatorname{coef}_M(e_PD_jr_{T_j}).
\]

For each increment, self-adjointness at its two ranks gives
\[
\operatorname{coef}_M(e_PD_jr_{T_j})
=
\operatorname{coef}_M\bigl((D_je_P)r_{T_j}\bigr).
\]
The propensity Hölder estimate and \(R_j=2R_{j-1}\) imply
\[
\begin{aligned}
|D_je_P|
&\le |e_P-\Pi_{R_j}e_P|
+|e_P-\Pi_{R_{j-1}}e_P|\\
&\le20R_j^{-\alpha}+20R_{j-1}^{-\alpha}
\le40R_{j-1}^{-\alpha}
=a_j.
\end{aligned}
\]
Together with the pooled tail bound, this gives
\[
|(D_je_P)r_{T_j}|
\le10a_jT_j^{1-p}\quad\text{almost everywhere},
\qquad
\|\operatorname{coef}_M(e_PD_jr_{T_j})\|_2
\le10a_jT_j^{1-p}.
\]
Hence the complete tail vector satisfies
\[
\left\|
\theta^{\mathrm{tail}}_{P,M}(T_0)
+\sum_{j=1}^L\operatorname{coef}_M(e_PD_jr_{T_j})
\right\|_2
\le
20T_0^{1-p}
+10\sum_{j=1}^La_jT_j^{1-p}.
\]
Combining this with the original-mean error bound proves
\[
\|\theta_P(c)-\theta^{\mathrm{wt}}_P(c)\|_2
\le
400K^{-S}+20T_0^{1-p}
+10\sum_{j=1}^La_jT_j^{1-p}
=B.
\]
When \(L=0\), \(K=M\), both increment sums are empty, and the same identities reduce to the single correction at rank \(M\). Thus, in either branch, we have established for every real \(c\)
\[
\left\|
\theta_P(c)
-
\int_0^1w_K(x)(\tau_P(x)-c)F_M(x)\,d\lambda(x)
\right\|_2
\le B.
\]
% lean: CausalSmith.Stat.FinitepHomogeneityDensegamma.ledger_final_rank_eq
% lean: CausalSmith.Stat.FinitepHomogeneityDensegamma.multiresMean_exact
% lean: CausalSmith.Stat.FinitepHomogeneityDensegamma.propensity_projection_increment_bound
% lean: CausalSmith.Stat.FinitepHomogeneityDensegamma.mean_increment_tail_norm_le
% lean: CausalSmith.Stat.FinitepHomogeneityDensegamma.multiresMeanTail_norm_le
% lean: CausalSmith.Stat.FinitepHomogeneityDensegamma.ledgerScore_weighted_mean_error_le

\item
Fix now \(c\in\mathbb R\) with \(|c|\le1/2\), and define
\[
f^{\mathrm{cell}}_{P,c}(x)
=
(\Pi_M\tau_P)(x)-c,
\qquad
\theta^{\mathrm{cell}}_P(c)
=
\operatorname{coef}_M\bigl(w_Kf^{\mathrm{cell}}_{P,c}\bigr).
\]
By \cref{obj:ass:overlap} and preservation of bounds under cell averaging,
\[
\frac14\le e_P(x),e_K(x)\le\frac34,
\qquad
0\le w_K(x)=e_P(x)(1-e_K(x))\le\frac9{16}.
\]
Since \(e_K\) is constant on each fine cell,
\[
(\Pi_Kw_K)(x)=e_K(x)(1-e_K(x)).
\]
For every real number in \([1/4,3/4]\), the product with its complement lies in \([3/16,1/4]\). Indeed,
\[
e_K(1-e_K)-\frac3{16}
=
\left(e_K-\frac14\right)\left(\frac34-e_K\right)\ge0,
\qquad
\frac14-e_K(1-e_K)
=
\left(e_K-\frac12\right)^2\ge0.
\]
Projection nesting and preservation of bounds now yield
\[
\Pi_Mw_K=\Pi_M(\Pi_Kw_K),
\qquad
\frac3{16}\le(\Pi_Mw_K)(x)\le\frac14.
\]

Define the coarse-cell average weights by
\[
\bar w_{j,M}
=
M\int_{I_{j,M}}w_K(x)\,d\lambda(x),
\qquad j\in\{1,\ldots,M\}.
\]
The function \(f^{\mathrm{cell}}_{P,c}\) is constant on each coarse cell, so the coordinate formula gives
\[
\bigl(\theta^{\mathrm{cell}}_P(c)\bigr)_j
=
\bar w_{j,M}
\bigl(\operatorname{coef}_M(f^{\mathrm{cell}}_{P,c})\bigr)_j,
\qquad
\bar w_{j,M}\ge\frac3{16}.
\]
Squaring and summing proves
\[
\|\theta^{\mathrm{cell}}_P(c)\|_2^2
\ge
\left(\frac3{16}\right)^2
\|\operatorname{coef}_M(f^{\mathrm{cell}}_{P,c})\|_2^2,
\]
and therefore
\[
\|\theta^{\mathrm{cell}}_P(c)\|_2
\ge
\frac3{16}
\|\operatorname{coef}_M(f^{\mathrm{cell}}_{P,c})\|_2.
\]
The effect cap in \cref{obj:ass:effect-cap} also gives the bounds used to integrate these vectors:
\[
|\Pi_M\tau_P|\le\frac12,
\qquad
|f^{\mathrm{cell}}_{P,c}|\le1,
\qquad
|\tau_P-c|\le1.
\]
% lean: CausalSmith.Stat.FinitepHomogeneityDensegamma.propensity_weight_abs_le
% lean: CausalSmith.Stat.FinitepHomogeneityDensegamma.coarse_propensity_weight_range
% lean: CausalSmith.Stat.FinitepHomogeneityDensegamma.weighted_histogram_coefficient_norm_ge

\item
To compare the unweighted coarse coefficients with heterogeneity, define
\[
f^{\mathrm{raw}}_{P,c}(x)=\tau_P(x)-c,
\qquad
\bar\tau_P=\int_0^1\tau_P(x)\,d\lambda(x).
\]
Preservation of constants and cell integrals gives
\[
\Pi_Mf^{\mathrm{raw}}_{P,c}=f^{\mathrm{cell}}_{P,c},
\qquad
\operatorname{coef}_M(f^{\mathrm{raw}}_{P,c})
=
\operatorname{coef}_M(f^{\mathrm{cell}}_{P,c}).
\]
The histogram coordinate formula reconstructs the projection:
\[
(\Pi_Mf^{\mathrm{raw}}_{P,c})(x)
=
\left\langle
\operatorname{coef}_M(f^{\mathrm{raw}}_{P,c}),F_M(x)
\right\rangle.
\]
Integrating this identity against \(f^{\mathrm{raw}}_{P,c}\) gives
\[
\int_0^1
f^{\mathrm{raw}}_{P,c}(x)(\Pi_Mf^{\mathrm{raw}}_{P,c})(x)\,d\lambda(x)
=
\|\operatorname{coef}_M(f^{\mathrm{raw}}_{P,c})\|_2^2.
\]
Disjointness of the cells and their measure \(1/M\) likewise give
\[
\int_0^1(\Pi_Mf^{\mathrm{raw}}_{P,c})(x)^2\,d\lambda(x)
=
\sum_{j=1}^M
\bigl(\operatorname{coef}_M(f^{\mathrm{raw}}_{P,c})\bigr)_j^2
=
\|\operatorname{coef}_M(f^{\mathrm{raw}}_{P,c})\|_2^2.
\]
Expanding the squared projection residual consequently proves the energy identity
\[
\int_0^1f^{\mathrm{raw}}_{P,c}(x)^2\,d\lambda(x)
=
\|\operatorname{coef}_M(f^{\mathrm{cell}}_{P,c})\|_2^2
+
\int_0^1
\bigl(f^{\mathrm{raw}}_{P,c}(x)-(\Pi_Mf^{\mathrm{raw}}_{P,c})(x)\bigr)^2
\,d\lambda(x).
\]

The Hölder estimate applied using
\cref{obj:ass:effect-smoothness} gives
\[
\left|f^{\mathrm{raw}}_{P,c}(x)-(\Pi_Mf^{\mathrm{raw}}_{P,c})(x)\right|
=
|\tau_P(x)-(\Pi_M\tau_P)(x)|
\le20M^{-\gamma}.
\]
Since \(\lambda\) is a probability measure,
\[
\int_0^1
\bigl(f^{\mathrm{raw}}_{P,c}-\Pi_Mf^{\mathrm{raw}}_{P,c}\bigr)^2
\,d\lambda
\le(20M^{-\gamma})^2.
\]
On the other hand, expansion around the mean gives
\[
\begin{aligned}
\int_0^1(\tau_P(x)-c)^2\,d\lambda(x)
&=\int_0^1\tau_P(x)^2\,d\lambda(x)-2c\bar\tau_P+c^2\\
&=d(P)^2+(\bar\tau_P-c)^2
\ge d(P)^2.
\end{aligned}
\]
Combining these displays yields
\[
d(P)^2
\le
\|\operatorname{coef}_M(f^{\mathrm{cell}}_{P,c})\|_2^2
+(20M^{-\gamma})^2
\le
\left(
\|\operatorname{coef}_M(f^{\mathrm{cell}}_{P,c})\|_2
+20M^{-\gamma}
\right)^2.
\]
Both sides inside the last square are nonnegative, so
\[
d(P)-20M^{-\gamma}
\le
\|\operatorname{coef}_M(f^{\mathrm{cell}}_{P,c})\|_2.
\]
% lean: CausalSmith.Stat.FinitepHomogeneityDensegamma.histogram_projection_energy
% lean: CausalSmith.Stat.FinitepHomogeneityDensegamma.effect_distance_sq_le_const_energy
% lean: CausalSmith.Stat.FinitepHomogeneityDensegamma.effect_histogram_coefficient_distance

\item
The difference between the weighted effect vector and its coarse-effect approximation is
\[
\theta^{\mathrm{wt}}_P(c)-\theta^{\mathrm{cell}}_P(c)
=
\operatorname{coef}_M\bigl(w_K(\tau_P-\Pi_M\tau_P)\bigr).
\]
The weight and Hölder bounds imply
\[
|w_K(x)(\tau_P(x)-(\Pi_M\tau_P)(x))|
\le
\frac9{16}\,20M^{-\gamma}
=
\frac{45}{4}M^{-\gamma}.
\]
Thus
\[
\|\theta^{\mathrm{wt}}_P(c)-\theta^{\mathrm{cell}}_P(c)\|_2
\le\frac{45}{4}M^{-\gamma}.
\]
The triangle inequality, together with the diagonal and distance bounds, gives
\[
\begin{aligned}
\frac3{16}\bigl(d(P)-20M^{-\gamma}\bigr)
&\le
\frac3{16}\|\operatorname{coef}_M(f^{\mathrm{cell}}_{P,c})\|_2\\
&\le\|\theta^{\mathrm{cell}}_P(c)\|_2\\
&\le
\|\theta^{\mathrm{wt}}_P(c)\|_2
+\frac{45}{4}M^{-\gamma}.
\end{aligned}
\]
Since
\[
\frac3{16}\,20+\frac{45}{4}=15,
\]
we obtain the weighted separation estimate
\[
\frac3{16}d(P)-15M^{-\gamma}
\le\|\theta^{\mathrm{wt}}_P(c)\|_2.
\]
% lean: CausalSmith.Stat.FinitepHomogeneityDensegamma.weighted_effect_coefficient_separation

\item
To transfer this estimate to the score mean, write
\[
\theta^{\mathrm{wt}}_P(c)
=
\theta_P(c)-\bigl(\theta_P(c)-\theta^{\mathrm{wt}}_P(c)\bigr).
\]
The triangle inequality and the established mean-error bound give
\[
\|\theta^{\mathrm{wt}}_P(c)\|_2
\le
\|\theta_P(c)\|_2
+\|\theta_P(c)-\theta^{\mathrm{wt}}_P(c)\|_2
\le
\|\theta_P(c)\|_2+B.
\]
Consequently,
\[
\|\theta_P(c)\|_2
\ge\frac3{16}d(P)-15M^{-\gamma}-B
\ge\frac3{16}d(P)-16M^{-\gamma}-B,
\]
where the last inequality uses \(M^{-\gamma}\ge0\).
% lean: CausalSmith.Stat.FinitepHomogeneityDensegamma.original_score_mean_ledger

\item
Finally, let \(P\in H_0(v)\), as specified in
\cref{obj:def:null}, and let \(c_0\in\mathbb R\) satisfy
\(\tau_P(x)=c_0\) for every \(x\in[0,1]\). Then
\[
\theta^{\mathrm{wt}}_P(c_0)
=
\int_0^1w_K(x)(\tau_P(x)-c_0)F_M(x)\,d\lambda(x)
=0.
\]
The mean-error bound was established for every real candidate constant. Applying it at \(c_0\) therefore gives
\[
\|\theta_P(c_0)\|_2
=
\|\theta_P(c_0)-\theta^{\mathrm{wt}}_P(c_0)\|_2
\le B.
\]
All calculations apply to either block, which completes the proof.
% lean: CausalSmith.Stat.FinitepHomogeneityDensegamma.ledgerTheta_null_norm_le
\end{enumerate}
\end{proof}

The mean bounds supply the population distinction used by the test. Under a constant effect, evaluating at that constant gives a score mean of norm at most \(B\). Under heterogeneity, every admissible candidate constant has a score mean bounded below by a multiple of \(d(P)\), after accounting for coarse approximation and original-mean bias. Uniformity in the candidate constant is what makes these bounds compatible with profiling.

\subsection{Singleton, pair, and cross-level covariance bounds}

The singleton projection in \cref{obj:def:mean-handle} combines the centered single-record contribution with twice the pair-kernel singleton projection. Its canonical pair projection retains the component centered with respect to either record. These two components have separate public bounds, \(L_1\) and \(L_2\), in \cref{obj:def:explicit-score-ledger}.

For the multiresolution construction, both bounds aggregate the main correction and every retained increment. In particular, \(L_1\) contains the sum of the increment-specific smoothness and clipping contributions, while \(L_2\) contains the sum of their rank-weighted second-moment bounds. The squared sums in
\[
\Lambda_{n,v}
=\frac{L_1^2}{s}
+\frac{2L_2^2}{s(s-1)}
\]
retain the cross-level contributions within each component. The covariance statement below controls both affine coefficients and their covariance in arbitrary, potentially different, directions.

\begin{lemmav}{lem:original-score-covariance-ledger}[Original score covariance bounds]
Let \(v=(p,\alpha,\beta,\gamma)\) be an admissible tuple of moment and smoothness exponents, and let \(P\in\mathcal M_v\), as specified in \cref{obj:def:model}. Let \(n\ge 2\) be an integer. Write the affine block score as
\[
Z_b(c)=Z_{b,0}-cZ_{b,A}.
\]
Under the independent-record sample law \(P^{\otimes n}\), for every evaluation block \(b\) and every \(u\in\mathbb R^M\), both \(u^T Z_{b,0}\) and \(u^T Z_{b,A}\) are square integrable and satisfy
\[
\operatorname{Var}_{P^{\otimes n}}(u^T Z_{b,0})
\le \Lambda_{n,v}\|u\|_2^2,
\qquad
\operatorname{Var}_{P^{\otimes n}}(u^T Z_{b,A})
\le \Lambda_{n,v}\|u\|_2^2.
\]
Moreover, for every evaluation block \(b\) and every \(u,z\in\mathbb R^M\),
\[
\left|
\operatorname{Cov}_{P^{\otimes n}}
\bigl(u^T Z_{b,0},z^T Z_{b,A}\bigr)
\right|
\le \Lambda_{n,v}\|u\|_2\|z\|_2.
\]
\end{lemmav}

\begin{proof}[Proof of \cref{obj:lem:original-score-covariance-ledger}]
\begin{enumerate}
\item Suppose first that \(n<4\). Since \(n\ge2\), the zero-score convention in \cref{obj:def:explicit-score-ledger,obj:def:calibration-handle} gives
\[
Z_{b,0}=Z_{b,A}=0,
\qquad
\Lambda_{n,v}=0.
\]
Both scalar coefficients are square integrable, and their variances and cross-covariance vanish. This proves all the asserted bounds in this case.
% lean: CausalSmith.Stat.FinitepHomogeneityDensegamma.covarianceLedger_small_sample

\item Henceforth \(n\ge4\). Use the ranks, cutoffs, and ledgers of \cref{obj:def:explicit-score-ledger}, and the blocks of \cref{obj:def:blocks}. In particular,
\[
s=s_n=\lfloor n/2\rfloor\ge2,
\qquad
M>0,\quad K>0,\quad M\mid K,
\qquad
T_0\ge1,
\]
and, in the multiresolution branch,
\[
R_j=2^jM,\qquad M\mid R_{j-1},\qquad T_j\ge1
\quad(1\le j\le L).
\]
The divisibility follows because the selected ranks are powers of two and \(K\ge M\).

Fix \(u\in\mathbb R^M\). Write
\[
\mathcal O=[0,1]\times\{0,1\}\times\mathbb R,
\qquad
O_i=(X_i,A_i,Y_i)\in\mathcal O,
\]
and let \(O,O'\) be independent records with law \(P\). For \(x\in[0,1]\), define the scalar histogram feature by
\[
f_u(x)=u^TF_M(x).
\]
The disjoint histogram cells in \cref{obj:def:projection}, together with the uniform covariate law, give
\[
\int_0^1 f_u(x)^2\,dx
=
\sum_{j=1}^{M}u_j^2
=
\|u\|_2^2,
\qquad
\mathbb E[f_u(X)^2]=\|u\|_2^2.
\]

Use a coefficient selector \(r\in\{0,1\}\), with
\[
Z_{b,r}^{\mathrm{coef}}
=
\begin{cases}
Z_{b,0},&r=0,\\
Z_{b,A},&r=1.
\end{cases}
\]
For every real cutoff \(T\ge1\), define
\[
\mathfrak m_{r,T}(o)
=
\begin{cases}
\ell_T(Y(o)),&r=0,\\
A(o),&r=1,
\end{cases}
\qquad
V_T=32T^{2-p},
\qquad o\in\mathcal O.
\]
Thus \(V_{T_0}=V_0\) and \(V_{T_j}=V_j\). For a pair kernel \(G\) from the score construction, define
\[
\mathscr E_{r,u}[G,T](o,o')
=
f_u(X(o))A(o)G(X(o),X(o'))\mathfrak m_{r,T}(o'),
\]
\[
\mathscr P_{r,u}[G,T](o,o')
=
\frac{
\mathscr E_{r,u}[G,T](o,o')
+
\mathscr E_{r,u}[G,T](o',o)
}{2},
\qquad o,o'\in\mathcal O.
\]
The selected scalar single-record and pair summands are
\[
h_{r,u}(o)=f_u(X(o))A(o)\mathfrak m_{r,T_0}(o),
\]
\[
g_{r,u}(o,o')
=
\begin{cases}
-\mathscr P_{r,u}[\Pi_K,T_0](o,o'),
&\text{single-correction branch},\\[1mm]
-\mathscr P_{r,u}[\Pi_M,T_0](o,o')
-\displaystyle\sum_{j=1}^{L}\mathscr P_{r,u}[D_j,T_j](o,o'),
&\text{multiresolution branch}.
\end{cases}
\]
Here \(A(o)^2=A(o)\), so the \(r=1\) single-record summand agrees with the treatment coefficient. Package these summands into
\[
k_{r,u}(o,o')
=
\frac{h_{r,u}(o)+h_{r,u}(o')}{2}
+
g_{r,u}(o,o').
\]
All these functions are measurable and bounded for the fixed finite tuning, hence square integrable.

For either block, symmetrization preserves the ordered-pair sum, and each index occurs \(s-1\) times in each position. Consequently, the coefficient definitions in \cref{obj:def:calibration-handle} and the score formulas in \cref{obj:def:score-family} give
\[
u^TZ_{b,r}^{\mathrm{coef}}
=
\frac1s\sum_{i\in\mathcal I_b}h_{r,u}(O_i)
+
\frac1{s(s-1)}
\sum_{\substack{i,j\in\mathcal I_b\\i\ne j}}
g_{r,u}(O_i,O_j)
=
\frac1{s(s-1)}
\sum_{\substack{i,j\in\mathcal I_b\\i\ne j}}
k_{r,u}(O_i,O_j).
\]
Empty increment sums have value zero throughout.
% lean: CausalSmith.Stat.FinitepHomogeneityDensegamma.ledgerScalarKernel_average

\item We first establish the variance identity for this pair-average representation. For any measurable symmetric square-integrable kernel
\(\mathscr H:\mathcal O^2\to\mathbb R\), define
\[
\overline{\mathscr H}=\mathbb E[\mathscr H(O,O')],
\qquad
\mathscr H^{(1)}(o)
=
\mathbb E[\mathscr H(o,O')]-\overline{\mathscr H},
\]
\[
\mathscr H^{(2)}(o,o')
=
\mathscr H(o,o')
-\overline{\mathscr H}
-\mathscr H^{(1)}(o)
-\mathscr H^{(1)}(o'),
\]
and
\[
\mathscr A_b(\mathscr H)
=
\frac1{s(s-1)}
\sum_{\substack{i,j\in\mathcal I_b\\i\ne j}}
\mathscr H(O_i,O_j).
\]
Conditional integration gives
\[
\mathbb E[\mathscr H^{(1)}(O)]=0,
\qquad
\mathbb E[\mathscr H^{(2)}(O,O')\mid O]
=
\mathbb E[\mathscr H^{(2)}(O,O')\mid O']
=
0.
\]
Expanding the square of
\(\mathscr H=\overline{\mathscr H}+\mathscr H^{(1)}+\mathscr H^{(1)}+\mathscr H^{(2)}\)
and integrating therefore yields
\[
\mathbb E[\mathscr H(O,O')^2]
=
\overline{\mathscr H}^{\,2}
+
2\mathbb E[\mathscr H^{(1)}(O)^2]
+
\mathbb E[\mathscr H^{(2)}(O,O')^2].
\]
Indeed, the two singleton terms have zero product expectation by independence, and their products with the canonical term have zero expectation by conditional integration.

Counting the occurrences of each singleton term gives
\[
\mathscr A_b(\mathscr H)-\overline{\mathscr H}
=
\frac2s\sum_{i\in\mathcal I_b}\mathscr H^{(1)}(O_i)
+
\mathscr A_b(\mathscr H^{(2)}).
\]
The singleton sum has variance
\[
\operatorname{Var}\!\left(
\frac2s\sum_{i\in\mathcal I_b}\mathscr H^{(1)}(O_i)
\right)
=
\frac4s\operatorname{Var}(\mathscr H^{(1)}(O)).
\]
For distinct indices within each pair, conditional centering and independence imply
\[
\operatorname{Cov}\!\left(
\mathscr H^{(2)}(O_i,O_j),
\mathscr H^{(2)}(O_{i'},O_{j'})
\right)
=
\begin{cases}
\mathbb E[\mathscr H^{(2)}(O,O')^2],
&(i',j')=(i,j)\text{ or }(j,i),\\
0,&\text{otherwise}.
\end{cases}
\]
The zero case includes pairs sharing exactly one index: integrate first in a coordinate that occurs in only one pair. There are \(s(s-1)\) ordered pairs and two contributing ordered pairs for each of them. Thus
\[
\operatorname{Var}(\mathscr A_b(\mathscr H^{(2)}))
=
\frac{2}{s(s-1)}
\mathbb E[\mathscr H^{(2)}(O,O')^2].
\]
The same conditional-centering argument gives
\[
\operatorname{Cov}\!\left(
\frac2s\sum_{i\in\mathcal I_b}\mathscr H^{(1)}(O_i),
\mathscr A_b(\mathscr H^{(2)})
\right)=0.
\]
Each evaluation block has marginal law \(P^{\otimes s}\) under \(P^{\otimes n}\), so
\[
\operatorname{Var}_{P^{\otimes n}}(\mathscr A_b(\mathscr H))
=
\frac4s\operatorname{Var}_P(\mathscr H^{(1)})
+
\frac{2}{s(s-1)}
\mathbb E[\mathscr H^{(2)}(O,O')^2].
\]
% lean: CausalSmith.Stat.FinitepHomogeneityDensegamma.evalBlockPairAverage_variance

\item We next bound the singleton channel. Admissibility gives \(1<p\le2\). For \(T\ge1\), clipping satisfies the pointwise bounds
\[
|y-\ell_T(y)|\le |y|^pT^{1-p},
\qquad
\ell_T(y)^2\le |y|^pT^{2-p}.
\]
Integrating under the conditional moment restriction in \cref{obj:ass:raw-moment} gives, for each treatment arm and almost every covariate value,
\[
\left|
\int\{y-\ell_T(y)\}\,Q_{a,P}(dy\mid x)
\right|
\le10T^{1-p},
\qquad
\int\ell_T(y)^2\,Q_{a,P}(dy\mid x)
\le10T^{2-p}.
\]
Define the conditional mark mean, for \(x\in[0,1]\), by
\[
\overline{\mathfrak m}_{r,T}(x)
=
e_P(x)\int\mathfrak m_{r,T}(x,1,y)\,Q_{1,P}(dy\mid x)
+
\{1-e_P(x)\}
\int\mathfrak m_{r,T}(x,0,y)\,Q_{0,P}(dy\mid x).
\]
The preceding bounds and \(0\le e_P\le1\) imply
\[
\mathbb E[\mathfrak m_{0,T}(O)^2\mid X=x]
\le10T^{2-p}\le V_T.
\]
For the treatment mark,
\[
\mathbb E[\mathfrak m_{1,T}(O)^2\mid X=x]
=e_P(x)\le1\le V_T,
\qquad
\overline{\mathfrak m}_{1,T}(x)=e_P(x).
\]
In particular, for both selectors,
\[
\mathbb E[\mathfrak m_{r,T}(O)^2\mid X=x]\le V_T,
\qquad
|\overline{\mathfrak m}_{r,T}(x)|\le\sqrt{V_T}
\quad\text{almost everywhere},
\]
where the second inequality is conditional Cauchy--Schwarz.

Since \(|A|\le1\), the single-record summand obeys
\[
h_{r,u}(O)^2
\le
2f_u(X)^2\mathfrak m_{r,T_0}(O)^2.
\]
Conditioning on \(X\) and using the feature isometry yields
\[
\mathbb E[h_{r,u}(O)^2]
\le2V_0\|u\|_2^2
\le B_h^2\|u\|_2^2,
\qquad
B_h=2\sqrt{V_0}.
\]
% lean: CausalSmith.Stat.FinitepHomogeneityDensegamma.multiresScalarSingle_energy_le

\item For any selected projection or difference kernel \(G\), and any integrable function \(\varphi:[0,1]\to\mathbb R\), use the operator notation
\[
(G\varphi)(x)=\int_0^1G(x,x')\varphi(x')\,dx'.
\]
Nested histogram cells imply
\[
G(x,x')f_u(x')=G(x,x')f_u(x)
\]
for \(G=\Pi_K,\Pi_M\), and for every \(G=D_j\). Integrating the two symmetrized legs separately therefore gives
\[
\mathbb E[\mathscr P_{r,u}[G,T](o,O')]
=
\frac12f_u(X(o))
\left[
A(o)(G\overline{\mathfrak m}_{r,T})(X(o))
+
\mathfrak m_{r,T}(o)(Ge_P)(X(o))
\right].
\]
For a positive histogram rank \(J_{\mathrm{aux}}\) divisible by \(M\), cell averaging gives
\[
|\Pi_{J_{\mathrm{aux}}}\overline{\mathfrak m}_{r,T}|
\le\sqrt{V_T},
\qquad
|\Pi_{J_{\mathrm{aux}}}e_P|\le1.
\]
Consequently,
\[
\left|
\mathbb E[\mathscr P_{r,u}[\Pi_{J_{\mathrm{aux}}},T](o,O')]
\right|
\le
\frac12|f_u(X(o))|
\left(\sqrt{V_T}+|\mathfrak m_{r,T}(o)|\right).
\]
Using
\[
\frac14\left(\sqrt{V_T}+|\mathfrak m_{r,T}(o)|\right)^2
\le
\frac12\left(V_T+\mathfrak m_{r,T}(o)^2\right)
\]
and then conditioning on \(X\), we obtain
\[
\mathbb E\!\left[
\left\{
\mathbb E[\mathscr P_{r,u}[\Pi_{J_{\mathrm{aux}}},T](O,O')\mid O]
\right\}^{\!2}
\right]
\le V_T\|u\|_2^2.
\]

The augmented kernel has singleton projection
\[
k_{r,u}^{(1)}(o)
=
\frac{h_{r,u}(o)-\mathbb E[h_{r,u}(O)]}{2}
+
g_{r,u}^{(1)}(o).
\]
Centering decreases second moment, so the triangle inequality in \(L_2(P)\) gives
\[
2\|k_{r,u}^{(1)}\|_{L_2(P)}
\le
\sqrt{\mathbb E[h_{r,u}(O)^2]}
+
2\sqrt{
\mathbb E\!\left[
\left\{\mathbb E[g_{r,u}(O,O')\mid O]\right\}^{2}
\right]}.
\]
Here
\[
\operatorname{Var}_P(k_{r,u}^{(1)})
=
\|k_{r,u}^{(1)}\|_{L_2(P)}^2
\]
because its mean is zero.

In the single-correction branch, the projection estimate with
\(J_{\mathrm{aux}}=K\) and \(T=T_0\) gives
\[
\mathbb E\!\left[
\left\{\mathbb E[g_{r,u}(O,O')\mid O]\right\}^{2}
\right]
\le B_{g,\mathrm{single}}^2\|u\|_2^2,
\qquad
B_{g,\mathrm{single}}=\sqrt{V_0}.
\]
It follows that
\[
4\operatorname{Var}_P(k_{r,u}^{(1)})
\le
(B_h+2B_{g,\mathrm{single}})^2\|u\|_2^2
=
(4\sqrt{V_0})^2\|u\|_2^2
\le
(16\sqrt{V_0})^2\|u\|_2^2
=
L_1^2\|u\|_2^2.
\]
% lean: CausalSmith.Stat.FinitepHomogeneityDensegamma.ledgerScalarKernel_singleton_variance_le

\item In the multiresolution branch, define
\[
s_0=\min\{\alpha,\beta,\gamma\}\in(0,1],
\qquad
f_{\mathrm{mean}}(x)=m_{0,P}(x)+e_P(x)\tau_P(x),
\qquad x\in[0,1].
\]
The smoothness and boundedness restrictions in
\cref{obj:ass:propensity-smoothness,obj:ass:baseline-smoothness,obj:ass:effect-smoothness,obj:ass:overlap,obj:ass:baseline-cap,obj:ass:effect-cap}
give
\[
|f_{\mathrm{mean}}(x)|\le2
\]
and
\[
\begin{aligned}
|f_{\mathrm{mean}}(x)-f_{\mathrm{mean}}(x')|
&\le
|m_{0,P}(x)-m_{0,P}(x')|
+
|e_P(x)|\,|\tau_P(x)-\tau_P(x')|\\
&\quad+
|e_P(x)-e_P(x')|\,|\tau_P(x')|\\
&\le
20|x-x'|^\beta
+
15|x-x'|^\gamma
+
20|x-x'|^\alpha\\
&\le55|x-x'|^{s_0}.
\end{aligned}
\]
For \(x=x'\) both sides vanish; for \(0<|x-x'|\le1\), the last inequality follows because each exponent is at least \(s_0\).

To apply the histogram approximation bound with constant \(20\), define
\[
\widetilde f_{\mathrm{mean}}(x)=\frac{20}{55}f_{\mathrm{mean}}(x).
\]
Then
\[
\|\widetilde f_{\mathrm{mean}}\|_\infty\le\frac{40}{55}\le20,
\qquad
|\widetilde f_{\mathrm{mean}}(x)-\widetilde f_{\mathrm{mean}}(x')|
\le20|x-x'|^{s_0}.
\]
Averaging over the histogram cell containing \(x\), whose diameter is at most \(J_{\mathrm{aux}}^{-1}\), gives, for every positive integer \(J_{\mathrm{aux}}\),
\[
|\Pi_{J_{\mathrm{aux}}}\widetilde f_{\mathrm{mean}}(x)
-\widetilde f_{\mathrm{mean}}(x)|
\le20J_{\mathrm{aux}}^{-s_0}.
\]
Thus, for \(1\le j\le L\),
\[
\begin{aligned}
|D_jf_{\mathrm{mean}}(x)|
&=
\frac{55}{20}
\left|
\Pi_{R_j}\widetilde f_{\mathrm{mean}}(x)
-
\Pi_{R_{j-1}}\widetilde f_{\mathrm{mean}}(x)
\right|\\
&\le
\frac{55}{20}
\left(20R_j^{-s_0}+20R_{j-1}^{-s_0}\right)
\le110R_{j-1}^{-s_0}.
\end{aligned}
\]

For \(T\ge1\), define on all of \([0,1]\)
\[
\mathfrak t_T(x)
=
f_{\mathrm{mean}}(x)-\overline{\mathfrak m}_{0,T}(x).
\]
The conditional mean versions in \cref{obj:def:model} imply, almost everywhere,
\[
\begin{aligned}
\mathfrak t_T(x)
&=
e_P(x)\int\{y-\ell_T(y)\}\,Q_{1,P}(dy\mid x)\\
&\quad+
\{1-e_P(x)\}
\int\{y-\ell_T(y)\}\,Q_{0,P}(dy\mid x).
\end{aligned}
\]
The clipping-tail estimate therefore gives
\[
\|\mathfrak t_T\|_{L_\infty(\lambda)}
\le10T^{1-p},
\qquad
|D_j\mathfrak t_{T_j}(x)|
\le20T_j^{1-p}.
\]
For the propensity, the same cell-averaging argument gives
\[
|\Pi_{J_{\mathrm{aux}}}e_P(x)-e_P(x)|
\le20J_{\mathrm{aux}}^{-\alpha},
\]
hence
\[
|D_je_P(x)|
\le20R_j^{-\alpha}+20R_{j-1}^{-\alpha}
\le40R_{j-1}^{-\alpha}=a_j.
\]
Combining the smoothness and tail bounds yields
\[
|D_j\overline{\mathfrak m}_{0,T_j}(x)|
\le110R_{j-1}^{-s_0}+20T_j^{1-p}=d_j.
\]

Substitution into the conditional pair formula gives the outcome-coefficient envelope
\[
\left|
\mathbb E[\mathscr P_{0,u}[D_j,T_j](o,O')]
\right|
\le
\frac12|f_u(X(o))|
\left(d_j+a_j|\mathfrak m_{0,T_j}(o)|\right).
\]
Since
\[
\frac14(d_j+a_j|\mathfrak m_{0,T_j}|)^2
\le
\frac12(d_j^2+a_j^2\mathfrak m_{0,T_j}^2),
\]
conditioning on \(X\) gives
\[
\mathbb E\!\left[
\left\{
\mathbb E[\mathscr P_{0,u}[D_j,T_j](O,O')\mid O]
\right\}^{\!2}
\right]
\le
\frac12(d_j^2+a_j^2V_j)\|u\|_2^2.
\]
For the treatment coefficient,
\(\overline{\mathfrak m}_{1,T_j}=e_P\) and
\(|\mathfrak m_{1,T_j}|\le1\), so
\[
\left|
\mathbb E[\mathscr P_{1,u}[D_j,T_j](o,O')]
\right|
\le a_j|f_u(X(o))|.
\]
Applying the same square-envelope inequality with constant term \(2a_j\) and zero mark term yields
\[
\mathbb E\!\left[
\left\{
\mathbb E[\mathscr P_{1,u}[D_j,T_j](O,O')\mid O]
\right\}^{\!2}
\right]
\le2a_j^2\|u\|_2^2.
\]
All \(a_j,d_j\) are nonnegative and \(V_j\ge1\). The two coefficient estimates fit the common bound because
\[
\frac12(d_j^2+a_j^2V_j)
\le\{2(d_j+a_j\sqrt{V_j})\}^2,
\qquad
2a_j^2
\le\{2(d_j+a_j\sqrt{V_j})\}^2.
\]
Define the nonnegative increment sum and conditional-pair budget by
\[
\mathcal S_{\mathrm{inc}}
=
\sum_{j=1}^{L}(d_j+a_j\sqrt{V_j}),
\qquad
B_{g,\mathrm{multi}}
=
\sqrt{V_0}+2\mathcal S_{\mathrm{inc}}.
\]
The initial projection estimate at rank \(M\), followed by the triangle inequality for the finite increment sum in \(L_2(P)\), gives
\[
\begin{aligned}
\left\|
\mathbb E[g_{r,u}(O,O')\mid O]
\right\|_{L_2(P)}
&\le
\sqrt{V_0}\|u\|_2
+
\sum_{j=1}^{L}2(d_j+a_j\sqrt{V_j})\|u\|_2\\
&=
B_{g,\mathrm{multi}}\|u\|_2.
\end{aligned}
\]
This includes \(L=0\), when the sum is zero. The singleton estimate for the augmented kernel now gives
\[
\begin{aligned}
4\operatorname{Var}_P(k_{r,u}^{(1)})
&\le
(B_h+2B_{g,\mathrm{multi}})^2\|u\|_2^2\\
&=
(4\sqrt{V_0}+4\mathcal S_{\mathrm{inc}})^2\|u\|_2^2\\
&\le
(16\sqrt{V_0}+8\mathcal S_{\mathrm{inc}})^2\|u\|_2^2
=
L_1^2\|u\|_2^2.
\end{aligned}
\]
% lean: CausalSmith.Stat.FinitepHomogeneityDensegamma.ledgerScalarKernel_singleton_variance_le

\item We next bound the canonical channel, again treating the two branches separately. The definitions of the projections give
\[
\overline{k_{r,u}}
=
\mathbb E[h_{r,u}(O)]+\overline{g_{r,u}},
\]
and cancellation of the single-record terms gives
\[
k_{r,u}^{(2)}(o,o')=g_{r,u}^{(2)}(o,o')
\quad\text{almost everywhere}.
\]
The pair-energy identity therefore implies
\[
\begin{aligned}
\mathbb E[k_{r,u}^{(2)}(O,O')^2]
&=
\mathbb E[g_{r,u}^{(2)}(O,O')^2]\\
&\le
\overline{g_{r,u}}^{\,2}
+
2\mathbb E[g_{r,u}^{(1)}(O)^2]
+
\mathbb E[g_{r,u}^{(2)}(O,O')^2]\\
&=
\mathbb E[g_{r,u}(O,O')^2].
\end{aligned}
\]

To bound this raw pair energy, the histogram kernels have exact row identities
\[
\int_0^1\Pi_{J_{\mathrm{aux}}}(x,x')^2\,dx'
=
J_{\mathrm{aux}}
\quad(J_{\mathrm{aux}}\in\mathbb N,\ J_{\mathrm{aux}}>0),
\]
\[
\int_0^1
\Pi_{R_j}(x,x')\Pi_{R_{j-1}}(x,x')\,dx'
=
R_{j-1},
\]
and hence
\[
\int_0^1D_j(x,x')^2\,dx'
=
R_j+R_{j-1}-2R_{j-1}
=
R_{j-1}
\quad(1\le j\le L).
\]
The first identity follows because a projection row equals its rank on one cell of reciprocal length. For the cross-product identity, the finer cell lies inside the coarser cell, so integrating their product gives the coarser rank.

Conditioning on the covariate of \(O'\), the mark bound gives, for every selected \(G,T\),
\[
\mathbb E[G(x,X(O'))^2\mathfrak m_{r,T}(O')^2]
\le
V_T\int_0^1G(x,x')^2\,dx'.
\]
Thus, integrating first in \(O'\) and using \(A(O)^2\le1\),
\[
\mathbb E[\mathscr E_{r,u}[\Pi_{J_{\mathrm{aux}}},T](O,O')^2]
\le
J_{\mathrm{aux}}V_T\,\mathbb E[f_u(X)^2]
=
J_{\mathrm{aux}}V_T\|u\|_2^2,
\]
\[
\mathbb E[\mathscr E_{r,u}[D_j,T_j](O,O')^2]
\le
R_{j-1}V_j\|u\|_2^2.
\]
Symmetrization satisfies
\[
\mathscr P_{r,u}[G,T](o,o')^2
\le
\frac{
\mathscr E_{r,u}[G,T](o,o')^2
+
\mathscr E_{r,u}[G,T](o',o)^2
}{2}.
\]
The two legs have equal integrated energy under \(P\otimes P\). Consequently,
\[
\mathbb E[\mathscr P_{r,u}[\Pi_{J_{\mathrm{aux}}},T](O,O')^2]
\le
J_{\mathrm{aux}}V_T\|u\|_2^2,
\]
\[
\mathbb E[\mathscr P_{r,u}[D_j,T_j](O,O')^2]
\le
R_{j-1}V_j\|u\|_2^2
\le
R_jV_j\|u\|_2^2.
\]

In the single-correction branch, this gives
\[
\begin{aligned}
\mathbb E[k_{r,u}^{(2)}(O,O')^2]
&\le
\mathbb E[g_{r,u}(O,O')^2]\\
&\le
KV_0\|u\|_2^2\\
&\le
(4\sqrt{KV_0})^2\|u\|_2^2
=
L_2^2\|u\|_2^2.
\end{aligned}
\]
In the multiresolution branch, the triangle inequality in \(L_2(P\otimes P)\) gives
\[
\|g_{r,u}\|_{L_2(P\otimes P)}
\le
\left(
\sqrt{MV_0}
+
\sum_{j=1}^{L}\sqrt{R_jV_j}
\right)\|u\|_2.
\]
Every summand is nonnegative, including the empty-sum case. Therefore
\[
\begin{aligned}
\mathbb E[k_{r,u}^{(2)}(O,O')^2]
&\le
\left(
\sqrt{MV_0}+\sum_{j=1}^{L}\sqrt{R_jV_j}
\right)^2\|u\|_2^2\\
&\le
\left[
4\left(
\sqrt{MV_0}+\sum_{j=1}^{L}\sqrt{R_jV_j}
\right)
\right]^2\|u\|_2^2\\
&=
L_2^2\|u\|_2^2.
\end{aligned}
\]
% lean: CausalSmith.Stat.FinitepHomogeneityDensegamma.ledgerScalarKernel_canonical_energy_le

\item Apply the pair-average variance identity to \(k_{r,u}\). Since \(s\ge2\), both sample factors are positive, and the two projection estimates imply
\[
\begin{aligned}
\operatorname{Var}_{P^{\otimes n}}(u^TZ_{b,r}^{\mathrm{coef}})
&=
\frac4s\operatorname{Var}_P(k_{r,u}^{(1)})
+
\frac{2}{s(s-1)}
\mathbb E[k_{r,u}^{(2)}(O,O')^2]\\
&\le
\frac{L_1^2}{s}\|u\|_2^2
+
\frac{2L_2^2}{s(s-1)}\|u\|_2^2\\
&=
\Lambda_{n,v}\|u\|_2^2.
\end{aligned}
\]
The finite-sum representation also gives
\[
u^TZ_{b,r}^{\mathrm{coef}}\in L_2(P^{\otimes n}).
\]
Taking \(r=0\) and \(r=1\) proves square integrability and the two variance bounds.
% lean: CausalSmith.Stat.FinitepHomogeneityDensegamma.original_score_covariance_ledger

\item Finally, the covariance ledger is nonnegative:
\[
\Lambda_{n,v}
=
\frac{L_1^2}{s}
+
\frac{2L_2^2}{s(s-1)}
\ge0
\qquad(n\ge4),
\]
and its value is zero for \(n=2,3\). For arbitrary \(u,z\in\mathbb R^M\), Cauchy--Schwarz for the centered square-integrable coefficients gives
\[
\begin{aligned}
\left|
\operatorname{Cov}_{P^{\otimes n}}
(u^TZ_{b,0},z^TZ_{b,A})
\right|
&\le
\sqrt{\operatorname{Var}_{P^{\otimes n}}(u^TZ_{b,0})}\,
\sqrt{\operatorname{Var}_{P^{\otimes n}}(z^TZ_{b,A})}\\
&\le
\sqrt{\Lambda_{n,v}\|u\|_2^2}\,
\sqrt{\Lambda_{n,v}\|z\|_2^2}\\
&=
(\sqrt{\Lambda_{n,v}})^2\|u\|_2\|z\|_2\\
&=
\Lambda_{n,v}\|u\|_2\|z\|_2.
\end{aligned}
\]
The square-root simplification applies also when either direction or the ledger is zero.
% lean: CausalSmith.Stat.FinitepHomogeneityDensegamma.ledger_cross_covariance_of_variance
\end{enumerate}
\end{proof}

The directional form gives a common covariance bound for the outcome coefficient and the constant-effect coefficient. The mixed bound retains their dependence within a block, including the dependence contributed by the different correction levels. Independence between the two evaluation blocks then provides the setting for the cross-block profiling statement.

\subsection{Simultaneous affine profiling and calibration}

The profiling multiplier \(C=128\) and the population mean \(\theta_P(c)\) were introduced with \cref{obj:thm:whole-null-calibration-resolved}. The following bound applies to the entire interval of candidate constants on one measurable event.

\begin{lemmav}{lem:uniform-affine-profile-bound}[Uniform affine profile bound]
Let \(v=(p,\alpha,\beta,\gamma)\) be an admissible public tuple as in \cref{obj:def:model}, and let \(n\ge 2\) be an integer. For every \(P\in\mathcal M_v\), consider the explicit affine score construction, with block scores \(Z_1(c)\) and \(Z_2(c)\), coarse histogram rank \(J_{\mathrm c}\), and public covariance ledger \(\Lambda_{n,v}\). Define the score inner product and its population target by
\[
W(c)=\langle Z_1(c),Z_2(c)\rangle,
\qquad
\theta_P(c)=\mathbb E_{P^{\otimes n}}[Z_1(c)].
\]
The event on which, simultaneously for every \(c\in[-1/2,1/2]\),
\[
\bigl|W(c)-\|\theta_P(c)\|_2^2\bigr|
\le
128\left\{
\sqrt{\Lambda_{n,v}}\,\|\theta_P(c)\|_2
+\sqrt{J_{\mathrm c}}\,\Lambda_{n,v}
\right\}
\]
is measurable and has probability at least \(0.9925\) under \(P^{\otimes n}\). Moreover, if \(\Lambda_{n,v}=0\), then \(P^{\otimes n}\)-almost surely, simultaneously for every \(c\in[-1/2,1/2]\),
\[
W(c)=\|\theta_P(c)\|_2^2.
\]
\end{lemmav}

\begin{proof}[Proof of \cref{obj:lem:uniform-affine-profile-bound}]
Fix an admissible \(v\), an integer \(n\ge2\), and \(P\in\mathcal M_v\). Write
\[
\mu=P^{\otimes n},
\qquad
M=J_{\mathrm c}.
\]
All expectations below are under \(\mu\).

\begin{enumerate}
\item
If \(n<4\), then \(n\in\{2,3\}\). The construction in \cref{obj:def:explicit-score-ledger} gives, for every sample and every \(c\in\mathbb R\),
\[
Z_1(c)=Z_2(c)=0,
\qquad
\theta_P(c)=0,
\qquad
W(c)=0,
\qquad
\Lambda_{n,v}=0.
\]
Consequently, the asserted profile event is the entire sample space, with probability one, and the equality assertion holds for every sample. Henceforth suppose \(n\ge4\).
% lean: CausalSmith.Stat.FinitepHomogeneityDensegamma.profileBound_small_sample

\item
Define the profile discrepancy and its simultaneous event by
\[
\mathsf F_{\mathrm{prof}}(c)
=
\bigl|W(c)-\|\theta_P(c)\|_2^2\bigr|
-
128\left(
\sqrt{\Lambda_{n,v}}\,\|\theta_P(c)\|_2
+\sqrt M\,\Lambda_{n,v}
\right),
\]
\[
\mathcal G_{\mathrm{prof}}
=
\left\{
\mathsf F_{\mathrm{prof}}(c)\le0
\text{ for every }c\in[-1/2,1/2]
\right\}.
\]
For fixed \(c\), the discrepancy is a measurable function of the sample. For each fixed sample, it is continuous in \(c\), because the scores and their population means are affine. Choose a countable dense set with
\[
\mathcal C_{\mathrm{dense}}\subseteq[-1/2,1/2],
\qquad
\overline{\mathcal C_{\mathrm{dense}}}=[-1/2,1/2].
\]
Continuity gives
\[
\mathcal G_{\mathrm{prof}}
=
\bigcap_{c\in\mathcal C_{\mathrm{dense}}}
\{\mathsf F_{\mathrm{prof}}(c)\le0\}.
\]
Indeed, for each sample, the set where the discrepancy is nonpositive is closed in the profiling interval. Thus \(\mathcal G_{\mathrm{prof}}\) is measurable.

We also record the positivity needed for the tail estimates and the zero-ledger assertion. The public formulas in \cref{obj:def:explicit-score-ledger} give
\[
s=s_n=\lfloor n/2\rfloor\ge2,
\qquad
M\ge1,
\qquad
T_0\ge1,
\qquad
V_0=32T_0^{2-p}>0.
\]
In the single-resolution branch, \(L_1=16\sqrt{V_0}\). In the multiresolution branch, the positive ranks and cutoffs give
\[
a_j\ge0,\qquad d_j\ge0,\qquad V_j>0
\quad(1\le j\le L),
\]
and hence
\[
\sum_{j=1}^{L}\bigl(d_j+a_j\sqrt{V_j}\bigr)\ge0,
\qquad
L_1
=
16\sqrt{V_0}
+8\sum_{j=1}^{L}\bigl(d_j+a_j\sqrt{V_j}\bigr)
\ge16\sqrt{V_0}>0.
\]
These inequalities include the empty-sum case \(L=0\). In either branch,
\[
\Lambda_{n,v}
=
\frac{L_1^2}{s}+\frac{2L_2^2}{s(s-1)}
\ge\frac{L_1^2}{s}>0.
\]
Together with the small-sample branch, this establishes, for admissible \(v\) and \(n\ge2\),
\[
\Lambda_{n,v}=0
\quad\Longleftrightarrow\quad
n<4.
\]
The zero-ledger assertion therefore follows from the pointwise equalities in the first case.
% lean: CausalSmith.Stat.FinitepHomogeneityDensegamma.measurableSet_profileEvent
% lean: CausalSmith.Stat.FinitepHomogeneityDensegamma.ledgerCov_pos_of_large_sample
% lean: CausalSmith.Stat.FinitepHomogeneityDensegamma.ledger_profile_zero_covariance

\item
For \(b\in\{1,2\}\), define the affine coefficients by
\[
Z_{b,0}=Z_b(0),
\qquad
Z_{b,A}=Z_b(0)-Z_b(1),
\qquad
Z_b(c)=Z_{b,0}-cZ_{b,A}
\quad(c\in\mathbb R).
\]
The common block-mean identity supplied by \cref{obj:lem:original-score-mean-ledger} is
\[
\mathbb E_\mu[Z_b(c)]=\theta_P(c)
\quad(b\in\{1,2\},\ c\in\mathbb R).
\]
Define the population coefficients and centered coefficient errors by
\[
\overline Z_0=\theta_P(0),
\qquad
\overline Z_A=\theta_P(0)-\theta_P(1),
\]
\[
\mathsf E_{b,\iota}=Z_{b,\iota}-\overline Z_\iota
\quad(b\in\{1,2\},\ \iota\in\{0,A\}),
\qquad
\mathsf E_b(c)=\mathsf E_{b,0}-c\mathsf E_{b,A}.
\]
Linearity of expectation gives
\[
\mathbb E_\mu[Z_{b,\iota}]=\overline Z_\iota,
\qquad
\mathbb E_\mu[\mathsf E_{b,\iota}]=0,
\]
\[
\theta_P(c)=\overline Z_0-c\overline Z_A,
\qquad
Z_b(c)=\theta_P(c)+\mathsf E_b(c).
\]
Introduce the fixed population subspace
\[
\mathcal S_P=\operatorname{span}\{\overline Z_0,\overline Z_A\}
\subseteq\mathbb R^M.
\]
Its orthogonal projection \(\operatorname{pr}_{\mathcal S_P}:\mathbb R^M\to\mathcal S_P\) is characterized by
\[
\operatorname{pr}_{\mathcal S_P}\mathsf v_{\mathrm{dir}}\in\mathcal S_P,
\qquad
\mathsf v_{\mathrm{dir}}-\operatorname{pr}_{\mathcal S_P}\mathsf v_{\mathrm{dir}}\in\mathcal S_P^\perp
\quad(\mathsf v_{\mathrm{dir}}\in\mathbb R^M).
\]

For \(b\in\{1,2\}\) and \(\iota,\jmath\in\{0,A\}\), define the bad events
\[
\mathcal P_{b,\iota}
=
\left\{
\|\operatorname{pr}_{\mathcal S_P}\mathsf E_{b,\iota}\|_2
>
40\sqrt{\Lambda_{n,v}}
\right\},
\]
\[
\mathcal Q_{\iota,\jmath}
=
\left\{
|\langle\mathsf E_{1,\iota},\mathsf E_{2,\jmath}\rangle|
>
40\sqrt M\,\Lambda_{n,v}
\right\},
\]
\[
\mathcal F_{\mathrm{bad}}
=
\left(
\bigcup_{b\in\{1,2\}}\ \bigcup_{\iota\in\{0,A\}}
\mathcal P_{b,\iota}
\right)
\cup
\left(
\bigcup_{\iota\in\{0,A\}}\ \bigcup_{\jmath\in\{0,A\}}
\mathcal Q_{\iota,\jmath}
\right).
\]
The coefficients are measurable, so continuity of projection, norm, and inner product makes all these events measurable.
% lean: CausalSmith.Stat.FinitepHomogeneityDensegamma.ledgerScore_eq_theta_add_error
% lean: CausalSmith.Stat.FinitepHomogeneityDensegamma.uniform_affine_profile_bound

\item
The square integrability and variance bounds in \cref{obj:lem:original-score-covariance-ledger}, together with the coefficient means just computed, imply
\[
\mathbb E_\mu\!\left[
\langle \mathsf v_{\mathrm{dir}},\mathsf E_{b,\iota}\rangle^2
\right]
=
\operatorname{Var}_\mu\!\left(\langle \mathsf v_{\mathrm{dir}},Z_{b,\iota}\rangle\right)
\le
\Lambda_{n,v}\|\mathsf v_{\mathrm{dir}}\|_2^2
\quad(\mathsf v_{\mathrm{dir}}\in\mathbb R^M).
\]
Because \(\mathcal S_P\) is spanned by two vectors,
\[
\dim\mathcal S_P\le2.
\]
Choose an orthonormal basis
\[
\mathcal B_P
=
\{\mathsf v^{\mathrm{sp}}_{\mathrm{dir},1},\ldots,\mathsf v^{\mathrm{sp}}_{\mathrm{dir},\dim\mathcal S_P}\}
\subseteq\mathcal S_P.
\]
Parseval's identity and the defining property of orthogonal projection give
\[
\begin{aligned}
\mathbb E_\mu
\|\operatorname{pr}_{\mathcal S_P}\mathsf E_{b,\iota}\|_2^2
&=
\sum_{\mathsf v_{\mathrm{dir}}\in\mathcal B_P}
\mathbb E_\mu
\langle \mathsf v_{\mathrm{dir}},\operatorname{pr}_{\mathcal S_P}\mathsf E_{b,\iota}\rangle^2\\
&=
\sum_{\mathsf v_{\mathrm{dir}}\in\mathcal B_P}
\mathbb E_\mu\langle \mathsf v_{\mathrm{dir}},\mathsf E_{b,\iota}\rangle^2\\
&\le
(\dim\mathcal S_P)\Lambda_{n,v}
\le2\Lambda_{n,v}.
\end{aligned}
\]
If \(\mathcal S_P=\{0\}\), the basis is empty and these sums equal zero. Markov's inequality applied to the squared projected norm now yields
\[
\mu(\mathcal P_{b,\iota})
\le
\frac{
\mathbb E_\mu
\|\operatorname{pr}_{\mathcal S_P}\mathsf E_{b,\iota}\|_2^2
}{
1600\Lambda_{n,v}
}
\le\frac1{800}.
\]
There are two blocks and two coefficients, so the finite union bound gives
\[
\mu\left(
\bigcup_{b\in\{1,2\}}\ \bigcup_{\iota\in\{0,A\}}
\mathcal P_{b,\iota}
\right)
\le
\sum_{b\in\{1,2\}}\ \sum_{\iota\in\{0,A\}}
\mu(\mathcal P_{b,\iota})
\le\frac4{800}.
\]
% lean: CausalSmith.Stat.FinitepHomogeneityDensegamma.ledgerCoeffError_inner_second_moment_le
% lean: CausalSmith.Stat.FinitepHomogeneityDensegamma.ledgerCoeffError_projection_tail

\item
The coefficients in the two blocks depend on disjoint sets of independent records, as specified in \cref{obj:def:blocks,obj:def:explicit-score-ledger}. Subtracting deterministic means preserves their independence. Define their probability laws on \(\mathbb R^M\) by
\[
\nu_{b,\iota}
=
\mu\circ\mathsf E_{b,\iota}^{-1}.
\]
For every \(\iota,\jmath\in\{0,A\}\), their joint law is therefore
\[
\mu\circ
(\mathsf E_{1,\iota},\mathsf E_{2,\jmath})^{-1}
=
\nu_{1,\iota}\otimes\nu_{2,\jmath}.
\]
Choose an orthonormal basis of the ambient space,
\[
\mathcal B_{\mathrm{amb}}
=
\{\mathsf v^{\mathrm{amb}}_{\mathrm{dir},1},\ldots,\mathsf v^{\mathrm{amb}}_{\mathrm{dir},M}\}
\subseteq\mathbb R^M.
\]
The directional second-moment bound from the preceding step gives
\[
\mathbb E_\mu\|\mathsf E_{1,\iota}\|_2^2
=
\sum_{\mathsf v_{\mathrm{dir}}\in\mathcal B_{\mathrm{amb}}}
\mathbb E_\mu\langle \mathsf v_{\mathrm{dir}},\mathsf E_{1,\iota}\rangle^2
\le M\Lambda_{n,v}.
\]
Square integrability and
\[
|\langle \mathsf v_{\mathrm{dir}},z\rangle|^2\le\|\mathsf v_{\mathrm{dir}}\|_2^2\|z\|_2^2
\quad(\mathsf v_{\mathrm{dir}},z\in\mathbb R^M)
\]
justify iterated integration under the product law. Applying the directional second-moment bound to the second coefficient, for each fixed \(\mathsf v_{\mathrm{dir}}\), gives
\[
\begin{aligned}
\mathbb E_\mu
\langle\mathsf E_{1,\iota},\mathsf E_{2,\jmath}\rangle^2
&=
\int_{\mathbb R^M}\int_{\mathbb R^M}
\langle \mathsf v_{\mathrm{dir}},z\rangle^2\,
d\nu_{2,\jmath}(z)\,d\nu_{1,\iota}(\mathsf v_{\mathrm{dir}})\\
&\le
\int_{\mathbb R^M}
\Lambda_{n,v}\|\mathsf v_{\mathrm{dir}}\|_2^2\,d\nu_{1,\iota}(\mathsf v_{\mathrm{dir}})\\
&=
\Lambda_{n,v}\,
\mathbb E_\mu\|\mathsf E_{1,\iota}\|_2^2\\
&\le M\Lambda_{n,v}^2.
\end{aligned}
\]
Since \(M\ge1\) and \(\Lambda_{n,v}>0\), Markov's inequality gives
\[
\mu(\mathcal Q_{\iota,\jmath})
\le
\frac{
\mathbb E_\mu
\langle\mathsf E_{1,\iota},\mathsf E_{2,\jmath}\rangle^2
}{
1600M\Lambda_{n,v}^2
}
\le\frac1{1600}.
\]
There are four coefficient pairs, whence
\[
\mu\left(
\bigcup_{\iota\in\{0,A\}}\ \bigcup_{\jmath\in\{0,A\}}
\mathcal Q_{\iota,\jmath}
\right)
\le
\sum_{\iota\in\{0,A\}}\ \sum_{\jmath\in\{0,A\}}
\mu(\mathcal Q_{\iota,\jmath})
\le\frac4{1600}.
\]
% lean: CausalSmith.Stat.FinitepHomogeneityDensegamma.ledgerCoeffError_indep
% lean: CausalSmith.Stat.FinitepHomogeneityDensegamma.ledgerCoeffError_cross_second_moment_le
% lean: CausalSmith.Stat.FinitepHomogeneityDensegamma.ledgerCoeffError_cross_tail

\item
Combining the two union bounds yields
\[
\mu(\mathcal F_{\mathrm{bad}})
\le
\frac4{800}+\frac4{1600}
=
\frac3{400}.
\]
% lean: CausalSmith.Stat.FinitepHomogeneityDensegamma.uniform_affine_profile_bound

\item
Fix a sample in \(\mathcal F_{\mathrm{bad}}^c\). By the definitions of the bad events,
\[
\|\operatorname{pr}_{\mathcal S_P}\mathsf E_{b,\iota}\|_2
\le40\sqrt{\Lambda_{n,v}},
\qquad
|\langle\mathsf E_{1,\iota},\mathsf E_{2,\jmath}\rangle|
\le40\sqrt M\,\Lambda_{n,v}.
\]
For every \(c\in\mathbb R\), the affine mean belongs to \(\mathcal S_P\). Orthogonal projection and Cauchy--Schwarz therefore give
\[
\begin{aligned}
|\langle\theta_P(c),\mathsf E_{b,\iota}\rangle|
&=
|\langle\theta_P(c),
\operatorname{pr}_{\mathcal S_P}\mathsf E_{b,\iota}\rangle|\\
&\le
\|\theta_P(c)\|_2\,
\|\operatorname{pr}_{\mathcal S_P}\mathsf E_{b,\iota}\|_2\\
&\le
40\sqrt{\Lambda_{n,v}}\,\|\theta_P(c)\|_2.
\end{aligned}
\]
Now take any \(|c|\le1/2\). The mean-plus-error expansion is
\[
W(c)-\|\theta_P(c)\|_2^2
=
\langle\theta_P(c),\mathsf E_1(c)+\mathsf E_2(c)\rangle
+
\langle\mathsf E_1(c),\mathsf E_2(c)\rangle.
\]
For the linear term, the triangle inequality yields
\[
\begin{aligned}
|\langle\theta_P(c),\mathsf E_1(c)+\mathsf E_2(c)\rangle|
&\le
\sum_{b=1}^{2}
\left(
|\langle\theta_P(c),\mathsf E_{b,0}\rangle|
+
|c|\,|\langle\theta_P(c),\mathsf E_{b,A}\rangle|
\right)\\
&\le
(2+2|c|)\,40\sqrt{\Lambda_{n,v}}\,\|\theta_P(c)\|_2\\
&\le
120\sqrt{\Lambda_{n,v}}\,\|\theta_P(c)\|_2.
\end{aligned}
\]
For the quadratic term, expansion of the two affine errors gives
\[
\begin{aligned}
|\langle\mathsf E_1(c),\mathsf E_2(c)\rangle|
&\le
|\langle\mathsf E_{1,0},\mathsf E_{2,0}\rangle|
+
|c|\,|\langle\mathsf E_{1,0},\mathsf E_{2,A}\rangle|\\
&\quad+
|c|\,|\langle\mathsf E_{1,A},\mathsf E_{2,0}\rangle|
+
|c|^2|\langle\mathsf E_{1,A},\mathsf E_{2,A}\rangle|\\
&\le
(1+2|c|+|c|^2)\,40\sqrt M\,\Lambda_{n,v}\\
&\le
\frac94\,40\sqrt M\,\Lambda_{n,v}
=
90\sqrt M\,\Lambda_{n,v}.
\end{aligned}
\]
Here \(|c|^2\le1/4\). Since both ledger terms are nonnegative,
\[
\begin{aligned}
|W(c)-\|\theta_P(c)\|_2^2|
&\le
120\sqrt{\Lambda_{n,v}}\,\|\theta_P(c)\|_2
+
90\sqrt M\,\Lambda_{n,v}\\
&\le
128\left(
\sqrt{\Lambda_{n,v}}\,\|\theta_P(c)\|_2
+
\sqrt M\,\Lambda_{n,v}
\right).
\end{aligned}
\]
The sample was fixed throughout and \(c\) was arbitrary in the profiling interval. Thus
\[
\mathcal F_{\mathrm{bad}}^c\subseteq\mathcal G_{\mathrm{prof}}.
\]
% lean: CausalSmith.Stat.FinitepHomogeneityDensegamma.ledgerTheta_inner_error_le_of_projection
% lean: CausalSmith.Stat.FinitepHomogeneityDensegamma.affine_profile_error_bound

\item
Finally, measurability of the bad event and the inclusion just proved imply
\[
\mu(\mathcal G_{\mathrm{prof}})
\ge
\mu(\mathcal F_{\mathrm{bad}}^c)
=
1-\mu(\mathcal F_{\mathrm{bad}})
\ge
1-\frac3{400}
=
\frac{397}{400}
=
0.9925.
\]
Together with the measurable-event argument and the zero-ledger case, this proves all the assertions.
% lean: CausalSmith.Stat.FinitepHomogeneityDensegamma.uniform_affine_profile_bound
\end{enumerate}
\end{proof}

Simultaneous control permits minimization over the candidate constant within the same event. Its fluctuation bound separates a term proportional to the population score norm from a term involving the coarse dimension and covariance ledger. When the covariance ledger vanishes, the simultaneous identity gives the population target exactly.

The public computation retains this affine structure explicitly. Using the one-based block convention of \cref{obj:def:calibration-handle}, set \(Z_{b,1}:=Z_{b,A}\) and
\[
w_{rr'}:=\langle Z_{1,r},Z_{2,r'}\rangle,
\qquad r,r'\in\{0,1\}.
\]
Then the computed quadratic is
\[
W(c)=w_{00}-c(w_{01}+w_{10})+c^2w_{11}.
\]
Its minimum over \([-1/2,1/2]\) is obtained by comparing the two endpoints and, when \(w_{11}>0\) and the vertex lies in the interval, the value at
\[
c_{\mathrm{vert}}:=\frac{w_{01}+w_{10}}{2w_{11}}.
\]
This formula preserves both mixed affine-coefficient contributions through their sum. The same calculation applies under the zero-based convention after translating the block indices.

Together, \cref{obj:lem:original-score-mean-ledger,obj:lem:original-score-covariance-ledger,obj:lem:uniform-affine-profile-bound} supply the population and fluctuation bounds used in \cref{obj:thm:whole-null-calibration-resolved}. The cutoff in \cref{obj:def:explicit-score-ledger} rounds
\[
2B^2+2^{16}\sqrt M\,\Lambda_{n,v}
\]
upward to a dyadic value. The calibration theorem gives size at most \(0.0075\) throughout the constant-effect null and type-II error at most \(0.0075\) at its calibrated separation, for \(n\ge4\) when that separation is below the model distance supremum. The zero rejection rule supplies the stated size guarantee at \(n=2,3\).

The calibrated separation \(R_{\mathrm{att}}(n,v)\) in \cref{obj:thm:whole-null-calibration-resolved} groups coarse approximation, the bias ledger, and the covariance contribution. The explicit attainment statement below expresses their combined requirement as a power of the sample size with a fully specified multiplier.

\subsection{Explicit attainment constants and separation}

Fix an admissible tuple \(v=(p,\alpha,\beta,\gamma)\) from \cref{obj:def:model}, and use the exponents \(q,S,t,E_0,E_4,E\) of \cref{obj:def:explicit-score-ledger}. Before stating attainment, define the minimum smoothness \(s_0\), the clipping allocation exponent \(\lambda_0\), the increment summability exponent \(\eta_0\), and the geometric-series factor \(A(x)\) by
\[
s_0:=\min\{\alpha,\beta,\gamma\},
\qquad
\lambda_0:=\frac{(p-1)^2}{4p},
\qquad
\eta_0:=\frac{(p-1)(p+6)}{4p},
\qquad
A(x):=(1-2^{-x})^{-1}\quad(x>0).
\]
Here \(A(x)\) denotes a scalar series factor; its argument distinguishes it from the treatment coordinate. Define the grouped bias, summation, singleton, and canonical-pair constants by
\[
\begin{aligned}
B_*&:=420+800\{A(\alpha)+A(\lambda_0)\},\\
D_*&:=1+\{\exp(1)S\log 2\}^{-1},\\
L_*&:=16+880A(s_0)+160D_*+960A(\alpha),\\
G_*&:=\sqrt{64}\{2^{t/2}A(1/2)+A(\eta_0/2)\}.
\end{aligned}
\]
The aggregate covariance constant \(A_*\) and attainment multiplier \(C^{\mathrm{att}}_v\) are
\[
A_*:=\sqrt{2}\{24576+64L_*^2+16384+256G_*^2\},
\qquad
C^{\mathrm{att}}_v:=16\{16+5B_*+1024\sqrt{A_*}\}.
\]
These definitions resolve every constant entering the public attainment multiplier.

Using the heterogeneity distance and model supremum defined in \cref{obj:thm:whole-null-calibration-resolved}, define the fixed witness distance \(d_0\) and public threshold \(N^{\mathrm{att}}_v\) by
\[
d_0:=\frac{1}{8\sqrt{2}},
\qquad
N^{\mathrm{att}}_v
:=\left\lceil\max\left\{4,
\left(\frac{2C^{\mathrm{att}}_v}{d_0}\right)^{1/E}\right\}\right\rceil.
\]
For the rejection expectation below, \(\mathcal D_n\) denotes \(n\) independent original records with law \(P\), and \(U\) denotes an independent uniform variable on \([0,1]\), as in \cref{obj:def:original-record-experiment}. Set
\[
\mathsf R_n(P,\phi):=\mathbb E[\phi(\mathcal D_n,U)].
\]

For an admissible bounded-model tuple \(w=(\alpha,\beta,\gamma)\) from \cref{obj:def:bounded-model}, define the matched exponent, multiplier, and distance supremum by
\[
E_{\mathrm b}
:=\min\left\{\frac{2\gamma}{4\gamma+1},
\frac{2(\alpha+\beta)}
{1+2(\alpha+\beta)+(\alpha+\beta)/(2\gamma)}\right\},
\qquad
C^{\mathrm{att}}_{(2,w)}
:=C^{\mathrm{att}}_{(2,\alpha,\beta,\gamma)},
\qquad
D_w^{\mathrm b}:=\sup_{P\in\mathcal M_w^{\mathrm b}}d(P).
\]
The following result gives attainment for the original model and its matched bounded submodel, together with public tuning bounds and uniform control of the constants on compact sets of admissible tuples.

\begin{theoremv}{thm:original-record-attainable-rate}[Original-record attainable separation rate]
For every admissible public tuple \(v=(p,\alpha,\beta,\gamma)\) of \cref{obj:def:model}, use the explicit rule \(\phi^*_{n,v}\) and its tuning quantities from \cref{obj:def:explicit-score-ledger,obj:def:sharp-frontier-scales}. Write
\[
q=\frac{p-1}{p},\qquad S=\alpha+\beta,\qquad
t=\frac{2-p}{p-1},\qquad
E_0=\frac{2\gamma q}{2\gamma+q},\qquad
E_4=\frac{2S}{1+2\alpha+\beta/q+S/(2\gamma)},\qquad
E=\min\{E_0,E_4\}.
\]
Define the minimum smoothness \(s_0\), increment summability exponent \(\eta_0\), geometric-series factor \(A(x)\), and attainment constants by
\[
\begin{aligned}
s_0&=\min\{\alpha,\beta,\gamma\},&
\lambda_0&=\frac{(p-1)^2}{4p},&
\eta_0&=\frac{(p-1)(p+6)}{4p},\\
A(x)&=(1-2^{-x})^{-1}\quad(x>0),\\
B_*&=420+800\{A(\alpha)+A(\lambda_0)\},\\
D_*&=1+\{\exp(1)S\log 2\}^{-1},\\
L_*&=16+880A(s_0)+160D_*+960A(\alpha),\\
G_*&=\sqrt{64}\{2^{t/2}A(1/2)+A(\eta_0/2)\},\\
A_*&=\sqrt{2}\{24576+64L_*^2+16384+256G_*^2\},\\
C^{\mathrm{att}}_v&=16\{16+5B_*+1024\sqrt{A_*}\}.
\end{aligned}
\]
Define the fixed witness distance \(d_0\) and public attainment threshold \(N^{\mathrm{att}}_v\) by
\[
d_0=\frac{1}{8\sqrt{2}},
\qquad
N^{\mathrm{att}}_v
=\left\lceil\max\left\{4,
\left(\frac{2C^{\mathrm{att}}_v}{d_0}\right)^{1/E}\right\}\right\rceil.
\]
Let \(\lambda\) be Lebesgue probability measure on \([0,1]\). Define the heterogeneity distance \(d(P)\) and its model supremum \(D_v\) by
\[
d(P)=
\left\{\int_{[0,1]}
\left(\tau_P(x)-\int_{[0,1]}\tau_P(z)\,d\lambda(z)\right)^2
\,d\lambda(x)\right\}^{1/2},
\qquad
D_v=\sup_{P\in\mathcal M_v}d(P).
\]
Let \(\mathcal D_n\) be a sample of \(n\) independent original records with law \(P\), and let \(U\) be an independent uniform randomization variable on \([0,1]\). For a test \(\phi\), define its rejection expectation by
\[
\mathsf R_n(P,\phi)=\mathbb E[\phi(\mathcal D_n,U)].
\]

Then \(C^{\mathrm{att}}_v>0\), and the following guarantees hold:
\begin{itemize}
\item \textbf{(Size and separation.)}
For every integer \(n\ge2\),
\[
\mathsf R_n(P,\phi^*_{n,v})\le0.0075
\qquad\text{for every }P\in H_0(v),
\]
where the null is given in \cref{obj:def:null}. Whenever
\(C^{\mathrm{att}}_v n^{-E}<D_v\),
\[
1-\mathsf R_n(P,\phi^*_{n,v})\le0.0075
\qquad
\text{for every }P\in\mathcal M_v
\text{ with }d(P)\ge C^{\mathrm{att}}_v n^{-E}.
\]
The capped critical radius of \cref{obj:def:critical-radius} satisfies
\[
r_n^*(v)\le C^{\mathrm{att}}_v n^{-E}.
\]

\item \textbf{(Witness threshold.)}
For every integer \(n\ge N^{\mathrm{att}}_v\),
\[
C^{\mathrm{att}}_v n^{-E}<d_0.
\]

\item \textbf{(Tuning bounds.)}
For every integer \(n\ge4\), the coarse rank \(M\), fine rank \(K\), block size \(s\), main clipping cutoff \(T_0\), and increment cutoffs \(T_j\) from \cref{obj:def:explicit-score-ledger,obj:def:blocks} satisfy
\[
M\le2s,\qquad K\le2s^2,\qquad
1\le T_0\le2s,\qquad
1\le T_j\le T_0\quad(1\le j\le L).
\]

\item \textbf{(Compact uniformity.)}
For every compact set \(K\) of public tuples in the product topology whose members are all admissible under \cref{obj:def:model}, there exist real constants \(C,N,e\), with \(e>0\), such that, for every \(v\in K\),
\[
C^{\mathrm{att}}_v\le C,\qquad
N^{\mathrm{att}}_v\le N,\qquad e\le E(v).
\]
Here \(E(v)\) denotes the displayed minimum exponent evaluated at \(v\).
\end{itemize}

For every admissible smoothness tuple \(w=(\alpha,\beta,\gamma)\) of \cref{obj:def:bounded-model}, define the bounded separation exponent \(E_{\mathrm b}\), matched attainment multiplier \(C^{\mathrm{att}}_{(2,w)}\), and bounded-model distance supremum \(D_w^{\mathrm b}\) by
\[
E_{\mathrm b}
=\min\left\{\frac{2\gamma}{4\gamma+1},
\frac{2(\alpha+\beta)}
{1+2(\alpha+\beta)+(\alpha+\beta)/(2\gamma)}\right\},
\qquad
C^{\mathrm{att}}_{(2,w)}
=C^{\mathrm{att}}_{(2,\alpha,\beta,\gamma)},
\qquad
D_w^{\mathrm b}=\sup_{P\in\mathcal M^{\mathrm b}_w}d(P).
\]
The original-model exponent evaluated at \((2,\alpha,\beta,\gamma)\) equals \(E_{\mathrm b}\). The bounded-model rule \(\phi^{*,\mathrm b}_{n,w}\) from \cref{obj:def:sharp-frontier-scales}, obtained by evaluating the same ledger rule at \((2,\alpha,\beta,\gamma)\), satisfies, for every integer \(n\ge2\),
\[
\mathsf R_n(P,\phi^{*,\mathrm b}_{n,w})\le0.0075
\qquad\text{for every }P\in H_0^{\mathrm b}(w),
\]
with the bounded null as in \cref{obj:def:bounded-null}. Whenever
\(C^{\mathrm{att}}_{(2,w)}n^{-E_{\mathrm b}}<D_w^{\mathrm b}\),
\[
1-\mathsf R_n(P,\phi^{*,\mathrm b}_{n,w})\le0.0075
\qquad
\text{for every }P\in\mathcal M^{\mathrm b}_w
\text{ with }d(P)\ge C^{\mathrm{att}}_{(2,w)}n^{-E_{\mathrm b}}.
\]
The bounded capped critical radius of \cref{obj:def:bounded-radius} satisfies
\[
r_n^{*,\mathrm b}(w)
\le C^{\mathrm{att}}_{(2,w)}n^{-E_{\mathrm b}}.
\]
\end{theoremv}

The attainment guarantee combines the mean separation in \cref{obj:lem:original-score-mean-ledger} with the simultaneous fluctuation bound in \cref{obj:lem:uniform-affine-profile-bound}. Its exponent is the minimum of the two public separation exponents, while its multiplier accounts explicitly for projection, clipping, and covariance contributions. Evaluating the same construction at moment order two gives the bounded-model exponent and multiplier. The power clauses apply when the stated separation is below the corresponding model distance supremum; the capped-radius bounds apply with the capping conventions in \cref{obj:def:critical-radius,obj:def:bounded-radius}.

The tuning bounds place every selected rank and outcome cutoff within the finite public grids of \cref{obj:def:explicit-score-ledger}. Compact uniformity supplies common bounds for the attainment constants and thresholds over compact sets of admissible tuples. Thus the explicit procedure links a finite original-record calculation to a separation guarantee whose dependence on the public moment and smoothness parameters is fully specified.


\section{Full-record lower bounds and tail essentiality}

The lower bounds complement the attainment guarantee in \cref{obj:thm:original-record-attainable-rate} by constructing null and alternative mixtures that remain close when the entire original record is observed. Two constructions supply the required comparisons: paired tents with constant propensity, and coupled continuous coefficient fields in the rough phase. Their separate bounded versions provide witnesses at moment order two.

We use total variation in the convention
\[
\operatorname{TV}(P,Q)=\sup_{E\ \mathrm{measurable}}|P(E)-Q(E)|.
\]
For probability laws \(P\ll Q\), the directed chi-square divergence is
\[
\chi^2(P\Vert Q)
=\int\left(\frac{dP}{dQ}-1\right)^2\,dQ
=\int\left(\frac{dP}{dQ}\right)^2\,dQ-1.
\]
These divergences compare probability laws on the same measurable space.

A finite prior draws one original-record law and then generates the independent sample from that law. Specifically, for
\[
\pi=\sum_{j=1}^{k}\omega_j\delta_{P_j},
\qquad
k\ge1,\quad \omega_j\ge0,\quad \sum_{j=1}^{k}\omega_j=1,
\]
write
\[
\operatorname{supp}_+(\pi)=\{P_j:\omega_j>0\},
\qquad
\mathbb P_{\pi,n}=\sum_{j=1}^{k}\omega_jP_j^{\otimes n}.
\]
Thus \(\mathbb P_{\pi,n}\) mixes product laws after the common latent law has been drawn. Repeated support laws have their weights combined. Normalization ensures that the positive-weight support is nonempty.

For the paired-tent comparison, let \(\lambda\) be Lebesgue probability measure on \([0,1]\), and let \(Q_{a,P}(dy\mid x)\) be the conditional outcome probability kernel in treatment arm \(a\). The following result supplies both rare-mark witnesses under a finite conditional moment and signed-binary witnesses in the matched bounded model.

\begin{lemmav}{lem:paired-tent-full-record-lower}[Paired-tent full-record lower bounds]
Let \(\lambda\) denote Lebesgue probability measure on \([0,1]\), and let \(Q_{a,P}(dy\mid x)\) denote the conditional outcome probability kernel under law \(P\) in treatment arm \(a\in\{0,1\}\). Define the heterogeneity distance, its model suprema, and the fixed witness distance by
\[
d(P)=
\left\{\int_0^1
\left(\tau_P(x)-\int_0^1\tau_P(z)\,\lambda(dz)\right)^2
\,\lambda(dx)\right\}^{1/2},
\qquad
D_v=\sup_{P\in\mathcal M_v}d(P),
\qquad
D_w^{\mathrm b}=\sup_{P\in\mathcal M^{\mathrm b}_w}d(P),
\qquad
d_0=\frac{1}{8\sqrt2},
\]
where the models are defined in \cref{obj:def:model,obj:def:bounded-model}.

Write a normalized finite probability prior as
\[
\pi=\sum_{j=1}^{k}\omega_j\delta_{P_j},
\qquad
\omega_j\ge0,
\qquad
\sum_{j=1}^{k}\omega_j=1.
\]
Define its positive-weight support and full original-record mixture by
\[
\operatorname{supp}_+(\pi)=\{P_j:\omega_j>0\},
\qquad
\mathbb P_{\pi,n}=\sum_{j=1}^{k}\omega_jP_j^{\otimes n}.
\]
For probability laws on the same measurable space, define total variation by
\[
\operatorname{TV}(P,Q)=
\sup_{E\ \mathrm{measurable}}|P(E)-Q(E)|.
\]

For every admissible moment and smoothness tuple \(v=(p,\alpha,\beta,\gamma)\) as in \cref{obj:def:model}, put
\[
q=\frac{p-1}{p},
\qquad
E_0=\frac{2\gamma q}{2\gamma+q},
\qquad
c^{\mathrm e}_v=\frac{4^{-\gamma}}{16\sqrt3}.
\]
There exists a sequence of mark magnitudes \(L:\mathbb N\to\mathbb R\) with
\[
\lim_{n\to\infty}L(n)=+\infty
\]
such that, for every integer \(n\ge2\), there exist a null law \(P_0\) and a normalized finite probability prior \(\pi_1\), with at least one strictly positive weight, satisfying:
\begin{itemize}
\item \textbf{(Null and propensity.)}
The law \(P_0\) belongs to \(H_0(v)\) from \cref{obj:def:null}, and
\[
e_{P_0}(x)=\frac12
\qquad\text{for every }x\in[0,1].
\]

\item \textbf{(Arm support and moments.)}
The magnitude \(L(n)\) is positive. For \(P=P_0\) and every \(P\in\operatorname{supp}_+(\pi_1)\), each arm \(a\in\{0,1\}\) satisfies, for \(\lambda\)-almost every \(x\),
\[
Q_{a,P}\bigl(\{0,-L(n),L(n)\}\mid x\bigr)=1,
\qquad
\int_{\mathbb R}|y|^p\,Q_{a,P}(dy\mid x)\le1.
\]

\item \textbf{(Alternative placement.)}
Every \(P\in\operatorname{supp}_+(\pi_1)\) belongs to \(\mathcal M_v\) from \cref{obj:def:model}, satisfies \(e_P(x)=1/2\) for every \(x\in[0,1]\), and obeys
\[
c^{\mathrm e}_v n^{-E_0}\le d(P)<d_0.
\]
Moreover, \(d_0\le D_v\).

\item \textbf{(Full-record comparison.)}
Define the alternative mixture by
\[
\overline P_1^{(n)}=\mathbb P_{\pi_1,n}.
\]
The same witnesses satisfy
\[
\operatorname{TV}\bigl(P_0^{\otimes n},\overline P_1^{(n)}\bigr)<\frac12.
\]
\end{itemize}
For every such \(n\), the supplied-propensity and original-record minimax testing risks in \cref{obj:def:oracle-risk,obj:def:testing-risk} satisfy
\[
B_n^{\mathrm e}\bigl(v,c^{\mathrm e}_v n^{-E_0}\bigr)\ge\frac25,
\qquad
B_n\bigl(v,c^{\mathrm e}_v n^{-E_0}\bigr)\ge\frac25.
\]

For every admissible smoothness tuple \(w=(\alpha,\beta,\gamma)\) as in \cref{obj:def:bounded-model}, put
\[
c_w^{\mathrm b}=\frac{4^{-\gamma}}{16\sqrt3}.
\]
For every integer \(n\ge2\), there exist a signed-binary null law \(P_0^{\mathrm b}\) and a normalized finite probability prior \(\pi_1^{\mathrm b}\), with at least one strictly positive weight, satisfying:
\begin{itemize}
\item \textbf{(Signed-binary null.)}
The law \(P_0^{\mathrm b}\) belongs to \(H_0^{\mathrm{bin}}(w)\) from \cref{obj:def:binary-null}, and
\[
e_{P_0^{\mathrm b}}(x)=\frac12
\qquad\text{for every }x\in[0,1].
\]

\item \textbf{(Signed-binary alternatives.)}
Every \(P\in\operatorname{supp}_+(\pi_1^{\mathrm b})\) belongs to \(\mathcal M^{\mathrm{bin}}_w\) from \cref{obj:def:binary-model}, satisfies \(e_P(x)=1/2\) for every \(x\in[0,1]\), and obeys
\[
c_w^{\mathrm b}n^{-2\gamma/(4\gamma+1)}
\le d(P)<d_0.
\]
Moreover, \(d_0\le D_w^{\mathrm b}\).

\item \textbf{(Full-record comparison.)}
Define the signed-binary alternative mixture by
\[
\overline P_{1,\mathrm b}^{(n)}
=\mathbb P_{\pi_1^{\mathrm b},n}.
\]
The same witnesses satisfy
\[
\operatorname{TV}\bigl((P_0^{\mathrm b})^{\otimes n},
\overline P_{1,\mathrm b}^{(n)}\bigr)<\frac12.
\]
\end{itemize}
For every such \(n\), the bounded-model minimax testing risk in \cref{obj:def:bounded-risk} satisfies
\[
B_n^{\mathrm b}\bigl(w,c_w^{\mathrm b}n^{-2\gamma/(4\gamma+1)}\bigr)
\ge\frac25.
\]
\end{lemmav}

\begin{proof}[Proof of \cref{obj:lem:paired-tent-full-record-lower}]
\begin{enumerate}
\item \emph{Choice of the mark sequence.}
Fix an admissible tuple \(v=(p,\alpha,\beta,\gamma)\). Define
\[
\vartheta_\triangle=\frac{2q}{2\gamma+q},
\qquad
\kappa_\triangle=\frac1{16},
\qquad
N_{\triangle,n}=2\left\lceil n^{\vartheta_\triangle}\right\rceil,
\qquad
h_{\triangle,n}=N_{\triangle,n}^{-1}
\quad(n\in\mathbb N),
\]
where the reciprocal at zero is defined to be zero. Choose the sequence
\[
L(n)=
\begin{cases}
0,&n=0,\\
N_{\triangle,n}^{\gamma/(p-1)},&n\ge1.
\end{cases}
\]
Admissibility gives \(p-1>0\), \(q>0\), and \(0<\gamma\le1\), hence
\[
\vartheta_\triangle>0,
\qquad
\frac{\gamma}{p-1}>0.
\]
For \(n\ge1\), the ceiling inequality gives
\[
N_{\triangle,n}\ge n^{\vartheta_\triangle}.
\]
Thus \(N_{\triangle,n}\to\infty\), and raising it to the positive power
\(\gamma/(p-1)\) proves \(L(n)\to\infty\).

Now fix \(n\ge2\). Since \(n^{\vartheta_\triangle}\ge1\),
\[
2n^{\vartheta_\triangle}
\le N_{\triangle,n}
<2\bigl(n^{\vartheta_\triangle}+1\bigr)
\le4n^{\vartheta_\triangle}.
\]
In particular, \(N_{\triangle,n}\) is even and at least two. Define
\[
\varepsilon_{\triangle,n}
=h_{\triangle,n}^{\gamma/q}.
\]
Then
\[
0<h_{\triangle,n}<1,
\qquad
0<\varepsilon_{\triangle,n}<1,
\qquad
L(n)=h_{\triangle,n}^{-\gamma/(p-1)}>0.
\]
The identities
\[
\frac{\gamma}{q}-\frac{\gamma}{p-1}=\gamma,
\qquad
\frac{\gamma}{q}-\frac{p\gamma}{p-1}=0
\]
give
\[
\varepsilon_{\triangle,n}L(n)=h_{\triangle,n}^{\gamma},
\qquad
\varepsilon_{\triangle,n}L(n)^p=1.
\]
% lean: CausalSmith.Stat.FinitepHomogeneityDensegamma.tentMagnitude_tendsto_atTop

\item \emph{Paired tents and the finite prior.}
Define the compactly supported tent on the whole real line by
\[
\psi_\triangle(t_\triangle)=
\begin{cases}
2\min\{t_\triangle,1-t_\triangle\},&0\le t_\triangle\le1,\\
0,&\text{otherwise},
\end{cases}
\qquad t_\triangle\in\mathbb R.
\]
Define the number of pairs, their index set, and their sign vectors by
\[
m_{\triangle,n}=\frac{N_{\triangle,n}}2,
\qquad
\mathscr J_{\triangle,n}=\{0,\ldots,m_{\triangle,n}-1\},
\qquad
\varsigma\in\{-1,1\}^{\mathscr J_{\triangle,n}}.
\]
For \(j_\triangle\in\mathscr J_{\triangle,n}\) and \(x\in\mathbb R\), put
\[
\psi_{\triangle,j_\triangle}^{\varsigma}(x)
=
\varsigma_{j_\triangle}
\left[
\psi_\triangle(N_{\triangle,n}x-2j_\triangle)
-\psi_\triangle(N_{\triangle,n}x-(2j_\triangle+1))
\right],
\]
\[
f_{\triangle,\varsigma}(x)
=
\sum_{j_\triangle\in\mathscr J_{\triangle,n}}
\psi_{\triangle,j_\triangle}^{\varsigma}(x).
\]
These functions are continuous. A nonzero paired bump satisfies
\[
\psi_{\triangle,j_\triangle}^{\varsigma}(x)\ne0
\quad\Longrightarrow\quad
2j_\triangle<N_{\triangle,n}x<2j_\triangle+2.
\]
The open intervals in this display are disjoint for distinct pair indices.
Within a pair, the two individual tents also have disjoint nonzero supports.
Consequently,
\[
\left|\psi_{\triangle,j_\triangle}^{\varsigma}(x)\right|\le1,
\qquad
|f_{\triangle,\varsigma}(x)|\le1.
\]

For the Hölder estimate, the identity
\[
\psi_\triangle(t_\triangle)
=\max\{0,2\min(t_\triangle,1-t_\triangle)\}
\]
and the Lipschitz inequalities for minimum and maximum imply
\[
|\psi_\triangle(t_\triangle)-\psi_\triangle(t_\triangle')|
\le2|t_\triangle-t_\triangle'|
\qquad(t_\triangle,t_\triangle'\in\mathbb R).
\]
Applying this to the two tents in a pair gives
\[
\left|
\psi_{\triangle,j_\triangle}^{\varsigma}(x)
-\psi_{\triangle,j_\triangle}^{\varsigma}(x')
\right|
\le4N_{\triangle,n}|x-x'|.
\]
Define the active pairs at a point by
\[
\mathscr A_{\triangle,\varsigma}(x)
=
\left\{
j_\triangle\in\mathscr J_{\triangle,n}:
\psi_{\triangle,j_\triangle}^{\varsigma}(x)\ne0
\right\}.
\]
The support calculation gives
\[
|\mathscr A_{\triangle,\varsigma}(x)|\le1,
\qquad
|\mathscr A_{\triangle,\varsigma}(x)
 \cup\mathscr A_{\triangle,\varsigma}(x')|\le2.
\]
Every summand outside this union vanishes at both points. Therefore
\[
\begin{aligned}
|f_{\triangle,\varsigma}(x)-f_{\triangle,\varsigma}(x')|
&\le
\sum_{j_\triangle\in
\mathscr A_{\triangle,\varsigma}(x)\cup
\mathscr A_{\triangle,\varsigma}(x')}
\left|
\psi_{\triangle,j_\triangle}^{\varsigma}(x)
-\psi_{\triangle,j_\triangle}^{\varsigma}(x')
\right|\\
&\le8N_{\triangle,n}|x-x'|.
\end{aligned}
\]
If \(N_{\triangle,n}|x-x'|\le1\), then
\[
N_{\triangle,n}|x-x'|
\le\bigl(N_{\triangle,n}|x-x'|\bigr)^\gamma.
\]
If \(N_{\triangle,n}|x-x'|>1\), the unit envelope instead gives
\[
|f_{\triangle,\varsigma}(x)-f_{\triangle,\varsigma}(x')|
\le2
\le8\bigl(N_{\triangle,n}|x-x'|\bigr)^\gamma.
\]
Both cases yield
\[
h_{\triangle,n}^{\gamma}
|f_{\triangle,\varsigma}(x)-f_{\triangle,\varsigma}(x')|
\le8|x-x'|^\gamma.
\]

Give every sign vector the same weight, and define
\[
\pi_1
=
2^{-m_{\triangle,n}}
\sum_{\varsigma\in\{-1,1\}^{\mathscr J_{\triangle,n}}}
\delta_{P_\varsigma},
\]
where the laws \(P_\varsigma\) are constructed below. There are exactly
\(2^{m_{\triangle,n}}\) sign vectors, so
\[
2^{-m_{\triangle,n}}>0,
\qquad
\sum_{\varsigma}2^{-m_{\triangle,n}}=1.
\]
The constant sign vector supplies an index, proving that the prior has a
strictly positive weight.
% lean: CausalSmith.Stat.FinitepHomogeneityDensegamma.tentEffect_holderBall

\item \emph{The null law and arm moments.}
Construct \(P_0\) and \(P_\varsigma\) with uniform covariate distribution and
propensity \(1/2\). Their common control kernel and their treated kernels are
\[
Q_{0,P_0}(\,\cdot\mid x)
=
Q_{0,P_\varsigma}(\,\cdot\mid x)=\delta_0,
\]
\[
Q_{1,P_0}(\,\cdot\mid x)
=
(1-\varepsilon_{\triangle,n})\delta_0
+\frac{\varepsilon_{\triangle,n}}2\delta_{L(n)}
+\frac{\varepsilon_{\triangle,n}}2\delta_{-L(n)},
\]
\[
\begin{aligned}
Q_{1,P_\varsigma}(\,\cdot\mid x)
={}&(1-\varepsilon_{\triangle,n})\delta_0\\
&+\frac{\varepsilon_{\triangle,n}}2
 \bigl(1+\kappa_\triangle f_{\triangle,\varsigma}(x)\bigr)\delta_{L(n)}\\
&+\frac{\varepsilon_{\triangle,n}}2
 \bigl(1-\kappa_\triangle f_{\triangle,\varsigma}(x)\bigr)\delta_{-L(n)}.
\end{aligned}
\]
Precisely, these kernels specify the record laws through
\[
P(dx,\{a\},dy)
=\frac12\,\lambda(dx)\,Q_{a,P}(dy\mid x),
\qquad
a\in\{0,1\},
\qquad
P\in\{P_0,P_\varsigma\}.
\]
The displayed masses sum to one. They are nonnegative because
\[
1\pm\kappa_\triangle f_{\triangle,\varsigma}(x)
\ge1-\kappa_\triangle=\frac{15}{16}>0,
\qquad
1-\varepsilon_{\triangle,n}>0.
\]
Continuity of the tents supplies the required kernel measurability.

For every \(x\in[0,1]\), direct integration gives
\[
e_{P_0}(x)=\frac12,
\qquad
m_{0,P_0}(x)=m_{1,P_0}(x)=\tau_{P_0}(x)=0,
\]
and, for \(P=P_0\) or \(P=P_\varsigma\),
\[
Q_{a,P}(\{0,-L(n),L(n)\}\mid x)=1,
\qquad
\int |y|^p\,Q_{a,P}(dy\mid x)
=
\begin{cases}
0,&a=0,\\
\varepsilon_{\triangle,n}L(n)^p=1,&a=1.
\end{cases}
\]
Thus both arm moments are at most one. For \(P_0\), the propensity lies in
\([1/4,3/4]\), its supremum norm is \(1/2\), and its Hölder differences vanish.
The baseline and effect have zero supremum norms and zero Hölder differences.
Together with uniform design and the moment bound \(1\le10\), these facts
verify all the restrictions in \cref{obj:def:model}. Its effect is the constant
zero, so \cref{obj:def:null} gives
\[
P_0\in H_0(v).
\]
% lean: CausalSmith.Stat.FinitepHomogeneityDensegamma.tentLaw_false_inNull

\item \emph{Alternative membership and separation.}
The alternative arm means are
\[
e_{P_\varsigma}(x)=\frac12,
\qquad
m_{0,P_\varsigma}(x)=0,
\]
\[
m_{1,P_\varsigma}(x)=\tau_{P_\varsigma}(x)
=
\varepsilon_{\triangle,n}L(n)\kappa_\triangle
f_{\triangle,\varsigma}(x)
=
\kappa_\triangle h_{\triangle,n}^{\gamma}
f_{\triangle,\varsigma}(x).
\]
The geometric estimates imply
\[
\|\tau_{P_\varsigma}\|_\infty
\le\kappa_\triangle h_{\triangle,n}^{\gamma}
\le\frac1{16}<\frac12,
\]
\[
|\tau_{P_\varsigma}(x)-\tau_{P_\varsigma}(x')|
\le8\kappa_\triangle|x-x'|^\gamma
=\frac12|x-x'|^\gamma
\le20|x-x'|^\gamma.
\]
Thus the effect satisfies its Hölder and supremum restrictions. The constant
propensity and zero baseline satisfy their respective smoothness and bound
restrictions, and the preceding arm calculation supplies the raw moments.
Hence
\[
P_\varsigma\in\mathcal M_v
\qquad\text{for every }\varsigma.
\]

To compute the centered distance, reflection about \(1/2\) gives
\[
\int_0^1\psi_\triangle(t_\triangle)\,dt_\triangle
=
2\int_0^{1/2}2t_\triangle\,dt_\triangle=\frac12,
\]
\[
\int_0^1\psi_\triangle(t_\triangle)^2\,dt_\triangle
=
2\int_0^{1/2}4t_\triangle^2\,dt_\triangle=\frac13.
\]
A change of variables, using the compact support of the tent, yields, for
every integer \(i_\triangle\) with \(0\le i_\triangle<N_{\triangle,n}\),
\[
\int_0^1
\psi_\triangle(N_{\triangle,n}x-i_\triangle)\,\lambda(dx)
=\frac{h_{\triangle,n}}2,
\]
\[
\int_0^1
\psi_\triangle(N_{\triangle,n}x-i_\triangle)^2\,\lambda(dx)
=\frac{h_{\triangle,n}}3.
\]
The opposite signs cancel the two equal areas in each pair. Disjoint nonzero
supports eliminate the cross terms in the square. Consequently,
\[
\int f_{\triangle,\varsigma}\,d\lambda
=
\sum_{j_\triangle}\varsigma_{j_\triangle}
\left(\frac{h_{\triangle,n}}2-\frac{h_{\triangle,n}}2\right)=0,
\]
\[
\int f_{\triangle,\varsigma}^2\,d\lambda
=
m_{\triangle,n}\frac{2h_{\triangle,n}}3=\frac13.
\]
It follows that
\[
\int\tau_{P_\varsigma}\,d\lambda=0,
\qquad
d(P_\varsigma)
=\frac{\kappa_\triangle h_{\triangle,n}^{\gamma}}{\sqrt3}.
\]
The upper rank bound supplies
\[
h_{\triangle,n}^{\gamma}
\ge(4n^{\vartheta_\triangle})^{-\gamma}
=4^{-\gamma}n^{-E_0},
\]
because \(\gamma\vartheta_\triangle=E_0\). Therefore
\[
d(P_\varsigma)\ge
\frac{4^{-\gamma}}{16\sqrt3}n^{-E_0}
=c_v^{\mathrm e}n^{-E_0}.
\]

For the strict upper comparison, expand the centered square and use the
effect envelope:
\[
\begin{aligned}
d(P_\varsigma)^2
&=
\int\tau_{P_\varsigma}^2\,d\lambda
-\left(\int\tau_{P_\varsigma}\,d\lambda\right)^2\\
&\le\int\tau_{P_\varsigma}^2\,d\lambda
\le\frac1{16^2}.
\end{aligned}
\]
Since \(\sqrt2<2\),
\[
d(P_\varsigma)\le\frac1{16}
<\frac1{8\sqrt2}=d_0.
\]
% lean: CausalSmith.Stat.FinitepHomogeneityDensegamma.paired_effect_distance_lower

\item \emph{Comparison of the complete-record laws.}
Define the one-record likelihood on
\([0,1]\times\{0,1\}\times\mathbb R\) by
\[
\ell_{\triangle,\varsigma}(x,a,y)=
\begin{cases}
1+\kappa_\triangle f_{\triangle,\varsigma}(x),
&a=1,\ y=L(n),\\
1-\kappa_\triangle f_{\triangle,\varsigma}(x),
&a=1,\ y=-L(n),\\
1,&\text{otherwise}.
\end{cases}
\]
Comparison of the conditional atom masses proves
\[
P_\varsigma\ll P_0,
\qquad
\frac{dP_\varsigma}{dP_0}=\ell_{\triangle,\varsigma}
\quad P_0\text{-almost surely},
\qquad
\frac{15}{16}\le\ell_{\triangle,\varsigma}\le\frac{17}{16}.
\]
At each covariate value, summing over the control record and the three
treated atoms gives
\[
\begin{aligned}
&\frac12+\frac{1-\varepsilon_{\triangle,n}}2\\
&\quad+\frac{\varepsilon_{\triangle,n}}4
\left[
(1+\kappa_\triangle f_{\triangle,\varsigma}(x))
(1+\kappa_\triangle f_{\triangle,\varsigma'}(x))
+
(1-\kappa_\triangle f_{\triangle,\varsigma}(x))
(1-\kappa_\triangle f_{\triangle,\varsigma'}(x))
\right]\\
&\hspace{2cm}
=1+\frac{\varepsilon_{\triangle,n}\kappa_\triangle^2}{2}
f_{\triangle,\varsigma}(x)f_{\triangle,\varsigma'}(x).
\end{aligned}
\]
The same support calculation used for the squared tents gives
\[
\begin{aligned}
\int f_{\triangle,\varsigma}f_{\triangle,\varsigma'}\,d\lambda
&=
\sum_{j_\triangle}
\varsigma_{j_\triangle}\varsigma'_{j_\triangle}
\left(\frac{h_{\triangle,n}}3+\frac{h_{\triangle,n}}3\right)\\
&=\frac{2h_{\triangle,n}}3
\sum_{j_\triangle}\varsigma_{j_\triangle}\varsigma'_{j_\triangle}.
\end{aligned}
\]
Define the overlap coefficient by
\[
\varrho_{\triangle,n}
=\frac{\varepsilon_{\triangle,n}\kappa_\triangle^2}{6}.
\]
Integration over the entire record now yields
\[
\int\ell_{\triangle,\varsigma}\ell_{\triangle,\varsigma'}\,dP_0
=
1+\frac{\varrho_{\triangle,n}}{m_{\triangle,n}}
\sum_{j_\triangle}\varsigma_{j_\triangle}\varsigma'_{j_\triangle}.
\]
Here \(m_{\triangle,n}\ge1\), \(0\le\varrho_{\triangle,n}\le1\), and
\[
\left|
\sum_{j_\triangle}\varsigma_{j_\triangle}\varsigma'_{j_\triangle}
\right|\le m_{\triangle,n}.
\]
Thus every overlap base is nonnegative.

For probability laws \(\mu_{\rm alt}\ll\mu_{\rm ref}\) with square-integrable
density, define the directed chi-square divergence by
\[
\chi^2(\mu_{\rm alt}\Vert\mu_{\rm ref})
=
\int
\left(\frac{d\mu_{\rm alt}}{d\mu_{\rm ref}}-1\right)^2
\,d\mu_{\rm ref}.
\]
The complete-record mixture and its density are
\[
\overline P_1^{(n)}
=
2^{-m_{\triangle,n}}\sum_\varsigma P_\varsigma^{\otimes n},
\]
\[
G_\triangle(\mathcal D_n)
=
\frac{d\overline P_1^{(n)}}{dP_0^{\otimes n}}(\mathcal D_n)
=
2^{-m_{\triangle,n}}\sum_\varsigma
\prod_{i_{\rm rec}=1}^n
\ell_{\triangle,\varsigma}(X_{i_{\rm rec}},A_{i_{\rm rec}},Y_{i_{\rm rec}}).
\]
These are bounded finite densities, so expansion of their square and
independence of the records give
\[
\begin{aligned}
1+\chi^2(\overline P_1^{(n)}\Vert P_0^{\otimes n})
&=\int G_\triangle^2\,dP_0^{\otimes n}\\
&=
2^{-2m_{\triangle,n}}
\sum_\varsigma\sum_{\varsigma'}
\left(
1+\frac{\varrho_{\triangle,n}}{m_{\triangle,n}}
\sum_{j_\triangle}\varsigma_{j_\triangle}\varsigma'_{j_\triangle}
\right)^n.
\end{aligned}
\]
Nonnegativity of the bases and \(1+t_\triangle\le e^{t_\triangle}\) imply
\[
\left(
1+\frac{\varrho_{\triangle,n}}{m_{\triangle,n}}
\sum_{j_\triangle}\varsigma_{j_\triangle}\varsigma'_{j_\triangle}
\right)^n
\le
\exp\left(
\frac{n\varrho_{\triangle,n}}{m_{\triangle,n}}
\sum_{j_\triangle}\varsigma_{j_\triangle}\varsigma'_{j_\triangle}
\right).
\]
For each fixed \(\varsigma\), the sum over \(\varsigma'\) factors coordinate
by coordinate:
\[
\begin{aligned}
&2^{-m_{\triangle,n}}\sum_{\varsigma'}
\exp\left(
\frac{n\varrho_{\triangle,n}}{m_{\triangle,n}}
\sum_{j_\triangle}\varsigma_{j_\triangle}\varsigma'_{j_\triangle}
\right)\\
&\qquad=
\prod_{j_\triangle}
\cosh\left(
\frac{n\varrho_{\triangle,n}}{m_{\triangle,n}}\varsigma_{j_\triangle}
\right)
\le
\exp\left(\frac{(n\varrho_{\triangle,n})^2}{2m_{\triangle,n}}\right).
\end{aligned}
\]
The last inequality follows from
\[
\cosh(t_\triangle)
=
\sum_{i_{\rm exp}=0}^{\infty}
\frac{t_\triangle^{2i_{\rm exp}}}{(2i_{\rm exp})!}
\le
\sum_{i_{\rm exp}=0}^{\infty}
\frac{(t_\triangle^2/2)^{i_{\rm exp}}}{i_{\rm exp}!}
=e^{t_\triangle^2/2},
\]
using \((2i_{\rm exp})!\ge2^{i_{\rm exp}}i_{\rm exp}!\).
Averaging also over \(\varsigma\) therefore gives
\[
1+\chi^2(\overline P_1^{(n)}\Vert P_0^{\otimes n})
\le
\exp\left(
\frac{\kappa_\triangle^4}{36}
n^2\varepsilon_{\triangle,n}^2h_{\triangle,n}
\right).
\]
Indeed, \(2m_{\triangle,n}=N_{\triangle,n}=h_{\triangle,n}^{-1}\).
Moreover,
\[
h_{\triangle,n}\le n^{-\vartheta_\triangle},
\qquad
\vartheta_\triangle\left(\frac{2\gamma}{q}+1\right)=2,
\]
so
\[
n^2\varepsilon_{\triangle,n}^2h_{\triangle,n}
=
n^2h_{\triangle,n}^{2\gamma/q+1}
\le n^2n^{-2}=1.
\]
Consequently,
\[
1+\chi^2(\overline P_1^{(n)}\Vert P_0^{\otimes n})
\le e^{\kappa_\triangle^4/36}
\le\frac1{1-\kappa_\triangle^4/36}<2.
\]
The middle inequality follows by comparing the exponential and geometric
series; the last uses \(\kappa_\triangle=1/16\).

Since \(\int(G_\triangle-1)\,dP_0^{\otimes n}=0\), its positive and negative
parts have equal integrals. The event \(\{G_\triangle\ge1\}\) then gives
\[
\operatorname{TV}(P_0^{\otimes n},\overline P_1^{(n)})
=
\frac12\int|G_\triangle-1|\,dP_0^{\otimes n}.
\]
Cauchy--Schwarz under the probability measure \(P_0^{\otimes n}\) yields
\[
\operatorname{TV}(P_0^{\otimes n},\overline P_1^{(n)})
\le
\frac12
\sqrt{\chi^2(\overline P_1^{(n)}\Vert P_0^{\otimes n})}
<\frac12.
\]
% lean: CausalSmith.Stat.FinitepHomogeneityDensegamma.tentMixture_tv_lt_half

\item \emph{Supplied-propensity and original-record risks.}
By \cref{obj:lem:causal-nonempty},
\[
d_0\le D_v.
\]
Define the separation used here by
\[
r_\triangle=c_v^{\mathrm e}n^{-E_0}.
\]
The positive prior weights and the established placement give
\[
0<r_\triangle\le d(P_\varsigma)<d_0\le D_v.
\]

Take any supplied-propensity test satisfying the size constraint in
\cref{obj:def:oracle-risk}. Define its section at the common propensity and
its seed average by
\[
\varpi_{\mathrm{seed}}\in[0,1],
\qquad
\phi_{1/2}(\mathcal D_n,\varpi_{\mathrm{seed}})
=
\phi^{\mathrm e}(\mathcal D_n,\varpi_{\mathrm{seed}};\,e\equiv1/2),
\qquad
\overline\phi_{1/2}(\mathcal D_n)
=
\int_0^1\phi_{1/2}(\mathcal D_n,\varpi_{\mathrm{seed}})\,\lambda(d\varpi_{\mathrm{seed}}).
\]
This is a measurable function with values in \([0,1]\). Integrating its level
sets shows that its expectation difference under two probability laws is
bounded by their total variation. Thus
\[
\left|
2^{-m_{\triangle,n}}\sum_\varsigma
\mathsf R_n^{\mathrm e}(P_\varsigma,\phi^{\mathrm e})
-
\mathsf R_n^{\mathrm e}(P_0,\phi^{\mathrm e})
\right|
\le
\operatorname{TV}(P_0^{\otimes n},\overline P_1^{(n)})
<\frac12.
\]
Since \(P_0\in H_0(v)\),
\[
2^{-m_{\triangle,n}}\sum_\varsigma
\mathsf R_n^{\mathrm e}(P_\varsigma,\phi^{\mathrm e})
<\frac1{10}+\frac12=\frac35.
\]
If every positive-weight alternative had type-II error strictly below
\(2/5\), normalization of the weights would instead imply
\[
2^{-m_{\triangle,n}}\sum_\varsigma
\mathsf R_n^{\mathrm e}(P_\varsigma,\phi^{\mathrm e})
\ge
2^{-m_{\triangle,n}}\sum_\varsigma\frac35
=\frac35.
\]
Hence some support law has type-II error at least \(2/5\). Its model
membership and separation imply
\[
\sup_{\substack{P\in\mathcal M_v\\d(P)\ge r_\triangle}}
\bigl\{1-\mathsf R_n^{\mathrm e}(P,\phi^{\mathrm e})\bigr\}
\ge\frac25.
\]
The constant-zero test makes the admissible test class nonempty, and all
rejection probabilities lie in \([0,1]\). Taking the infimum therefore gives
\[
B_n^{\mathrm e}(v,r_\triangle)\ge\frac25.
\]

For an original-record test \(\phi\), use instead
\[
\overline\phi(\mathcal D_n)
=\int_0^1\phi(\mathcal D_n,\varpi_{\mathrm{seed}})\,\lambda(d\varpi_{\mathrm{seed}}).
\]
The identical expectation comparison gives
\[
2^{-m_{\triangle,n}}\sum_\varsigma
\mathsf R_n(P_\varsigma,\phi)<\frac35.
\]
The same weighted-sum contradiction supplies a support law with type-II
error at least \(2/5\). Taking the supremum and then the infimum in
\cref{obj:def:testing-risk} proves
\[
B_n(v,r_\triangle)\ge\frac25.
\]
This completes the first witness package.
% lean: CausalSmith.Stat.FinitepHomogeneityDensegamma.tent_canonical_lower_receipt

\item \emph{The signed-binary construction.}
Fix an admissible \(w=(\alpha,\beta,\gamma)\) and \(n\ge2\). Define
\[
\vartheta_\triangle^{\mathrm b}=\frac2{4\gamma+1},
\qquad
N_{\triangle,n}^{\mathrm b}
=2\left\lceil n^{\vartheta_\triangle^{\mathrm b}}\right\rceil,
\qquad
h_{\triangle,n}^{\mathrm b}=(N_{\triangle,n}^{\mathrm b})^{-1},
\]
\[
m_{\triangle,n}^{\mathrm b}=\frac{N_{\triangle,n}^{\mathrm b}}2,
\qquad
\mathscr J_{\triangle,n}^{\mathrm b}
=\{0,\ldots,m_{\triangle,n}^{\mathrm b}-1\},
\qquad
\delta_{\triangle,n}^{\mathrm b}
=\kappa_\triangle(h_{\triangle,n}^{\mathrm b})^\gamma.
\]
For \(\varsigma\in\{-1,1\}^{\mathscr J_{\triangle,n}^{\mathrm b}}\), define
\[
f_{\triangle,\varsigma}^{\mathrm b}(x)
=
\sum_{j_\triangle\in\mathscr J_{\triangle,n}^{\mathrm b}}
\varsigma_{j_\triangle}
\left[
\psi_\triangle(N_{\triangle,n}^{\mathrm b}x-2j_\triangle)
-\psi_\triangle(N_{\triangle,n}^{\mathrm b}x-(2j_\triangle+1))
\right],
\qquad x\in\mathbb R.
\]
These are the preceding paired-tent constructions evaluated at moment order
two, for which \(q=1/2\). In particular,
\[
2n^{\vartheta_\triangle^{\mathrm b}}
\le N_{\triangle,n}^{\mathrm b}
\le4n^{\vartheta_\triangle^{\mathrm b}},
\qquad
0<h_{\triangle,n}^{\mathrm b}<1,
\]
\[
|f_{\triangle,\varsigma}^{\mathrm b}(x)|\le1,
\qquad
(h_{\triangle,n}^{\mathrm b})^\gamma
|f_{\triangle,\varsigma}^{\mathrm b}(x)
-f_{\triangle,\varsigma}^{\mathrm b}(x')|
\le8|x-x'|^\gamma.
\]

To establish legality through the signed-binary conversion, first construct
deterministic record laws with uniform design, propensity \(1/2\), baseline
zero, and outcomes
\[
Y=0
\quad\text{under }P_{\mathrm{det},0}^{\mathrm b},
\qquad
Y=A\delta_{\triangle,n}^{\mathrm b}
f_{\triangle,\varsigma}^{\mathrm b}(X)
\quad\text{under }P_{\mathrm{det},\varsigma}^{\mathrm b}.
\]
Their effects are respectively zero and
\(\delta_{\triangle,n}^{\mathrm b}f_{\triangle,\varsigma}^{\mathrm b}\).
The displayed geometric estimates give
\[
\left\|
\delta_{\triangle,n}^{\mathrm b}f_{\triangle,\varsigma}^{\mathrm b}
\right\|_\infty\le\frac1{16},
\]
\[
\left|
\delta_{\triangle,n}^{\mathrm b}f_{\triangle,\varsigma}^{\mathrm b}(x)
-\delta_{\triangle,n}^{\mathrm b}f_{\triangle,\varsigma}^{\mathrm b}(x')
\right|
\le\frac12|x-x'|^\gamma.
\]
The control second moments are zero, and the treated second moments satisfy
\[
\bigl|\delta_{\triangle,n}^{\mathrm b}
f_{\triangle,\varsigma}^{\mathrm b}(x)\bigr|^2
\le\frac1{16^2}\le10.
\]
Together with the constant propensity and zero baseline, these inequalities
put both deterministic laws in
\(\mathcal M_{(2,\alpha,\beta,\gamma)}\); the first has constant zero effect.

Apply the conversion in \cref{obj:prop:bounded-submodel-comparison}. It
preserves the design, propensity, means, effect, and null constancy. The
resulting laws \(P_0^{\mathrm b}\) and \(P_\varsigma^{\mathrm b}\) have kernels
\[
Q_{a,P_0^{\mathrm b}}(\,\cdot\mid x)
=\frac12\delta_1+\frac12\delta_{-1}
\qquad(a\in\{0,1\}),
\]
\[
Q_{0,P_\varsigma^{\mathrm b}}(\,\cdot\mid x)
=\frac12\delta_1+\frac12\delta_{-1},
\]
\[
Q_{1,P_\varsigma^{\mathrm b}}(\,\cdot\mid x)
=
\frac{1+\delta_{\triangle,n}^{\mathrm b}
f_{\triangle,\varsigma}^{\mathrm b}(x)}2\delta_1
+
\frac{1-\delta_{\triangle,n}^{\mathrm b}
f_{\triangle,\varsigma}^{\mathrm b}(x)}2\delta_{-1}.
\]
Their propensity is identically \(1/2\), and
\[
P_0^{\mathrm b}\in H_0^{\mathrm{bin}}(w),
\qquad
P_\varsigma^{\mathrm b}\in\mathcal M_w^{\mathrm{bin}},
\qquad
\tau_{P_\varsigma^{\mathrm b}}
=\delta_{\triangle,n}^{\mathrm b}f_{\triangle,\varsigma}^{\mathrm b}.
\]
Define the binary prior by
\[
\pi_1^{\mathrm b}
=
2^{-m_{\triangle,n}^{\mathrm b}}
\sum_{\varsigma\in\{-1,1\}^{\mathscr J_{\triangle,n}^{\mathrm b}}}
\delta_{P_\varsigma^{\mathrm b}}.
\]
Again,
\[
2^{-m_{\triangle,n}^{\mathrm b}}>0,
\qquad
\sum_\varsigma2^{-m_{\triangle,n}^{\mathrm b}}=1,
\]
and the constant sign vector gives a positive-weight index.
% lean: CausalSmith.Stat.FinitepHomogeneityDensegamma.binaryTentLaw_inBinaryModel

\item \emph{Signed-binary separation.}
The equal-area cancellation and the disjoint-support square calculation give
\[
\int f_{\triangle,\varsigma}^{\mathrm b}\,d\lambda=0,
\qquad
\int(f_{\triangle,\varsigma}^{\mathrm b})^2\,d\lambda=\frac13.
\]
Consequently,
\[
d(P_\varsigma^{\mathrm b})
=\frac{\kappa_\triangle(h_{\triangle,n}^{\mathrm b})^\gamma}{\sqrt3}.
\]
Using the upper rank bound and
\(\gamma\vartheta_\triangle^{\mathrm b}=2\gamma/(4\gamma+1)\),
\[
\begin{aligned}
d(P_\varsigma^{\mathrm b})
&\ge
\frac{\kappa_\triangle}{\sqrt3}
(4n^{\vartheta_\triangle^{\mathrm b}})^{-\gamma}\\
&=
c_w^{\mathrm b}n^{-2\gamma/(4\gamma+1)}.
\end{aligned}
\]
The effect envelope also gives
\[
\begin{aligned}
d(P_\varsigma^{\mathrm b})^2
&=
\int\tau_{P_\varsigma^{\mathrm b}}^2\,d\lambda
-\left(\int\tau_{P_\varsigma^{\mathrm b}}\,d\lambda\right)^2\\
&\le\frac1{16^2},
\end{aligned}
\qquad
d(P_\varsigma^{\mathrm b})\le\frac1{16}<d_0.
\]
% lean: CausalSmith.Stat.FinitepHomogeneityDensegamma.binaryTentLaw_true_distance_lt_d0

\item \emph{Signed-binary complete-record comparison.}
Define, on the whole record space,
\[
\ell_{\triangle,\varsigma}^{\mathrm b}(x,a,y)=
\begin{cases}
1+\delta_{\triangle,n}^{\mathrm b}
f_{\triangle,\varsigma}^{\mathrm b}(x),&a=1,\ y=1,\\
1-\delta_{\triangle,n}^{\mathrm b}
f_{\triangle,\varsigma}^{\mathrm b}(x),&a=1,\ y=-1,\\
1,&\text{otherwise}.
\end{cases}
\]
The kernel probabilities imply
\[
P_\varsigma^{\mathrm b}\ll P_0^{\mathrm b},
\qquad
\frac{dP_\varsigma^{\mathrm b}}{dP_0^{\mathrm b}}
=\ell_{\triangle,\varsigma}^{\mathrm b}
\quad P_0^{\mathrm b}\text{-almost surely}.
\]
Under the null, each treatment--sign combination has conditional probability
\(1/4\). Summing the two control terms and the two treated terms yields
\[
\begin{aligned}
&\frac12+\frac14
\left[
(1+\delta_{\triangle,n}^{\mathrm b}f_{\triangle,\varsigma}^{\mathrm b}(x))
(1+\delta_{\triangle,n}^{\mathrm b}f_{\triangle,\varsigma'}^{\mathrm b}(x))
\right.\\[-2pt]
&\hspace{3cm}\left.
+
(1-\delta_{\triangle,n}^{\mathrm b}f_{\triangle,\varsigma}^{\mathrm b}(x))
(1-\delta_{\triangle,n}^{\mathrm b}f_{\triangle,\varsigma'}^{\mathrm b}(x))
\right]\\
&\qquad=
1+\frac{(\delta_{\triangle,n}^{\mathrm b})^2}{2}
f_{\triangle,\varsigma}^{\mathrm b}(x)
f_{\triangle,\varsigma'}^{\mathrm b}(x).
\end{aligned}
\]
Define
\[
\varepsilon_{\triangle,n}^{\mathrm b}
=(h_{\triangle,n}^{\mathrm b})^{2\gamma},
\qquad
\varrho_{\triangle,n}^{\mathrm b}
=\frac{\varepsilon_{\triangle,n}^{\mathrm b}\kappa_\triangle^2}{6}.
\]
Then
\[
(\delta_{\triangle,n}^{\mathrm b})^2
=\kappa_\triangle^2\varepsilon_{\triangle,n}^{\mathrm b},
\]
and integration of the paired squared tents gives
\[
\int f_{\triangle,\varsigma}^{\mathrm b}
f_{\triangle,\varsigma'}^{\mathrm b}\,d\lambda
=
\frac{2h_{\triangle,n}^{\mathrm b}}3
\sum_{j_\triangle}
\varsigma_{j_\triangle}\varsigma'_{j_\triangle}.
\]
Thus the complete one-record overlap is
\[
\int
\ell_{\triangle,\varsigma}^{\mathrm b}
\ell_{\triangle,\varsigma'}^{\mathrm b}\,dP_0^{\mathrm b}
=
1+\frac{\varrho_{\triangle,n}^{\mathrm b}}{m_{\triangle,n}^{\mathrm b}}
\sum_{j_\triangle}\varsigma_{j_\triangle}\varsigma'_{j_\triangle}.
\]
These bases are nonnegative, since
\[
0\le\varrho_{\triangle,n}^{\mathrm b}\le1,
\qquad
\left|\sum_{j_\triangle}
\varsigma_{j_\triangle}\varsigma'_{j_\triangle}\right|
\le m_{\triangle,n}^{\mathrm b}.
\]

Define the mixture and its density by
\[
\overline P_{1,\mathrm b}^{(n)}
=
2^{-m_{\triangle,n}^{\mathrm b}}
\sum_\varsigma(P_\varsigma^{\mathrm b})^{\otimes n},
\]
\[
G_\triangle^{\mathrm b}(\mathcal D_n)
=
\frac{d\overline P_{1,\mathrm b}^{(n)}}
 {d(P_0^{\mathrm b})^{\otimes n}}(\mathcal D_n)
=
2^{-m_{\triangle,n}^{\mathrm b}}
\sum_\varsigma\prod_{i_{\rm rec}=1}^n
\ell_{\triangle,\varsigma}^{\mathrm b}
(X_{i_{\rm rec}},A_{i_{\rm rec}},Y_{i_{\rm rec}}).
\]
Expansion and independence give
\[
\begin{aligned}
1+\chi^2\!\left(
\overline P_{1,\mathrm b}^{(n)}
\Vert(P_0^{\mathrm b})^{\otimes n}\right)
={}&2^{-2m_{\triangle,n}^{\mathrm b}}
\sum_\varsigma\sum_{\varsigma'}
\left(
1+\frac{\varrho_{\triangle,n}^{\mathrm b}}{m_{\triangle,n}^{\mathrm b}}
\sum_{j_\triangle}
\varsigma_{j_\triangle}\varsigma'_{j_\triangle}
\right)^n.
\end{aligned}
\]
For fixed \(\varsigma\), the exponential comparison and sign factorization
already established give
\[
\begin{aligned}
&2^{-m_{\triangle,n}^{\mathrm b}}\sum_{\varsigma'}
\left(
1+\frac{\varrho_{\triangle,n}^{\mathrm b}}{m_{\triangle,n}^{\mathrm b}}
\sum_{j_\triangle}
\varsigma_{j_\triangle}\varsigma'_{j_\triangle}
\right)^n\\
&\quad\le
\prod_{j_\triangle}
\cosh\left(
\frac{n\varrho_{\triangle,n}^{\mathrm b}}
 {m_{\triangle,n}^{\mathrm b}}\varsigma_{j_\triangle}
\right)
\le
\exp\left(
\frac{(n\varrho_{\triangle,n}^{\mathrm b})^2}
 {2m_{\triangle,n}^{\mathrm b}}\right).
\end{aligned}
\]
Averaging over the remaining signs and simplifying the exponent yields
\[
1+\chi^2\!\left(
\overline P_{1,\mathrm b}^{(n)}
\Vert(P_0^{\mathrm b})^{\otimes n}\right)
\le
\exp\left(
\frac{\kappa_\triangle^4}{36}
n^2(h_{\triangle,n}^{\mathrm b})^{4\gamma+1}
\right).
\]
Since
\[
h_{\triangle,n}^{\mathrm b}
\le n^{-\vartheta_\triangle^{\mathrm b}},
\qquad
\vartheta_\triangle^{\mathrm b}(4\gamma+1)=2,
\]
we have
\[
n^2(h_{\triangle,n}^{\mathrm b})^{4\gamma+1}\le1,
\]
and therefore
\[
1+\chi^2\!\left(
\overline P_{1,\mathrm b}^{(n)}
\Vert(P_0^{\mathrm b})^{\otimes n}\right)
\le e^{\kappa_\triangle^4/36}<2.
\]
Applying the zero-integral density-deviation identity and
Cauchy--Schwarz to \(G_\triangle^{\mathrm b}\) gives
\[
\begin{aligned}
\operatorname{TV}\!\left(
(P_0^{\mathrm b})^{\otimes n},\overline P_{1,\mathrm b}^{(n)}
\right)
&=
\frac12\int|G_\triangle^{\mathrm b}-1|\,d(P_0^{\mathrm b})^{\otimes n}\\
&\le
\frac12\sqrt{\chi^2\!\left(
\overline P_{1,\mathrm b}^{(n)}
\Vert(P_0^{\mathrm b})^{\otimes n}\right)}
<\frac12.
\end{aligned}
\]
% lean: CausalSmith.Stat.FinitepHomogeneityDensegamma.binaryTentMixture_tv_lt_half

\item \emph{Bounded-model placement and risk.}
The inclusion supplied by \cref{obj:prop:bounded-submodel-comparison} gives
\[
P_\varsigma^{\mathrm b}\in\mathcal M_w^{\mathrm b},
\qquad
P_0^{\mathrm b}\in H_0^{\mathrm b}(w).
\]
For the null, this uses its bounded-model membership together with the
preserved constant zero effect. The same comparison, together with
\cref{obj:lem:causal-nonempty} at moment order two, supplies
\[
D_w^{\mathrm b}=D_{(2,\alpha,\beta,\gamma)}\ge d_0.
\]
Define
\[
r_\triangle^{\mathrm b}
=c_w^{\mathrm b}n^{-2\gamma/(4\gamma+1)}.
\]
The binary alternatives satisfy
\[
0<r_\triangle^{\mathrm b}
\le d(P_\varsigma^{\mathrm b})<d_0\le D_w^{\mathrm b}.
\]

For any test satisfying the bounded-null size constraint in
\cref{obj:def:bounded-risk}, average its public seed as above. The
complete-record total-variation bound gives
\[
\left|
2^{-m_{\triangle,n}^{\mathrm b}}\sum_\varsigma
\mathsf R_n(P_\varsigma^{\mathrm b},\phi)
-\mathsf R_n(P_0^{\mathrm b},\phi)
\right|<\frac12,
\]
hence
\[
2^{-m_{\triangle,n}^{\mathrm b}}\sum_\varsigma
\mathsf R_n(P_\varsigma^{\mathrm b},\phi)<\frac35.
\]
If all positive-weight support laws had type-II error below \(2/5\), their
weighted rejection sum would be at least
\[
2^{-m_{\triangle,n}^{\mathrm b}}\sum_\varsigma\frac35=\frac35,
\]
a contradiction. Thus a support law has type-II error at least \(2/5\).
All support laws belong to the bounded model at separation
\(r_\triangle^{\mathrm b}\), so
\[
\sup_{\substack{P\in\mathcal M_w^{\mathrm b}\\
d(P)\ge r_\triangle^{\mathrm b}}}
\{1-\mathsf R_n(P,\phi)\}\ge\frac25.
\]
The constant-zero test is admissible, and taking the infimum over the
bounded-null level class proves
\[
B_n^{\mathrm b}(w,r_\triangle^{\mathrm b})\ge\frac25.
\]
The constructed signed-binary null, normalized positive-weight prior,
placement, and complete-record comparison establish the second witness
package.
% lean: CausalSmith.Stat.FinitepHomogeneityDensegamma.binary_tent_canonical_lower_receipt
\end{enumerate}
\end{proof}

The fixed propensity in \cref{obj:lem:paired-tent-full-record-lower} makes the comparison applicable to both observation regimes. Every alternative support law meets the separation requirement, while the mixture remains close to the null in the complete original-record experiment. The growing mark magnitudes coexist with a conditional \(p\)-moment bound of one. The signed-binary package supplies a separate bounded comparison at the moment-order-two exponent.

\subsection{Marked records and continuous coefficient fields}

The rough-phase construction uses a normalized marked table \leanref{sym:Cone}{\ensuremath{\mathsf T}} to specify treatment, outcome sign, and mark occurrence jointly. Keeping its six categories explicit preserves the assignment information and the zero outcomes in the likelihood comparison.

\begin{definitionv}{def:marked-table}[Six-category record table \(\mathsf T\)]
The conditional original-record table is
\[
\begin{aligned}
\mathsf T\bigl(A=(1+\ell)/2,Y=hL\mid X=x\bigr)
&=\frac{\varepsilon}{4}
\{1+\ell\xi(x)+h\upsilon(x)+\ell h\zeta(x)\},
&&(\ell,h)\in\{-1,1\}^2,\\
\mathsf T\bigl(A=(1+\ell)/2,Y=0\mid X=x\bigr)
&=\frac{1-\varepsilon}{2}\{1+\ell\xi(x)\},
&&\ell\in\{-1,1\}.
\end{aligned}
\]
Here \(\varepsilon\) is the rare-mark probability, \(L>0\) is the mark magnitude, and \(\xi,\upsilon,\zeta\) may be Borel functions of the covariate. The table is evaluated on inputs for which every displayed probability is nonnegative. Likelihood products include all six record categories.
\end{definitionv}

In \cref{obj:def:marked-table}, the treatment sign \leanref{sym:treatmentSign}{\ensuremath{\ell}} and marked-outcome sign \leanref{sym:markSign}{\ensuremath{h}} index the four nonzero-outcome categories. The coordinate \leanref{sym:xi}{\ensuremath{\xi}} governs treatment probabilities, \leanref{sym:upsilon}{\ensuremath{\upsilon}} governs marked-outcome signs, and \leanref{sym:zeta}{\ensuremath{\zeta}} governs their interaction. Summing over the categories gives total probability one and nonzero-outcome probability \leanref{sym:epsilon}{\ensuremath{\varepsilon}}; the positive mark magnitude is \leanref{sym:Lmark}{\ensuremath{L}}. The remaining two categories retain the treatment coordinate when the outcome is zero.

Continuous coefficient fields are formed from the finite-overlap frame \leanref{sym:Frame}{\ensuremath{(f_i)_{i=0}^{K}}} below, where \(K\) is the fine rank. Its adjacent-cell support gives a continuous construction on the fine partition.

\begin{definitionv}{def:frame-handle}[Finite-overlap frame \(f_i\)]
The finite-overlap frame at rank \(K=J_{\mathrm f}\), the fine histogram rank, consists of the functions
\[
f_i(x)=
\begin{cases}
\cos\!\bigl(\pi(Kx-i)/2\bigr),
&x\in[i/K,(i+1)/K],\\
\sin\!\bigl(\pi(Kx-(i-1))/2\bigr),
&i>0,\ x\in[(i-1)/K,i/K],\\
0,&\text{otherwise},
\end{cases}
\qquad x\in[0,1],\quad 0\le i\le K.
\]
For a coefficient vector
\[
a=(a_0,\ldots,a_K),
\]
the associated frame field is
\[
x\longmapsto\sum_{i=0}^{K}a_i f_i(x).
\]
\end{definitionv}

On each fine cell, the two active frame coordinates have squared values summing to one. The prior construction uses these squares to interpolate coarse paired tents, while the frame coordinates themselves carry the random propensity and outcome coefficients. The following algorithm specifies both branches through a finite collection of sign draws.

\begin{algorithmv}{def:copula-frame-priors}[Copula-frame priors \(\pi^{K,M}_{\nu;v,a,u,\varepsilon,L}\)]
Inputs are the public tuple \(v\), dyadic ranks \(K=J_{\mathrm f}\) and \(M=J_{\mathrm c}\) satisfying \(M\ge2\) and \(K\ge16M\), amplitudes \(0<a,u\le1/16\), a rare-mark probability \(0<\varepsilon<1\), and a mark magnitude \(L>0\).
\begin{enumerate}
\item Set the coupling amplitude and triangular tent to
\[
\kappa=\frac1{16},
\qquad
b_\triangle(t)=2\min\{t,1-t\},
\quad t\in[0,1].
\]
Draw independent uniform signs
\[
\Pr(\sigma_j=1)=\Pr(\sigma_j=-1)=\frac12,
\qquad 1\le j\le M/2.
\]
Define the opposite paired coarse tents by
\[
g_\sigma(x)=
\begin{cases}
\sigma_j b_\triangle(Mx-(2j-2)),
&x\in[(2j-2)/M,(2j-1)/M],\\
-\sigma_j b_\triangle(Mx-(2j-1)),
&x\in[(2j-1)/M,2j/M],
\end{cases}
\quad 1\le j\le M/2.
\]
Their values at coarse-cell boundaries are zero.
\item On each fine cell define the two nonzero frame coordinates by
\[
f_i(x)=\cos\!\bigl(\pi(Kx-i)/2\bigr),
\qquad
f_{i+1}(x)=\sin\!\bigl(\pi(Kx-i)/2\bigr),
\]
\[
x\in[i/K,(i+1)/K],
\qquad 0\le i<K.
\]
Set the other coordinates to zero and form the squared-frame interpolation
\[
\widetilde g_\sigma(x)
=
\sum_{i=0}^{K}g_\sigma(i/K)f_i(x)^2.
\]
\item For each branch \(\nu\in\{0,1\}\), draw coefficient pairs independently conditional on \(\sigma\), with
\[
\Pr_\nu(\lambda_i=l,\eta_i=h\mid\sigma)
=
\frac{1+\nu\kappa g_\sigma(i/K)lh}{4},
\qquad
l,h\in\{-1,1\},
\quad 0\le i\le K.
\]
\item For each coefficient draw, define the conditional-table coordinates
\[
\xi(x)=a\sum_{i=0}^{K}f_i(x)\lambda_i,
\qquad
\upsilon(x)=u\sum_{i=0}^{K}f_i(x)\eta_i,
\]
\[
t_\nu(x)=
-\frac{\nu au\kappa}{1-a^2}\widetilde g_\sigma(x),
\qquad
\zeta(x)=
\xi(x)\upsilon(x)+t_\nu(x)\bigl(1-\xi(x)^2\bigr).
\]
Let \(P\) be the original-record law with uniform covariate \(X\) and the normalized six-category marked conditional table evaluated at these coordinates, rare-mark probability \(\varepsilon\), and mark magnitude \(L\).
\item Output the finite prior obtained by pushing the coefficient distribution forward to original-record laws:
\[
\pi^{K,M}_{\nu;v,a,u,\varepsilon,L}
=
\mathcal L_\nu(P),
\qquad \nu\in\{0,1\}.
\]
\end{enumerate}
\end{algorithmv}

The copula-frame prior \leanref{sym:copulaPrior}{\ensuremath{\pi^{K,M}_{\nu;v,a,u,\varepsilon,L}}} in \cref{obj:def:copula-frame-priors} changes the dependence between the fine coefficient signs across branches. Its propensity-coordinate amplitude is \leanref{sym:aamp}{\ensuremath{a}}, and its marked-outcome-coordinate amplitude is \leanref{sym:uamp}{\ensuremath{u}}. Each coefficient draw determines one original-record law, so the pushforward is a finite probability prior even when several draws determine the same law.

For the model tuning, write \(b\) for the baseline amplitude, and denote the resulting null and alternative priors by \(\pi_0\) and \(\pi_1\). The next result gives their support placement under the phase and rank conditions, together with a separate bounded tuning.

\begin{lemmav}{lem:copula-frame-legality}[Copula-frame legality and moments]
Consider an admissible tuple \(v=(p,\alpha,\beta,\gamma)\) as stipulated in \cref{obj:def:model}. Put
\[
q=\frac{p-1}{p},\qquad S=\alpha+\beta,\qquad
F(p)=\frac{2\alpha}{p-1}+\frac{\beta}{q}+\frac{S}{2\gamma}.
\]
Suppose the following conditions hold:
\begin{itemize}
\item \textbf{(Phase.)} \(F(p)<1\).
\item \textbf{(Dyadic ranks.)} \(K\) and \(M\) are powers of two, with \(16M\le K\) and \(2\le M\).
\item \textbf{(Rank compatibility.)} \(M\le K^{S/\gamma}\).
\end{itemize}
Define the construction amplitudes by
\[
a=\frac{K^{-\alpha}}{16},\qquad b=\frac{K^{-\beta}}{256}.
\]
Let \(\lambda\) denote Lebesgue probability measure on \([0,1]\), and write \(Q_{a,P}(dy\mid x)\) for the conditional outcome kernel in treatment arm \(a\) under \(P\). Define the population effect distance by
\[
d(P)=
\left(
\int_0^1
\left|\tau_P(x)-\int_0^1\tau_P(z)\,\lambda(dz)\right|^2
\,\lambda(dx)
\right)^{1/2},
\]
and let \(d_0\) be the fixed positive distance supplied by the legal model witness.

Use the copula-frame priors of \cref{obj:def:copula-frame-priors} with
\[
u=\frac1{16},\qquad
\varepsilon=(16b)^{1/q},\qquad L=\varepsilon^{-1/p},
\qquad
\pi_0=\pi^{K,M}_{0;v,a,u,\varepsilon,L},\qquad
\pi_1=\pi^{K,M}_{1;v,a,u,\varepsilon,L}.
\]
For every realization of the \(M/2\) coarse signs and the \(K+1\) fine coefficient-sign pairs, the null-branch law belongs to \(H_0(v)\) of \cref{obj:def:null}, and the alternative-branch law belongs to \(\mathcal M_v\) of \cref{obj:def:model}. Each alternative-branch law satisfies
\[
\tau_P(x)=-\frac{2ab\kappa}{1-a^2}\widetilde g_\sigma(x)
\quad\text{for every }x\in[0,1],
\qquad
\int_0^1\tau_P(x)\,\lambda(dx)=0,
\]
\[
2^{-17}K^{-S}\le d(P)\le 2^{-14}K^{-S},
\qquad d(P)<d_0.
\]
Every law in either branch satisfies
\[
\int_{\mathbb R}|y|^p\,Q_{a,P}(dy\mid x)=1
\quad\text{for every }a\in\{0,1\}\text{ and }x\in[0,1].
\]

For the bounded construction, let \(w=(\alpha,\beta,\gamma)\) be an admissible tuple as stipulated in \cref{obj:def:bounded-model}, set \(v=(2,\alpha,\beta,\gamma)\), and impose the same phase, dyadic-rank, and rank-compatibility conditions with \(p=2\). Retain the displayed amplitudes \(a,b\), and use the priors of \cref{obj:def:copula-frame-priors} with
\[
u=2b,\qquad \varepsilon=\frac12,\qquad L=1,
\qquad
\pi_0=\pi^{K,M}_{0;v,a,u,\varepsilon,L},\qquad
\pi_1=\pi^{K,M}_{1;v,a,u,\varepsilon,L}.
\]
For every realization of the coarse signs and fine coefficient-sign pairs, the null-branch law belongs to \(H_0^{\mathrm b}(w)\) of \cref{obj:def:bounded-null}, and the alternative-branch law belongs to \(\mathcal M^{\mathrm b}_w\) of \cref{obj:def:bounded-model}. Each alternative-branch law satisfies the displayed effect identity, centering identity, and distance bounds. Every law in either bounded branch satisfies
\[
\int_{\mathbb R}|y|^2\,Q_{a,P}(dy\mid x)=\frac12
\quad\text{for every }a\in\{0,1\}\text{ and }x\in[0,1].
\]
\end{lemmav}

\begin{proof}[Proof of \cref{obj:lem:copula-frame-legality}]
Fix an admissible tuple and ranks satisfying the hypotheses, and fix an arbitrary realization of all coefficient signs. Admissibility gives
\[
1<p\le2,\qquad
0<\alpha,\beta\le1,\qquad
\frac14\le\gamma\le1.
\]
We first treat the finite-moment construction.

\begin{enumerate}
\item \textit{Tuning and construction inputs.}
The rank assumptions imply \(K\ge32\), so the displayed amplitudes satisfy
\[
0<a\le\frac1{16},\qquad
0<b\le\frac1{256},\qquad
aK^\alpha=\frac1{16},\qquad
bK^\beta=\frac1{256}.
\]
Indeed, \(K^{-\alpha},K^{-\beta}\in(0,1]\), and multiplication by the corresponding positive powers of \(K\) cancels the exponents. Also \(q>0\) and \(0<16b\le1/16<1\), whence
\[
0<\varepsilon=(16b)^{1/q}<1,\qquad
L=\varepsilon^{-1/p}>0.
\]
The tuning identities are
\[
\varepsilon L^p
=\varepsilon\varepsilon^{-1}=1,
\qquad
\varepsilon Lu
=\frac{\varepsilon^{\,1-1/p}}{16}
=\frac{\varepsilon^q}{16}
=b.
\]
Thus the finite-moment tuning meets the input conditions of
\cref{obj:def:copula-frame-priors}. The second tuning also meets those conditions, since
\[
0<2b\le\frac1{128}\le\frac1{16},
\qquad
0<\frac12<1,\qquad 1>0.
\]
For that tuning,
\[
\varepsilon Lu=\frac12\cdot1\cdot2b=b,
\qquad
\varepsilon L^2=\frac12.
\]
% lean: CausalSmith.Stat.FinitepHomogeneityDensegamma.legality_copula_domains

\item \textit{Validity of the conditional table and exact moments.}
Use the coordinates and signs of \cref{obj:def:copula-frame-priors}, and define the two coefficient fields by
\[
\mathcal F_\lambda(x)=\sum_{i=0}^{K}\lambda_i f_i(x),
\qquad
\mathcal F_\eta(x)=\sum_{i=0}^{K}\eta_i f_i(x),
\qquad x\in[0,1].
\]
On a fine cell, write
\[
\theta_i(x)=\frac{\pi(Kx-i)}2,
\qquad
0\le i<K,\quad x\in[i/K,(i+1)/K].
\]
There the only possibly nonzero coordinates are
\(f_i(x)=\cos\theta_i(x)\) and
\(f_{i+1}(x)=\sin\theta_i(x)\). Consequently,
\[
\sum_{i=0}^{K}f_i(x)^2=1,\qquad
|\mathcal F_\lambda(x)|,\ |\mathcal F_\eta(x)|\le2.
\]
The coarse tents have absolute value at most one. Hence their squared-frame interpolation satisfies
\[
|\widetilde g_\sigma(x)|
\le\sum_{i=0}^{K}|g_\sigma(i/K)|f_i(x)^2
\le\sum_{i=0}^{K}f_i(x)^2=1.
\]
For either valid tuning and either branch,
\[
|\xi(x)|\le2a\le\frac18,\qquad
|\upsilon(x)|\le2u\le\frac18,
\qquad
a^2\le\frac1{256},\qquad au\le\frac1{256}.
\]
In the null branch \(t_0=0\). In the alternative branch, positivity of \(1-a^2\) and the preceding bounds give
\[
|t_1(x)|
\le\frac{au\kappa}{1-a^2}
\le\frac1{128},
\]
where the last inequality follows from
\(au\le1/256\), \(\kappa=1/16\), and
\(1-a^2\ge255/256\). Since
\(0\le\xi(x)^2\le1/64\), we also have
\[
|1-\xi(x)^2|\le1,
\qquad
|\zeta(x)|
\le|\xi(x)\upsilon(x)|
   +|t_\nu(x)|\,|1-\xi(x)^2|
\le\frac1{64}+\frac1{128}\le\frac18.
\]
Thus every unmarked and marked bracket in
\cref{obj:def:marked-table} is positive:
\[
1+\ell\xi(x)\ge\frac78,
\qquad
1+\ell\xi(x)+h\upsilon(x)+\ell h\zeta(x)
\ge\frac58,
\qquad \ell,h\in\{-1,1\}.
\]
The frame coordinates are continuous at their cell endpoints, and all three table coordinates are continuous finite expressions in them. The table therefore supplies Borel conditional probability kernels.

Summing the marked probabilities over the outcome sign gives
\[
\sum_{h\in\{-1,1\}}
\frac{\varepsilon}{4}
\{1+\ell\xi(x)+h\upsilon(x)+\ell h\zeta(x)\}
=\frac{\varepsilon}{2}\{1+\ell\xi(x)\}.
\]
Adding the unmarked probability gives
\((1+\ell\xi(x))/2\), and summing over \(\ell\) gives one. In particular, the conditional probability of a mark within either treatment arm is exactly \(\varepsilon\). The marked outcomes have absolute value \(L\), and the unmarked outcome is zero. Therefore
\[
\int_{\mathbb R}|y|^{r_{\mathrm{mom}}}
\,Q_{(1+\ell)/2,P}(dy\mid x)
=\varepsilon L^{r_{\mathrm{mom}}},
\qquad
r_{\mathrm{mom}}\in(0,\infty),\
\ell\in\{-1,1\},\ x\in[0,1].
\]
The tuning identities from the preceding step give the required value \(1\) at order \(p\), and the value \(1/2\) at order two for the bounded tuning.

Summing the signed marked outcomes within each arm also gives
\[
\begin{aligned}
e_P(x)&=\frac{1+\xi(x)}2,\\
m_{0,P}(x)
&=\varepsilon L\frac{\upsilon(x)-\zeta(x)}{1-\xi(x)}
 =\varepsilon L\{\upsilon(x)-t_\nu(x)(1+\xi(x))\},\\
m_{1,P}(x)
&=\varepsilon L\frac{\upsilon(x)+\zeta(x)}{1+\xi(x)}
 =\varepsilon L\{\upsilon(x)+t_\nu(x)(1-\xi(x))\},\\
\tau_P(x)&=2\varepsilon L t_\nu(x).
\end{aligned}
\]
All denominators are positive because \(|\xi(x)|\le1/8\).
% lean: CausalSmith.Stat.FinitepHomogeneityDensegamma.legality_arm_moments

\item \textit{Null membership.}
We record the frame estimates needed for the primitive restrictions. For either coefficient field,
\[
|\mathcal F_\lambda(x)|,\ |\mathcal F_\eta(x)|
\le\sqrt2.
\]
Indeed, on a fine cell its absolute value is at most
\(|\cos\theta_i(x)|+|\sin\theta_i(x)|\), whose square is at most two by
\(\cos^2\theta_i(x)+\sin^2\theta_i(x)=1\) and
\((|\cos\theta_i(x)|-|\sin\theta_i(x)|)^2\ge0\).

For either field, denoted below by
\(\mathcal F\in\{\mathcal F_\lambda,\mathcal F_\eta\}\), the sine and cosine Lipschitz inequalities give, when \(x,z\) belong to the same fine cell,
\[
|\mathcal F(x)-\mathcal F(z)|
\le2|\theta_i(x)-\theta_i(z)|
=\pi K|x-z|.
\]
If \(K|x-z|<1\), the points lie in the same cell or in adjacent cells. In the latter case, split at the common endpoint and add the two cellwise bounds; the two distances sum to \(|x-z|\). Equality of the points gives zero, and reversing their order leaves the bound unchanged. Thus
\[
|\mathcal F(x)-\mathcal F(z)|
\le\pi K|x-z|
\qquad\text{when }K|x-z|<1.
\]
For \(0<s_{\mathrm H}\le1\), the near case uses
\(K|x-z|\le(K|x-z|)^{s_{\mathrm H}}\) and \(\pi<4\); the complementary case uses
\(|\mathcal F(x)-\mathcal F(z)|\le2\sqrt2\le4\) and
\((K|x-z|)^{s_{\mathrm H}}\ge1\). Together these give
\[
|\mathcal F(x)-\mathcal F(z)|
\le4K^{s_{\mathrm H}}|x-z|^{s_{\mathrm H}},
\qquad
0<s_{\mathrm H}\le1,\quad x,z\in[0,1].
\]

The membership calculation applies to valid construction inputs satisfying
\[
\varepsilon Lu=b,\qquad \varepsilon L^p\le10.
\]
In the null branch, the table identities reduce to
\[
e_P(x)=\frac{1+a\mathcal F_\lambda(x)}2,
\qquad
m_{0,P}(x)=b\mathcal F_\eta(x),
\qquad
\tau_P(x)=0.
\]
The frame envelope and amplitude bounds imply
\[
|\xi(x)|\le a\sqrt2\le\frac18,\qquad
\frac14\le e_P(x)\le\frac34,\qquad
|m_{0,P}(x)|\le b\sqrt2\le\frac12.
\]
The frame Hölder estimates, together with the amplitude-power identities, give
\[
\begin{aligned}
|e_P(x)-e_P(z)|
&\le2aK^\alpha|x-z|^\alpha
=\frac18|x-z|^\alpha
\le20|x-z|^\alpha,\\
|m_{0,P}(x)-m_{0,P}(z)|
&\le4bK^\beta|x-z|^\beta
=\frac1{64}|x-z|^\beta
\le20|x-z|^\beta.
\end{aligned}
\]
The effect is continuous, has supremum norm zero, and has zero Hölder differences. These estimates establish
\cref{obj:ass:overlap,obj:ass:propensity-smoothness,obj:ass:baseline-smoothness,obj:ass:effect-smoothness,obj:ass:baseline-cap,obj:ass:effect-cap}.
The normalized conditional table is integrated against \(\lambda\), so
\(P_X=\lambda\), as required by \cref{obj:ass:uniform}. The exact conditional-moment formula yields
\[
\int_{\mathbb R}|y|^p Q_{(1+\ell)/2,P}(dy\mid x)
=\varepsilon L^p\le10,
\]
and hence \cref{obj:ass:raw-moment}. Finally, \(\tau_P=0\) supplies
\cref{obj:ass:null-constancy} with constant zero. Therefore
\(P\in H_0(v)\) by \cref{obj:def:null}. The finite-moment tuning satisfies the two displayed input conditions by the first step.
% lean: CausalSmith.Stat.FinitepHomogeneityDensegamma.legality_null_inNull

\item \textit{Alternative membership.}
Define the positive correction amplitude by
\[
C_{\mathrm{cop}}=\frac{ab\kappa}{1-a^2}>0.
\]
The amplitude bounds give
\[
1-a^2\ge\frac12,\qquad
C_{\mathrm{cop}}\le2ab\kappa\le\frac1{32768}.
\]
To control its smoothness factors, first observe that \(p-1\le1\) and \(q\le1/2\). Consequently,
\[
2\alpha\le\frac{2\alpha}{p-1},
\qquad
2\beta\le\frac{\beta}{q},
\qquad
2S+\frac{S}{2\gamma}\le F(p)<1.
\]
If \(S\ge\gamma\), the left side would be at least
\(2\gamma+1/2\ge1\). Thus
\[
S<\gamma,\qquad \beta\le\gamma.
\]
Rank compatibility and positivity of \(\gamma\) give
\(M^\gamma\le K^S\). Using \(M\ge1\) and the amplitude identities, we obtain
\[
\begin{aligned}
ab&=\frac{K^{-S}}{4096},\\
abM^\gamma&\le\frac{K^{-S}K^S}{4096}=\frac1{4096},\\
abM^\beta&\le abM^\gamma\le\frac1{4096},\\
a^2bK^\beta&=\frac{a^2}{256}\le\frac1{65536}.
\end{aligned}
\]
Multiplication by \(2\kappa=1/8\) now gives
\[
C_{\mathrm{cop}}M^\gamma\le\frac1{32768},
\qquad
C_{\mathrm{cop}}M^\beta\le\frac1{32768},
\qquad
C_{\mathrm{cop}}aK^\beta\le\frac1{524288}.
\]
The frame estimates also give
\[
|\xi(x)|\le\frac18,\qquad
|1+\xi(x)|\le2,\qquad
|\xi(x)-\xi(z)|
\le4aK^\beta|x-z|^\beta.
\]

For valid inputs with \(\varepsilon Lu=b\), the alternative table identities become
\[
m_{0,P}(x)
=b\mathcal F_\eta(x)
 +C_{\mathrm{cop}}\widetilde g_\sigma(x)(1+\xi(x)),
\qquad
\tau_P(x)=-2C_{\mathrm{cop}}\widetilde g_\sigma(x).
\]
Their uniform bounds are
\[
|m_{0,P}(x)|
\le2b+2C_{\mathrm{cop}}\le\frac12,
\qquad
|\tau_P(x)|\le2C_{\mathrm{cop}}\le\frac12.
\]
The design remains \(\lambda\), because the alternative conditional table is normalized. Its propensity is the same function
\((1+a\mathcal F_\lambda)/2\) as in the null branch, so the preceding overlap and propensity-smoothness estimates apply.

For the remaining smoothness estimates, extend the triangular tent by
\[
b_{\triangle,\mathrm{ext}}(r_\triangle)
=\max\{0,\,2\min(r_\triangle,1-r_\triangle)\},
\qquad r_\triangle\in\mathbb R.
\]
This agrees with \(b_\triangle\) on \([0,1]\) and vanishes outside that interval. The elementary Lipschitz inequalities for minimum and maximum give
\[
|b_{\triangle,\mathrm{ext}}(r_\triangle)
 -b_{\triangle,\mathrm{ext}}(s_\triangle)|
\le2|r_\triangle-s_\triangle|,
\qquad r_\triangle,s_\triangle\in\mathbb R.
\]
Define each signed pair contribution and its active index set by
\[
\begin{aligned}
\mathfrak g_j(x)
&=\sigma_j\bigl\{
b_{\triangle,\mathrm{ext}}(Mx-(2j-2))
-b_{\triangle,\mathrm{ext}}(Mx-(2j-1))
\bigr\},\\
\mathcal J_\sigma(x)
&=\{j\in\{1,\ldots,M/2\}:\mathfrak g_j(x)\ne0\},
\qquad x\in[0,1].
\end{aligned}
\]
Thus \(g_\sigma=\sum_{j=1}^{M/2}\mathfrak g_j\).
Each contribution satisfies
\[
|\mathfrak g_j(x)-\mathfrak g_j(z)|
\le4M|x-z|.
\]
The interiors of the pair supports are disjoint, so
\(|\mathcal J_\sigma(x)|\le1\). Only indices in
\(\mathcal J_\sigma(x)\cup\mathcal J_\sigma(z)\) contribute to a difference, and this union has cardinality at most two. Hence
\[
|g_\sigma(x)-g_\sigma(z)|\le8M|x-z|.
\]

On a fine cell, the squared-frame formula is
\[
\widetilde g_\sigma(x)
=g_\sigma(i/K)\cos^2\theta_i(x)
 +g_\sigma((i+1)/K)\sin^2\theta_i(x).
\]
For \(x,z\) in that cell, subtraction and
\(\cos^2+\sin^2=1\) give
\[
\begin{aligned}
\widetilde g_\sigma(x)-\widetilde g_\sigma(z)
={}&\bigl(g_\sigma((i+1)/K)-g_\sigma(i/K)\bigr)\\
&\quad{}\times
\bigl(\sin\theta_i(x)-\sin\theta_i(z)\bigr)
\bigl(\sin\theta_i(x)+\sin\theta_i(z)\bigr).
\end{aligned}
\]
The sine Lipschitz inequality and its unit envelope therefore imply
\[
|\widetilde g_\sigma(x)-\widetilde g_\sigma(z)|
\le\frac{8M}{K}\,\pi K|x-z|
\le32M|x-z|.
\]
Adjacent cell formulas agree at their common endpoint. Splitting any ordered interval at the intervening fine-grid endpoints and summing the cellwise estimates proves this bound for all \(x,z\in[0,1]\); equal points give zero, and reversing their order preserves the bound.

If \(M|x-z|\le1\), use
\(M|x-z|\le(M|x-z|)^{s_{\mathrm H}}\). If \(M|x-z|>1\), use the envelope
\(|\widetilde g_\sigma(x)-\widetilde g_\sigma(z)|\le2\). Thus
\[
|\widetilde g_\sigma(x)-\widetilde g_\sigma(z)|
\le32M^{s_{\mathrm H}}|x-z|^{s_{\mathrm H}},
\qquad 0<s_{\mathrm H}\le1.
\]
The same inequality holds for \(s_{\mathrm H}=0\), with the convention
\(r^0=1\) for \(r\ge0\), since the left side is at most two.

The baseline product difference decomposes as
\[
\begin{aligned}
m_{0,P}(x)-m_{0,P}(z)
={}&b\{\mathcal F_\eta(x)-\mathcal F_\eta(z)\}\\
&+C_{\mathrm{cop}}\bigl[
 \{\widetilde g_\sigma(x)-\widetilde g_\sigma(z)\}(1+\xi(x))
 +\widetilde g_\sigma(z)\{\xi(x)-\xi(z)\}
\bigr].
\end{aligned}
\]
Applying the displayed estimates to these three terms gives
\[
\begin{aligned}
|m_{0,P}(x)-m_{0,P}(z)|
&\le
\bigl(4bK^\beta
 +64C_{\mathrm{cop}}M^\beta
 +4C_{\mathrm{cop}}aK^\beta\bigr)|x-z|^\beta\\
&\le
\left(\frac4{256}+\frac{64}{32768}
                  +\frac4{524288}\right)|x-z|^\beta\\
&\le20|x-z|^\beta.
\end{aligned}
\]
Similarly,
\[
\begin{aligned}
|\tau_P(x)-\tau_P(z)|
&\le64C_{\mathrm{cop}}M^\gamma|x-z|^\gamma\\
&\le\frac{64}{32768}|x-z|^\gamma
\le20|x-z|^\gamma.
\end{aligned}
\]
All these representatives are continuous. Their uniform bounds and Hölder estimates establish the baseline and effect restrictions in
\cref{obj:ass:baseline-smoothness,obj:ass:effect-smoothness,obj:ass:baseline-cap,obj:ass:effect-cap}.
The conditional-moment calculation supplies
\cref{obj:ass:raw-moment} whenever \(\varepsilon L^p\le10\).
Together with the design, overlap, and propensity restrictions already established, this proves
\(P\in\mathcal M_v\) by \cref{obj:def:model}. In particular it applies to the finite-moment tuning.
% lean: CausalSmith.Stat.FinitepHomogeneityDensegamma.legality_alternative_inModel

\item \textit{Centering of the interpolated tents.}
Choose the dyadic exponents and the refinement factor so that
\[
K=2^{k_{\mathrm{rank}}},\qquad
M=2^{m_{\mathrm{rank}}},
\qquad
k_{\mathrm{rank}},m_{\mathrm{rank}}\in\mathbb N,
\qquad
1\le m_{\mathrm{rank}}\le k_{\mathrm{rank}},
\]
\[
r_{\mathrm{grid}}=2^{k_{\mathrm{rank}}-m_{\mathrm{rank}}}
=\frac KM\in\mathbb N,\qquad r_{\mathrm{grid}}\ge16.
\]
The inequality between the dyadic exponents follows from \(M\le K\), and \(M\ge2\) implies that \(M\) is even.

Every coarse cell has the same fine-grid sample sum of its unsigned tent. More precisely, support of the extended tent and the change of index
\(i=i_{\mathrm c}r_{\mathrm{grid}}+h_{\mathrm{grid}}\) give
\[
\sum_{i=0}^{K}
 b_{\triangle,\mathrm{ext}}(Mi/K-i_{\mathrm c})
=
\sum_{h_{\mathrm{grid}}=0}^{r_{\mathrm{grid}}}
 b_\triangle(h_{\mathrm{grid}}/r_{\mathrm{grid}}),
\qquad
i_{\mathrm c}\in\{0,\ldots,M-1\}.
\]
The two cells in each signed pair therefore have equal sample sums and opposite signs. Consequently,
\[
\sum_{i=0}^{K}g_\sigma(i/K)=0,
\qquad g_\sigma(0)=g_\sigma(1)=0.
\]
In particular,
\[
\sum_{i=0}^{K-1}g_\sigma(i/K)=0,
\qquad
\sum_{i=0}^{K-1}g_\sigma((i+1)/K)=0.
\]

Integrating the sine and cosine squares over one fine cell, by an affine change of variable, gives
\[
\int_{i/K}^{(i+1)/K}\cos^2\theta_i(x)\,dx
=
\int_{i/K}^{(i+1)/K}\sin^2\theta_i(x)\,dx
=\frac1{2K}.
\]
The cell formula therefore yields
\[
\begin{aligned}
\int_0^1\widetilde g_\sigma(x)\,\lambda(dx)
&=\sum_{i=0}^{K-1}
 \frac{g_\sigma(i/K)+g_\sigma((i+1)/K)}{2K}\\
&=0.
\end{aligned}
\]
% lean: CausalSmith.Stat.FinitepHomogeneityDensegamma.smoothedTent_integral_zero_of_legalityRanks

\item \textit{Second-moment bounds.}
We first sharpen the coarse-field difference estimate used for interpolation error. Define the midpoint between the two cells of a pair by
\[
x_{j,\mathrm{mid}}=\frac{2j-1}{M},
\qquad j\in\{1,\ldots,M/2\}.
\]
On either side of this point, one of the two extended tents in
\(\mathfrak g_j\) vanishes. The tent Lipschitz estimate thus gives a constant \(2M\) when both arguments are on the same side. If
\(x\le x_{j,\mathrm{mid}}\le z\), both tents vanish at the midpoint, and
\[
\begin{aligned}
|\mathfrak g_j(x)-\mathfrak g_j(z)|
&\le|\mathfrak g_j(x)|+|\mathfrak g_j(z)|\\
&\le2M(x_{j,\mathrm{mid}}-x)
     +2M(z-x_{j,\mathrm{mid}})
=2M|x-z|.
\end{aligned}
\]
Reversing the arguments covers the other crossing case. Summing over the at most two active pair indices gives
\[
|g_\sigma(x)-g_\sigma(z)|\le4M|x-z|,
\qquad x,z\in[0,1].
\]

The partition identity gives the exact interpolation-error expression
\[
\widetilde g_\sigma(x)-g_\sigma(x)
=
\sum_{i=0}^{K}
 \{g_\sigma(i/K)-g_\sigma(x)\}f_i(x)^2.
\]
Whenever \(f_i(x)\ne0\), its adjacent-cell support implies
\(|i/K-x|\le1/K\); terms with \(f_i(x)=0\) vanish. Hence
\[
\begin{aligned}
|\widetilde g_\sigma(x)-g_\sigma(x)|
&\le\sum_{i=0}^{K}\frac{4M}{K}f_i(x)^2\\
&=\frac{4M}{K}\le\frac14.
\end{aligned}
\]

For the exact coarse second moment, direct integration gives
\[
\int_0^1 b_\triangle(r_\triangle)^2\,dr_\triangle
=2\int_0^{1/2}(2r_\triangle)^2\,dr_\triangle
=\frac13.
\]
Each coarse cell therefore contributes \(1/(3M)\). The cell interiors are disjoint, the signs square to one, and the even rank supplies all \(M\) cells, so
\[
\int_0^1g_\sigma(x)^2\,\lambda(dx)=\frac13.
\]
Pointwise, the error estimate and the elementary square inequality give
\[
\begin{aligned}
g_\sigma(x)^2
&\le2\widetilde g_\sigma(x)^2
   +2\{\widetilde g_\sigma(x)-g_\sigma(x)\}^2\\
&\le2\widetilde g_\sigma(x)^2+\frac18.
\end{aligned}
\]
Integration against the probability measure \(\lambda\) yields
\[
\frac13
\le2\int_0^1\widetilde g_\sigma(x)^2\,\lambda(dx)+\frac18,
\]
which implies the lower bound below. The upper bound follows from the unit envelope:
\[
\frac1{16}
\le\int_0^1\widetilde g_\sigma(x)^2\,\lambda(dx)
\le\int_0^1 1\,\lambda(dx)=1.
\]
% lean: CausalSmith.Stat.FinitepHomogeneityDensegamma.smoothedTent_second_moment_lower_of_legalityRanks

\item \textit{Effect identity and distance calibration.}
Define the effect multiplier by
\[
A_{\mathrm{eff}}
=2C_{\mathrm{cop}}
=\frac{2ab\kappa}{1-a^2}>0.
\]
The alternative table calculation gives, pointwise,
\[
\tau_P(x)=-A_{\mathrm{eff}}\widetilde g_\sigma(x)
=-\frac{2ab\kappa}{1-a^2}\widetilde g_\sigma(x).
\]
Since the interpolated tent is centered,
\[
\int_0^1\tau_P(x)\,\lambda(dx)
=-A_{\mathrm{eff}}
 \int_0^1\widetilde g_\sigma(x)\,\lambda(dx)=0.
\]
Substitution into the definition of \(d(P)\), followed by factoring the nonnegative multiplier out of the square root, gives
\[
d(P)
=A_{\mathrm{eff}}
 \left(\int_0^1\widetilde g_\sigma(x)^2\,\lambda(dx)\right)^{1/2}.
\]

The amplitude product and \(\kappa=1/16\) give
\[
2ab\kappa=2^{-15}K^{-S}.
\]
Because \(0\le a^2\le1/256\), the denominator is positive and lies between \(1/2\) and one. Therefore
\[
2^{-15}K^{-S}
\le A_{\mathrm{eff}}
\le2^{-14}K^{-S}.
\]
Taking square roots of the second-moment bounds gives
\[
\frac14
\le
\left(\int_0^1\widetilde g_\sigma(x)^2\,\lambda(dx)\right)^{1/2}
\le1.
\]
Multiplication of the two intervals proves
\[
2^{-17}K^{-S}\le d(P)\le2^{-14}K^{-S}.
\]

For the strict comparison, \cref{obj:lem:causal-nonempty} supplies
\(d_0=1/(8\sqrt2)\). Since \(S>0\) and \(K\ge1\), we have
\(K^{-S}\le1\). Also \(\sqrt2<2\), so
\[
d(P)\le2^{-14}K^{-S}
\le\frac1{16384}
<\frac1{8\sqrt2}=d_0.
\]
This establishes the effect, centering, and distance conclusions for every alternative realization in the finite-moment construction.
% lean: CausalSmith.Stat.FinitepHomogeneityDensegamma.legality_separation_of_moments

\item \textit{Bounded construction.}
Now take an admissible \(w\) and \(v=(2,\alpha,\beta,\gamma)\). This \(v\) is admissible, and the assumed phase and rank conditions are exactly those used above. Under
\[
u=2b,\qquad \varepsilon=\frac12,\qquad L=1,
\]
the input verification in the first step applies, and
\[
\varepsilon Lu=b,\qquad
\varepsilon L^p=\varepsilon L^2=\frac12\le10.
\]
Thus the null-membership calculation gives all restrictions of
\(H_0(v)\), and the alternative-membership calculation gives all restrictions of \(\mathcal M_v\).

In either branch, the normalized table has outcome support
\[
Y\in\{-1,0,1\},
\qquad P(|Y|\le1)=1.
\]
Indeed, each conditional table atom has absolute outcome at most one, so its conditional probability of this event is one; integrating over \(\lambda\) preserves that value. This supplies
\cref{obj:ass:bounded-outcome}. Together with the previously established restrictions, it gives bounded null and alternative membership by
\cref{obj:def:bounded-null,obj:def:bounded-model}.

The table effect calculation uses \(\varepsilon Lu=b\), so every bounded alternative has the same displayed effect formula. The tent centering and second-moment estimates depend on the ranks and coarse signs, and the distance calibration uses the retained amplitudes. Hence the centering identity, both distance bounds, and the strict comparison with \(d_0\) also apply. Finally, the exact arm-moment calculation gives
\[
\int_{\mathbb R}|y|^2
\,Q_{(1+\ell)/2,P}(dy\mid x)
=\varepsilon L^2=\frac12,
\qquad \ell\in\{-1,1\},\quad x\in[0,1].
\]
The coefficient realization was arbitrary throughout, so all conclusions hold for every realization in both priors.
% lean: CausalSmith.Stat.FinitepHomogeneityDensegamma.copula_frame_legality
\end{enumerate}
\end{proof}

The support conclusion in \cref{obj:lem:copula-frame-legality} applies to every coefficient realization. The rank-compatibility condition links the coarse effect resolution to the product of the two fine-scale amplitudes. The displayed effect is centered and has distance of order \(K^{-S}\). In the finite-moment tuning, mark rarity and magnitude yield conditional \(p\)-moment exactly one; in the bounded tuning, unit mark magnitude and mark probability one half yield conditional second moment one half. Both tunings retain the same effect identity and distance bounds.

\subsection{Augmentation and complete-record likelihood budgets}

The component comparison retains the complete record through a common augmentation. Write \(B_i=\mathbf1\{Y_i\ne0\}\) for the mark indicator, and let \(\mathcal D_{\partial}\) disclose coefficient pairs at the coarse-cell boundary nodes. The augmentation consists of the covariates, these indicators, and this boundary disclosure. Its common probability law is denoted by \(\mu\).

Conditional on the augmentation, use \(P_0\) and \(P_1\) for the two branch label laws. These are local conditional laws in the component comparison. Each label contains treatment and an outcome sign, including a retained sign when its mark indicator is zero. The intermediate conditional label law is denoted by \(Q\), so that the intermediate augmented probability law is \(\mu\otimes Q\). Recovering the observed outcome from the sign and mark indicator gives the original-record mixtures.

Conditioning on the augmentation fixes the allocation of covariates to fine and coarse cells and identifies which records carry nonzero marks. Boundary disclosure is part of this conditioning. The first and second conditional likelihood budgets, denoted by \(\mathcal A\) and \(\mathcal B\), bound the comparisons from \(P_0\) to \(Q\) and from \(Q\) to \(P_1\), respectively. The following result averages these budgets over the same augmentation law and gives the resulting complete-record comparison.

\begin{lemmav}{lem:full-record-copula-component-bound}[Full-record copula component bound]
Let \(v=(p,\alpha,\beta,\gamma)\) be any real tuple, and suppose:
\begin{itemize}
\item \textbf{(Sample size and ranks.)} \(n,K,M\) are natural numbers, \(n\ge2\), \(K\) and \(M\) are powers of two, \(M\ge2\), \(16M\le K\), and \(2^{40}n\le K\).
\item \textbf{(Amplitudes.)} \(0<a\le1/16\) and \(0<u\le1/16\).
\item \textbf{(Mark parameters.)} \(0<\varepsilon<1\) and \(L>0\).
\end{itemize}
Write \(\pi_0\) and \(\pi_1\) for the two copula-frame priors
\[
\pi_\nu=\pi^{K,M}_{\nu;v,a,u,\varepsilon,L},
\qquad \nu\in\{0,1\},
\]
and define their complete-record mixtures by
\[
\mathbb P_\nu=\int P^{\otimes n}\,\pi_\nu(dP),
\qquad \nu\in\{0,1\}.
\]
Let \(\lambda\) denote Lebesgue probability measure on \([0,1]\). Total variation, denoted by \(\operatorname{TV}\), is the supremum of the absolute difference of event probabilities. The directed chi-square divergence \(\chi^2\) is the integral of the squared Radon--Nikodym density minus one, for probability laws with the stated absolute continuity.

There exist a common augmentation law \(\mu\), conditional label kernels \(P_0,P_1,Q\), and nonnegative measurable, integrable conditional budgets \(\mathcal A,\mathcal B\) with the following properties. The augmentation coordinate is
\[
\left((X_i)_{i=1}^n,(B_i)_{i=1}^n,\mathcal D_{\partial}\right),
\qquad B_i=\mathbf1\{Y_i\ne0\},
\]
where \(\mathcal D_{\partial}\) records the sampled coefficient pairs
\((\lambda_j,\eta_j)\) exactly at nodes satisfying
\[
jM\bmod K=0,\qquad 0\le j\le K.
\]
Its covariate marginal is \(\lambda^{\otimes n}\), and its flag marginal is
\(\operatorname{Bernoulli}(\varepsilon)^{\otimes n}\).

Each label consists of a binary treatment and a binary outcome sign, retaining the sign also when the mark flag is zero. Under branch \(\nu\), the joint law of the augmentation and these labels induced by the latent copula construction is \(\mu\otimes P_\nu\); the disclosure equals the boundary restriction of the sampled coefficients almost surely. Under each branch, the \(M/2\) coarse signs are independent fair signs, and their joint law with the augmentation is the product of their fair-sign law and \(\mu\).

For each branch, recovering the records by retaining \(X_i\) and the treatment label and setting the outcome to the signed magnitude \(L\) when \(B_i=1\), and to zero when \(B_i=0\), sends \(\mu\otimes P_\nu\) to \(\mathbb P_\nu\). Both augmented branch laws have augmentation marginal \(\mu\).

The intermediate augmented law \(\mu\otimes Q\) is a probability law. For \(\mu\)-almost every augmentation, \(P_0,P_1,Q\) are probability laws on the labels and satisfy
\[
Q\ll P_0,\qquad P_1\ll Q,\qquad
\mathcal A\ge0,\qquad \mathcal B\ge0,
\]
\[
1+\chi^2(Q\Vert P_0)\le \exp(\mathcal A),
\qquad
1+\chi^2(P_1\Vert Q)\le \exp(\mathcal B).
\]
Here all three conditional laws and both budgets are evaluated at the full augmentation triple. Their expectations satisfy
\[
\int \mathcal A\,d\mu
\le \frac{2^{38}a^4u^4\varepsilon^2 n^2}{K},
\qquad
\int \mathcal B\,d\mu
\le \frac{2^{56}a^4u^4\varepsilon^2 n^4}{MK^2}.
\]
If
\[
\frac{2^{38}a^4u^4\varepsilon^2 n^2}{K}
+
\frac{2^{56}a^4u^4\varepsilon^2 n^4}{MK^2}
\le 2^{-16},
\]
then
\[
\operatorname{TV}(\mathbb P_0,\mathbb P_1)<\frac18.
\]
\end{lemmav}

\begin{proof}[Proof of \cref{obj:lem:full-record-copula-component-bound}]
\begin{enumerate}
\item
We first construct the common augmentation and the three label kernels. For \(J\in\{K,M\}\), define the cell index, including its value at the right endpoint, by
\[
\operatorname{cell}_J(x)=\min\{J-1,\lfloor Jx\rfloor\},
\qquad x\in[0,1].
\]
Here \(K>0\), since \(16M\le K\) and \(M\ge2\). Define the boundary-node set and the disclosure map by
\[
\mathsf N_{\partial}
 =\{j\in\{0,\ldots,K\}:jM\bmod K=0\},
\qquad
\mathsf D(\boldsymbol\lambda,\boldsymbol\eta)_j
 =
 \begin{cases}
 (\lambda_j,\eta_j),&j\in\mathsf N_{\partial},\\
 \bot,&j\notin\mathsf N_{\partial}.
 \end{cases}
\]
The symbol \(\bot\) denotes an undisclosed pair. Thus the space of possible disclosures is
\[
\mathsf\Delta_K
 =\bigl(\{\bot\}\cup\{-1,1\}^2\bigr)^{\{0,\ldots,K\}},
\]
and the augmentation space and a generic augmentation are
\[
\mathsf\Omega_{\mathrm{aug}}
 =[0,1]^n\times\{0,1\}^n\times\mathsf\Delta_K,
\qquad
\omega=(\mathbf X,\mathbf B,\delta).
\]

Connect observations \(i\) and \(i'\) when
\[
\begin{aligned}
&\operatorname{cell}_M(X_i)=\operatorname{cell}_M(X_{i'}),\\
&\operatorname{cell}_K(X_i)=\operatorname{cell}_K(X_{i'})
\quad\text{or}\quad
\left[
\begin{array}{c}
|\operatorname{cell}_K(X_i)-\operatorname{cell}_K(X_{i'})|=1,\\
\max\{\operatorname{cell}_K(X_i),\operatorname{cell}_K(X_{i'})\}
 \notin\mathsf N_{\partial}.
\end{array}
\right]
\end{aligned}
\]
Let \(\mathscr C(\omega)\) be the collection of connected components of this graph, with connectivity including paths of length zero. It partitions \(\{1,\ldots,n\}\). For every subset \(\mathsf c\subseteq\{1,\ldots,n\}\), set
\[
\mathsf n_{\mathsf c}=|\mathsf c|,
\qquad
\mathsf k_{\mathsf c}=\sum_{i\in\mathsf c}B_i,
\qquad
\mathsf j(\mathsf c)
 =\max\left(
 \left\{\left\lfloor\frac{\operatorname{cell}_M(X_i)}2\right\rfloor:
 i\in\mathsf c\right\}\cup\{0\}\right).
\]
The last convention defines \(\mathsf j(\varnothing)=0\).

We represent the binary labels by their signs:
\[
\mathbf z=(z_i^{\mathrm t},z_i^{\mathrm y})_{i=1}^n
 \in\mathsf Z_n:=\bigl(\{-1,1\}^2\bigr)^n,
\qquad
\mathsf F_n
 =4^{-n}\sum_{\mathbf z\in\mathsf Z_n}
       \operatorname{Dirac}_{\mathbf z}.
\]
The treatment label is \((1+z_i^{\mathrm t})/2\); the outcome-sign label is retained for both values of \(B_i\).

For \(\nu\in\{0,1\}\), a coarse-sign vector
\(\sigma\in\{-1,1\}^{M/2}\), and a disclosure \(\delta\), define the conditional coefficient weights by
\[
\mathsf w_{\nu,\sigma,\delta}
 (\boldsymbol\lambda,\boldsymbol\eta)
 =
 \prod_{j=0}^{K}
 \begin{cases}
 \mathbf1\{(\lambda_j,\eta_j)=\delta_j\},&\delta_j\ne\bot,\\[2pt]
 \displaystyle
 \frac{1+\nu\kappa g_\sigma(j/K)\lambda_j\eta_j}{4},
 &\delta_j=\bot,
 \end{cases}
\qquad \kappa=\frac1{16}.
\]
For every real function \(\Psi\) on the finite coefficient space, write
\[
\mathbb E_{\nu,\sigma,\delta}[\Psi]
 =
 \sum_{\boldsymbol\lambda,\boldsymbol\eta\in\{-1,1\}^{K+1}}
 \mathsf w_{\nu,\sigma,\delta}
 (\boldsymbol\lambda,\boldsymbol\eta)\,
 \Psi(\boldsymbol\lambda,\boldsymbol\eta).
\]
Each node weight is nonnegative and sums to one: for an undisclosed node this follows from
\[
|\kappa g_\sigma(j/K)|\le\frac1{16},
\qquad
\sum_{\lambda_j,\eta_j\in\{-1,1\}}\lambda_j\eta_j=0.
\]
Consequently,
\[
\mathsf w_{\nu,\sigma,\delta}\ge0,
\qquad
\sum_{\boldsymbol\lambda,\boldsymbol\eta}
 \mathsf w_{\nu,\sigma,\delta}
 (\boldsymbol\lambda,\boldsymbol\eta)=1.
\]

Use the coordinate fields of \cref{obj:def:copula-frame-priors}, and define the record-label density by
\[
\mathsf r_{\nu,i}
 =
 1+z_i^{\mathrm t}\xi(X_i)
 +B_i\left\{
 z_i^{\mathrm y}\upsilon(X_i)
 +z_i^{\mathrm t}z_i^{\mathrm y}\zeta(X_i)
 \right\}.
\]
Its dependence on the coefficient draw and \(\sigma\) is through those coordinate fields. On each fine cell, the two active frame coordinates satisfy
\[
\sum_{j=0}^{K}f_j(x)^2=1,
\qquad
\sum_{j=0}^{K}|f_j(x)|\le\sqrt2\le2.
\]
The same identities hold at cell endpoints. They imply
\[
|\xi(x)|\le\frac18,
\qquad
|\upsilon(x)|\le\frac18,
\qquad
|\widetilde g_\sigma(x)|\le1,
\qquad
|t_\nu(x)|
 \le\frac{au\kappa}{1-a^2}\le\frac1{128}.
\]
Indeed, \(a^2\le1/256\), and hence \(1-a^2\ge255/256>0\). Moreover,
\[
|1-\xi(x)^2|\le1,
\qquad
|\zeta(x)|
 \le\frac1{64}+\frac1{128}\le\frac18.
\]
Thus, for every augmentation, coefficient draw, and label vector,
\[
\frac12\le\mathsf r_{\nu,i}\le2,
\qquad
\frac14\sum_{z_i^{\mathrm t},z_i^{\mathrm y}\in\{-1,1\}}
 \mathsf r_{\nu,i}=1.
\]

For \(s_{\mathrm{coarse}}\in\{-1,1\}\), let
\(\sigma^{s_{\mathrm{coarse}}}\) denote the constant coarse-sign vector:
\[
\sigma^{s_{\mathrm{coarse}}}_j=s_{\mathrm{coarse}}
\quad(1\le j\le M/2).
\]
Define, for every subset \(\mathsf c\),
\[
\mathsf d_{\mathsf c,\nu,s_{\mathrm{coarse}}}
 (\omega,\mathbf z)
 =
 \mathbb E_{\nu,\sigma^{s_{\mathrm{coarse}}},\delta}
 \left[\prod_{i\in\mathsf c}\mathsf r_{\nu,i}\right],
\]
and abbreviate
\[
\mathsf d_{\mathsf c,0}
 =\mathsf d_{\mathsf c,0,-1},
\qquad
\mathsf d_{\mathsf c,+}
 =\mathsf d_{\mathsf c,1,1},
\qquad
\mathsf d_{\mathsf c,-}
 =\mathsf d_{\mathsf c,1,-1},
\qquad
\mathsf d_{\mathsf c,\mathrm{av}}
 =\frac{\mathsf d_{\mathsf c,+}+\mathsf d_{\mathsf c,-}}2.
\]
The null density is independent of the coarse signs. Empty products equal one, so all these component densities equal one when \(\mathsf c=\varnothing\). Finite averaging and label normalization give
\[
2^{-\mathsf n_{\mathsf c}}
 \le\mathsf d_{\mathsf c,\nu,s_{\mathrm{coarse}}}
 \le2^{\mathsf n_{\mathsf c}},
\qquad
\int\mathsf d_{\mathsf c,\nu,s_{\mathrm{coarse}}}\,d\mathsf F_n=1,
\]
and the same bounds and normalization hold for
\(\mathsf d_{\mathsf c,\mathrm{av}}\).

Now define the three full-label densities and kernels by
\[
\begin{aligned}
\mathsf f_0(\omega,\mathbf z)
 &=\prod_{\mathsf c\in\mathscr C(\omega)}
       \mathsf d_{\mathsf c,0}(\omega,\mathbf z),\\
\mathsf f_1(\omega,\mathbf z)
 &=2^{-M/2}\sum_{\sigma\in\{-1,1\}^{M/2}}
   \mathbb E_{1,\sigma,\delta}
       \left[\prod_{i=1}^n\mathsf r_{1,i}\right],\\
\mathsf f_{\mathrm{av}}(\omega,\mathbf z)
 &=\prod_{\mathsf c\in\mathscr C(\omega)}
       \mathsf d_{\mathsf c,\mathrm{av}}(\omega,\mathbf z),\\
P_\nu(\omega,d\mathbf z)
 &=\mathsf f_\nu(\omega,\mathbf z)\,\mathsf F_n(d\mathbf z),
 \qquad \nu\in\{0,1\},\\
Q(\omega,d\mathbf z)
 &=\mathsf f_{\mathrm{av}}(\omega,\mathbf z)\,
       \mathsf F_n(d\mathbf z).
\end{aligned}
\]
These definitions apply to every augmentation, including disclosures outside the support of the augmentation law. Component densities depend only on labels in their component. Since the components partition the observations, their normalized products are normalized. The full coefficient average defining \(\mathsf f_1\) is normalized by the record-label normalization above. Therefore \(P_0(\omega),P_1(\omega),Q(\omega)\) are probability laws for every \(\omega\).

Let the disclosure law be the pushforward of independent fair coefficient pairs:
\[
\mu_{\partial}
 =
 4^{-(K+1)}
 \sum_{\boldsymbol\lambda,\boldsymbol\eta\in\{-1,1\}^{K+1}}
 \operatorname{Dirac}_{\mathsf D(\boldsymbol\lambda,\boldsymbol\eta)},
\qquad
\mu
 =\lambda^{\otimes n}
   \otimes\operatorname{Bernoulli}(\varepsilon)^{\otimes n}
   \otimes\mu_{\partial}.
\]
It is a probability law with the required covariate and flag marginals, and
\[
\delta_j\ne\bot\ \Longleftrightarrow\ j\in\mathsf N_{\partial}
\qquad\text{for }\mu\text{-almost every }\omega.
\]
The kernels are measurable: the graph is determined by finitely many measurable cell assignments, and all coefficient and label averages are finite. Their normalization gives
\[
(\mu\otimes P_\nu)\circ\operatorname{pr}_{\mathrm{aug}}^{-1}
 =\mu,\qquad \nu\in\{0,1\},
\qquad
(\mu\otimes Q)(\mathsf\Omega_{\mathrm{aug}}\times\mathsf Z_n)=1.
\]
% lean: CausalSmith.Stat.FinitepHomogeneityDensegamma.copula_common_augmentation_certificate
% lean: CausalSmith.Stat.FinitepHomogeneityDensegamma.intermediateConditionalLaw_isProbabilityMeasure

\item\label{proof:full-record-copula-component-bound:conditional-comparisons}
We establish the conditional comparisons. Write
\[
K=2^{k_{\mathrm f}},\qquad M=2^{k_{\mathrm c}},
\qquad k_{\mathrm f},k_{\mathrm c}\in\mathbb N.
\]
Since \(M\le K\) and \(M\ge2\),
\[
k_{\mathrm c}\le k_{\mathrm f},
\qquad k_{\mathrm c}\ge1,
\qquad M\mid K,\qquad 2\mid M.
\]
In particular,
\[
\mathsf N_{\partial}
 =\{jK/M:0\le j\le M\}.
\]

Fix an augmentation in the support of \(\mu\). Its disclosed coefficients are fixed. The undisclosed coefficient nodes used by a component are
\[
\mathsf N_{\mathsf c}
 =
 \{j\notin\mathsf N_{\partial}:
       f_j(X_i)\ne0\text{ for some }i\in\mathsf c\}.
\]
These node sets are disjoint for distinct components. Indeed, two observations sharing an undisclosed active node lie in the same fine interval or in neighboring fine intervals joined through that node, and therefore are connected. Each component lies in one coarse cell. Its coefficient coupling and its correction consequently depend on the single sign indexed by \(\mathsf j(\mathsf c)\). Conditional coefficient independence then gives
\[
\mathbb E_{\nu,\sigma,\delta}
 \left[\prod_{i=1}^n\mathsf r_{\nu,i}\right]
 =
 \prod_{\mathsf c\in\mathscr C(\omega)}
 \mathsf d_{\mathsf c,\nu,\sigma_{\mathsf j(\mathsf c)+1}}.
\]
Define the occupied coarse-pair indices and their component groups by
\[
\mathscr J(\omega)
 =\{\mathsf j(\mathsf c):\mathsf c\in\mathscr C(\omega)\},
\qquad
\mathscr C_j(\omega)
 =\{\mathsf c\in\mathscr C(\omega):\mathsf j(\mathsf c)=j\}.
\]
The index \(\mathsf j(\mathsf c)\) belongs to \(\{0,\ldots,M/2-1\}\). Averaging the independent fair coarse signs yields
\[
\mathsf f_1
 =
 \prod_{j\in\mathscr J(\omega)}
 \frac12\left\{
   \prod_{\mathsf c\in\mathscr C_j(\omega)}
      \mathsf d_{\mathsf c,+}
   +
   \prod_{\mathsf c\in\mathscr C_j(\omega)}
      \mathsf d_{\mathsf c,-}
 \right\}.
\]

For every augmentation and every subset \(\mathsf c\), define
\[
\Delta_{\mathsf c,\mathrm e}
 =\mathsf d_{\mathsf c,\mathrm{av}}-\mathsf d_{\mathsf c,0},
\qquad
\Delta_{\mathsf c,\mathrm o}
 =\frac{\mathsf d_{\mathsf c,+}-\mathsf d_{\mathsf c,-}}2,
\]
\[
\mathfrak e_{\mathsf c}
 =\int\frac{\Delta_{\mathsf c,\mathrm e}^2}
                 {\mathsf d_{\mathsf c,0}}\,d\mathsf F_n,
\qquad
\mathfrak o_{\mathsf c}
 =\int\frac{\Delta_{\mathsf c,\mathrm o}^2}
                 {\mathsf d_{\mathsf c,\mathrm{av}}}\,d\mathsf F_n.
\]
The denominators are strictly positive. Component normalization gives
\[
\int\Delta_{\mathsf c,\mathrm e}\,d\mathsf F_n
 =\int\Delta_{\mathsf c,\mathrm o}\,d\mathsf F_n=0,
\qquad
\mathfrak e_{\mathsf c}\ge0,\qquad
\mathfrak o_{\mathsf c}\ge0.
\]
Choose the conditional budgets
\[
\mathcal A(\omega)
 =\sum_{\mathsf c\in\mathscr C(\omega)}
       \mathfrak e_{\mathsf c},
\qquad
\mathcal B(\omega)
 =\frac12
   \sum_{\mathsf c,\mathsf c'\in\mathscr C(\omega)}
   \mathbf1\{\mathsf c\ne\mathsf c',\
             \mathsf j(\mathsf c)=\mathsf j(\mathsf c')\}
       \mathfrak o_{\mathsf c}\mathfrak o_{\mathsf c'}.
\]

For the first comparison, cancellation of the even discrepancy gives
\[
\int
 \frac{\mathsf d_{\mathsf c,\mathrm{av}}^2}
      {\mathsf d_{\mathsf c,0}}\,d\mathsf F_n
 =1+\mathfrak e_{\mathsf c}.
\]
Integration over disjoint label blocks therefore gives
\[
1+\chi^2(Q(\omega)\Vert P_0(\omega))
 =\prod_{\mathsf c\in\mathscr C(\omega)}
      (1+\mathfrak e_{\mathsf c})
 \le\exp\!\left(
       \sum_{\mathsf c\in\mathscr C(\omega)}
          \mathfrak e_{\mathsf c}\right)
 =\exp(\mathcal A(\omega)).
\]

For the second comparison, define
\[
\vartheta_{\mathsf c}
 =\frac{\Delta_{\mathsf c,\mathrm o}}
        {\mathsf d_{\mathsf c,\mathrm{av}}}.
\]
Under the component law with density
\(\mathsf d_{\mathsf c,\mathrm{av}}\),
\[
\int\vartheta_{\mathsf c}\,
       \mathsf d_{\mathsf c,\mathrm{av}}\,d\mathsf F_n=0,
\qquad
\int\vartheta_{\mathsf c}^{\,2}\,
       \mathsf d_{\mathsf c,\mathrm{av}}\,d\mathsf F_n
 =\mathfrak o_{\mathsf c}.
\]
For two signs \(s_{\mathrm{coarse}},s'_{\mathrm{coarse}}\in\{-1,1\}\), it follows that
\[
\int
 (1+s_{\mathrm{coarse}}\vartheta_{\mathsf c})
 (1+s'_{\mathrm{coarse}}\vartheta_{\mathsf c})
 \mathsf d_{\mathsf c,\mathrm{av}}\,d\mathsf F_n
 =
 1+s_{\mathrm{coarse}}s'_{\mathrm{coarse}}
       \mathfrak o_{\mathsf c}.
\]
Using two independent copies of each fair coarse sign when squaring the likelihood ratio gives the exact identity
\[
1+\chi^2(P_1(\omega)\Vert Q(\omega))
 =
 \prod_{j\in\mathscr J(\omega)}
 \frac12\left\{
 \prod_{\mathsf c\in\mathscr C_j(\omega)}
      (1+\mathfrak o_{\mathsf c})
 +
 \prod_{\mathsf c\in\mathscr C_j(\omega)}
      (1-\mathfrak o_{\mathsf c})
 \right\}.
\]

Here is the exponential bound for each factor. For any finite collection
\(\mathscr U\) of components, define
\[
\begin{aligned}
\mathsf H_{\mathrm e}(\mathscr U)
 &=\frac12\left\{
   \prod_{\mathsf c\in\mathscr U}(1+\mathfrak o_{\mathsf c})
   +\prod_{\mathsf c\in\mathscr U}(1-\mathfrak o_{\mathsf c})
   \right\},\\
\mathsf H_{\mathrm o}(\mathscr U)
 &=\frac12\left\{
   \prod_{\mathsf c\in\mathscr U}(1+\mathfrak o_{\mathsf c})
   -\prod_{\mathsf c\in\mathscr U}(1-\mathfrak o_{\mathsf c})
   \right\},\\
\mathsf D_{\mathrm o}(\mathscr U)
 &=\sum_{\mathsf c\in\mathscr U}\mathfrak o_{\mathsf c},\\
\mathsf T_{\mathrm o}(\mathscr U)
 &=\frac12\sum_{\substack{\mathsf c,\mathsf c'\in\mathscr U\\
                         \mathsf c\ne\mathsf c'}}
      \mathfrak o_{\mathsf c}\mathfrak o_{\mathsf c'}.
\end{aligned}
\]
For the empty collection,
\[
\mathsf H_{\mathrm e}=1,\qquad
\mathsf H_{\mathrm o}=\mathsf D_{\mathrm o}
 =\mathsf T_{\mathrm o}=0.
\]
Upon adjoining \(\mathsf c_*\notin\mathscr U\), the recurrences are
\[
\begin{aligned}
\mathsf H_{\mathrm e}(\mathscr U\cup\{\mathsf c_*\})
 &=\mathsf H_{\mathrm e}(\mathscr U)
   +\mathfrak o_{\mathsf c_*}\mathsf H_{\mathrm o}(\mathscr U),\\
\mathsf H_{\mathrm o}(\mathscr U\cup\{\mathsf c_*\})
 &=\mathsf H_{\mathrm o}(\mathscr U)
   +\mathfrak o_{\mathsf c_*}\mathsf H_{\mathrm e}(\mathscr U),\\
\mathsf T_{\mathrm o}(\mathscr U\cup\{\mathsf c_*\})
 &=\mathsf T_{\mathrm o}(\mathscr U)
   +\mathfrak o_{\mathsf c_*}\mathsf D_{\mathrm o}(\mathscr U).
\end{aligned}
\]
These recurrences preserve
\[
\mathsf H_{\mathrm e}\ge0,\qquad
\mathsf H_{\mathrm o}\ge0,\qquad
\mathsf H_{\mathrm o}
 \le\mathsf D_{\mathrm o}\mathsf H_{\mathrm e}.
\]
For the last inequality, the old bound gives
\[
\mathsf H_{\mathrm o}
 +\mathfrak o_{\mathsf c_*}\mathsf H_{\mathrm e}
 \le
 (\mathsf D_{\mathrm o}+\mathfrak o_{\mathsf c_*})
 \mathsf H_{\mathrm e}
 \le
 (\mathsf D_{\mathrm o}+\mathfrak o_{\mathsf c_*})
 (\mathsf H_{\mathrm e}
   +\mathfrak o_{\mathsf c_*}\mathsf H_{\mathrm o}).
\]
Induction, using \(1+x\le e^x\), now gives
\[
\begin{aligned}
\mathsf H_{\mathrm e}(\mathscr U\cup\{\mathsf c_*\})
 &\le
 \mathsf H_{\mathrm e}(\mathscr U)
 \bigl(1+\mathfrak o_{\mathsf c_*}\mathsf D_{\mathrm o}(\mathscr U)\bigr)\\
 &\le
 \exp\!\left(
 \mathsf T_{\mathrm o}(\mathscr U)
 +\mathfrak o_{\mathsf c_*}\mathsf D_{\mathrm o}(\mathscr U)\right).
\end{aligned}
\]
Applying this to each \(\mathscr C_j(\omega)\) proves
\[
1+\chi^2(P_1(\omega)\Vert Q(\omega))
 \le\exp(\mathcal B(\omega)).
\]
Finally, \(\mathsf f_0\) and \(\mathsf f_{\mathrm{av}}\) are strictly positive on the finite label alphabet. Hence
\[
Q(\omega)\ll P_0(\omega),
\qquad
P_1(\omega)\ll Q(\omega).
\]
All these comparisons are conditional on the full triple \(\omega\).
% lean: CausalSmith.Stat.FinitepHomogeneityDensegamma.copula_conditional_comparison_certificate

\item
We derive the component estimates used for occupancy counting. In this step, the finite density expressions are considered as algebraic functions of
\[
a'\in(-1,1),\qquad u'\in\mathbb R,
\]
with the augmentation and coefficient weights fixed. Quantitative derivative bounds will be used on
\([0,1/16]^2\). Write \(\Delta_{\mathsf c,\mathrm e}(a',u')\) and
\(\Delta_{\mathsf c,\mathrm o}(a',u')\) for the resulting discrepancies.

The entire propensity-coefficient marginal and the entire outcome-coefficient marginal are common to both branches, with disclosed coordinates pinned. Therefore
\[
\Delta_{\mathsf c,\mathrm e}(a',0)
 =\Delta_{\mathsf c,\mathrm o}(a',0)
 =\Delta_{\mathsf c,\mathrm e}(0,u')
 =\Delta_{\mathsf c,\mathrm o}(0,u')=0.
\]
We also need the two first-order even cancellations.

To see them, the conditional coefficient means at an undisclosed node are
\[
\mathbb E_{\nu,\sigma,\delta}
 [\eta_j\mid\boldsymbol\lambda]
 =\nu\kappa g_\sigma(j/K)\lambda_j,
\qquad
\mathbb E_{\nu,\sigma,\delta}
 [\lambda_j\mid\boldsymbol\eta]
 =\nu\kappa g_\sigma(j/K)\eta_j.
\]
At a disclosed node, these means are the corresponding pinned values. For constant coarse sign, the undisclosed means change sign when the coarse sign changes.

If \(\mathsf c\) has exactly one marked observation \(i\), its record product equals
\[
\begin{aligned}
\prod_{i'\in\mathsf c}\mathsf r_{\nu,i'}
={}&
\prod_{i'\in\mathsf c\setminus\{i\}}
 (1+z_{i'}^{\mathrm t}\xi(X_{i'}))\\
&\times
\left[
 (1+z_i^{\mathrm t}\xi(X_i))
 (1+z_i^{\mathrm y}\upsilon(X_i))
 +z_i^{\mathrm t}z_i^{\mathrm y}
   t_\nu(X_i)(1-\xi(X_i)^2)
\right].
\end{aligned}
\]
Conditioning on the propensity coefficients, its outcome-field term is affine in the outcome coefficients. The branch-dependent conditional means of the undisclosed outcome coefficients, and the correction \(t_1\), are linear in the constant coarse sign. Averaging its two values cancels these changes. Disclosed-coordinate terms are common. Thus
\[
\mathsf d_{\mathsf c,+}+\mathsf d_{\mathsf c,-}
 =2\mathsf d_{\mathsf c,0},
\qquad
\Delta_{\mathsf c,\mathrm e}=0
\quad\text{if }\mathsf k_{\mathsf c}\le1.
\]
For zero marks, every record factor uses only the propensity coefficients, so
\[
\Delta_{\mathsf c,\mathrm o}=0
\quad\text{if }\mathsf k_{\mathsf c}=0.
\]

At \(a'=0\), all undifferentiated record factors depend only on the outcome coefficients. Differentiating one factor produces a term linear in the propensity coefficients and a correction term linear in the coarse sign. Conditioning on the outcome coefficients and averaging the two constant coarse signs, using the conditional means above, therefore gives
\[
\partial_{a'}\Delta_{\mathsf c,\mathrm e}(0,u')=0.
\]
For the outcome derivative, define the modified augmentation
\[
\omega^{(i)}
 =\bigl(\mathbf X,(B_{i'}\mathbf1\{i'=i\})_{i'=1}^n,\delta\bigr).
\]
Each record factor is affine in \(u'\), so the product rule at \(u'=0\) gives
\[
\partial_{u'}\Delta_{\mathsf c,\mathrm e}(a',0;\omega)
 =
 \sum_{i\in\mathsf c}
 \left\{
 \Delta_{\mathsf c,\mathrm e}(a',1;\omega^{(i)})
 -\Delta_{\mathsf c,\mathrm e}(a',0;\omega)
 \right\}=0.
\]
Here evaluation at unit outcome amplitude extracts the coefficient of the affine marked factor; the algebraic functions are defined there. Each modified augmentation has at most one mark, so the preceding cancellation applies.

For the derivative bounds, define the coefficient fields and the rational correction factor by
\[
\mathsf F_\lambda(x)=\sum_{j=0}^{K}f_j(x)\lambda_j,
\qquad
\mathsf F_\eta(x)=\sum_{j=0}^{K}f_j(x)\eta_j,
\qquad
\mathsf c_x(a')
 =\frac{a'(1-(a')^2\mathsf F_\lambda(x)^2)}
        {1-(a')^2}.
\]
The frame identities give
\[
|\mathsf F_\lambda|,|\mathsf F_\eta|\le\sqrt2\le2,
\qquad
|\mathsf F_\lambda\mathsf F_\eta|\le2.
\]
Direct differentiation gives
\[
\begin{aligned}
\mathsf c_x'(a')
 &=1+
 (1-\mathsf F_\lambda(x)^2)
 \frac{(a')^2(3-(a')^2)}{(1-(a')^2)^2},\\
\mathsf c_x''(a')
 &=(1-\mathsf F_\lambda(x)^2)
 \frac{2a'(3+(a')^2)}{(1-(a')^2)^3}.
\end{aligned}
\]
On \(0\le a'\le1/16\),
\[
1-(a')^2\ge\frac{15}{16},
\qquad
|1-\mathsf F_\lambda^2|\le1,
\]
and hence
\[
|\mathsf c_x|\le\frac1{15},
\qquad
|\mathsf c_x'|
 \le1+\frac{3/256}{(15/16)^2}\le2,
\qquad
|\mathsf c_x''|
 \le\frac{1/2}{(15/16)^3}\le1.
\]
In this notation the interaction coordinate is
\[
\zeta(x)
 =u'\left\{
 a'\mathsf F_\lambda(x)\mathsf F_\eta(x)
 -\nu\kappa\widetilde g_\sigma(x)\mathsf c_x(a')
 \right\}.
\]
The estimates needed for its derivatives are
\[
\begin{aligned}
|\mathsf F_\lambda\mathsf F_\eta
 -\nu\kappa\widetilde g_\sigma\mathsf c_x'|
 &\le 4+\frac18=\frac{33}{8},\\
|\mathsf F_\lambda\mathsf F_\eta
 -\nu\kappa\widetilde g_\sigma\mathsf c_x'|
 &\le 2+\frac18=\frac{17}{8},\\
|a'\mathsf F_\lambda\mathsf F_\eta
 -\nu\kappa\widetilde g_\sigma\mathsf c_x|
 &\le\frac18+\frac1{240}=\frac{31}{240}.
\end{aligned}
\]
The first, coarser bound is sufficient for the first propensity derivative. For every record factor on the amplitude rectangle,
\[
\begin{aligned}
|\partial_{a'}\mathsf r_{\nu,i}|
 &\le2+\frac1{16}\frac{33}{8}\le4,\\
|\partial_{a'}^2\mathsf r_{\nu,i}|
 &\le\frac1{256}\le4,\\
|\partial_{u'}\mathsf r_{\nu,i}|
 &\le2+\frac{31}{240}\le4,\\
|\partial_{a'}\partial_{u'}\mathsf r_{\nu,i}|
 &\le\frac{17}{8}\le4,\\
|\partial_{a'}^2\partial_{u'}\mathsf r_{\nu,i}|
 &\le\frac1{16}\le4,
\qquad
\partial_{u'}^2\mathsf r_{\nu,i}=0.
\end{aligned}
\]
Also \(|\mathsf r_{\nu,i}|\le2\) throughout this closed rectangle.

Assigning the four labelled differentiations to record factors gives at most
\(\mathsf n_{\mathsf c}^4\) terms in the fourth mixed product derivative. Each differentiated factor is bounded by four, and every remaining factor by two. Finite averaging therefore gives
\[
\left|
\partial_{a'}^2\partial_{u'}^2
 \mathsf d_{\mathsf c,\nu,s_{\mathrm{coarse}}}
\right|
 \le256\,\mathsf n_{\mathsf c}^4\,2^{\mathsf n_{\mathsf c}},
\]
and the two-differentiation product rule similarly gives
\[
\left|
\partial_{a'}\partial_{u'}
 \mathsf d_{\mathsf c,\nu,s_{\mathrm{coarse}}}
\right|
 \le16\,\mathsf n_{\mathsf c}^2\,2^{\mathsf n_{\mathsf c}}.
\]
These statements include the empty product, whose derivatives vanish. Taking the even and odd contrasts yields
\[
\left|
\partial_{a'}^2\partial_{u'}^2
 \Delta_{\mathsf c,\mathrm e}
\right|
 \le512\,\mathsf n_{\mathsf c}^4\,2^{\mathsf n_{\mathsf c}},
\qquad
\left|
\partial_{a'}\partial_{u'}
 \Delta_{\mathsf c,\mathrm o}
\right|
 \le16\,\mathsf n_{\mathsf c}^2\,2^{\mathsf n_{\mathsf c}}.
\]

The axis identities and first-order cancellations imply
\[
\partial_{a'}^2\Delta_{\mathsf c,\mathrm e}(a',0)
 =
\partial_{a'}^2\partial_{u'}
 \Delta_{\mathsf c,\mathrm e}(a',0)=0.
\]
Applying the second-order remainder formula first in outcome amplitude and then in propensity amplitude gives
\[
\begin{aligned}
|\Delta_{\mathsf c,\mathrm e}(a,u)|
&\le
512\,\mathsf n_{\mathsf c}^4\,2^{\mathsf n_{\mathsf c}}
 \int_0^a(a-a')\,da'
 \int_0^u(u-u')\,du'\\
&=
128a^2u^2\mathsf n_{\mathsf c}^4\,2^{\mathsf n_{\mathsf c}}.
\end{aligned}
\]
For the odd discrepancy,
\(\partial_{a'}\Delta_{\mathsf c,\mathrm o}(a',0)=0\).
The mixed derivative bound and the mean-value inequality give, successively,
\[
|\partial_{a'}\Delta_{\mathsf c,\mathrm o}(a',u)|
 \le16u\,\mathsf n_{\mathsf c}^2\,2^{\mathsf n_{\mathsf c}},
\qquad
|\Delta_{\mathsf c,\mathrm o}(a,u)|
 \le16au\,\mathsf n_{\mathsf c}^2\,2^{\mathsf n_{\mathsf c}}.
\]

It remains to check singleton cancellation under the actual disclosure. Distinct boundary nodes are separated by \(K/M\ge16\) fine intervals, so at most one active frame node at any covariate is disclosed. Distinct undisclosed coefficient pairs are independent, with each individual coefficient having mean zero. Thus, for every fixed coarse-sign vector,
\[
\mathbb E_{1,\sigma,\delta}[\xi(x)^2]=a^2,
\]
\[
\mathbb E_{1,\sigma,\delta}[\xi(x)\upsilon(x)]
 -\mathbb E_{0,\sigma,\delta}[\xi(x)\upsilon(x)]
 =au\kappa\widetilde g_\sigma(x).
\]
For the first identity, the diagonal terms sum to \(a^2\sum_jf_j(x)^2=a^2\), and every cross term has an undisclosed mean-zero factor. For the second, the undisclosed same-node correlation is
\(\kappa g_\sigma(j/K)\); disclosed same-node terms are common, and
\(g_\sigma(j/K)=0\) at disclosed nodes. Cross terms again vanish. The first moments of \(\xi\) and \(\upsilon\) are common to both branches, and
\[
au\kappa\widetilde g_\sigma(x)
 +t_1(x)(1-a^2)=0.
\]
Consequently every singleton label density agrees with its null density, for either value of its coarse sign. These identities also hold at endpoints by the frame formulas. For the empty subset, normalization of the coefficient weights gives all component densities equal to one. Hence
\[
\Delta_{\mathsf c,\mathrm e}
 =\Delta_{\mathsf c,\mathrm o}=0
\quad\text{if }\mathsf n_{\mathsf c}\le1.
\]

Squaring the discrepancy bounds and using the denominator lower bounds now gives, for \(\mu\)-almost every augmentation and every component,
\[
\mathfrak e_{\mathsf c}
 \le
 2^{14}a^4u^4\mathsf n_{\mathsf c}^{\,8}
       8^{\mathsf n_{\mathsf c}}
 \mathbf1\{\mathsf n_{\mathsf c}\ge2,\
           \mathsf k_{\mathsf c}\ge2\},
\]
\[
\mathfrak o_{\mathsf c}
 \le
 2^8a^2u^2\mathsf n_{\mathsf c}^{\,4}
       8^{\mathsf n_{\mathsf c}}
 \mathbf1\{\mathsf n_{\mathsf c}\ge2,\
           \mathsf k_{\mathsf c}\ge1\}.
\]
% lean: CausalSmith.Stat.FinitepHomogeneityDensegamma.component_discrepancy_bounds
% lean: CausalSmith.Stat.FinitepHomogeneityDensegamma.component_activity_envelopes

\item
We bound the expectations of the budgets. Define the two occupancy scores by
\[
\mathsf S_{\mathrm{one}}(\omega)
 =
 \sum_{\mathsf c\in\mathscr C(\omega)}
 \mathsf n_{\mathsf c}^{\,8}8^{\mathsf n_{\mathsf c}}
 \mathbf1\{\mathsf n_{\mathsf c}\ge2,\
           \mathsf k_{\mathsf c}\ge2\},
\]
\[
\begin{aligned}
\mathsf S_{\mathrm{pair}}(\omega)
 =\frac12
 \sum_{\mathsf c,\mathsf c'\in\mathscr C(\omega)}
 &\mathsf n_{\mathsf c}^{\,4}\mathsf n_{\mathsf c'}^{\,4}
   8^{\mathsf n_{\mathsf c}+\mathsf n_{\mathsf c'}}\\
 &\times
 \mathbf1\left\{
 \begin{array}{c}
 \mathsf c\ne\mathsf c',\
 \mathsf j(\mathsf c)=\mathsf j(\mathsf c'),\\
 \mathsf n_{\mathsf c},\mathsf n_{\mathsf c'}\ge2,\
 \mathsf k_{\mathsf c},\mathsf k_{\mathsf c'}\ge1
 \end{array}\right\}.
\end{aligned}
\]
The component estimates imply
\[
\mathcal A\le2^{14}a^4u^4\mathsf S_{\mathrm{one}},
\qquad
\mathcal B\le2^{16}a^4u^4\mathsf S_{\mathrm{pair}}
\quad\mu\text{-almost surely}.
\]

Under \(\mu\), the fine-cell indices are independent uniform draws from
\(\{0,\ldots,K-1\}\), and the flags are independent Bernoulli draws, independent of those indices. A connected component of size \(r_{\mathrm{occ}}\) occupies consecutive fine cells spanning at most \(r_{\mathrm{occ}}\) cells. For a specified observation subset of that size and a specified starting cell, the necessary occupancy event has probability at most
\[
\left(\frac{r_{\mathrm{occ}}}{K}\right)^{r_{\mathrm{occ}}}.
\]
A window extending past the final cell has fewer available cells and obeys the same bound. The probability of at least two marks in that subset is at most
\(r_{\mathrm{occ}}^2\varepsilon^2\). There are at most \(K\) starting cells and
\(\binom n{r_{\mathrm{occ}}}\) observation subsets. Overcounting these necessary events yields
\[
\int\mathsf S_{\mathrm{one}}\,d\mu
 \le
 \varepsilon^2K
 \sum_{r_{\mathrm{occ}}=2}^{n}
 \binom n{r_{\mathrm{occ}}}
 r_{\mathrm{occ}}^{10}8^{r_{\mathrm{occ}}}
 \left(\frac{r_{\mathrm{occ}}}{K}\right)^{r_{\mathrm{occ}}}.
\]

For two distinct components, their observation subsets are disjoint. Each requires a mark, so for sizes
\(r_{\mathrm{occ}},r'_{\mathrm{occ}}\) the joint flag probability is at most
\(r_{\mathrm{occ}}r'_{\mathrm{occ}}\varepsilon^2\). A coarse pair contains
\(2K/M\) fine cells. Thus the number of ordered starting-cell choices in a common coarse pair is at most
\[
\frac M2\left(\frac{2K}{M}\right)^2=\frac{2K^2}{M}.
\]
The number of disjoint subset choices is bounded by
\(\binom n{r_{\mathrm{occ}}}\binom n{r'_{\mathrm{occ}}}\).
Using independence on the disjoint subsets, and overcounting by ordered pairs, gives
\[
\int\mathsf S_{\mathrm{pair}}\,d\mu
 \le
 \frac{2\varepsilon^2K^2}{M}
 \left\{
 \sum_{r_{\mathrm{occ}}=2}^{n}
 \binom n{r_{\mathrm{occ}}}
 r_{\mathrm{occ}}^{5}8^{r_{\mathrm{occ}}}
 \left(\frac{r_{\mathrm{occ}}}{K}\right)^{r_{\mathrm{occ}}}
 \right\}^{2}.
\]
When the collection of eligible pairs is empty, its score is zero and this bound still applies.

Define the occupancy ratio by
\[
\zeta_{\mathrm{occ}}=8e\,\frac nK.
\]
The factorial and binomial bounds
\[
r_{\mathrm{occ}}!\ge
 \left(\frac{r_{\mathrm{occ}}}{e}\right)^{r_{\mathrm{occ}}},
\qquad
\binom n{r_{\mathrm{occ}}}
 \le\frac{n^{r_{\mathrm{occ}}}}{r_{\mathrm{occ}}!}
\]
imply, for every \(r_{\mathrm{occ}}\ge2\) and every nonnegative integer
\(d_{\mathrm{occ}}\),
\[
\binom n{r_{\mathrm{occ}}}
 r_{\mathrm{occ}}^{d_{\mathrm{occ}}}8^{r_{\mathrm{occ}}}
 \left(\frac{r_{\mathrm{occ}}}{K}\right)^{r_{\mathrm{occ}}}
 \le
 r_{\mathrm{occ}}^{d_{\mathrm{occ}}}
 \zeta_{\mathrm{occ}}^{r_{\mathrm{occ}}}
 \le
 r_{\mathrm{occ}}^{d_{\mathrm{occ}}+1}
 \zeta_{\mathrm{occ}}^{r_{\mathrm{occ}}}.
\]
The factorial inequality follows from the usual Stirling lower bound, whose square-root factor is at least one for positive integers. Extending the nonnegative finite sums gives
\[
\int\mathcal A\,d\mu
 \le
 2^{14}a^4u^4\varepsilon^2K
 \sum_{r_{\mathrm{occ}}=2}^{\infty}
 r_{\mathrm{occ}}^{11}
 \zeta_{\mathrm{occ}}^{r_{\mathrm{occ}}},
\]
\[
\int\mathcal B\,d\mu
 \le
 \frac{2^{17}a^4u^4\varepsilon^2K^2}{M}
 \left\{
 \sum_{r_{\mathrm{occ}}=2}^{\infty}
 r_{\mathrm{occ}}^{6}
 \zeta_{\mathrm{occ}}^{r_{\mathrm{occ}}}
 \right\}^{2}.
\]

For \(d_{\mathrm{ser}}\in\{6,11\}\), the elementary inequality
\(r_{\mathrm{occ}}\le2^{r_{\mathrm{occ}}-1}\) gives
\[
\begin{aligned}
\sum_{r_{\mathrm{occ}}=2}^{\infty}
 r_{\mathrm{occ}}^{d_{\mathrm{ser}}}
 \zeta_{\mathrm{occ}}^{r_{\mathrm{occ}}}
 &\le
 2^{d_{\mathrm{ser}}}\zeta_{\mathrm{occ}}^2
 \sum_{r_{\mathrm{occ}}=2}^{\infty}
 (2^{d_{\mathrm{ser}}}\zeta_{\mathrm{occ}})^{r_{\mathrm{occ}}-2}\\
 &=
 \frac{2^{d_{\mathrm{ser}}}\zeta_{\mathrm{occ}}^2}
      {1-2^{d_{\mathrm{ser}}}\zeta_{\mathrm{occ}}}
 \le2^{d_{\mathrm{ser}}+1}\zeta_{\mathrm{occ}}^2.
\end{aligned}
\]
The geometric series converges because
\[
0\le\zeta_{\mathrm{occ}}\le24\,\frac nK,
\qquad
2^{11}\zeta_{\mathrm{occ}}
 \le2^{11}\cdot24\cdot2^{-40}<\frac12.
\]
Here we used \(e<3\) and \(2^{40}n\le K\). Consequently,
\[
\begin{aligned}
\int\mathcal A\,d\mu
 &\le
 2^{26}24^2\,
 \frac{a^4u^4\varepsilon^2n^2}{K}
 \le
 \frac{2^{38}a^4u^4\varepsilon^2n^2}{K},\\
\int\mathcal B\,d\mu
 &\le
 2^{31}24^4\,
 \frac{a^4u^4\varepsilon^2n^4}{MK^2}
 \le
 \frac{2^{56}a^4u^4\varepsilon^2n^4}{MK^2},
\end{aligned}
\]
since \(2^{26}24^2<2^{38}\) and \(2^{31}24^4<2^{56}\).
The budgets are measurable finite sums. The uniform component denominator bounds and finite label alphabet make them bounded for fixed parameters, hence integrable under the probability law \(\mu\).
% lean: CausalSmith.Stat.FinitepHomogeneityDensegamma.component_occupancy_budgets

\item\label{proof:full-record-copula-component-bound:record-projections}
We identify the original-record projections. Define the recovery map by
\[
\mathsf O_L(\omega,\mathbf z)
 =
 \left(
 X_i,\frac{1+z_i^{\mathrm t}}2,B_i z_i^{\mathrm y}L
 \right)_{i=1}^{n}.
\]
This map is measurable.

For a fixed \(\sigma\), the coefficient law in
\cref{obj:def:copula-frame-priors} assigns mass
\[
\prod_{j=0}^{K}
 \frac{1+\nu\kappa g_\sigma(j/K)\lambda_j\eta_j}{4}
\]
to a coefficient draw. At boundary nodes, \(g_\sigma(j/K)=0\). Its disclosure law is therefore \(\mu_{\partial}\), and conditioning on a disclosure gives precisely the weights
\(\mathsf w_{\nu,\sigma,\delta}\). In particular, for every function \(\Psi\) on the coefficient space,
\[
\int
 \mathbb E_{\nu,\sigma,\delta}[\Psi]\,
 \mu_{\partial}(d\delta)
 =
 \sum_{\boldsymbol\lambda,\boldsymbol\eta}
 \left\{
 \prod_{j=0}^{K}
 \frac{1+\nu\kappa g_\sigma(j/K)\lambda_j\eta_j}{4}
 \right\}
 \Psi(\boldsymbol\lambda,\boldsymbol\eta).
\]
For the null branch, the component factorization in \cref{proof:full-record-copula-component-bound:conditional-comparisons} identifies
\(\mathsf f_0\) with the conditional full-record label average. For the alternative branch, that average is the definition of \(\mathsf f_1\). Integrating the disclosure thus recovers the respective original coefficient averages.

For a fixed latent draw and covariate \(x\), the flag and fair-label weights give
\[
\Pr_\nu\left(
 A=\frac{1+z^{\mathrm t}}2,\,
 Y=z^{\mathrm y}L\mid x
 \right)
 =
 \frac{\varepsilon}{4}
 \left\{
 1+z^{\mathrm t}\xi(x)+z^{\mathrm y}\upsilon(x)
   +z^{\mathrm t}z^{\mathrm y}\zeta(x)
 \right\},
\]
and, summing the retained outcome-sign label when the flag is zero,
\[
\Pr_\nu\left(
 A=\frac{1+z^{\mathrm t}}2,\,
 Y=0\mid x
 \right)
 =
 \sum_{z^{\mathrm y}\in\{-1,1\}}
 \frac{1-\varepsilon}{4}(1+z^{\mathrm t}\xi(x))
 =
 \frac{1-\varepsilon}{2}(1+z^{\mathrm t}\xi(x)).
\]
These are exactly the categories of \cref{obj:def:marked-table}. Independent record generation conditional on the latent draw, followed by its prior average, therefore proves
\[
(\mu\otimes P_0)\circ\mathsf O_L^{-1}=\mathbb P_0,
\qquad
(\mu\otimes P_1)\circ\mathsf O_L^{-1}=\mathbb P_1.
\]
% lean: CausalSmith.Stat.FinitepHomogeneityDensegamma.null_augmented_projection_label_bind
% lean: CausalSmith.Stat.FinitepHomogeneityDensegamma.alternative_augmented_projection_label_bind
% lean: CausalSmith.Stat.FinitepHomogeneityDensegamma.prior_label_design_bind_eq_mixture

\item
Assume the sum of the two stated public budget bounds is at most \(2^{-16}\). We first obtain an envelope for conditional total variation.

On a finite alphabet, absolute continuity and Cauchy--Schwarz give
\[
\operatorname{TV}(Q(\omega),P_0(\omega))
 \le\frac12\sqrt{\chi^2(Q(\omega)\Vert P_0(\omega))},
\]
\[
\operatorname{TV}(P_1(\omega),Q(\omega))
 \le\frac12\sqrt{\chi^2(P_1(\omega)\Vert Q(\omega))}.
\]
Indeed, total variation is half the sum of absolute differences of the masses; writing those differences relative to the dominating law and applying Cauchy--Schwarz gives the displayed bounds.

For \(0\le x_{\mathrm{exp}}\le1/256\), comparison of the exponential series with the geometric series gives
\[
e^{x_{\mathrm{exp}}}
 \le\frac1{1-x_{\mathrm{exp}}}
 \le1+2x_{\mathrm{exp}},
\]
where the second inequality follows after multiplication by
\(1-x_{\mathrm{exp}}>0\). If
\(\mathcal A(\omega)+\mathcal B(\omega)\le1/256\), the conditional divergence bounds and the triangle inequality therefore yield
\[
\begin{aligned}
\operatorname{TV}(P_0(\omega),P_1(\omega))
 &\le
 \frac12\sqrt{2\mathcal A(\omega)}
 +\frac12\sqrt{2\mathcal B(\omega)}\\
 &\le\sqrt{\mathcal A(\omega)+\mathcal B(\omega)}
 \le\frac1{16}.
\end{aligned}
\]
The middle inequality follows by squaring and using
\(2\sqrt{\mathcal A\mathcal B}\le\mathcal A+\mathcal B\).
In the complementary case,
\(\mathcal A(\omega)+\mathcal B(\omega)>1/256\), total variation between probability laws is at most one. Both cases consequently give
\[
\operatorname{TV}(P_0(\omega),P_1(\omega))
 \le\frac1{16}
      +256\bigl(\mathcal A(\omega)+\mathcal B(\omega)\bigr)
\quad\mu\text{-almost surely}.
\]

For any measurable augmented event
\(\mathsf E_{\mathrm{aug}}\subseteq
 \mathsf\Omega_{\mathrm{aug}}\times\mathsf Z_n\),
define its label section by
\[
\mathsf E_{\mathrm{aug},\omega}
 =\{\mathbf z:(\omega,\mathbf z)\in\mathsf E_{\mathrm{aug}}\}.
\]
The common augmentation marginal implies
\[
\begin{aligned}
&|(\mu\otimes P_0)(\mathsf E_{\mathrm{aug}})
  -(\mu\otimes P_1)(\mathsf E_{\mathrm{aug}})|\\
&\qquad\le
 \int
 |P_0(\omega)(\mathsf E_{\mathrm{aug},\omega})
  -P_1(\omega)(\mathsf E_{\mathrm{aug},\omega})|\,\mu(d\omega)\\
&\qquad\le
 \frac1{16}
 +256\left(\int\mathcal A\,d\mu+\int\mathcal B\,d\mu\right)
 \le\frac1{16}+256\cdot2^{-16}
 =\frac{17}{256}.
\end{aligned}
\]
Taking the supremum over augmented events proves
\[
\operatorname{TV}(\mu\otimes P_0,\mu\otimes P_1)
 \le\frac{17}{256}.
\]
% lean: CausalSmith.Stat.FinitepHomogeneityDensegamma.copula_augmented_tv_budget

For every measurable original-record event
\(\mathsf E_{\mathrm{rec}}\), its inverse image under \(\mathsf O_L\) is an augmented event. The projection identities from \cref{proof:full-record-copula-component-bound:record-projections} therefore give
\[
\operatorname{TV}(\mathbb P_0,\mathbb P_1)
 \le
 \operatorname{TV}(\mu\otimes P_0,\mu\otimes P_1)
 \le\frac{17}{256}<\frac18.
\]
% lean: CausalSmith.Stat.FinitepHomogeneityDensegamma.copula_tv_map_le

\item
Finally, we verify that the augmentation is induced by the stipulated latent construction. Define its latent coefficient law by the masses
\[
\mathsf W_\nu
 \bigl(\{\sigma,\boldsymbol\lambda,\boldsymbol\eta\}\bigr)
 =
 2^{-M/2}
 \prod_{j=0}^{K}
 \frac{1+\nu\kappa g_\sigma(j/K)\lambda_j\eta_j}{4}.
\]
Define the joint latent augmented law by
\[
\begin{aligned}
&\mathsf L_\nu
 (d\sigma,d\boldsymbol\lambda,d\boldsymbol\eta,
  d\mathbf X,d\mathbf B,d\delta,d\mathbf z)\\
&\quad=
 \mathsf W_\nu(d\sigma,d\boldsymbol\lambda,d\boldsymbol\eta)\,
 \lambda^{\otimes n}(d\mathbf X)\,
 \operatorname{Bernoulli}(\varepsilon)^{\otimes n}(d\mathbf B)\,
 \operatorname{Dirac}_{\mathsf D(\boldsymbol\lambda,\boldsymbol\eta)}
       (d\delta)\,
 \left(\prod_{i=1}^{n}\mathsf r_{\nu,i}\right)
 \mathsf F_n(d\mathbf z).
\end{aligned}
\]
The coefficient weights sum to one, and the conditional label product integrates to one. Thus
\[
\mathsf L_\nu(\text{entire space})=1,
\qquad
\mathsf L_\nu\{\delta=\mathsf D(\boldsymbol\lambda,\boldsymbol\eta)\}=1.
\]
Conditioning its finite coefficient average on the disclosure, and then using the component factorization for branch zero and the full coefficient average for branch one, gives
\[
\mathsf L_\nu\circ
 \bigl((\sigma,\boldsymbol\lambda,\boldsymbol\eta,\omega,\mathbf z)
       \mapsto(\omega,\mathbf z)\bigr)^{-1}
 =\mu\otimes P_\nu,
\qquad \nu\in\{0,1\}.
\]

There are \(M+1\) disclosed nodes, and the coarse tents vanish at all of them. Their coefficient pairs are therefore independent fair pairs conditional on every coarse-sign vector. Summing the undisclosed coefficient weights gives, for each
\(\sigma_*\in\{-1,1\}^{M/2}\) and each actual boundary disclosure \(\delta\),
\[
\Pr_\nu\{\sigma=\sigma_*,\mathcal D_{\partial}=\delta\}
 =2^{-M/2}4^{-(M+1)}.
\]
Define the fair coarse-sign law by
\[
\mathsf U_\sigma
 =2^{-M/2}\sum_{\sigma\in\{-1,1\}^{M/2}}
       \operatorname{Dirac}_{\sigma}.
\]
The covariates and flags are drawn independently of these latent variables, and integrating the normalized labels leaves their law unchanged. Hence
\[
\mathsf L_\nu\circ
 \bigl((\sigma,\boldsymbol\lambda,\boldsymbol\eta,\omega,\mathbf z)
       \mapsto(\sigma,\omega)\bigr)^{-1}
 =\mathsf U_\sigma\otimes\mu.
\]
This proves both the independent fair-sign assertion and its product relation with the full augmentation. Since \(L>0\), recovery also satisfies
\(B_i=\mathbf1\{Y_i\ne0\}\). Together with the kernel, budget, projection, and total-variation properties already established, these are all the required conclusions.
% lean: CausalSmith.Stat.FinitepHomogeneityDensegamma.copulaLatentAugmentedLaw_map_conditional
% lean: CausalSmith.Stat.FinitepHomogeneityDensegamma.copulaLatentAugmentedLaw_sign_independent
% lean: CausalSmith.Stat.FinitepHomogeneityDensegamma.full_record_copula_component_bound
\end{enumerate}
\end{proof}

\begin{algorithmv}{def:converse-handle}[Branchwise converse priors \(\pi_{0,n,v},\pi_{1,n,v}\)]
Inputs are the sample size \(n\) and the public tuple \(v=(p,\alpha,\beta,\gamma)\); \(q=(p-1)/p\), \(F(p)\) is the exponent-comparison scalar, \(D_p\) is the rough-confounding exponent denominator, \(H_{\mathrm{fine}}\) is the fine-rank multiplier, and \(N_{\mathrm{low}}(v)\) is the rough-construction threshold.
\begin{enumerate}
\item If \(F(p)<1\) and \(n\ge N_{\mathrm{low}}(v)\), select
\[
K=2^{\left\lceil\log_2\!\left(H_{\mathrm{fine}}n^{2/D_p}\right)\right\rceil},
\qquad
M=2^{\left\lfloor\log_2\!\left(K^{(\alpha+\beta)/\gamma}/16\right)\right\rfloor},
\]
and set
\[
a=\frac{K^{-\alpha}}{16},
\qquad
u=\frac1{16},
\qquad
\varepsilon_{n,v}=16^{-1/q}K^{-\beta/q},
\qquad
L_{n,v}=\varepsilon_{n,v}^{-1/p}.
\]
Select the finite copula-frame priors
\[
\pi_{\nu,n,v}
=
\pi^{K,M}_{\nu;v,a,u,\varepsilon_{n,v},L_{n,v}},
\qquad \nu\in\{0,1\}.
\]
\item Otherwise, set
\[
N=2\left\lceil n^{2q/(2\gamma+q)}\right\rceil,
\qquad
h=\frac1N,
\qquad
\varepsilon_{n,v}=h^{\gamma/q},
\qquad
L_{n,v}=h^{-\gamma/(p-1)}.
\]
Select the opposite paired-tent null and alternative priors at this rank, cell width, rare-mark probability, and mark magnitude, denoted respectively by
\[
\pi_{0,n,v},\qquad \pi_{1,n,v}.
\]
\item Output the selected priors and their complete independent-record mixtures
\[
\bigl(\pi_{0,n,v},\pi_{1,n,v}\bigr),
\qquad
\mathbb P_{\nu,n,v}
=
\int P^{\otimes n}\,\pi_{\nu,n,v}(dP),
\quad \nu\in\{0,1\}.
\]
The selected construction satisfies the stipulated balancing, support-law membership, true null-mixture denominator, complete occupancy and transition accounting, and bounded-scale comparison properties.
\end{enumerate}
\end{algorithmv}

The common augmentation in \cref{obj:lem:full-record-copula-component-bound} includes the disclosed boundary coefficients, and the coarse signs remain independent of it under both branches. All three conditional label laws are normalized probability laws. The first expected budget has sample-size and rank factor \(n^2/K\); the second has factor \(n^4/(MK^2)\). Both include the amplitude and mark factor \(a^4u^4\varepsilon^2\). Their stated sum condition gives total variation below one eighth for the recovered complete-record mixtures, including the treatment labels and zero outcomes.

\subsection{Selected priors and the rough lower bound}

The selected converse construction in \cref{obj:def:converse-handle} now has its ingredients. For \(F(p)<1\) and sample size at least the public lower-construction threshold, it selects the copula-frame priors at its specified ranks and finite-moment tuning. The other branch uses the paired-tent package. This selection combines the rough-phase construction with witnesses available for every integer sample size at least two.

At moment order two, the bounded companion uses the bounded copula tuning in the large-sample rough branch and the signed-binary paired-tent package otherwise. Let \(\mathcal V\) denote the admissible tuple domain stipulated in \cref{obj:def:model}. The following definition fixes the bounded prior pair and its complete-record mixtures.

\begin{definitionv}{synth_1}[Bounded converse priors]
Fix \(v=(2,\alpha,\beta,\gamma)\in\mathcal V\), write \(w=(\alpha,\beta,\gamma)\), and let \(n\ge2\). Use the phase functional and public constants of \cref{obj:def:sharp-frontier-scales}. For \(\nu\in\{0,1\}\), define the bounded prior pair \(\pi^{\mathrm b}_{\nu,n,v}\) as follows.

If \(F(2)<1\) and \(n\ge N_{\mathrm{low}}(v)\), put \(S=\alpha+\beta\), let \(K\) be the least power of two at least \(H_{\mathrm{fine}}n^{2/D_2}\), and set
\[
 M=2^{\lfloor\log_2(K^{S/\gamma}/16)\rfloor},
 \qquad a=\frac{K^{-\alpha}}{16},
 \qquad b=\frac{K^{-\beta}}{256}.
\]
Define
\[
 \pi^{\mathrm b}_{\nu,n,v}
   =\pi^{K,M}_{\nu;v,a,\,2b,\,1/2,\,1},
 \qquad \nu\in\{0,1\},
\]
using the finite coefficient draws and record-law pushforward of \cref{obj:def:copula-frame-priors}. This is the bounded tuning of \cref{obj:lem:copula-frame-legality}: it uses the same ranks and amplitudes \(a,b\) as the rough branch of \cref{obj:def:converse-handle}, with outcome-coordinate amplitude \(u=2b\), mark probability \(\varepsilon=1/2\), and mark magnitude \(L=1\).

If \(F(2)\ge1\), or if \(F(2)<1\) and \(n<N_{\mathrm{low}}(v)\), take the signed-binary paired-tent package of \cref{obj:lem:paired-tent-full-record-lower} at sample size \(n\) and tuple \(w\). Writing \(P_0^{\mathrm b}\) for its selected null law and \(\pi_1^{\mathrm b}\) for its selected finite alternative prior, define
\[
 \pi^{\mathrm b}_{0,n,v}=\delta_{P_0^{\mathrm b}},
 \qquad
 \pi^{\mathrm b}_{1,n,v}=\pi_1^{\mathrm b}.
\]
This supplies the signed-binary counterpart of the paired-tent branch and small-sample fallback in \cref{obj:def:converse-handle}.

In either branch, the associated complete-record mixtures are
\[
 \mathbb P^{\mathrm b}_{\nu,n,v}
   =\int P^{\otimes n}\,\pi^{\mathrm b}_{\nu,n,v}(dP),
 \qquad \nu\in\{0,1\}.
\]
\end{definitionv}

The bounded prior pair \(\pi^{\mathrm b}_{\nu,n,v}\) in \cref{obj:synth_1} is defined for every admissible moment-order-two tuple and every \(n\ge2\). Its complete-record mixture \(\mathbb P^{\mathrm b}_{\nu,n,v}\) uses the same draw-once convention as \(\mathbb P_{\pi,n}\). The next result assembles support placement, likelihood control, and the testing-risk lower bound in the rough phase.

\begin{lemmav}{lem:rough-full-record-sharp-lower}[Sharp rough-phase record lower bound]
Use the selected converse priors and complete-record mixtures of \cref{obj:def:converse-handle}, and the phase functional, exponents, multipliers, and scales of \cref{obj:def:sharp-frontier-scales}. Let \(\lambda\) denote Lebesgue probability measure on \([0,1]\). Define the heterogeneity distance, model supremum, and witness distance by
\[
d(P)=\left(\int_{[0,1]}
\left(\tau_P(x)-\int_{[0,1]}\tau_P(z)\,d\lambda(z)\right)^2
\,d\lambda(x)\right)^{1/2},
\qquad
D_v=\sup_{P\in\mathcal M_v}d(P),
\qquad
d_0=\frac{1}{8\sqrt2}.
\]
For probability laws, \(\operatorname{TV}\) denotes the supremum of the absolute differences between their probabilities of the same measurable event.

Suppose that:
\begin{itemize}
\item \textbf{(Admissibility.)} \(v=(p,\alpha,\beta,\gamma)\) is admissible as in \cref{obj:def:model}.
\item \textbf{(Rough phase.)} \(F(p)<1\).
\item \textbf{(Sample size.)} \(n\) is an integer with \(n\ge2\).
\end{itemize}
The finite priors \(\pi_{0,n,v}\) and \(\pi_{1,n,v}\) have nonnegative weights summing to one. Every positive-weight null support law belongs to \(H_0(v)\) of \cref{obj:def:null}, and every positive-weight alternative support law satisfies
\[
P\in\mathcal M_v,
\qquad
c_vn^{-E_4(p)}\le d(P)<d_0.
\]
Moreover,
\[
d_0\le D_v,
\qquad
\operatorname{TV}(\mathbb P_{0,n,v},\mathbb P_{1,n,v})<\frac12,
\qquad
B_n\!\left(v,c_vn^{-E_4(p)}\right)\ge\frac25,
\]
with testing risk as in \cref{obj:def:testing-risk}.

If \(n\ge N_{\mathrm{low}}(v)\), the selected fine and coarse ranks \(K,M\) of \cref{obj:def:converse-handle} satisfy the rank conditions of \cref{obj:lem:copula-frame-legality}. In particular, they are powers of two and satisfy
\[
16M\le K,
\qquad
2\le M\le K^{S/\gamma},
\qquad
2^{40}n\le K.
\]
With the selected rare-mark probability \(\varepsilon_{n,v}\) of \cref{obj:def:converse-handle}, the deterministic budget bounds are
\[
\begin{aligned}
\frac{2^{38}(K^{-\alpha}/16)^4(1/16)^4
      \varepsilon_{n,v}^{\,2}n^2}{K}
&\le \frac{2^{-10}}{H_{\mathrm{fine}}},\\
\frac{2^{56}(K^{-\alpha}/16)^4(1/16)^4
      \varepsilon_{n,v}^{\,2}n^4}{MK^2}
&\le \frac{2^{13}}{H_{\mathrm{fine}}^2}.
\end{aligned}
\]
The same complete-record mixtures then satisfy
\[
\operatorname{TV}(\mathbb P_{0,n,v},\mathbb P_{1,n,v})<\frac18.
\]

For the bounded experiment, suppose that:
\begin{itemize}
\item \textbf{(Admissibility.)} \(w=(\alpha,\beta,\gamma)\) is admissible as in \cref{obj:def:bounded-model}; write \(v=(2,\alpha,\beta,\gamma)\).
\item \textbf{(Rough phase.)} \(F(2)<1\).
\item \textbf{(Sample size.)} \(n\) is an integer with \(n\ge2\).
\end{itemize}
Let \(\pi^{\mathrm b}_{0,n,v}\) and \(\pi^{\mathrm b}_{1,n,v}\) be the selected bounded converse priors of \cref{obj:synth_1}, and let
\(\mathbb P^{\mathrm b}_{0,n,v}\) and \(\mathbb P^{\mathrm b}_{1,n,v}\) denote their respective complete-record mixtures. Both finite priors have nonnegative weights summing to one. Every positive-weight null support law belongs to \(H_0^{\mathrm b}(w)\) of \cref{obj:def:bounded-null}, and every positive-weight alternative support law satisfies
\[
P\in\mathcal M^{\mathrm b}_w,
\qquad
c_w^{\mathrm b}n^{-E_4(2)}\le d(P)<d_0.
\]
Furthermore,
\[
\operatorname{TV}
\bigl(\mathbb P^{\mathrm b}_{0,n,v},\mathbb P^{\mathrm b}_{1,n,v}\bigr)
<\frac12,
\qquad
B_n^{\mathrm b}\!\left(w,c_w^{\mathrm b}n^{-E_4(2)}\right)\ge\frac25,
\]
with bounded testing risk as in \cref{obj:def:bounded-risk}.
\end{lemmav}

\begin{proof}[Proof of \cref{obj:lem:rough-full-record-sharp-lower}]
Fix an admissible tuple \(v=(p,\alpha,\beta,\gamma)\) with \(F(p)<1\), and an integer \(n\ge2\). Use the quantities of \cref{obj:def:sharp-frontier-scales}, in particular
\[
q=\frac{p-1}{p},
\qquad S=\alpha+\beta,
\qquad
D_p=1+2\alpha+\frac{\beta}{q}+\frac{S}{2\gamma}.
\]
Admissibility gives
\[
1<p\le2,\qquad \alpha,\beta>0,\qquad \gamma\ge\frac14,
\]
so
\[
0<q\le\frac12,\qquad S>0,\qquad
2\gamma+q>0,\qquad D_p>1.
\]
We first establish the original-model conclusions, splitting at \(N_{\mathrm{low}}(v)\).

\begin{enumerate}
\item \textbf{The small-sample branch.}
Suppose \(n<N_{\mathrm{low}}(v)\). Subtracting the two exponent formulas gives
\[
\begin{aligned}
E_4(p)-E_0(p)
&=
\frac{2S(2\gamma+q)-2\gamma qD_p}
     {(2\gamma+q)D_p}\\
&=
\frac{2\gamma q}{(2\gamma+q)D_p}
\left(
\frac{2S}{q}-2\alpha-\frac{\beta}{q}
+\frac{S}{2\gamma}-1
\right)\\
&=
\frac{2\gamma q}{(2\gamma+q)D_p}\{F(p)-1\}<0.
\end{aligned}
\]
The last identity uses \(q^{-1}-1=(p-1)^{-1}\). Thus \(E_0(p)-E_4(p)>0\). The definition of \(c_v\), together with \(n\le N_{\mathrm{low}}(v)\), yields
\[
\begin{aligned}
c_vn^{-E_4(p)}
&\le
c_0^{\mathrm e}
N_{\mathrm{low}}(v)^{-(E_0(p)-E_4(p))}
n^{-E_4(p)}\\
&\le
c_0^{\mathrm e}
n^{-(E_0(p)-E_4(p))}n^{-E_4(p)}
=
c_0^{\mathrm e}n^{-E_0(p)}.
\end{aligned}
\]

The selected priors are the paired-tent priors of
\cref{obj:def:converse-handle}, at the even rank
\[
N=2\left\lceil n^{2q/(2\gamma+q)}\right\rceil,
\qquad h=\frac1N,
\qquad \varepsilon_{n,v}=h^{\gamma/q},
\qquad L_{n,v}=h^{-\gamma/(p-1)}.
\]
We identify their canonical support laws explicitly. For
\(\sigma\in\{-1,1\}^{N/2}\), define the paired-tent field by
\[
g^{\mathrm{tent}}_\sigma(x)=
\begin{cases}
\sigma_j b_\triangle(Nx-(2j-2)),
&x\in[(2j-2)/N,(2j-1)/N],\\
-\sigma_j b_\triangle(Nx-(2j-1)),
&x\in[(2j-1)/N,2j/N],
\end{cases}
\qquad 1\le j\le N/2,
\]
where \(b_\triangle\) is the triangular tent of
\cref{obj:def:copula-frame-priors}; boundary values are zero.
Let \(P^{\mathrm{tent}}_\sigma\) be the law with uniform covariate,
propensity \(1/2\), and conditional outcome kernels
\[
Q_{0,P^{\mathrm{tent}}_\sigma}(\cdot\mid x)=\delta_0,
\]
\[
\begin{aligned}
Q_{1,P^{\mathrm{tent}}_\sigma}(\cdot\mid x)
={}&(1-\varepsilon_{n,v})\delta_0\\
&+\frac{\varepsilon_{n,v}}2
\left(1+\frac{g^{\mathrm{tent}}_\sigma(x)}{16}\right)\delta_{L_{n,v}}
+\frac{\varepsilon_{n,v}}2
\left(1-\frac{g^{\mathrm{tent}}_\sigma(x)}{16}\right)\delta_{-L_{n,v}}.
\end{aligned}
\]
Here \(\delta_y\) denotes a point mass at \(y\). Define \(P_0\) by the same kernels with the tent field replaced by zero. These are exactly the null and alternative laws selected by the paired-tent branch: its treated mean is
\[
\tau_{P^{\mathrm{tent}}_\sigma}(x)
=\frac{\varepsilon_{n,v}L_{n,v}}{16}g^{\mathrm{tent}}_\sigma(x)
=\frac{h^\gamma}{16}g^{\mathrm{tent}}_\sigma(x),
\qquad \tau_{P_0}(x)=0,
\]
since \(q^{-1}-(p-1)^{-1}=1\). Thus the selected alternative prior and its complete-record mixture are
\[
\pi_{1,n,v}
=2^{-N/2}\sum_{\sigma\in\{-1,1\}^{N/2}}
\delta_{P^{\mathrm{tent}}_\sigma},
\qquad
\overline P_1^{(n)}
=2^{-N/2}\sum_{\sigma\in\{-1,1\}^{N/2}}
(P^{\mathrm{tent}}_\sigma)^{\otimes n}.
\]
Before pushing forward to laws, each of the \(2^{N/2}\) sign configurations has weight \(2^{-N/2}\); hence
\[
2^{-N/2}>0,
\qquad
2^{N/2}2^{-N/2}=1.
\]
Every null sign configuration produces the same zero-effect law \(P_0\) just defined. Consequently,
\[
\pi_{0,n,v}=\delta_{P_0},
\qquad
\mathbb P_{0,n,v}=P_0^{\otimes n},
\qquad
\mathbb P_{1,n,v}=\overline P_1^{(n)}.
\]
Here the equality of priors is as probability measures on laws; the finite sign enumeration may retain repeated copies of \(P_0\).

We now certify these canonical laws directly. Ceiling bounds give
\[
2n^{2q/(2\gamma+q)}\le N\le4n^{2q/(2\gamma+q)},
\qquad N\ge2,
\qquad 0<h\le\frac12.
\]
Hence \(0<\varepsilon_{n,v}<1\) and \(L_{n,v}>0\). The displayed kernels have nonnegative masses summing to one, and their arm moments satisfy
\[
\int |y|^p\,Q_{0,P}(dy\mid x)=0,
\qquad
\int |y|^p\,Q_{1,P}(dy\mid x)
=\varepsilon_{n,v}L_{n,v}^p=1
\]
for \(P=P_0\) and every \(P=P^{\mathrm{tent}}_\sigma\).
The propensity and baseline mean are constant, while
\[
\|g^{\mathrm{tent}}_\sigma\|_\infty\le1,
\qquad
|g^{\mathrm{tent}}_\sigma(x)-g^{\mathrm{tent}}_\sigma(z)|
\le8N|x-z|.
\]
For the latter bound, each paired bump has Lipschitz constant at most \(4N\), and at most two paired bumps contribute at the two endpoints. Combining this bound with the envelope bound gives, for \(0<\gamma\le1\),
\[
h^\gamma
|g^{\mathrm{tent}}_\sigma(x)-g^{\mathrm{tent}}_\sigma(z)|
\le8|x-z|^\gamma.
\]
Indeed, when \(N|x-z|\le1\), use
\(N|x-z|\le(N|x-z|)^\gamma\); when \(N|x-z|>1\), use the envelope bound \(2\le8(N|x-z|)^\gamma\).
Consequently,
\[
\|\tau_{P^{\mathrm{tent}}_\sigma}\|_\infty\le\frac1{16},
\qquad
|\tau_{P^{\mathrm{tent}}_\sigma}(x)
-\tau_{P^{\mathrm{tent}}_\sigma}(z)|
\le\frac12|x-z|^\gamma.
\]
These calculations verify all restrictions of
\cref{obj:def:model,obj:def:null}, so
\[
P_0\in H_0(v),
\qquad P^{\mathrm{tent}}_\sigma\in\mathcal M_v.
\]
Opposite pairing and integration of the squared triangular tent give
\[
\int g^{\mathrm{tent}}_\sigma\,d\lambda=0,
\qquad
\int (g^{\mathrm{tent}}_\sigma)^2\,d\lambda=\frac13.
\]
Thus the distance of each canonical alternative is
\[
d(P^{\mathrm{tent}}_\sigma)
=\frac{h^\gamma}{16\sqrt3}
\ge\frac{4^{-\gamma}}{16\sqrt3}
 n^{-2\gamma q/(2\gamma+q)}
=c_0^{\mathrm e}n^{-E_0(p)},
\qquad
d(P^{\mathrm{tent}}_\sigma)\le\frac1{16}<d_0.
\]

To compare the canonical complete-record mixtures, define the one-record likelihood ratio on the full record space by
\[
\ell^{\mathrm{tent}}_\sigma(X,A,Y)
=1+\frac{g^{\mathrm{tent}}_\sigma(X)}{16}
\begin{cases}
1,&A=1,\ Y=L_{n,v},\\
-1,&A=1,\ Y=-L_{n,v},\\
0,&\text{otherwise}.
\end{cases}
\]
The conditional masses show directly that
\[
\frac{dP^{\mathrm{tent}}_\sigma}{dP_0}
=\ell^{\mathrm{tent}}_\sigma
\quad P_0\text{-almost surely}.
\]
For two sign configurations \(\sigma,\sigma'\), integration over all record categories and then over the paired cells gives
\[
\int g^{\mathrm{tent}}_\sigma
 g^{\mathrm{tent}}_{\sigma'}\,d\lambda
=\frac{2}{3N}\sum_{j=1}^{N/2}\sigma_j\sigma'_j,
\]
\[
\int\ell^{\mathrm{tent}}_\sigma
\ell^{\mathrm{tent}}_{\sigma'}\,dP_0
=1+\frac{\varepsilon_{n,v}}{768N}
\sum_{j=1}^{N/2}\sigma_j\sigma'_j.
\]
Write \(\chi^2(Q\Vert P)=\int(dQ/dP-1)^2\,dP\) for the directed chi-square divergence. Independence of the records and finite-mixture expansion yield
\[
1+\chi^2(\overline P_1^{(n)}\Vert P_0^{\otimes n})
=\mathbb E_{\sigma,\sigma'}
\left[
\left(1+\frac{\varepsilon_{n,v}}{768N}
\sum_{j=1}^{N/2}\sigma_j\sigma'_j\right)^n
\right],
\]
where both sign configurations are independent and uniform.
The expression inside the power is positive. Using
\((1+\vartheta_{\mathrm{exp}})^n\le e^{n\vartheta_{\mathrm{exp}}}\) for \(\vartheta_{\mathrm{exp}}>-1\), independence of the sign products, and
\(\cosh(\vartheta_{\mathrm{exp}})\le e^{\vartheta_{\mathrm{exp}}^2/2}\) for \(\vartheta_{\mathrm{exp}}\in\mathbb R\), we obtain
\[
\begin{aligned}
1+\chi^2(\overline P_1^{(n)}\Vert P_0^{\otimes n})
&\le
\left[\cosh\!\left(\frac{n\varepsilon_{n,v}}{768N}\right)\right]^{N/2}\\
&\le
\exp\!\left(\frac{n^2\varepsilon_{n,v}^2h}{36\cdot16^4}\right).
\end{aligned}
\]
The selected rank calibrates this exponent because
\[
n^2\varepsilon_{n,v}^2h
=n^2N^{-(1+2\gamma/q)}\le1.
\]
Here the lower rank bound implies
\(N\ge n^{2q/(2\gamma+q)}\), and
\(\frac{2q}{2\gamma+q}(1+2\gamma/q)=2\).
Since \(e^{\vartheta_{\mathrm{exp}}}\le(1-\vartheta_{\mathrm{exp}})^{-1}\) for \(0\le\vartheta_{\mathrm{exp}}<1\),
\[
1+\chi^2(\overline P_1^{(n)}\Vert P_0^{\otimes n})
\le\exp\!\left(\frac1{36\cdot16^4}\right)<2.
\]
Cauchy--Schwarz therefore gives
\[
\operatorname{TV}(\mathbb P_{0,n,v},\mathbb P_{1,n,v})
\le\frac12
\sqrt{\chi^2(\overline P_1^{(n)}\Vert P_0^{\otimes n})}
<\frac12.
\]
The model-supremum bound is supplied by
\cref{obj:lem:paired-tent-full-record-lower}:
\[
d_0\le D_v.
\]
Together with the multiplier comparison, these canonical calculations give, for every positive-weight alternative support law,
\[
P\in\mathcal M_v,
\qquad
c_vn^{-E_4(p)}
\le c_0^{\mathrm e}n^{-E_0(p)}
\le d(P)<d_0.
\]
Finally, for every size-valid test \(\phi\), averaging over its independent uniform seed and using the total-variation expectation bound gives
\[
\begin{aligned}
\sup_{\substack{P\in\mathcal M_v\\
d(P)\ge c_0^{\mathrm e}n^{-E_0(p)}}}
\{1-\mathsf R_n(P,\phi)\}
&\ge\int\{1-\mathsf R_n(P,\phi)\}\,\pi_{1,n,v}(dP)\\
&\ge\frac9{10}
-\operatorname{TV}(P_0^{\otimes n},\overline P_1^{(n)})
>\frac25.
\end{aligned}
\]
Taking the infimum over size-valid tests proves the canonical risk bound
\[
B_n\!\left(v,c_0^{\mathrm e}n^{-E_0(p)}\right)\ge\frac25.
\]
% lean: CausalSmith.Stat.FinitepHomogeneityDensegamma.tent_canonical_lower_receipt

For each test satisfying the size constraint, decreasing the separation enlarges the alternative set, so
\[
\sup_{\substack{P\in\mathcal M_v\\
d(P)\ge c_0^{\mathrm e}n^{-E_0(p)}}}
\{1-\mathsf R_n(P,\phi)\}
\le
\sup_{\substack{P\in\mathcal M_v\\
d(P)\ge c_vn^{-E_4(p)}}}
\{1-\mathsf R_n(P,\phi)\}.
\]
The positive-weight alternative supplied by the paired-tent construction makes both sets nonempty. Taking infima over the same size-valid tests and using its risk bound gives
\[
\frac25
\le B_n\!\left(v,c_0^{\mathrm e}n^{-E_0(p)}\right)
\le B_n\!\left(v,c_vn^{-E_4(p)}\right).
\]
Finally, \(n<N_{\mathrm{low}}(v)\) excludes the hypothesis
\(n\ge N_{\mathrm{low}}(v)\) of the additional rank and activity conclusions.
% lean: CausalSmith.Stat.FinitepHomogeneityDensegamma.rough_small_sample_lower

\item \textbf{Ranks in the large-sample branch.}
Suppose now \(n\ge N_{\mathrm{low}}(v)\). Define the selected fine and coarse ranks by
\[
K=
2^{\left\lceil\log_2(H_{\mathrm{fine}}n^{2/D_p})\right\rceil},
\qquad
M=
2^{\left\lfloor\log_2(K^{S/\gamma}/16)\right\rfloor}.
\]
Since \(p-1\le1\) and \(q\le1/2\),
\[
2S+\frac{S}{2\gamma}\le F(p)<1,
\qquad
S<\frac{2\gamma}{4\gamma+1}\le\gamma.
\]
Also,
\[
D_p-1
=
2\alpha+\frac{\beta}{q}+\frac{S}{2\gamma}
\le F(p)<1,
\qquad
1<D_p<2.
\]
In particular \(2/D_p\ge1\), and therefore \(n^{2/D_p}\ge n\).
The ceiling inequalities
\[
\log_2(H_{\mathrm{fine}}n^{2/D_p})
\le
\left\lceil\log_2(H_{\mathrm{fine}}n^{2/D_p})\right\rceil
\le
\log_2(H_{\mathrm{fine}}n^{2/D_p})+1
\]
give
\[
H_{\mathrm{fine}}n^{2/D_p}\le K
\le2H_{\mathrm{fine}}n^{2/D_p},
\qquad
2^{40}n\le K.
\]

The definition of the threshold gives
\[
n\ge32^{D_p\gamma/(2S)},
\qquad
K^{S/\gamma}\ge n^{2S/(D_p\gamma)}\ge32.
\]
Thus the coarse target \(K^{S/\gamma}/16\) is at least \(2\).
Applying the floor inequalities to its base-two logarithm gives
\[
\frac{K^{S/\gamma}}{32}\le M
\le\frac{K^{S/\gamma}}{16},
\qquad M\ge2.
\]
Both ranks are powers of two. Since \(K\ge1\) and \(S/\gamma<1\),
\[
16M\le K^{S/\gamma}\le K,
\qquad
M\le K^{S/\gamma}.
\]
These are all the required rank-legality conditions.
% lean: CausalSmith.Stat.FinitepHomogeneityDensegamma.rough_rank_legality

\item \textbf{Normalization, support, and separation.}
At these ranks, define the legality tuning and its copula priors by
\[
a_{\mathrm{leg}}=\frac{K^{-\alpha}}{16},
\qquad
b_{\mathrm{leg}}=\frac{K^{-\beta}}{256},
\qquad
u=\frac1{16},
\qquad
\varepsilon=(16b_{\mathrm{leg}})^{1/q},
\qquad
L_{\mathrm{mark}}=\varepsilon^{-1/p},
\]
\[
\pi_\nu=\pi^{K,M}_{\nu;v,a_{\mathrm{leg}},u,\varepsilon,L_{\mathrm{mark}}},
\qquad
\mathbb P_\nu=\int P^{\otimes n}\,\pi_\nu(dP),
\qquad \nu\in\{0,1\}.
\]
Because \(K\ge1\) and \(\alpha,\beta>0\),
\[
0<a_{\mathrm{leg}}\le\frac1{16},
\qquad
0<b_{\mathrm{leg}}\le\frac1{256},
\qquad
0<u\le\frac1{16},
\qquad
0<\varepsilon<1,
\qquad L_{\mathrm{mark}}>0.
\]

For completeness, the coefficient weights in
\cref{obj:def:copula-frame-priors} are normalized as follows.
The triangular tent satisfies
\[
0\le b_\triangle(t_{\triangle})=2\min\{t_{\triangle},1-t_{\triangle}\}\le1
\qquad(t_{\triangle}\in[0,1]),
\]
so \(|g_\sigma(x)|\le1\) on \([0,1]\). With \(\kappa=1/16\), each node-pair probability satisfies
\[
\frac{1+\nu\kappa g_\sigma(i/K)\lambda_i\eta_i}{4}
\ge\frac{1-\kappa}{4}>0,
\]
and cancellation of the sign product gives
\[
\sum_{(\lambda_i,\eta_i)\in\{-1,1\}^2}
\frac{1+\nu\kappa g_\sigma(i/K)\lambda_i\eta_i}{4}
=1.
\]
Conditional independence therefore makes the sum of the fine-pair product weights equal to one for every coarse-sign configuration. Summing over the fair coarse signs then gives
\[
\sum_{\sigma\in\{-1,1\}^{M/2}}2^{-M/2}=1.
\]
Thus both finite coefficient distributions, and their pushforwards \(\pi_0,\pi_1\), have nonnegative weights summing to one.
% lean: CausalSmith.Stat.FinitepHomogeneityDensegamma.copula_prior_normalized

The rank conditions established above allow us to apply
\cref{obj:lem:copula-frame-legality}. Every null support law belongs to \(H_0(v)\), and every alternative support law satisfies
\[
P\in\mathcal M_v,
\qquad
2^{-17}K^{-S}\le d(P)<d_0.
\]
The upper rounding bound and \(S>0\) imply
\[
\begin{aligned}
c_vn^{-E_4(p)}
&\le k_vn^{-E_4(p)}\\
&=2^{-17}(2H_{\mathrm{fine}})^{-S}n^{-2S/D_p}\\
&=2^{-17}(2H_{\mathrm{fine}}n^{2/D_p})^{-S}\\
&\le2^{-17}K^{-S}\le d(P).
\end{aligned}
\]
This proves the required alternative placement for the copula pair.
% lean: CausalSmith.Stat.FinitepHomogeneityDensegamma.rough_rank_separation

\item \textbf{Activity calibration.}
Define the deterministic first- and second-component activity bounds by
\[
\mathcal A_{\mathrm{det}}=
\frac{2^{38}a_{\mathrm{leg}}^4u^4\varepsilon^2n^2}{K},
\qquad
\mathcal B_{\mathrm{det}}=
\frac{2^{56}a_{\mathrm{leg}}^4u^4\varepsilon^2n^4}{MK^2}.
\]
In \cref{obj:lem:full-record-copula-component-bound}, let \(\mu\) be the common augmentation law and write \(\mathcal A_{\mathrm{cond}},\mathcal B_{\mathrm{cond}}\) for its conditional chi-square exponent budgets. The expectation bounds supplied there are
\[
\int\mathcal A_{\mathrm{cond}}\,d\mu
\le\mathcal A_{\mathrm{det}},
\qquad
\int\mathcal B_{\mathrm{cond}}\,d\mu
\le\mathcal B_{\mathrm{det}}.
\]
We calibrate the deterministic quantities on the right.

First, raising the lower fine-rank bound to \(D_p>0\) gives
\[
K^{D_p}\ge H_{\mathrm{fine}}^{D_p}n^2
\ge H_{\mathrm{fine}}n^2.
\]
Consequently,
\[
n^2K^{-D_p}\le\frac1{H_{\mathrm{fine}}},
\qquad
n^4K^{-2D_p}
=\bigl(n^2K^{-D_p}\bigr)^2
\le\frac1{H_{\mathrm{fine}}^2}.
\]

The tuning gives the exact amplitude product
\[
\varepsilon=16^{-1/q}K^{-\beta/q},
\qquad
a_{\mathrm{leg}}^4u^4\varepsilon^2
=
16^{-8-2/q}K^{-(4\alpha+2\beta/q)}
\le
2^{-48}K^{-(4\alpha+2\beta/q)},
\]
where the last inequality follows from \(q\le1/2\).
Furthermore,
\[
\frac{S}{2\gamma}\le2S
=2\alpha+2\beta
\le2\alpha+\frac{\beta}{q},
\]
and hence
\[
D_p\le1+4\alpha+\frac{2\beta}{q}.
\]
As \(K\ge1\), these inequalities yield
\[
\begin{aligned}
\mathcal A_{\mathrm{det}}
&\le
2^{-10}n^2K^{-(1+4\alpha+2\beta/q)}\\
&\le2^{-10}n^2K^{-D_p}
\le\frac{2^{-10}}{H_{\mathrm{fine}}}.
\end{aligned}
\]
For the second budget, use the lower coarse-rank bound:
\[
\begin{aligned}
\mathcal B_{\mathrm{det}}
&\le
\frac{2^8n^4K^{-(2+4\alpha+2\beta/q)}}{M}\\
&\le
2^{13}n^4
K^{-(2+4\alpha+2\beta/q+S/\gamma)}\\
&=
2^{13}n^4K^{-2D_p}
\le\frac{2^{13}}{H_{\mathrm{fine}}^2}.
\end{aligned}
\]
The exponent identity in the last line is
\[
2D_p=2+4\alpha+\frac{2\beta}{q}+\frac{S}{\gamma}.
\]
Since \(H_{\mathrm{fine}}=2^{40}\),
\[
\mathcal A_{\mathrm{det}}+\mathcal B_{\mathrm{det}}
\le
\frac{2^{-10}}{H_{\mathrm{fine}}}
+\frac{2^{13}}{H_{\mathrm{fine}}^2}
=
2^{-50}+2^{-67}<2^{-16}.
\]
All sample-size, rank, amplitude, and mark hypotheses of
\cref{obj:lem:full-record-copula-component-bound} have now been discharged. Its complete-record comparison gives
\[
\operatorname{TV}(\mathbb P_0,\mathbb P_1)
<\frac18<\frac12.
\]
% lean: CausalSmith.Stat.FinitepHomogeneityDensegamma.tail_rough_activity_budgets

\item \textbf{Testing risk and identification of the selected pair.}
Let \(\phi\) be any randomized test satisfying the size constraint in
\cref{obj:def:testing-risk}. Average over the independent public uniform seed \(U\) of
\cref{obj:def:original-record-experiment}, defining
\[
\overline\phi:
\bigl([0,1]\times\{0,1\}\times\mathbb R\bigr)^n
\longrightarrow[0,1],
\qquad
\overline\phi(\mathcal D_n)
=
\mathbb E_U[\phi(\mathcal D_n,U)].
\]
Finite-mixture linearity and null support give
\[
\int\overline\phi\,d\mathbb P_0
=
\int\mathsf R_n(P,\phi)\,\pi_0(dP)
\le\frac1{10}.
\]
For a measurable function taking values in \([0,1]\), the expectation difference is bounded by total variation. Thus
\[
\left|
\int\overline\phi\,d\mathbb P_1
-
\int\overline\phi\,d\mathbb P_0
\right|
\le\operatorname{TV}(\mathbb P_0,\mathbb P_1).
\]
Using alternative support and normalization,
\[
\begin{aligned}
\sup_{\substack{P\in\mathcal M_v\\
d(P)\ge c_vn^{-E_4(p)}}}
\{1-\mathsf R_n(P,\phi)\}
&\ge
\int\{1-\mathsf R_n(P,\phi)\}\,\pi_1(dP)\\
&=
1-\int\overline\phi\,d\mathbb P_1\\
&\ge
\frac9{10}-\operatorname{TV}(\mathbb P_0,\mathbb P_1).
\end{aligned}
\]
Normalization ensures that at least one alternative weight is positive, so the alternative set is nonempty; zero-weight terms contribute zero. The zero test satisfies the size constraint, so the infimum is also over a nonempty set. Taking that infimum gives
\[
B_n\!\left(v,c_vn^{-E_4(p)}\right)
\ge
\frac9{10}-\operatorname{TV}(\mathbb P_0,\mathbb P_1)
>\frac25.
\]
% lean: CausalSmith.Stat.FinitepHomogeneityDensegamma.finite_two_prior_testing_lower

It remains to identify this copula pair with the selected pair. In the present branch, the formulas in
\cref{obj:def:converse-handle} satisfy
\[
\varepsilon_{n,v}
=
16^{-1/q}K^{-\beta/q}
=
\left(16\frac{K^{-\beta}}{256}\right)^{1/q}
=\varepsilon,
\qquad
L_{n,v}=\varepsilon_{n,v}^{-1/p}=L_{\mathrm{mark}}.
\]
Therefore
\[
\pi_{\nu,n,v}=\pi_\nu,
\qquad
\mathbb P_{\nu,n,v}=\mathbb P_\nu,
\qquad \nu\in\{0,1\}.
\]
For the selected pair, the conditional budgets supplied by
\cref{obj:lem:full-record-copula-component-bound} consequently satisfy
\[
\int\mathcal A_{\mathrm{cond}}\,d\mu
\le\mathcal A_{\mathrm{det}}
\le\frac{2^{-10}}{H_{\mathrm{fine}}},
\qquad
\int\mathcal B_{\mathrm{cond}}\,d\mu
\le\mathcal B_{\mathrm{det}}
\le\frac{2^{13}}{H_{\mathrm{fine}}^2}.
\]
The calibrated activity bounds apply to the deterministic right-hand sides of the component expectation estimates.
Finally, the model witness supplied by
\cref{obj:lem:paired-tent-full-record-lower} gives
\[
d_0\le D_v.
\]
Together with the preceding calculations, this establishes normalization, support, separation, the model-supremum bound, testing risk, total variation, rank legality, the fine-rank bound, and the deterministic activity calibration in the large-sample branch.
% lean: CausalSmith.Stat.FinitepHomogeneityDensegamma.tail_rough_large_sample_lower

\item \textbf{The bounded small-sample branch.}
Fix an admissible \(w=(\alpha,\beta,\gamma)\), put
\[
v=(2,\alpha,\beta,\gamma),
\qquad
c_w^{\mathrm b}=c_v,
\qquad
E_0(2)=\frac{2\gamma}{4\gamma+1},
\]
and suppose \(F(2)<1\). This \(v\) is admissible. If \(n<N_{\mathrm{low}}(v)\), the exponent and multiplier calculation at \(p=2\) gives
\[
c_w^{\mathrm b}n^{-E_4(2)}
\le c_0^{\mathrm e}n^{-E_0(2)}.
\]
The bounded fallback in \cref{obj:synth_1} selects the signed-binary paired-tent package of
\cref{obj:lem:paired-tent-full-record-lower}:
\[
\pi^{\mathrm b}_{0,n,v}=\delta_{P_0^{\mathrm b}},
\qquad
\pi^{\mathrm b}_{1,n,v}=\pi_1^{\mathrm b},
\]
\[
\mathbb P^{\mathrm b}_{0,n,v}
=(P_0^{\mathrm b})^{\otimes n},
\qquad
\mathbb P^{\mathrm b}_{1,n,v}
=\overline P_{1,\mathrm b}^{(n)}.
\]
The point mass and the finite alternative prior are normalized. Signed-binary support implies \(|Y|\le1\), so their support laws satisfy
\[
P_0^{\mathrm b}\in H_0^{\mathrm b}(w),
\]
and, on every positive-weight alternative,
\[
P\in\mathcal M_w^{\mathrm b},
\qquad
c_w^{\mathrm b}n^{-E_4(2)}
\le c_0^{\mathrm e}n^{-E_0(2)}
\le d(P)<d_0.
\]
The same paired-tent bound supplies
\[
\operatorname{TV}
\bigl(\mathbb P^{\mathrm b}_{0,n,v},
      \mathbb P^{\mathrm b}_{1,n,v}\bigr)<\frac12,
\qquad
B_n^{\mathrm b}\!\left(w,c_0^{\mathrm e}n^{-E_0(2)}\right)
\ge\frac25.
\]
For every bounded-null size-valid test, the alternative-set inclusion gives
\[
\sup_{\substack{P\in\mathcal M_w^{\mathrm b}\\
d(P)\ge c_0^{\mathrm e}n^{-E_0(2)}}}
\{1-\mathsf R_n(P,\phi)\}
\le
\sup_{\substack{P\in\mathcal M_w^{\mathrm b}\\
d(P)\ge c_w^{\mathrm b}n^{-E_4(2)}}}
\{1-\mathsf R_n(P,\phi)\}.
\]
The paired-tent prior has a positive weight, ensuring nonempty alternative sets. Taking infima yields
\[
B_n^{\mathrm b}\!\left(w,c_w^{\mathrm b}n^{-E_4(2)}\right)
\ge
B_n^{\mathrm b}\!\left(w,c_0^{\mathrm e}n^{-E_0(2)}\right)
\ge\frac25.
\]
% lean: CausalSmith.Stat.FinitepHomogeneityDensegamma.bounded_rough_small_sample_lower

\item \textbf{The bounded large-sample branch.}
Suppose \(n\ge N_{\mathrm{low}}(v)\). Evaluate the rank formulas and amplitudes \(a_{\mathrm{leg}},b_{\mathrm{leg}}\) at \(v=(2,\alpha,\beta,\gamma)\). The rank calculation gives
\[
K=
2^{\left\lceil\log_2(H_{\mathrm{fine}}n^{2/D_2})\right\rceil},
\qquad
M=
2^{\left\lfloor\log_2(K^{S/\gamma}/16)\right\rfloor},
\]
\[
H_{\mathrm{fine}}n^{2/D_2}\le K
\le2H_{\mathrm{fine}}n^{2/D_2},
\qquad
2^{40}n\le K,
\]
\[
\frac{K^{S/\gamma}}{32}\le M
\le\frac{K^{S/\gamma}}{16},
\qquad
M\ge2,\qquad16M\le K,\qquad M\le K^{S/\gamma}.
\]
Use the separate bounded tuning
\[
a_{\mathrm{leg}}=\frac{K^{-\alpha}}{16},
\qquad
b_{\mathrm{leg}}=\frac{K^{-\beta}}{256},
\qquad
u_{\mathrm b}=2b_{\mathrm{leg}}=\frac{K^{-\beta}}{128},
\qquad
\varepsilon_{\mathrm b}=\frac12,
\qquad
L_{\mathrm b}=1,
\]
and define its priors and complete-record mixtures by
\[
\pi_\nu^{\mathrm{b,frame}}
=
\pi^{K,M}_{\nu;v,a_{\mathrm{leg}},u_{\mathrm b},\varepsilon_{\mathrm b},L_{\mathrm b}},
\qquad
\mathbb P_\nu^{\mathrm{b,frame}}
=
\int P^{\otimes n}\,\pi_\nu^{\mathrm{b,frame}}(dP),
\qquad \nu\in\{0,1\}.
\]
Here
\[
0<a_{\mathrm{leg}}\le\frac1{16},
\qquad
0<u_{\mathrm b}\le\frac1{128}\le\frac1{16},
\qquad
0<\varepsilon_{\mathrm b}<1,
\qquad L_{\mathrm b}>0.
\]
The coefficient weights depend on the ranks and the copula coupling, so the product-weight calculation proves normalization for both bounded priors as well. The bounded clause of
\cref{obj:lem:copula-frame-legality} supplies
\[
P\in H_0^{\mathrm b}(w)
\quad\text{on null support},
\]
\[
P\in\mathcal M_w^{\mathrm b},
\qquad
2^{-17}K^{-S}\le d(P)<d_0
\quad\text{on alternative support}.
\]
The upper fine-rank bound gives
\[
c_w^{\mathrm b}n^{-E_4(2)}
\le k_vn^{-E_4(2)}
=
2^{-17}(2H_{\mathrm{fine}}n^{2/D_2})^{-S}
\le2^{-17}K^{-S}.
\]

For this tuning, the exact amplitude product is
\[
a_{\mathrm{leg}}^4u_{\mathrm b}^4\varepsilon_{\mathrm b}^2
=2^{-46}K^{-4S}.
\]
Thus the deterministic component budgets, denoted by
\(\mathcal A_{\mathrm b}\) and \(\mathcal B_{\mathrm b}\), are
\[
\mathcal A_{\mathrm b}
=2^{-8}n^2K^{-(1+4S)},
\qquad
\mathcal B_{\mathrm b}
=\frac{2^{10}n^4K^{-(2+4S)}}{M}.
\]
At \(p=2\),
\[
D_2=1+2S+\frac{S}{2\gamma}
\le1+4S,
\]
because \(\gamma\ge1/4\). Raising the fine-rank bound to \(D_2\), and then squaring the resulting reciprocal bound, gives
\[
n^2K^{-D_2}\le\frac1{H_{\mathrm{fine}}},
\qquad
n^4K^{-2D_2}\le\frac1{H_{\mathrm{fine}}^2}.
\]
It follows that
\[
\mathcal A_{\mathrm b}
\le2^{-8}n^2K^{-D_2}
\le\frac{2^{-8}}{H_{\mathrm{fine}}},
\]
and, using \(M\ge K^{S/\gamma}/32\),
\[
\begin{aligned}
\mathcal B_{\mathrm b}
&\le
2^{15}n^4K^{-(2+4S+S/\gamma)}\\
&=
2^{15}n^4K^{-2D_2}
\le\frac{2^{15}}{H_{\mathrm{fine}}^2}.
\end{aligned}
\]
Consequently,
\[
\mathcal A_{\mathrm b}+\mathcal B_{\mathrm b}
\le
\frac{2^{-8}}{H_{\mathrm{fine}}}
+\frac{2^{15}}{H_{\mathrm{fine}}^2}
=
2^{-48}+2^{-65}<2^{-16}.
\]
Applying \cref{obj:lem:full-record-copula-component-bound}, with the verified bounded parameters, gives
\[
\operatorname{TV}
\bigl(\mathbb P_0^{\mathrm{b,frame}},
      \mathbb P_1^{\mathrm{b,frame}}\bigr)
<\frac18<\frac12.
\]
% lean: CausalSmith.Stat.FinitepHomogeneityDensegamma.bounded_rough_activity_budgets

For a test satisfying the bounded-null size constraint, seed averaging, null support, and the total-variation expectation bound give
\[
\int\overline\phi\,d\mathbb P_0^{\mathrm{b,frame}}\le\frac1{10},
\]
\[
\begin{aligned}
\sup_{\substack{P\in\mathcal M_w^{\mathrm b}\\
d(P)\ge c_w^{\mathrm b}n^{-E_4(2)}}}
\{1-\mathsf R_n(P,\phi)\}
&\ge
1-\int\overline\phi\,d\mathbb P_1^{\mathrm{b,frame}}\\
&\ge
\frac9{10}
-\operatorname{TV}
\bigl(\mathbb P_0^{\mathrm{b,frame}},
      \mathbb P_1^{\mathrm{b,frame}}\bigr).
\end{aligned}
\]
Both priors are normalized, so their alternative support is nonempty; the zero test again satisfies the size constraint. Taking the infimum therefore yields
\[
B_n^{\mathrm b}\!\left(w,c_w^{\mathrm b}n^{-E_4(2)}\right)
\ge
\frac9{10}
-\operatorname{TV}
\bigl(\mathbb P_0^{\mathrm{b,frame}},
      \mathbb P_1^{\mathrm{b,frame}}\bigr)
>\frac25.
\]
Finally, the selection in \cref{obj:synth_1} identifies
\[
\pi^{\mathrm b}_{\nu,n,v}=\pi_\nu^{\mathrm{b,frame}},
\qquad
\mathbb P^{\mathrm b}_{\nu,n,v}
=\mathbb P_\nu^{\mathrm{b,frame}},
\qquad \nu\in\{0,1\}.
\]
This transfers normalization, support, separation, total variation, and the testing-risk bound to the selected bounded pair, completing the bounded clause.
% lean: CausalSmith.Stat.FinitepHomogeneityDensegamma.bounded_rough_large_sample_lower
\end{enumerate}
% lean: CausalSmith.Stat.FinitepHomogeneityDensegamma.rough_full_record_sharp_lower
\end{proof}

Under \(F(p)<1\), \cref{obj:lem:rough-full-record-sharp-lower} gives the rough separation exponent for every \(n\ge2\). Above the public lower-construction threshold, the selected ranks also meet the component-comparison conditions and give the stronger total-variation bound of one eighth. The deterministic budget bounds control the upper bounds on the expectations of \(\mathcal A\) and \(\mathcal B\) in \cref{obj:lem:full-record-copula-component-bound}. The bounded conclusion uses its own outcome tuning at \(p=2\), preserving the rough exponent and the finite-sample risk guarantee.

\subsection{Tail essentiality and the moment-order-two face}

The final result combines the rough and paired-tent branches with attainment. Its bounded-prior comparison quantifies over normalized finite priors in the convention fixed above: all support restrictions apply to positive-weight laws, and each mixture draws its latent law once for the complete sample. For each fixed admissible tuple with moment order below two, the comparison has a tuple-dependent eventual threshold. On the moment-order-two face, the selected bounded constructions give witnesses for every \(n\ge2\).

\begin{theoremv}{thm:tail-essential-full-record-converse}[Tail essentiality and bounded witnesses]
Fix an admissible tuple \(v=(p,\alpha,\beta,\gamma)\) as in \cref{obj:def:model}, and write \(w=(\alpha,\beta,\gamma)\). Use the scales, constants, thresholds, and total tests of \cref{obj:def:sharp-frontier-scales}, and the selected priors, mark magnitudes, and complete-record mixtures of \cref{obj:def:converse-handle}.

Let \(\lambda\) denote Lebesgue probability measure on \([0,1]\). Define the heterogeneity distance and its model supremum by
\[
d(P)=
\left\{\int_{[0,1]}
\left(\tau_P(x)-\int_{[0,1]}\tau_P(z)\,\lambda(dz)\right)^2
\,\lambda(dx)\right\}^{1/2},
\qquad
D_v=\sup_{P\in\mathcal M_v}d(P).
\]
For probability laws, \(\operatorname{TV}\) denotes the supremum of the absolute differences of their probabilities over measurable events.

A normalized finite prior on original-record laws is
\[
\pi=\sum_{j=1}^{k}\omega_j\delta_{P_j},
\qquad
k\ge1,\quad \omega_j\ge0,\quad \sum_{j=1}^{k}\omega_j=1.
\]
Its positive-weight support and complete-record mixture are defined by
\[
\operatorname{supp}_+(\pi)=\{P_j:\omega_j>0\},
\qquad
\mathbb P_{\pi,n}=\sum_{j=1}^{k}\omega_jP_j^{\otimes n}.
\]

For every integer \(n\ge2\), both \(\pi_{0,n,v}\) and \(\pi_{1,n,v}\) are normalized finite priors, and
\[
\begin{gathered}
\operatorname{supp}_+(\pi_{0,n,v})\subseteq H_0(v),\\
\operatorname{supp}_+(\pi_{1,n,v})
\subseteq
\{P\in\mathcal M_v:c_v\rho_n(v)\le d(P)<D_v\},\\
0<c_v\rho_n(v)<D_v,\\
\operatorname{TV}(\mathbb P_{0,n,v},\mathbb P_{1,n,v})<\frac12,
\qquad
B_n(v,c_v\rho_n(v))\ge\frac25,
\end{gathered}
\]
where the classes and testing risk are those of
\cref{obj:def:null,obj:def:model,obj:def:testing-risk}.
For every positive-weight law in either prior, each conditional outcome arm assigns probability one to
\(\{0,-L_{n,v},L_{n,v}\}\) for \(\lambda\)-almost every covariate value. Moreover,
\[
L_{n,v}\longrightarrow\infty
\qquad(n\longrightarrow\infty).
\]

Use the matched bounded model \(\mathcal M^{\mathrm b}_w\) and null
\(H_0^{\mathrm b}(w)\) of
\cref{obj:def:bounded-model,obj:def:bounded-null}.
If \(p<2\), then
\[
\frac{\rho_n(v)}{\rho_n^{\mathrm b}(w)}
\longrightarrow\infty
\qquad(n\longrightarrow\infty).
\]
For each such fixed \(v\), there exists an integer \(N(v)\ge2\) such that, for every integer \(n\ge N(v)\) and every pair of normalized finite priors \(\pi_0,\pi_1\),
\[
\left.
\begin{gathered}
\operatorname{supp}_+(\pi_0)\subseteq H_0^{\mathrm b}(w),\\
\operatorname{supp}_+(\pi_1)\subseteq
\{P\in\mathcal M^{\mathrm b}_w:c_v\rho_n(v)\le d(P)\}
\end{gathered}
\right\}
\quad\Longrightarrow\quad
\operatorname{TV}(\mathbb P_{\pi_0,n},\mathbb P_{\pi_1,n})
\ge\frac12.
\]

If \(p=2\), the following conclusions hold for every integer \(n\ge2\).
Let \(\pi^{\mathrm b}_{\nu,n,v}\), for \(\nu\in\{0,1\}\), denote the selected bounded witness priors: the bounded rough priors of
\cref{obj:lem:rough-full-record-sharp-lower} when
\(F(2)<1\) and \(n\ge N_{\mathrm{low}}(v)\), and the signed-binary paired-tent priors of
\cref{obj:lem:paired-tent-full-record-lower} otherwise. Define
\[
\mathbb P^{\mathrm b}_{\nu,n,v}
=\mathbb P_{\pi^{\mathrm b}_{\nu,n,v},n},
\qquad \nu\in\{0,1\}.
\]
Both priors are normalized and have nonempty positive-weight supports. Their positive-weight masses satisfy
\[
\sum_{P\in\operatorname{supp}_+(\pi^{\mathrm b}_{\nu,n,v})}
\pi^{\mathrm b}_{\nu,n,v}(\{P\})=1,
\qquad \nu\in\{0,1\}.
\]
Let \(d_0>0\) be the fixed witness distance of
\cref{obj:lem:causal-nonempty}. The support, total-variation, and risk conclusions are
\[
\begin{gathered}
\operatorname{supp}_+(\pi^{\mathrm b}_{0,n,v})
\subseteq H_0^{\mathrm b}(w),
\qquad
\operatorname{supp}_+(\pi^{\mathrm b}_{0,n,v})
\subseteq H_0(v),\\
\operatorname{supp}_+(\pi^{\mathrm b}_{1,n,v})
\subseteq
\{P\in\mathcal M^{\mathrm b}_w:c_v\rho_n(v)\le d(P)<d_0\},\\
\operatorname{supp}_+(\pi^{\mathrm b}_{1,n,v})
\subseteq
\{P\in\mathcal M_v:c_v\rho_n(v)\le d(P)\},\\
\operatorname{TV}(\mathbb P^{\mathrm b}_{0,n,v},
                  \mathbb P^{\mathrm b}_{1,n,v})<\frac12,\\
\frac25\le B_n^{\mathrm b}(w,c_v\rho_n(v))
\le B_n(v,c_v\rho_n(v)),
\end{gathered}
\]
with bounded risk defined in \cref{obj:def:bounded-risk}.

Define the bounded-model distance supremum by
\[
D_w^{\mathrm b}=\sup_{P\in\mathcal M^{\mathrm b}_w}d(P).
\]
On this face, the scales, multipliers, thresholds, and distance suprema agree:
\[
\begin{gathered}
\rho_n(v)=\rho_n^{\mathrm b}(w),
\qquad c_v=c_w^{\mathrm b},
\qquad C_v=C_w^{\mathrm b},\\
N_v=N_w^{\mathrm b},
\qquad D_v=D_w^{\mathrm b}.
\end{gathered}
\]
The capped critical radii of
\cref{obj:def:critical-radius,obj:def:bounded-radius,obj:def:binary-radius}
satisfy
\[
\begin{gathered}
c_v\rho_n(v)\le r_n^*(v)\le C_v\rho_n(v),\\
c_w^{\mathrm b}\rho_n^{\mathrm b}(w)
\le r_n^{*,\mathrm b}(w)=r_n^{*,\mathrm{bin}}(w)
\le C_w^{\mathrm b}\rho_n^{\mathrm b}(w).
\end{gathered}
\]

Let \(\mathsf R_n(P,\phi)\) denote the expected rejection probability under \(n\) independent original records and the public randomization of
\cref{obj:def:original-record-experiment}. The total tests satisfy, for every law in the indicated null class,
\[
\begin{aligned}
\mathsf R_n(P,\phi^*_{n,v})&\le0.0075
&&\quad(P\in H_0(v)),\\
\mathsf R_n(P,\phi^{*,\mathrm b}_{n,w})&\le0.0075
&&\quad(P\in H_0^{\mathrm b}(w)).
\end{aligned}
\]
Whenever \(C_v\rho_n(v)<D_v\), both separated alternative classes are nonempty:
\[
\begin{gathered}
\{P\in\mathcal M_v:C_v\rho_n(v)\le d(P)\}\ne\varnothing,\\
\{P\in\mathcal M^{\mathrm b}_w:
C_w^{\mathrm b}\rho_n^{\mathrm b}(w)\le d(P)\}\ne\varnothing.
\end{gathered}
\]
Under this same condition, every law in the corresponding separated class satisfies
\[
\begin{aligned}
1-\mathsf R_n(P,\phi^*_{n,v})&\le0.0075
&&\quad(P\in\mathcal M_v,\ C_v\rho_n(v)\le d(P)),\\
1-\mathsf R_n(P,\phi^{*,\mathrm b}_{n,w})&\le0.0075
&&\quad(P\in\mathcal M^{\mathrm b}_w,\\
C_w^{\mathrm b}\rho_n^{\mathrm b}(w)\le d(P)).
\end{aligned}
\]
For \(n\ge N_v\), the common upper separation satisfies
\[
\begin{gathered}
C_v\rho_n(v)\le\frac{d_0}{2},\\
C_v\rho_n(v)<D_v,
\qquad
C_w^{\mathrm b}\rho_n^{\mathrm b}(w)<D_w^{\mathrm b}.
\end{gathered}
\]
At \(n=2,3\), both total tests have zero rejection probability at every sample:
\[
\phi^*_{n,v}(z)=0,
\qquad
\phi^{*,\mathrm b}_{n,w}(z)=0
\qquad\text{for every }z.
\]
\end{theoremv}

The distinction in \cref{obj:thm:tail-essential-full-record-converse} concerns the outcome distributions that sustain close full-record comparisons at the finite-moment separation scale. The selected finite-moment witnesses have growing mark magnitudes and total variation below one half for every \(n\ge2\). For each fixed tuple with \(p<2\), the bounded scale is asymptotically smaller, and every normalized finite bounded-prior pair satisfying the stated support restrictions has total variation at least one half once \(n\ge N(v)\). This threshold depends on the fixed tuple.

At \(p=2\), the bounded copula construction and the signed-binary paired-tent construction together supply normalized, nonempty bounded witnesses for every sample size. The common scales and constants then connect these lower comparisons to the attained upper bounds. The power conclusions apply when the common upper separation is below the model distance supremum, a condition guaranteed by the public threshold \(N_v\); the capped-radius conclusions retain the conventions of \cref{obj:def:critical-radius,obj:def:bounded-radius,obj:def:binary-radius}.


\section{Supplied-propensity proofs and comparison algebra}

The supplied-propensity results combine inverse weighting, removal of the constant effect, and a quadratic comparison of independent blocks. The construction in \cref{obj:def:oracle-handle} applies over the continuous-carrier class of \cref{obj:def:oracle-broad-model}. Its upper bounds and the common-propensity lower witnesses give the matching separation guarantees in \cref{obj:thm:oracle-frontier-broad-class,obj:thm:oracle-frontier-resolved}. The comparison algebra then distinguishes the value of supplied assignment information on the original model from the cost of outcome tails at matched primitive smoothness.

\subsection{Inverse weighting and the public error bounds}

Fix an admissible tuple \(v=(p,\alpha,\beta,\gamma)\) and a law \(P\in\mathcal M_v^{\mathrm{e,br}}\), as defined in \cref{obj:def:oracle-broad-model}. Throughout this section, \(\lambda\) denotes Lebesgue probability measure on \([0,1]\). The conditional-mean identity in \cref{obj:thm:oracle-frontier-broad-class} identifies the mean effect through the untruncated inverse-propensity score. The treatment probability cancels its inverse weight separately in each arm, leaving the difference between the treated and baseline conditional means. This applies the inverse-probability weighting principle of \citet[pp.~665--666]{Horvitz1952} to the conditional mean contrast.

Outcome clipping makes this identity usable under the conditional moment envelope. The score \(R_i\), cutoff \(T\), and public bounds \(b,\Lambda\) are those of \cref{obj:def:oracle-handle}:
\[
b=20T^{1-p},
\qquad
\Lambda=\frac{80T^{2-p}}{s},
\qquad
s=\lfloor n/2\rfloor.
\]
The bias bound accounts for clipping in both treatment arms. The covariance bound combines the clipped second-moment envelope with the bounded inverse weights and averaging over an evaluation block. These calculations use uniform design, overlap, and the conditional moment restriction in
\cref{obj:ass:uniform,obj:ass:overlap,obj:ass:raw-moment}.

The statistical roles of the two bounds are distinct. Increasing the cutoff reduces clipping bias through \(T^{1-p}\), while the covariance allowance grows through \(T^{2-p}\). At \(p=2\), the latter factor equals one. Effect smoothness in \cref{obj:ass:effect-smoothness} controls histogram approximation of the mean contrast. Thus the broader-class attainment result depends on the moment order and effect smoothness, with continuous propensity and arm-mean representatives supplied by \cref{obj:def:oracle-broad-model}.

\subsection{Removing the unknown constant}

The histogram representation in \cref{obj:def:projection} makes the constant-effect direction explicit. Let
\[
u=J^{-1/2}(1,\ldots,1)^T,
\qquad
C_J=I_J-uu^T,
\]
where \(I_J\) is the \(J\times J\) identity matrix and \(C_J\) is the orthogonal projection onto the complement of the constant direction. In the notation of \cref{obj:def:oracle-handle}, \(u\) is a unit vector and
\[
C_JF_J(x)=F_J(x)-u.
\]
Consequently, the centered block vectors and their inner product have the equivalent representations
\[
V_b=\frac1s\sum_{i\in\mathcal I_b}C_JF_J(X_i)R_i,
\qquad b\in\{1,2\},
\]
and
\[
\langle V_1,V_2\rangle
=
\frac1{s^2}
\sum_{i\in\mathcal I_1}\sum_{k\in\mathcal I_2}
\bigl\{\Pi_J(X_i,X_k)-1\bigr\}R_iR_k.
\]
These are representations of the statistic already specified in \cref{obj:def:oracle-handle}.

Constant centering targets the distance of the effect from its average. Every unknown constant in the null of \cref{obj:def:oracle-broad-null} occupies the same removed direction. Clipping bias remains covered by the public bound \(b\), while effect smoothness governs the approximation of the remaining heterogeneity. The two disjoint evaluation blocks provide independent score vectors, so their quadratic comparison measures the squared population signal subject to the stated bias and covariance allowances.

The construction also specifies the test on its full supplied-function domain. For an arbitrary supplied continuous function \(f\), define its overlap-clipped version by
\[
\widetilde f(x)=\max\{1/4,\min\{f(x),3/4\}\}.
\]
The score \(R_i(f)\) and block vector \(V_b(f)\) denote the expressions in \cref{obj:def:oracle-handle} evaluated with \(\widetilde f\) in place of \(\widetilde e\). This convention bounds the denominators throughout the test domain. For a legal true supplied propensity, it agrees with the propensity itself, as required by the rejection expectation in \cref{obj:def:oracle-broad-risk}.

\subsection{Calibration, attainment, and common-propensity witnesses}

The threshold in \cref{obj:def:oracle-handle} is
\[
2b^2+1024\sqrt J\,\Lambda.
\]
Its first term accommodates the population discrepancy caused by clipping under a constant effect. The second supplies the quadratic fluctuation allowance. The resulting calibration in \cref{obj:thm:oracle-frontier-broad-class} controls rejection probability by \(0.01\) uniformly over the entire broader-class null for every \(n\ge2\). At the smallest sample sizes, the public tuning gives \(J=1\), for which constant centering makes the statistic identically zero.

The tuning coordinates clipping, projection, and quadratic variability. With the target scale \(a=s^{-E_0}\) from \cref{obj:def:oracle-handle}, the cutoff and rank are upper dyadic roundings of
\[
a^{-1/(p-1)}
\quad\text{and}\quad
a^{-1/\gamma},
\]
respectively. Their common separation exponent is
\[
E_0(p)=\frac{2\gamma q}{2\gamma+q},
\qquad q=\frac{p-1}{p},
\]
as defined in \cref{obj:def:sharp-frontier-scales}. The attained scale is \(\rho_n^{\mathrm e}(v)=n^{-E_0(p)}\). At moment order two,
\[
E_0(2)=\frac{2\gamma}{4\gamma+1}.
\]
This agrees numerically with the familiar one-dimensional quadratic-testing exponent in the Gaussian Sobolev setting of \citet[Section~6.1, pp.~249--250]{IngsterSapatinas2009GoodnessOfFit}, where the stated smoothness condition is strictly above \(1/4\). The supplied-propensity guarantees here have the conditional-moment and continuous-carrier scope specified in \cref{obj:thm:oracle-frontier-broad-class}, including \(\gamma=1/4\).

The lower and upper multipliers \(c_v^{\mathrm e},C_v^{\mathrm e}\) and threshold \(N_v^{\mathrm e}\) are defined in \cref{obj:thm:oracle-frontier-broad-class}. Their radius bounds hold for every \(n\ge2\). The power guarantee applies when
\[
C_v^{\mathrm e}\rho_n^{\mathrm e}(v)<D_v^{\mathrm{e,br}}.
\]
The public threshold places this upper separation at most \(d_0/2\), where \(d_0=1/(8\sqrt2)\) is the positive witness distance in that result. The lower bound \(D_v^{\mathrm{e,br}}\ge d_0\) then supplies a nonempty separation range.

The common-propensity witnesses connect this attainment result to both supplied experiments. In \cref{obj:thm:oracle-frontier-broad-class}, the null law and every positive-weight alternative law belong to the original and broader classes and share the propensity \(x\mapsto1/2\). Their complete-record distributions have total variation below one half. Supplying that same function preserves the comparison underlying the lower error bound.

Class inclusion transfers the broader-class calibration and power guarantees to the original model:
\[
\mathcal M_v\subseteq\mathcal M_v^{\mathrm{e,br}},
\qquad
H_0(v)\subseteq H_0^{\mathrm{e,br}}(v),
\]
using \cref{obj:def:model,obj:def:null,obj:def:oracle-broad-model,obj:def:oracle-broad-null}. The witnesses also lie within the original class, where \cref{obj:thm:oracle-frontier-resolved} supplies the corresponding necessity statement. Its multipliers \(c^{\mathrm e}_v,C^{\mathrm e}_v\) have the same displayed values as their broader-class counterparts; its power condition uses the original supremum \(D_v\).

\subsection{The preservation boundary on the original class}

Hold \(w=(\alpha,\beta,\gamma)\) fixed when evaluating the exponents in \cref{obj:def:sharp-frontier-scales}. The comparison identity in \cref{obj:thm:heavy-tail-comparison} is
\[
E_4(p)-E_0(p)
=
\frac{2\gamma q}{(2\gamma+q)D_p}\{F(p)-1\}.
\]
Its positive prefactor makes \(F(p)=1\) the equality boundary between the two separation exponents. The original-model exponent is their minimum.

The nonnegative comparison exponent \(G(v)\), defined in \cref{obj:thm:exact-propensity-preservation-boundary}, therefore has the representations
\[
G(v)=E_0(p)-E(p)
=
\frac{2\gamma q}{(2\gamma+q)D_p}
\max\{0,1-F(p)\}.
\]
The same-class radius ratio is bounded on both sides by positive multiples of \(n^{G(v)}\), with the exact multipliers given in that result. For \(F(p)\ge1\), this exponent equals zero and the ratio remains bounded over all \(n\ge2\). Equality points \(F(p)=1\) are included. For \(F(p)<1\), the positive exponent quantifies the polynomial gain from supplied assignment information. This yields precisely the preservation region in \cref{obj:def:preservation-region,obj:thm:exact-propensity-preservation-boundary}.

\subsection{The finite-moment and bounded comparison}

The bounded comparison evaluates the original-record exponent at moment order two while preserving \(w\). For presentation, let \(\delta_0(p)\) and \(\delta_4(p)\) denote the within-mechanism exponent differences. The formulas in \cref{obj:def:sharp-frontier-scales} give
\[
\delta_0(p)
:=E_0(2)-E_0(p)
=
\frac{4\gamma^2(2-p)}
{p(4\gamma+1)(2\gamma+q)},
\]
and
\[
\delta_4(p)
:=E_4(2)-E_4(p)
=
\frac{2S\beta(1/q-2)}{D_2D_p},
\qquad
D_2=1+2S+\frac{S}{2\gamma}.
\]
Both differences are strictly positive for \(p<2\) and vanish at \(p=2\).

At fixed smoothness, \(F(p)\ge F(2)\). The selected exponent difference consequently has three compatible representations:
\[
E(2)-E(p)=
\begin{cases}
\delta_0(p), & F(2)\ge1,\\[3pt]
\displaystyle
\delta_4(p)+
\frac{2\gamma q}{(2\gamma+q)D_p}\{F(p)-1\},
& F(p)\ge1>F(2),\\[8pt]
\delta_4(p), & F(p)<1.
\end{cases}
\]
At \(F(p)=1\), the additional term in the middle expression vanishes. At \(F(2)=1\), the bounded exponents coincide, so the adjoining expressions agree. These equality points preserve the exponent comparison across the change in the rate-determining mechanism.

The matched capped-radius inequalities in \cref{obj:thm:heavy-tail-comparison} attach the factor \(n^{E(2)-E(p)}\) to \(\Delta_n(v)\), the ratio defined in \cref{obj:def:comparison-handle}. Strict positivity for \(p<2\) gives the diverging ratio characterized by \cref{obj:def:tail-region}. At \(p=2\), the exponent difference is zero and the ratio remains uniformly bounded over sample sizes, as characterized by \cref{obj:def:equality-region}. The phase equality \(F(p)=1\) and the moment-order equality \(p=2\) thus describe distinct comparisons.

All these statements retain the finite-sample multipliers and saturation conventions of the cited radius results. Their local uniformity is expressed through common positive lower multipliers, finite upper multipliers, and finite bounds on the public power thresholds over compact subsets of the admissible parameter domains, as stipulated in \cref{obj:thm:heavy-tail-comparison,obj:thm:oracle-frontier-broad-class,obj:thm:oracle-frontier-resolved,obj:thm:exact-propensity-preservation-boundary}. These uniform constants accompany the parameter-dependent powers, including at the stated equality boundaries.

\begin{theoremv}{thm:heavy-tail-comparison}[Heavy-tail critical-radius comparison]
Let the admissible exponent domains be
\[
\begin{aligned}
\mathcal V
&=\{(p,\alpha,\beta,\gamma)\in\mathbb R^4:
1<p\le2,\ 0<\alpha\le1,\ 0<\beta\le1,\ 1/4\le\gamma\le1\},\\
\mathcal W
&=\{(\alpha,\beta,\gamma)\in\mathbb R^3:
0<\alpha\le1,\ 0<\beta\le1,\ 1/4\le\gamma\le1\}.
\end{aligned}
\]
For \(v=(p,\alpha,\beta,\gamma)\in\mathcal V\), write
\(w=(\alpha,\beta,\gamma)\). Use the public scales, multipliers, thresholds, and original-record tests defined in \cref{obj:def:sharp-frontier-scales}, and the models, null classes, risks, and capped radii defined in \cref{obj:def:model,obj:def:null,obj:def:bounded-model,obj:def:bounded-null,obj:def:testing-risk,obj:def:bounded-risk,obj:def:critical-radius,obj:def:bounded-radius,obj:def:binary-radius}. The matched radius ratio from \cref{obj:def:comparison-handle} is
\[
\Delta_n(v)=\frac{r_n^*(v)}{r_n^{*,\mathrm b}(w)}.
\]
Here \(\lambda\) is Lebesgue probability measure on \([0,1]\), and the population distance of the mean effect from its average is
\[
d(P)=
\left(
\int_0^1
\left|\tau_P(x)-\int_0^1\tau_P(z)\,d\lambda(z)\right|^2
\,d\lambda(x)
\right)^{1/2}.
\]
The maximal distances over the respective models are
\[
D_v=\sup_{P\in\mathcal M_v}d(P),
\qquad
D_w^{\mathrm b}=\sup_{P\in\mathcal M^{\mathrm b}_w}d(P),
\]
and \(d_0>0\) is the fixed witness distance supplied by \cref{obj:lem:causal-nonempty}. For the original-record experiment in \cref{obj:def:original-record-experiment}, write the rejection expectation as
\[
\mathsf R_n(P,\phi)=\mathbb E_{\mathsf S_{n,P}}[\phi].
\]

The following assertions hold.
\begin{itemize}
\item \textbf{(Matched radii and lower separations.)}
For every \(v\in\mathcal V\), \(c_v>0\) and \(C_v>0\). For every integer \(n\ge2\), both scales are positive and
\[
\begin{aligned}
c_v\rho_n(v)
&\le r_n^*(v)\le C_v\rho_n(v),\\
c_w^{\mathrm b}\rho_n^{\mathrm b}(w)
&\le r_n^{*,\mathrm b}(w)
=r_n^{*,\mathrm{bin}}(w)
\le C_w^{\mathrm b}\rho_n^{\mathrm b}(w).
\end{aligned}
\]
The lower separations satisfy
\[
0<c_v\rho_n(v)<D_v,
\qquad
0<c_w^{\mathrm b}\rho_n^{\mathrm b}(w)<D_w^{\mathrm b},
\]
and the corresponding minimax type-II errors satisfy
\[
B_n\!\left(v,c_v\rho_n(v)\right)>\frac1{10},
\qquad
B_n^{\mathrm b}\!\left(w,c_w^{\mathrm b}\rho_n^{\mathrm b}(w)\right)>\frac1{10}.
\]

\item \textbf{(Size and conditional power.)}
For every \(v\in\mathcal V\) and every integer \(n\ge2\),
\[
\begin{aligned}
\mathsf R_n(P,\phi^*_{n,v})&\le\frac1{10}
&&\text{for every }P\in H_0(v),\\
\mathsf R_n(P,\phi^{*,\mathrm b}_{n,w})&\le\frac1{10}
&&\text{for every }P\in H_0^{\mathrm b}(w).
\end{aligned}
\]
If \(C_v\rho_n(v)<D_v\), then
\[
1-\mathsf R_n(P,\phi^*_{n,v})\le\frac1{10}
\]
for every \(P\in\mathcal M_v\) satisfying \(d(P)\ge C_v\rho_n(v)\).
If \(C_w^{\mathrm b}\rho_n^{\mathrm b}(w)<D_w^{\mathrm b}\), then
\[
1-\mathsf R_n(P,\phi^{*,\mathrm b}_{n,w})\le\frac1{10}
\]
for every \(P\in\mathcal M^{\mathrm b}_w\) satisfying
\(d(P)\ge C_w^{\mathrm b}\rho_n^{\mathrm b}(w)\).

\item \textbf{(Vanishing scales and public thresholds.)}
For every \(v\in\mathcal V\),
\[
\rho_n(v)\longrightarrow0,
\qquad
\rho_n^{\mathrm b}(w)\longrightarrow0
\quad\text{as }n\longrightarrow\infty.
\]
For every natural number \(n\),
\[
\begin{aligned}
n\ge N_v
&\ \Longrightarrow\ C_v\rho_n(v)<d_0,\\
n\ge N_w^{\mathrm b}
&\ \Longrightarrow\ C_w^{\mathrm b}\rho_n^{\mathrm b}(w)<d_0.
\end{aligned}
\]

\item \textbf{(Phase identities and radius comparison.)}
Use the supplied-propensity exponent \(E_0(p)\), rough-confounding exponent \(E_4(p)\), phase-comparison scalar \(F(p)\), and denominator \(D_p\) from \cref{obj:def:sharp-frontier-scales}, with the smoothness tuple \(w\) fixed. With
\[
q=\frac{p-1}{p},
\qquad
E(p)=\min\{E_0(p),E_4(p)\},
\]
every \(v\in\mathcal V\) satisfies
\[
E_4(p)-E_0(p)
=
\frac{2\gamma q}{(2\gamma+q)D_p}\bigl(F(p)-1\bigr),
\]
and
\[
F(p)\ge1\ \Longrightarrow\ E(p)=E_0(p),
\qquad
F(p)<1\ \Longrightarrow\ E(p)=E_4(p).
\]
With
\[
D_2=1+2S+\frac{S}{2\gamma},
\qquad
F_2=2S+\frac{S}{2\gamma},
\]
the exponent difference, with evaluation at \(2\) preserving \(w\), is
\[
E(2)-E(p)=
\begin{cases}
\displaystyle
\frac{4\gamma^2(2-p)}{(4\gamma+1)(2\gamma p+p-1)},
&F_2\ge1,\\[6pt]
\displaystyle E_4(2)-E_0(p),
&F_2<1\le F(p),\\[4pt]
\displaystyle
\frac{2S\beta(2-p)}{(p-1)D_pD_2},
&F(p)<1.
\end{cases}
\]
For every integer \(n\ge2\),
\[
\frac{c_v}{C_w^{\mathrm b}}\,n^{E(2)-E(p)}
\le\Delta_n(v)
\le
\frac{C_v}{c_w^{\mathrm b}}\,n^{E(2)-E(p)},
\qquad
1\le\Delta_n(v).
\]

\item \textbf{(Strict and equality regions.)}
The strict-enlargement and bounded-ratio regions of \cref{obj:def:tail-region,obj:def:equality-region} are exactly
\[
\mathcal T=\{v\in\mathcal V:p<2\},
\qquad
\mathcal E=\{v\in\mathcal V:p=2\}.
\]
There exists a set \(U\subseteq\mathbb R^4\), open in the product topology of the four public exponents, such that
\[
U\cap\mathcal V\ne\varnothing,
\qquad
U\cap\mathcal V\subseteq\mathcal T.
\]
For every \(v\in\mathcal V\) with \(p<2\),
\[
E(p)<E(2),
\qquad
\frac{\rho_n(v)}{\rho_n^{\mathrm b}(w)}
\longrightarrow\infty
\quad\text{as }n\longrightarrow\infty.
\]

\item \textbf{(Tail-essential complete-record priors.)}
For every \(v\in\mathcal V\) with \(p<2\), the selected mark magnitudes from \cref{obj:def:converse-handle} satisfy
\[
L_{n,v}\longrightarrow\infty
\quad\text{as }n\longrightarrow\infty.
\]
For every integer \(n\ge2\), the selected priors
\(\pi_{0,n,v}\) and \(\pi_{1,n,v}\) from \cref{obj:def:converse-handle} are normalized finite probability priors. Every law receiving positive null-prior weight belongs to \(H_0(v)\), and every law receiving positive alternative-prior weight belongs to \(\mathcal M_v\) and satisfies
\[
c_v\rho_n(v)\le d(P)<D_v.
\]
For every law receiving positive weight in either prior, each conditional arm law assigns probability one to the three-point mark set: writing \(Q_{a,P}(dy\mid x)\) for the conditional outcome probability kernel in arm \(a\),
\[
Q_{a,P}\!\left(\{0,-L_{n,v},L_{n,v}\}\mid x\right)=1
\quad
\text{for }a\in\{0,1\}\text{ and }\lambda\text{-almost every }x.
\]
The complete-record mixtures defined in \cref{obj:def:converse-handle} satisfy
\[
\operatorname{TV}\!\left(\mathbb P_{0,n,v},\mathbb P_{1,n,v}\right)<\frac12,
\]
where total variation is the supremum of the absolute difference in event probabilities.

For each such \(v\), there exists a natural-number threshold \(N(v)\) such that, for every \(n\ge N(v)\), every pair of normalized finite probability priors supported respectively on
\[
H_0^{\mathrm b}(w)
\quad\text{and}\quad
\{P\in\mathcal M^{\mathrm b}_w:c_v\rho_n(v)\le d(P)\}
\]
has complete-record mixture total variation at least \(1/2\), in the sense of \cref{obj:thm:tail-essential-full-record-converse}.

\item \textbf{(Compact uniformity.)}
For every compact set of admissible public tuples, there exist real constants \(c,C,N\), with \(c>0\), such that
\[
c\le c_v,\qquad C_v\le C,\qquad N_v\le N
\]
for every tuple in that set. Likewise, for every compact set of admissible smoothness tuples, there exist real constants \(c,C,N\), with \(c>0\), such that
\[
c\le c_w^{\mathrm b},
\qquad
C_w^{\mathrm b}\le C,
\qquad
N_w^{\mathrm b}\le N
\]
for every smoothness tuple in that set.
\end{itemize}
\end{theoremv}

\begin{theoremv}{thm:oracle-frontier-broad-class}[Supplied-propensity broader-class frontier]
For every admissible public tuple \(v=(p,\alpha,\beta,\gamma)\) of \cref{obj:def:model}, consider the broader model and its constant-effect null in \cref{obj:def:oracle-broad-model,obj:def:oracle-broad-null}, with testing risk and capped critical radius as in \cref{obj:def:oracle-broad-risk,obj:def:oracle-broad-radius}. Put
\[
q=\frac{p-1}{p},\qquad
E_0=\frac{2\gamma q}{2\gamma+q},\qquad
\rho_n^{\mathrm e}(v)=n^{-E_0},\qquad
d_0=\frac{1}{8\sqrt2},
\]
and define the lower multiplier, upper multiplier, and public sample-size threshold by
\[
c_v^{\mathrm e}=\frac{4^{-\gamma}}{16\sqrt3},\qquad
C_v^{\mathrm e}=3\bigl(120+128\sqrt{160\sqrt2}\bigr),\qquad
N_v^{\mathrm e}
=\left\lceil\max\left\{2,
\left(\frac{2C_v^{\mathrm e}}{d_0}\right)^{1/E_0}
\right\}\right\rceil .
\]
Let \(\lambda\) denote Lebesgue probability measure on \([0,1]\), and write the heterogeneity distance as
\[
d(P)=
\left\{\int_{[0,1]}
\left(\tau_P(x)-\int_{[0,1]}\tau_P(z)\,\lambda(dz)\right)^2
\,\lambda(dx)\right\}^{1/2}.
\]
Then
\[
d_0\le D_v^{\mathrm{e,br}}\le\frac12,\qquad
\frac{1}{64\sqrt3}\le c_v^{\mathrm e},\qquad
0<C_v^{\mathrm e}.
\]

For every integer \(n\ge2\), the following guarantees hold.
\begin{itemize}
\item \textbf{(Critical radius.)}
\[
c_v^{\mathrm e}\rho_n^{\mathrm e}(v)
\le r_n^{*,\mathrm{e,br}}(v)
\le C_v^{\mathrm e}\rho_n^{\mathrm e}(v),
\qquad
B_n^{\mathrm{e,br}}
\bigl(v,c_v^{\mathrm e}\rho_n^{\mathrm e}(v)\bigr)
\ge\frac25.
\]

\item \textbf{(Calibration and power.)}
The supplied-propensity test \(\phi^{*,\mathrm e}_{n,v}\) of
\cref{obj:def:oracle-handle} satisfies
\[
\mathsf R_n^{\mathrm e}
(P,\phi^{*,\mathrm e}_{n,v})\le0.01
\qquad
\text{for every }P\in H_0^{\mathrm{e,br}}(v).
\]
Here \(\mathsf R_n^{\mathrm e}(P,\phi^{\mathrm e})\) is the expected rejection probability when the test receives the entire true continuous propensity \(e_P\), together with the independent original-record sample and public randomization. If
\[
C_v^{\mathrm e}\rho_n^{\mathrm e}(v)<D_v^{\mathrm{e,br}},
\]
then every \(P\in\mathcal M_v^{\mathrm{e,br}}\) with
\(d(P)\ge C_v^{\mathrm e}\rho_n^{\mathrm e}(v)\) satisfies
\[
1-\mathsf R_n^{\mathrm e}
(P,\phi^{*,\mathrm e}_{n,v})\le0.01.
\]

\item \textbf{(Finite lower-bound witnesses.)}
There exist a continuous function \(f\), a null law \(P_0\), and a normalized finite probability prior
\[
f(x)=\frac12\quad(x\in[0,1]),\qquad
\pi_1=\sum_{j=1}^{k}\omega_j\delta_{P_j},\qquad
\omega_j\ge0,\qquad
\sum_{j=1}^{k}\omega_j=1,
\]
with at least one strictly positive weight, such that
\[
P_0\in H_0(v)\cap H_0^{\mathrm{e,br}}(v),
\qquad e_{P_0}=f.
\]
Here \(H_0(v)\) is the original-model null of \cref{obj:def:null}.
For every \(j\) with \(\omega_j>0\),
\[
P_j\in\mathcal M_v\cap\mathcal M_v^{\mathrm{e,br}},
\qquad e_{P_j}=f,\qquad
c_v^{\mathrm e}\rho_n^{\mathrm e}(v)\le d(P_j)<d_0,
\]
where \(\mathcal M_v\) is defined in \cref{obj:def:model}.
Define the complete-record alternative mixture by
\[
\overline P_1^{(n)}
=\sum_{j=1}^{k}\omega_j P_j^{\otimes n}.
\]
With total variation defined as the supremum of absolute event-probability differences,
\[
\operatorname{TV}
\bigl(P_0^{\otimes n},\overline P_1^{(n)}\bigr)<\frac12.
\]
For every allowed randomized supplied-propensity test \(\phi^{\mathrm e}\) satisfying
\[
\mathsf R_n^{\mathrm e}(P,\phi^{\mathrm e})\le\frac1{10}
\qquad
\text{for every }P\in H_0^{\mathrm{e,br}}(v),
\]
there is an index \(j\) with \(\omega_j>0\) such that
\[
1-\mathsf R_n^{\mathrm e}(P_j,\phi^{\mathrm e})
\ge\frac25.
\]
\end{itemize}

For every integer \(n\ge N_v^{\mathrm e}\),
\[
C_v^{\mathrm e}\rho_n^{\mathrm e}(v)\le\frac{d_0}{2}.
\]

For each \(P\in\mathcal M_v^{\mathrm{e,br}}\), let
\(Q_{a,P}(dy\mid x)\) denote the conditional outcome probability kernel in treatment arm \(a\). The untruncated inverse-propensity score has conditional mean equal to the canonical effect: for \(\lambda\)-almost every \(x\),
\[
e_P(x)\int_{\mathbb R}\frac{y}{e_P(x)}\,Q_{1,P}(dy\mid x)
+
\bigl(1-e_P(x)\bigr)
\int_{\mathbb R}\frac{-y}{1-e_P(x)}\,Q_{0,P}(dy\mid x)
=\tau_P(x).
\]

The constants and exponent are locally uniform on compact subsets of the admissible public parameter domain, equipped with the product topology on \((p,\alpha,\beta,\gamma)\).
\end{theoremv}

\begin{theoremv}{thm:oracle-frontier-resolved}[Supplied-propensity frontier and attainment]
Let \(\mathcal V\) denote the admissible domain of public tuples \(v=(p,\alpha,\beta,\gamma)\) in \cref{obj:def:model}, equipped with the topology inherited from \(\mathbb R^4\). For \(v\in\mathcal V\), define
\[
q=\frac{p-1}{p},\qquad
E_0=\frac{2\gamma q}{2\gamma+q},\qquad
\rho_n^{\mathrm e}(v)=n^{-E_0},\qquad
c^{\mathrm e}_v=\frac{4^{-\gamma}}{16\sqrt3},\qquad
C^{\mathrm e}_v=3\left(120+128\sqrt{160\sqrt2}\right).
\]
Use the supplied-propensity rule \(\phi^{*,\mathrm e}_{n,v}\) of \cref{obj:def:oracle-handle}, and the risk and capped radius of \cref{obj:def:oracle-risk,obj:def:oracle-radius}. Write \(\mathsf R_n^{\mathrm e}(P,\phi^{\mathrm e})\) for the expected rejection probability when the entire true continuous propensity \(e_P\) is supplied to the test.

Let \(\lambda\) be Lebesgue probability measure on \([0,1]\), and define
\[
d(P)=
\left\{\int_0^1
\left(\tau_P(x)-\int_0^1\tau_P(z)\,d\lambda(z)\right)^2
\,d\lambda(x)\right\}^{1/2},
\qquad
D_v=\sup_{P\in\mathcal M_v}d(P).
\]
Let \(d_0>0\) be the fixed distance supplied by the legal model witness, and set
\[
N^{\mathrm e}_v=
\left\lceil
\max\left\{2,\left(\frac{2C^{\mathrm e}_v}{d_0}\right)^{1/E_0}\right\}
\right\rceil.
\]
Then
\[
\frac{1}{64\sqrt3}\le c^{\mathrm e}_v,
\qquad
0<C^{\mathrm e}_v.
\]
For every integer \(n\ge2\),
\[
c^{\mathrm e}_v\rho_n^{\mathrm e}(v)
\le r_n^{*,\mathrm e}(v)
\le C^{\mathrm e}_v\rho_n^{\mathrm e}(v),
\qquad
B_n^{\mathrm e}\!\left(v,c^{\mathrm e}_v\rho_n^{\mathrm e}(v)\right)
\ge\frac25.
\]
The attaining rule satisfies
\[
\mathsf R_n^{\mathrm e}\!\left(P,\phi^{*,\mathrm e}_{n,v}\right)
\le0.01
\qquad(P\in H_0(v)).
\]
Whenever \(C^{\mathrm e}_v\rho_n^{\mathrm e}(v)<D_v\), it also satisfies
\[
1-\mathsf R_n^{\mathrm e}\!\left(P,\phi^{*,\mathrm e}_{n,v}\right)
\le0.01
\qquad
\left(P\in\mathcal M_v,\quad
C^{\mathrm e}_v\rho_n^{\mathrm e}(v)\le d(P)\right).
\]

For each such \(v,n\), there exist a null law \(P_0\in H_0(v)\) and a normalized finite prior
\[
\pi_1=\sum_{j=1}^{k}\omega_j\delta_{P_j},
\qquad
\omega_j\ge0,\qquad
\sum_{j=1}^{k}\omega_j=1,\qquad
\exists j:\ \omega_j>0,
\]
such that every positive-weight support law satisfies
\[
P_j\in\mathcal M_v,\qquad
e_{P_j}=e_{P_0},\qquad
c^{\mathrm e}_v\rho_n^{\mathrm e}(v)\le d(P_j)<D_v.
\]
Define its complete-record mixture by
\[
\mathbb P_{\pi_1,n}=\sum_{j=1}^{k}\omega_jP_j^{\otimes n},
\]
and let \(\operatorname{TV}\) denote the supremum of absolute event-probability differences. These laws satisfy
\[
\operatorname{TV}\!\left(P_0^{\otimes n},\mathbb P_{\pi_1,n}\right)<\frac12.
\]
For every supplied-propensity test \(\phi^{\mathrm e}\) satisfying
\[
\mathsf R_n^{\mathrm e}(P,\phi^{\mathrm e})\le\frac1{10}
\qquad(P\in H_0(v)),
\]
there is an index \(j\) with
\[
\omega_j>0,\qquad
1-\mathsf R_n^{\mathrm e}(P_j,\phi^{\mathrm e})\ge\frac25.
\]
For every integer \(n\ge N^{\mathrm e}_v\),
\[
C^{\mathrm e}_v\rho_n^{\mathrm e}(v)\le\frac{d_0}{2}.
\]

For \(P\in\mathcal M_v\), let \(Q_{a,P}(dy\mid x)\) denote the conditional outcome probability kernel in treatment arm \(a\). The untruncated inverse-propensity score identifies the original mean effect: for \(\lambda\)-almost every \(x\),
\[
e_P(x)\int_{\mathbb R}\frac{y}{e_P(x)}\,Q_{1,P}(dy\mid x)
+
\bigl(1-e_P(x)\bigr)
\int_{\mathbb R}\frac{-y}{1-e_P(x)}\,Q_{0,P}(dy\mid x)
=\tau_P(x).
\]
The frontier constants and nonempty-power thresholds are locally uniform over admissible tuples: on every compact subset of \(\mathcal V\), the lower multipliers admit a common positive lower bound, the upper multipliers admit a common finite upper bound, and the thresholds admit a common finite upper bound, including at \(p=2\) and \(\gamma=1/4\).
\end{theoremv}

\begin{theoremv}{thm:exact-propensity-preservation-boundary}[Exact propensity preservation boundary]
Let \(\mathcal V\) be the domain of public tuples \(v=(p,\alpha,\beta,\gamma)\) satisfying:
\begin{itemize}
\item \textbf{(Moment order.)} \(1<p\le2\).
\item \textbf{(Propensity smoothness.)} \(0<\alpha\le1\).
\item \textbf{(Baseline smoothness.)} \(0<\beta\le1\).
\item \textbf{(Effect smoothness.)} \(1/4\le\gamma\le1\).
\end{itemize}
For every \(v\in\mathcal V\), use the quantities and original-model radius multipliers \(c_v,C_v\) of \cref{obj:def:sharp-frontier-scales}, with
\[
q=\frac{p-1}{p},\qquad S=\alpha+\beta,\qquad
D_p=1+2\alpha+\frac{\beta}{q}+\frac{S}{2\gamma},
\]
\[
E_0(p)=\frac{2\gamma q}{2\gamma+q},\qquad
E_4(p)=\frac{2S}{D_p},\qquad
F(p)=\frac{2\alpha}{p-1}+\frac{\beta}{q}+\frac{S}{2\gamma}.
\]
Define the comparison exponent and the supplied-propensity radius multipliers by
\[
G(v)=\max\{0,E_0(p)-E_4(p)\},\qquad
c^{\mathrm e}_v=\frac{4^{-\gamma}}{16\sqrt3},\qquad
C^{\mathrm e}_v=3\bigl(120+128\sqrt{160\sqrt2}\bigr).
\]
Then
\[
G(v)=
\frac{2\gamma q}{(2\gamma+q)D_p}\max\{0,1-F(p)\}.
\]
For the unknown- and supplied-propensity experiments on the same class \(\mathcal M_v\) of \cref{obj:def:model}, the capped critical radii of \cref{obj:def:critical-radius,obj:def:oracle-radius} satisfy, for every integer \(n\ge2\),
\[
\frac{c_v}{C^{\mathrm e}_v}n^{G(v)}
\le
\frac{r_n^*(v)}{r_n^{*,\mathrm e}(v)}
\le
\frac{C_v}{c^{\mathrm e}_v}n^{G(v)}.
\]
If \(F(p)=1\), then
\[
0<\frac{c_v}{C^{\mathrm e}_v},
\qquad
0<\frac{C_v}{c^{\mathrm e}_v},
\]
and, for every integer \(n\ge2\),
\[
\frac{c_v}{C^{\mathrm e}_v}
\le
\frac{r_n^*(v)}{r_n^{*,\mathrm e}(v)}
\le
\frac{C_v}{c^{\mathrm e}_v}.
\]
If \(F(p)<1\), then
\[
\lim_{n\to\infty}\frac{r_n^*(v)}{r_n^{*,\mathrm e}(v)}=+\infty,
\]
and, for every integer \(n\ge2\),
\[
\frac{c_v}{C^{\mathrm e}_v}n^{E_0(p)-E_4(p)}
\le
\frac{r_n^*(v)}{r_n^{*,\mathrm e}(v)}
\le
\frac{C_v}{c^{\mathrm e}_v}n^{E_0(p)-E_4(p)}.
\]
The preservation region of \cref{obj:def:preservation-region} is exactly
\[
\mathcal R=
\left\{v\in\mathcal V:
\frac{2\alpha}{p-1}+\frac{\beta}{q}
+\frac{\alpha+\beta}{2\gamma}\ge1\right\}.
\]
Moreover, for every compact subset \(K\) of \(\mathcal V\), in the product topology on the four real coordinates, there exist real constants \(c,C\), with \(c>0\), such that, for every \(v\in K\),
\[
c\le\frac{c_v}{C^{\mathrm e}_v},
\qquad
64\sqrt3\,C_v\le C.
\]
\end{theoremv}


\section{Verification note}

The theorem statements and their proofs are machine-checked in Lean 4. The anchors attached to mathematical environments connect the displayed definitions and results to their corresponding formal declarations. Checking establishes the stated conclusions from the hypotheses and dependencies of each result, preserving its parameter domain, quantifiers, constants, and sample-size conditions.

The restrictions defining the statistical models, including those in \cref{obj:def:model,obj:def:bounded-model,obj:def:binary-model,obj:def:oracle-broad-model}, are stipulated hypotheses. The verified conclusions concern deductions under these restrictions, including calibration, attainment, lower bounds, and separation-radius comparisons. Where a result uses a published conclusion as a dependency, its theorem-local verification-scope footnote identifies that conclusion and records whether the source proof was formalized. These local disclosures specify the dependency boundary alongside the checked statements and proofs.


\section{Proofs of the main results}\label{sec:deferred-proofs}

\begin{proof}[Proof of \cref{obj:prop:bounded-submodel-comparison}]
Fix the admissible tuple \(v=(p,\alpha,\beta,\gamma)\) and its matched smoothness tuple \(w=(\alpha,\beta,\gamma)\).

\begin{enumerate}
\item We first establish the integrated conversion identity. Write
\[
\mathcal O=[0,1]\times\{0,1\}\times\mathbb R
\]
for the Borel record space. For a law \(P\) in either
\(\mathcal M_v\) or \(\mathcal M_{(2,\alpha,\beta,\gamma)}\), let
\(Q_{a_{\mathrm{arm}},P}(\,\cdot\mid x)\) be its conditional outcome probability kernel in arm \(a_{\mathrm{arm}}\in\{0,1\}\). Thus
\[
\begin{split}
P(dx,\{a_{\mathrm{arm}}\},dy)
&=\lambda(dx)e_P(x)^{a_{\mathrm{arm}}}
  (1-e_P(x))^{1-a_{\mathrm{arm}}}
  Q_{a_{\mathrm{arm}},P}(dy\mid x),\\
m_{a_{\mathrm{arm}},P}(x)
&=m_{0,P}(x)+a_{\mathrm{arm}}\tau_P(x)
 =\int y\,Q_{a_{\mathrm{arm}},P}(dy\mid x)
 \qquad\text{for \(\lambda\)-almost every \(x\)}.
\end{split}
\]
The mean bounds give
\[
|m_{0,P}(x)|\le\frac12,\qquad
|m_{1,P}(x)|
\le |m_{0,P}(x)|+|\tau_P(x)|\le1.
\]
Define the signed-binary conversion \(\mathcal C_{\mathrm{bin}}(P)\) by
\[
\begin{split}
Q^{\mathrm{bin}}_{a_{\mathrm{arm}},P}(dy\mid x)
&=\frac{1+m_{a_{\mathrm{arm}},P}(x)}2\,\delta_1(dy)
 +\frac{1-m_{a_{\mathrm{arm}},P}(x)}2\,\delta_{-1}(dy),\\
\mathcal C_{\mathrm{bin}}(P)(dx,\{a_{\mathrm{arm}}\},dy)
&=\lambda(dx)e_P(x)^{a_{\mathrm{arm}}}
 (1-e_P(x))^{1-a_{\mathrm{arm}}}
 Q^{\mathrm{bin}}_{a_{\mathrm{arm}},P}(dy\mid x).
\end{split}
\]
The two arm weights are nonnegative and sum to one, and
\[
\int y\,Q^{\mathrm{bin}}_{a_{\mathrm{arm}},P}(dy\mid x)
=m_{a_{\mathrm{arm}},P}(x),\qquad
\int |y|^2\,Q^{\mathrm{bin}}_{a_{\mathrm{arm}},P}(dy\mid x)=1.
\]
Consequently the conversion has signed-binary support and retains the uniform design, overlap, smoothness restrictions, and mean bounds. In particular,
\[
\begin{gathered}
\mathcal C_{\mathrm{bin}}(P)\in\mathcal M_w^{\mathrm{bin}},\qquad
\bigl(\mathcal C_{\mathrm{bin}}(P)\bigr)_{X,A}=P_{X,A},\\
e_{\mathcal C_{\mathrm{bin}}(P)}=e_P,\qquad
m_{0,\mathcal C_{\mathrm{bin}}(P)}=m_{0,P},\qquad
\tau_{\mathcal C_{\mathrm{bin}}(P)}=\tau_P.
\end{gathered}
\]
Substitution into the definitions of distance and constancy yields
\[
d(\mathcal C_{\mathrm{bin}}(P))=d(P),
\]
and
\[
\begin{split}
&\exists c\in[-1/2,1/2]\ 
  \forall x,\ \tau_{\mathcal C_{\mathrm{bin}}(P)}(x)=c\\
&\hspace{35mm}\Longleftrightarrow
\exists c\in[-1/2,1/2]\ 
  \forall x,\ \tau_P(x)=c.
\end{split}
\]
% lean: CausalSmith.Stat.FinitepHomogeneityDensegamma.binaryConversion_primitives
% lean: CausalSmith.Stat.FinitepHomogeneityDensegamma.binaryConversion_inModel
% lean: CausalSmith.Stat.FinitepHomogeneityDensegamma.binaryConversion_preserves_effect

To define record randomization on the whole record space, set
\[
\ell_1(y)=\max\{-1,\min\{1,y\}\}\qquad(y\in\mathbb R)
\]
and define the probability kernel
\[
\mathsf K_{\mathrm{sign}}((x,a_{\mathrm{obs}},y),\cdot)
=
\frac{1+\ell_1(y)}2\,\delta_{(x,a_{\mathrm{obs}},1)}
+\frac{1-\ell_1(y)}2\,\delta_{(x,a_{\mathrm{obs}},-1)},
\qquad (x,a_{\mathrm{obs}},y)\in\mathcal O.
\]
For \(P\in\mathcal M_w^{\mathrm b}\), uniform design and fixed overlap transfer the observed outcome bound to both conditional arm kernels:
\[
Q_{a_{\mathrm{arm}},P}([-1,1]\mid x)=1
\qquad
\text{for \(\lambda\)-almost every \(x\), for both arms}.
\]
Indeed, disintegrating \(P(|Y|>1)=0\) gives
\[
\int
\left[
(1-e_P(x))Q_{0,P}(\{|y|>1\}\mid x)
+e_P(x)Q_{1,P}(\{|y|>1\}\mid x)
\right]\lambda(dx)=0.
\]
Both summands are nonnegative, and both arm weights are at least \(1/4\). Each conditional probability therefore vanishes almost everywhere. Since \(\ell_1(y)=y\) on \([-1,1]\), for \(b_{\mathrm{sign}}\in\{-1,1\}\),
\[
\int
\frac{1+b_{\mathrm{sign}}\ell_1(y)}2\,
Q_{a_{\mathrm{arm}},P}(dy\mid x)
=
\frac{1+b_{\mathrm{sign}}m_{a_{\mathrm{arm}},P}(x)}2
\qquad\text{for \(\lambda\)-almost every \(x\)}.
\]
Integrating the two atomic weights and then the treatment mixture proves the one-record identity
\[
\int_{\mathcal O}\mathsf K_{\mathrm{sign}}(o,\cdot)\,P(do)
=\mathcal C_{\mathrm{bin}}(P).
\]
% lean: CausalSmith.Stat.FinitepHomogeneityDensegamma.binaryRecordKernel_comp

Now let \(n\in\mathbb N\), let \(\phi\) be a measurable \([0,1]\)-valued test, and write its input as
\[
\boldsymbol o=(o_i)_{i=1}^n\in\mathcal O^n,\qquad
o_i=(x_i,a_{\mathrm{obs},i},y_i),\qquad
u_{\mathrm{seed}}\in[0,1].
\]
Define its integrated rejection map by
\[
\widetilde\phi(\boldsymbol o,u_{\mathrm{seed}})
=
\sum_{\boldsymbol b_{\mathrm{sign}}\in\{-1,1\}^n}
\phi\!\left(
 ((x_i,a_{\mathrm{obs},i},b_{\mathrm{sign},i}))_{i=1}^n,
 u_{\mathrm{seed}}
\right)
\prod_{i=1}^n
\frac{1+b_{\mathrm{sign},i}\ell_1(y_i)}2,
\qquad
\boldsymbol b_{\mathrm{sign}}
=(b_{\mathrm{sign},i})_{i=1}^n.
\]
Its weights are nonnegative and satisfy
\[
\sum_{\boldsymbol b_{\mathrm{sign}}\in\{-1,1\}^n}
\prod_{i=1}^n
\frac{1+b_{\mathrm{sign},i}\ell_1(y_i)}2
=
\prod_{i=1}^n
\left(
\frac{1+\ell_1(y_i)}2+\frac{1-\ell_1(y_i)}2
\right)=1.
\]
Thus \(\widetilde\phi\) is measurable and takes values in \([0,1]\). For \(n=0\), the sign-vector set contains the single empty vector and the empty product is one, so the same definition applies.

The finite sum is the integral of \(\phi\) against the independent product of the record kernels. Independence and the one-record identity imply
\[
\int_{\mathcal O^n}
\bigotimes_{i=1}^n\mathsf K_{\mathrm{sign}}(o_i,\cdot)\,
P^{\otimes n}(d\boldsymbol o)
=
\mathcal C_{\mathrm{bin}}(P)^{\otimes n}.
\]
This follows by factorization on measurable rectangles; the empty product gives the identity for \(n=0\). Hence, for every fixed \(u_{\mathrm{seed}}\),
\[
\int_{\mathcal O^n}
\widetilde\phi(\boldsymbol o,u_{\mathrm{seed}})
\,P^{\otimes n}(d\boldsymbol o)
=
\int_{\mathcal O^n}
\phi(\boldsymbol o,u_{\mathrm{seed}})
\,\mathcal C_{\mathrm{bin}}(P)^{\otimes n}(d\boldsymbol o).
\]
Integrating in the public seed and using the experiment of
\cref{obj:def:original-record-experiment} gives
\[
\mathsf R_n(P,\widetilde\phi)
=
\mathsf R_n(\mathcal C_{\mathrm{bin}}(P),\phi).
\]
Boundedness of the tests justifies the integral exchanges. On bounded records, the kernel uses exactly the probabilities \((1+Y)/2\) and \((1-Y)/2\) stated in the proposition.
% lean: CausalSmith.Stat.FinitepHomogeneityDensegamma.binaryRandomization_rejectProb

\item We next compare the binary and bounded testing risks. In the infimum--supremum comparisons, a supremum of type-II errors over an empty alternative set is assigned the value zero. The test
\[
\phi_{\mathrm{zero}}(\boldsymbol o,u_{\mathrm{seed}})=0
\]
is level-valid for every null class. The sets of level-valid tests are therefore nonempty, and
\[
0\le 1-\mathsf R_n(P,\phi)\le1
\]
bounds every individual type-II error and every corresponding worst-case error.

Signed-binary support implies bounded support. Keeping the remaining restrictions and the constancy condition gives
\[
\mathcal M_w^{\mathrm{bin}}\subseteq\mathcal M_w^{\mathrm b},
\qquad
H_0^{\mathrm{bin}}(w)\subseteq H_0^{\mathrm b}(w).
\]
Every bounded-null level-valid test is consequently binary-null level-valid, with
\[
\sup_{\substack{P\in\mathcal M_w^{\mathrm{bin}}\\d(P)\ge r}}
\{1-\mathsf R_n(P,\phi)\}
\le
\sup_{\substack{P\in\mathcal M_w^{\mathrm b}\\d(P)\ge r}}
\{1-\mathsf R_n(P,\phi)\}.
\]
For a nonempty binary alternative set, this is the supremum comparison under inclusion; for an empty binary alternative set, it follows from the zero convention and nonnegativity. Taking infima gives
\[
B_n^{\mathrm{bin}}(w,r)\le B_n^{\mathrm b}(w,r).
\]
% lean: CausalSmith.Stat.FinitepHomogeneityDensegamma.binaryTestingRisk_le_boundedTestingRisk

Conversely, take any binary-null level-valid test \(\phi\). For every \(P\in H_0^{\mathrm b}(w)\), the conversion belongs to \(H_0^{\mathrm{bin}}(w)\), so the preceding identity gives
\[
\mathsf R_n(P,\widetilde\phi)
=
\mathsf R_n(\mathcal C_{\mathrm{bin}}(P),\phi)
\le\frac1{10}.
\]
For a bounded alternative \(P\) with \(d(P)\ge r\), its conversion is a binary alternative with the same distance, and
\[
1-\mathsf R_n(P,\widetilde\phi)
=
1-\mathsf R_n(\mathcal C_{\mathrm{bin}}(P),\phi)
\le
\sup_{\substack{P\in\mathcal M_w^{\mathrm{bin}}\\d(P)\ge r}}
\{1-\mathsf R_n(P,\phi)\}.
\]
If the bounded alternative set is nonempty, taking its supremum preserves this inequality; the binary supremum is bounded above by one. If the bounded alternative set is empty, its worst-case error is zero and the binary worst-case error is nonnegative. Thus in either case,
\[
\sup_{\substack{P\in\mathcal M_w^{\mathrm b}\\d(P)\ge r}}
\{1-\mathsf R_n(P,\widetilde\phi)\}
\le
\sup_{\substack{P\in\mathcal M_w^{\mathrm{bin}}\\d(P)\ge r}}
\{1-\mathsf R_n(P,\phi)\}.
\]
The bounded-model infimum is at most the left-hand side for each such \(\phi\). Taking the binary-model infimum proves
\[
B_n^{\mathrm b}(w,r)\le B_n^{\mathrm{bin}}(w,r),
\qquad
B_n^{\mathrm{bin}}(w,r)=B_n^{\mathrm b}(w,r).
\]
% lean: CausalSmith.Stat.FinitepHomogeneityDensegamma.boundedTestingRisk_eq_binary_of_conversion

\item For the remaining model inclusion, take \(P\in\mathcal M_w^{\mathrm b}\). The conditional support established above and the nonnegative moment exponent supplied by admissibility give
\[
\int |y|^p\,Q_{a_{\mathrm{arm}},P}(dy\mid x)
\le
\int 1\,Q_{a_{\mathrm{arm}},P}(dy\mid x)
=1\le10
\qquad\text{for \(\lambda\)-almost every \(x\)}.
\]
All the other primitive restrictions are identical at the matched smoothness tuple. Therefore
\[
\mathcal M_w^{\mathrm{bin}}
\subseteq\mathcal M_w^{\mathrm b}
\subseteq\mathcal M_v,
\qquad
H_0^{\mathrm b}(w)\subseteq H_0(v).
\]
% lean: CausalSmith.Stat.FinitepHomogeneityDensegamma.boundedModel_inModel

\item The distance suprema are equal because the attained distance sets are equal. The binary-to-bounded inclusion gives one direction, and converting each bounded law gives the other:
\[
\{d(P):P\in\mathcal M_w^{\mathrm{bin}}\}
=
\{d(P):P\in\mathcal M_w^{\mathrm b}\}.
\]
For comparison with the full model, convert each bounded law. Its converted arm kernels also satisfy
\[
\int |y|^p\,Q^{\mathrm{bin}}_{a_{\mathrm{arm}},P}(dy\mid x)
=1\le10,
\]
so the conversion belongs to \(\mathcal M_v\) and has the same distance. Conversely, the conversion of every \(P\in\mathcal M_v\) belongs to the binary model, hence to the bounded model, and retains \(d(P)\). Thus
\[
\{d(P):P\in\mathcal M_w^{\mathrm b}\}
=
\{d(P):P\in\mathcal M_v\}.
\]
Taking suprema yields
\[
D_w^{\mathrm{bin}}=D_w^{\mathrm b}=D_v.
\]
Finally, \cref{obj:lem:causal-nonempty} supplies the boundedness of the full-model distance set and the witness lower bound
\[
d_0\le D_v\le\frac12.
\]
% lean: CausalSmith.Stat.FinitepHomogeneityDensegamma.maxDistBinary_eq_maxDistBounded
% lean: CausalSmith.Stat.FinitepHomogeneityDensegamma.maxDistBounded_eq_maxDist
% lean: CausalSmith.Stat.FinitepHomogeneityDensegamma.model_distance_lower

\item A test level-valid for \(H_0(v)\) is level-valid for \(H_0^{\mathrm b}(w)\), and model inclusion gives
\[
\sup_{\substack{P\in\mathcal M_w^{\mathrm b}\\d(P)\ge r}}
\{1-\mathsf R_n(P,\phi)\}
\le
\sup_{\substack{P\in\mathcal M_v\\d(P)\ge r}}
\{1-\mathsf R_n(P,\phi)\}.
\]
For a nonempty bounded alternative set, this follows by inclusion; for an empty set, it follows from the zero convention and nonnegativity of the full-model error. Taking infima over full-null level-valid tests, and allowing all bounded-null level-valid tests on the left, proves
\[
B_n^{\mathrm b}(w,r)\le B_n(v,r).
\]
Together with the equality already established, this gives, for \(n\ge2\) and \(0<r<D_v\),
\[
B_n^{\mathrm{bin}}(w,r)=B_n^{\mathrm b}(w,r)\le B_n(v,r).
\]
% lean: CausalSmith.Stat.FinitepHomogeneityDensegamma.boundedTestingRisk_le_testingRisk

\item We establish strict positivity of the binary capped radius by a small cosine alternative. For \(t_{\mathrm{cos}}\in[0,1]\), define the deterministic law and its conversion by
\[
\begin{gathered}
P^{\mathrm{det}}_{\mathrm{cos},t_{\mathrm{cos}}}:
\quad X\sim\lambda,\qquad
\Pr(A=1\mid X=x)=\frac12,\qquad
Y=A\,\frac{t_{\mathrm{cos}}}{8}\cos(2\pi X),\\
P_{\mathrm{cos},t_{\mathrm{cos}}}
=\mathcal C_{\mathrm{bin}}
 \bigl(P^{\mathrm{det}}_{\mathrm{cos},t_{\mathrm{cos}}}\bigr).
\end{gathered}
\]
Their primitive functions are
\[
e_{P_{\mathrm{cos},t_{\mathrm{cos}}}}(x)=\frac12,\qquad
m_{0,P_{\mathrm{cos},t_{\mathrm{cos}}}}(x)=0,\qquad
\tau_{P_{\mathrm{cos},t_{\mathrm{cos}}}}(x)
=\frac{t_{\mathrm{cos}}}{8}\cos(2\pi x).
\]
The propensity and baseline are constant. For the effect, the cosine increment bound and \(|x-z|\le1\) imply
\[
\begin{split}
|\tau_{P_{\mathrm{cos},t_{\mathrm{cos}}}}(x)|
&\le\frac{t_{\mathrm{cos}}}{8}\le\frac12,\\
|\tau_{P_{\mathrm{cos},t_{\mathrm{cos}}}}(x)
 -\tau_{P_{\mathrm{cos},t_{\mathrm{cos}}}}(z)|
&\le t_{\mathrm{cos}}\frac{2\pi}{8}|x-z|
 \le t_{\mathrm{cos}}\,20|x-z|^\gamma
 \le20|x-z|^\gamma.
\end{split}
\]
The deterministic outcomes have absolute value at most \(1/8\), so their conditional second moments are at most one. These bounds verify membership of the deterministic law in
\(\mathcal M_{(2,\alpha,\beta,\gamma)}\); conversion then gives
\[
P_{\mathrm{cos},t_{\mathrm{cos}}}\in\mathcal M_w^{\mathrm{bin}},
\qquad
P_{\mathrm{cos},0}\in H_0^{\mathrm{bin}}(w).
\]
Moreover,
\[
\int\tau_{P_{\mathrm{cos},t_{\mathrm{cos}}}}\,d\lambda=0,\qquad
\int\tau_{P_{\mathrm{cos},t_{\mathrm{cos}}}}^2\,d\lambda
=\frac{t_{\mathrm{cos}}^2}{128},
\qquad
d(P_{\mathrm{cos},t_{\mathrm{cos}}})=t_{\mathrm{cos}}d_0.
\]
% lean: CausalSmith.Stat.FinitepHomogeneityDensegamma.smallCosineLaw_inBinaryModel
% lean: CausalSmith.Stat.FinitepHomogeneityDensegamma.smallCosineLaw_distance

For probability measures \(\mu_{\mathrm{tv}},\nu_{\mathrm{tv}}\) on a common measurable space \(\mathcal Z_{\mathrm{tv}}\), use
\[
\operatorname{TV}(\mu_{\mathrm{tv}},\nu_{\mathrm{tv}})
=
\sup_{E_{\mathrm{tv}}\text{ measurable in }\mathcal Z_{\mathrm{tv}}}
|\mu_{\mathrm{tv}}(E_{\mathrm{tv}})
 -\nu_{\mathrm{tv}}(E_{\mathrm{tv}})|.
\]
The conditional distribution of \((A,Y)\) under the cosine law is
\[
\begin{split}
\kappa_{\mathrm{cos},t_{\mathrm{cos}}}(x,\cdot)
={}&\frac{1+\tau_{P_{\mathrm{cos},t_{\mathrm{cos}}}}(x)}4
       \delta_{(1,1)}
+\frac{1-\tau_{P_{\mathrm{cos},t_{\mathrm{cos}}}}(x)}4
       \delta_{(1,-1)}\\
&+\frac14\delta_{(0,1)}+\frac14\delta_{(0,-1)}.
\end{split}
\]
For a measurable record event \(E_{\mathrm{tv}}\subseteq\mathcal O\), define its section by
\[
E_{\mathrm{tv},x}
=\{(a_{\mathrm{arm}},y):
  (x,a_{\mathrm{arm}},y)\in E_{\mathrm{tv}}\}.
\]
Subtracting the four displayed masses gives
\[
\begin{split}
&\kappa_{\mathrm{cos},0}(x,E_{\mathrm{tv},x})
-\kappa_{\mathrm{cos},t_{\mathrm{cos}}}(x,E_{\mathrm{tv},x})\\
&\quad=
-\frac{\tau_{P_{\mathrm{cos},t_{\mathrm{cos}}}}(x)}4
\left[
\mathbf1_{E_{\mathrm{tv},x}}(1,1)
-\mathbf1_{E_{\mathrm{tv},x}}(1,-1)
\right].
\end{split}
\]
The absolute difference of these two indicators is at most one. Therefore
\[
\left|
\kappa_{\mathrm{cos},0}(x,E_{\mathrm{tv},x})
-\kappa_{\mathrm{cos},t_{\mathrm{cos}}}(x,E_{\mathrm{tv},x})
\right|
\le
\frac{|\tau_{P_{\mathrm{cos},t_{\mathrm{cos}}}}(x)|}{4}
\le\frac{t_{\mathrm{cos}}}{16}.
\]
Integration over the common uniform design gives
\[
\begin{split}
|P_{\mathrm{cos},0}(E_{\mathrm{tv}})
-P_{\mathrm{cos},t_{\mathrm{cos}}}(E_{\mathrm{tv}})|
&\le
\int
\left|
\kappa_{\mathrm{cos},0}(x,E_{\mathrm{tv},x})
-\kappa_{\mathrm{cos},t_{\mathrm{cos}}}(x,E_{\mathrm{tv},x})
\right|\lambda(dx)\\
&\le\frac{t_{\mathrm{cos}}}{16},
\end{split}
\]
and hence
\[
\operatorname{TV}
(P_{\mathrm{cos},0},P_{\mathrm{cos},t_{\mathrm{cos}}})
\le\frac{t_{\mathrm{cos}}}{16}.
\]
% lean: CausalSmith.Stat.FinitepHomogeneityDensegamma.smallCosineLaw_tv

We use the expectation bound for any measurable
\(f_{\mathrm{tv}}:\mathcal Z_{\mathrm{tv}}\to[0,1]\). The layer-cake identity is
\[
\int f_{\mathrm{tv}}\,d\eta_{\mathrm{tv}}
=
\int_0^1
\eta_{\mathrm{tv}}
\{z_{\mathrm{tv}}:f_{\mathrm{tv}}(z_{\mathrm{tv}})\ge h_{\mathrm{tv}}\}
\,dh_{\mathrm{tv}},
\qquad
\eta_{\mathrm{tv}}\in\{\mu_{\mathrm{tv}},\nu_{\mathrm{tv}}\}.
\]
Subtracting these identities, taking absolute values, and bounding each event-probability difference gives
\[
\left|
\int f_{\mathrm{tv}}\,d\mu_{\mathrm{tv}}
-\int f_{\mathrm{tv}}\,d\nu_{\mathrm{tv}}
\right|
\le\operatorname{TV}(\mu_{\mathrm{tv}},\nu_{\mathrm{tv}}).
\]
Inserting the mixed product
\(P_{\mathrm{cos},0}\otimes
P_{\mathrm{cos},t_{\mathrm{cos}}}^{\otimes n_{\mathrm{prod}}}\)
between the two product laws, applying the triangle inequality, and integrating event sections gives
\[
\begin{split}
&\operatorname{TV}
\left(
P_{\mathrm{cos},0}^{\otimes(n_{\mathrm{prod}}+1)},
P_{\mathrm{cos},t_{\mathrm{cos}}}^{\otimes(n_{\mathrm{prod}}+1)}
\right)\\
&\quad\le
\operatorname{TV}(P_{\mathrm{cos},0},P_{\mathrm{cos},t_{\mathrm{cos}}})
+
\operatorname{TV}
\left(
P_{\mathrm{cos},0}^{\otimes n_{\mathrm{prod}}},
P_{\mathrm{cos},t_{\mathrm{cos}}}^{\otimes n_{\mathrm{prod}}}
\right),
\qquad n_{\mathrm{prod}}\in\mathbb N.
\end{split}
\]
For the first term, the integrated section probability is a \([0,1]\)-valued function, so the expectation bound applies; for the second, each section difference is bounded by the product total variation. At \(n_{\mathrm{prod}}=0\), both product laws are the point mass on the empty sample. Induction therefore yields
\[
\operatorname{TV}
\left(
P_{\mathrm{cos},0}^{\otimes n},
P_{\mathrm{cos},t_{\mathrm{cos}}}^{\otimes n}
\right)
\le
n\,\operatorname{TV}
(P_{\mathrm{cos},0},P_{\mathrm{cos},t_{\mathrm{cos}}})
\le\frac{nt_{\mathrm{cos}}}{16}.
\]
% lean: Causalean.Stat.Minimax.MomentMatchedMixture.tvDist_pi_iid_le

For the fixed \(n\ge2\), choose
\[
t_{\mathrm{small}}=\frac8{100n},
\qquad
0<t_{\mathrm{small}}\le1.
\]
The preceding product bound becomes
\[
\operatorname{TV}
\left(
P_{\mathrm{cos},0}^{\otimes n},
P_{\mathrm{cos},t_{\mathrm{small}}}^{\otimes n}
\right)
\le\frac{nt_{\mathrm{small}}}{16}
=\frac1{200}.
\]
For a binary-null level-valid test \(\phi\), define its seed average by
\[
\overline\phi(\boldsymbol o)
=\int_{[0,1]}\phi(\boldsymbol o,u_{\mathrm{seed}})
\,\lambda(du_{\mathrm{seed}}).
\]
It is measurable, satisfies \(0\le\overline\phi\le1\), and obeys
\[
\mathsf R_n(P,\phi)
=\int_{\mathcal O^n}\overline\phi\,dP^{\otimes n}.
\]
Applying the expectation bound to this average and using null validity gives
\[
\begin{split}
1-\mathsf R_n(P_{\mathrm{cos},t_{\mathrm{small}}},\phi)
&\ge
\frac9{10}
-\operatorname{TV}
\left(
P_{\mathrm{cos},0}^{\otimes n},
P_{\mathrm{cos},t_{\mathrm{small}}}^{\otimes n}
\right)\\
&\ge\frac{179}{200}>\frac1{10}.
\end{split}
\]
Whenever \(0<r\le t_{\mathrm{small}}d_0\), this cosine law belongs to the separated alternative set. Taking its error as a lower bound for the worst-case error, and then taking the infimum over level-valid tests, proves
\[
B_n^{\mathrm{bin}}(w,r)\ge\frac{179}{200}>\frac1{10}
\qquad(0<r\le t_{\mathrm{small}}d_0).
\]
% lean: CausalSmith.Stat.FinitepHomogeneityDensegamma.testingRiskOn_ge_of_two_point

Thus every successful separation in the binary capped-radius set exceeds \(t_{\mathrm{small}}d_0\). Its sentinel also satisfies
\[
t_{\mathrm{small}}d_0\le d_0\le D_w^{\mathrm{bin}}.
\]
The set is nonempty because it contains the sentinel. Taking its infimum, as defined in \cref{obj:def:binary-radius}, gives
\[
0<t_{\mathrm{small}}d_0
\le r_n^{*,\mathrm{bin}}(w).
\]
% lean: CausalSmith.Stat.FinitepHomogeneityDensegamma.binaryCriticalRadius_pos

\item Finally, the common cap and the risk equality give identical sets in the definitions of the binary and bounded radii:
\[
\begin{split}
&\left\{r\in(0,D_w^{\mathrm{bin}}):
 B_n^{\mathrm{bin}}(w,r)\le\frac1{10}\right\}
 \cup\{D_w^{\mathrm{bin}}\}\\
&\qquad=
\left\{r\in(0,D_w^{\mathrm b}):
 B_n^{\mathrm b}(w,r)\le\frac1{10}\right\}
 \cup\{D_w^{\mathrm b}\}.
\end{split}
\]
Their infima are therefore equal:
\[
r_n^{*,\mathrm{bin}}(w)=r_n^{*,\mathrm b}(w).
\]
The risk ordering also gives
\[
\begin{split}
&\left\{r\in(0,D_v):B_n(v,r)\le\frac1{10}\right\}\cup\{D_v\}\\
&\qquad\subseteq
\left\{r\in(0,D_v):B_n^{\mathrm b}(w,r)\le\frac1{10}\right\}
\cup\{D_v\}.
\end{split}
\]
Both sets contain \(D_v\) and are bounded below by zero, since \(D_v\ge d_0>0\). Taking infima reverses this inclusion. By
\cref{obj:def:bounded-radius,obj:def:critical-radius}, it follows that
\[
r_n^{*,\mathrm b}(w)\le r_n^*(v).
\]
Combining this with the positivity just proved establishes
\[
0<r_n^{*,\mathrm{bin}}(w)
=r_n^{*,\mathrm b}(w)\le r_n^*(v).
\]
The conversion properties and the identity for every \(n\in\mathbb N\) were established in the first step.
% lean: CausalSmith.Stat.FinitepHomogeneityDensegamma.cappedRadius_eq_of_risk_eq
% lean: CausalSmith.Stat.FinitepHomogeneityDensegamma.boundedCriticalRadius_le_criticalRadius
\end{enumerate}
\end{proof}

\begin{proof}[Proof of \cref{obj:thm:whole-null-calibration-resolved}]
Fix an admissible tuple \(v\) and an integer \(n\ge2\). Throughout, write
\[
M:=J_{\mathrm c},
\qquad
B:=\mathfrak b_{n,v}.
\]
These quantities have the values prescribed in \cref{obj:def:explicit-score-ledger}, including its small-sample branch.

\begin{enumerate}
\item First,
\[
B\ge0,
\qquad
M\ge1.
\]
For \(n<4\), these follow from \(B=0\) and \(M=1\). For \(n\ge4\), every summand in the formula for \(B\) is nonnegative: the clipping cutoffs and ranks are positive, and the increment weights \(a_j\) are nonnegative. This also covers an empty increment sum. The coarse rank is a nonnegative integer power of two, hence is at least one.
% lean: CausalSmith.Stat.FinitepHomogeneityDensegamma.ledgerBias_nonneg
% lean: CausalSmith.Stat.FinitepHomogeneityDensegamma.ledgerM_ge_one

\item Let \(P\in\mathcal M_v\). By \cref{obj:lem:original-score-covariance-ledger}, the affine coefficients in
\[
Z_b(c)=Z_{b,0}-cZ_{b,A},
\qquad b\in\{1,2\},
\]
have square-integrable scalar projections, and, for all \(u,z\in\mathbb R^M\),
\[
\begin{aligned}
\operatorname{Var}(u^T Z_{b,0})
&\le \Lambda_{n,v}\|u\|_2^2,\\
\operatorname{Var}(u^T Z_{b,A})
&\le \Lambda_{n,v}\|u\|_2^2,\\
\left|\operatorname{Cov}(u^T Z_{b,0},z^T Z_{b,A})\right|
&\le \Lambda_{n,v}\|u\|_2\|z\|_2.
\end{aligned}
\]
Taking the coordinate directions and summing their second moments proves square integrability of each coefficient vector.

Define the sample event
\[
\mathcal E_{\mathrm{prof}}(P)
:=
\left\{
\text{sample}:
\forall |c|\le\tfrac12,\ 
\left|W(c)-\|\theta_P(c)\|_2^2\right|
\le
128\left(
\sqrt{\Lambda_{n,v}}\|\theta_P(c)\|_2
+\sqrt M\,\Lambda_{n,v}
\right)
\right\}.
\]
By \cref{obj:lem:uniform-affine-profile-bound}, this event is measurable and
\[
P^{\otimes n}\!\left(\mathcal E_{\mathrm{prof}}(P)\right)\ge0.9925.
\]
The same result supplies, when \(\Lambda_{n,v}=0\), the simultaneous almost-sure identity
\[
W(c)=\|\theta_P(c)\|_2^2
\qquad\text{for every }|c|\le\tfrac12.
\]
% lean: CausalSmith.Stat.FinitepHomogeneityDensegamma.original_score_covariance_ledger
% lean: CausalSmith.Stat.FinitepHomogeneityDensegamma.uniform_affine_profile_bound

\item Suppose \(P\in H_0(v)\). By \cref{obj:def:null}, there is a constant \(c_0\) satisfying
\[
c_0\in[-\tfrac12,\tfrac12],
\qquad
\tau_P(x)=c_0\quad\text{for every }x\in[0,1].
\]
If \(n\ge4\), \cref{obj:lem:original-score-mean-ledger} gives
\[
\|\theta_P(c_0)\|_2\le B.
\]
If \(n<4\), the zero-score branch of \cref{obj:def:explicit-score-ledger} gives
\[
\theta_P(c)=0\quad\text{for every }c\in\mathbb R,
\qquad B=0,
\]
so the same bound holds. Since norms are nonnegative and \(c_0\) is an admissible candidate,
\[
\inf_{|c|\le1/2}\|\theta_P(c)\|_2
\le \|\theta_P(c_0)\|_2
\le B.
\]
% lean: CausalSmith.Stat.FinitepHomogeneityDensegamma.whole_null_calibration_resolved

\item We next establish the size bound for \(n\ge4\). The covariance ledger is strictly positive. Indeed, the construction gives
\[
s_n=\lfloor n/2\rfloor\ge2,
\qquad
T_0\ge1,
\qquad
V_0=32T_0^{2-p}>0,
\qquad
L_1\ge16\sqrt{V_0}>0.
\]
Consequently,
\[
\Lambda_{n,v}
=
\frac{L_1^2}{s_n}
+\frac{2L_2^2}{s_n(s_n-1)}
>0.
\]
% lean: CausalSmith.Stat.FinitepHomogeneityDensegamma.ledgerCov_pos

On \(\mathcal E_{\mathrm{prof}}(P)\), define
\[
u_{c_0}:=\|\theta_P(c_0)\|_2.
\]
Then \(0\le u_{c_0}\le B\), and the profiling inequality implies
\[
W(c_0)
\le
u_{c_0}^2
+128\sqrt{\Lambda_{n,v}}\,u_{c_0}
+128\sqrt M\,\Lambda_{n,v}.
\]
Completing the square gives
\[
0\le\left(B-64\sqrt{\Lambda_{n,v}}\right)^2,
\qquad
128\sqrt{\Lambda_{n,v}}\,B
\le B^2+4096\Lambda_{n,v}.
\]
Also, \(M\ge1\) implies
\[
\sqrt M\ge1,
\qquad
\Lambda_{n,v}\le\sqrt M\,\Lambda_{n,v}.
\]
Thus
\[
\begin{aligned}
W(c_0)
&\le B^2+128\sqrt{\Lambda_{n,v}}\,B
       +128\sqrt M\,\Lambda_{n,v}\\
&\le2B^2+4096\Lambda_{n,v}
       +128\sqrt M\,\Lambda_{n,v}\\
&\le2B^2+4224\sqrt M\,\Lambda_{n,v}\\
&\le2B^2+65536\sqrt M\,\Lambda_{n,v}.
\end{aligned}
\]
% lean: CausalSmith.Stat.FinitepHomogeneityDensegamma.profile_null_scalar

\item We justify the comparison with the computed profile minimum and the rounded cutoff. At a fixed sample, define
\[
Z_{b,0}:=Z_b(0),
\qquad
Z_{b,A}:=Z_b(0)-Z_b(1),
\qquad b\in\{1,2\}.
\]
The affine score construction gives \(Z_b(c)=Z_{b,0}-cZ_{b,A}\). Define its quadratic coefficients by
\[
\begin{aligned}
\omega_0&:=\langle Z_{1,0},Z_{2,0}\rangle,\\
\omega_1&:=-\langle Z_{1,A},Z_{2,0}\rangle
           -\langle Z_{1,0},Z_{2,A}\rangle,\\
\omega_2&:=\langle Z_{1,A},Z_{2,A}\rangle.
\end{aligned}
\]
Expansion yields
\[
W(c)=\omega_0+\omega_1c+\omega_2c^2.
\]
Define the smaller endpoint value and the computed profile value by
\[
\omega_{\mathrm{end}}
:=
\min\left\{
\omega_0-\frac{\omega_1}{2}+\frac{\omega_2}{4},
\omega_0+\frac{\omega_1}{2}+\frac{\omega_2}{4}
\right\},
\]
\[
\mathsf W_{\min}
:=
\begin{cases}
\displaystyle
\min\left\{\omega_{\mathrm{end}},
\omega_0-\frac{\omega_1^2}{4\omega_2}\right\},
&
\displaystyle
\omega_2>0,\quad
-\frac12\le-\frac{\omega_1}{2\omega_2}\le\frac12,\\[2mm]
\omega_{\mathrm{end}},&\text{otherwise}.
\end{cases}
\]
The quotient in the first branch is used under \(\omega_2>0\).

For every \(|c|\le1/2\),
\[
\mathsf W_{\min}\le W(c).
\]
To see this in the first branch, complete the square:
\[
W(c)-\left(\omega_0-\frac{\omega_1^2}{4\omega_2}\right)
=
\omega_2\left(c+\frac{\omega_1}{2\omega_2}\right)^2
\ge0.
\]
If \(\omega_2>0\) and the vertex lies to the left of the interval, then \(\omega_1>\omega_2\), and
\[
W(c)-W(-\tfrac12)
=
(c+\tfrac12)\bigl[\omega_1+\omega_2(c-\tfrac12)\bigr]
\ge0.
\]
If the vertex lies to the right, then \(\omega_1<-\omega_2\), and
\[
W(c)-W(\tfrac12)
=
(\tfrac12-c)\bigl[-\omega_1-\omega_2(c+\tfrac12)\bigr]
\ge0.
\]
Finally, if \(\omega_2\le0\), then \(c^2\le1/4\). For \(\omega_1\ge0\),
\[
W(c)-W(-\tfrac12)
=
\omega_1(c+\tfrac12)+\omega_2(c^2-\tfrac14)
\ge0,
\]
whereas for \(\omega_1<0\),
\[
W(c)-W(\tfrac12)
=
\omega_1(c-\tfrac12)+\omega_2(c^2-\tfrac14)
\ge0.
\]
These cases include a constant or linear statistic.
% lean: CausalSmith.Stat.FinitepHomogeneityDensegamma.ledger_profiledMin_le

For the cutoff, define
\[
x_{\mathrm{cut}}
:=2B^2+65536\sqrt M\,\Lambda_{n,v}>0,
\qquad
k_{\mathrm{cut}}
:=\left\lceil\log_2 x_{\mathrm{cut}}\right\rceil.
\]
The defining ceiling inequalities give
\[
\log_2 x_{\mathrm{cut}}
\le k_{\mathrm{cut}}
<\log_2 x_{\mathrm{cut}}+1.
\]
Exponentiating and using the cutoff formula therefore yields
\[
2B^2+65536\sqrt M\,\Lambda_{n,v}
\le h_{n,v}
\le4B^2+131072\sqrt M\,\Lambda_{n,v}.
\]
% lean: CausalSmith.Stat.FinitepHomogeneityDensegamma.ledgerCutoff_bounds

Combining these bounds with the preceding step, on \(\mathcal E_{\mathrm{prof}}(P)\),
\[
\mathsf W_{\min}
\le W(c_0)
\le2B^2+65536\sqrt M\,\Lambda_{n,v}
\le h_{n,v}.
\]
The strict rejection rule in \cref{obj:def:explicit-score-ledger} consequently gives
\[
\psi_{n,v}(\text{sample},\text{seed})=0
\quad\text{on }\mathcal E_{\mathrm{prof}}(P)\times[0,1].
\]

Define the experiment measure and its good event by
\[
\mathbb P_{\mathrm{exp}}:=P^{\otimes n}\otimes\lambda,
\qquad
\mathcal E_{\mathrm{exp}}(P)
:=\mathcal E_{\mathrm{prof}}(P)\times[0,1].
\]
Independence and the unit mass of \(\lambda\) imply
\[
\mathbb P_{\mathrm{exp}}\!\left(\mathcal E_{\mathrm{exp}}(P)\right)
=
P^{\otimes n}\!\left(\mathcal E_{\mathrm{prof}}(P)\right).
\]
The test takes values in \([0,1]\), so pointwise
\[
0\le\psi_{n,v}
\le\mathbf1_{\mathcal E_{\mathrm{exp}}(P)^c}.
\]
Integrating this inequality gives
\[
\begin{aligned}
\mathsf R_n(P,\psi_{n,v})
&\le
\mathbb P_{\mathrm{exp}}\!\left(\mathcal E_{\mathrm{exp}}(P)^c\right)\\
&=1-P^{\otimes n}\!\left(\mathcal E_{\mathrm{prof}}(P)\right)\\
&\le0.0075.
\end{aligned}
\]
For \(n<4\), the zero-test branch gives directly
\[
\mathsf R_n(P,\psi_{n,v})=0\le0.0075.
\]
% lean: CausalSmith.Stat.FinitepHomogeneityDensegamma.integral_le_bad_event
% lean: CausalSmith.Stat.FinitepHomogeneityDensegamma.whole_null_calibration_resolved

\item Now assume
\[
n\ge4,
\qquad
R_{\mathrm{att}}(n,v)<D_v,
\qquad
P\in\mathcal M_v,
\qquad
d(P)\ge R_{\mathrm{att}}(n,v).
\]
Work on \(\mathcal E_{\mathrm{prof}}(P)\). The computed profile value is attained: there exists
\[
c_{\min}\in[-\tfrac12,\tfrac12],
\qquad
W(c_{\min})=\mathsf W_{\min}.
\]
Indeed, the smaller endpoint value is attained at one of the two endpoints. In the vertex branch, choose that endpoint if its value is at most the vertex value; otherwise choose
\[
c_{\mathrm{vert}}:=-\frac{\omega_1}{2\omega_2},
\qquad
W(c_{\mathrm{vert}})
=\omega_0-\frac{\omega_1^2}{4\omega_2}.
\]
Here \(\omega_2>0\) and the branch condition places \(c_{\mathrm{vert}}\) in the candidate interval. In the remaining branch, choose the endpoint attaining \(\omega_{\mathrm{end}}\).
% lean: CausalSmith.Stat.FinitepHomogeneityDensegamma.ledger_profiledMin_attained

Define
\[
u_{\min}:=\|\theta_P(c_{\min})\|_2.
\]
By \cref{obj:lem:original-score-mean-ledger},
\[
u_{\min}
\ge\frac{3}{16}d(P)-16M^{-\gamma}-B.
\]
Substituting the assumed separation and its displayed definition gives
\[
\begin{aligned}
u_{\min}
&\ge
\frac{3}{16}R_{\mathrm{att}}(n,v)-16M^{-\gamma}-B\\
&=4B+1024M^{1/4}\sqrt{\Lambda_{n,v}}.
\end{aligned}
\]

We derive the scalar power bound. Since \(M\ge1\) and \(\Lambda_{n,v}>0\),
\[
M^{1/4}\ge1,
\qquad
\left(M^{1/4}\sqrt{\Lambda_{n,v}}\right)^2
=\sqrt M\,\Lambda_{n,v}>0.
\]
Thus
\[
u_{\min}\ge256\sqrt{\Lambda_{n,v}}\ge0,
\qquad
128\sqrt{\Lambda_{n,v}}\,u_{\min}
\le\frac12u_{\min}^2.
\]
Squaring the signal bound, with all its terms nonnegative, also gives
\[
\begin{aligned}
u_{\min}^2
&\ge
\left(4B+1024M^{1/4}\sqrt{\Lambda_{n,v}}\right)^2\\
&\ge16B^2+1048576\sqrt M\,\Lambda_{n,v}.
\end{aligned}
\]
The lower side of the profiling inequality now yields
\[
\begin{aligned}
W(c_{\min})
&\ge u_{\min}^2
      -128\sqrt{\Lambda_{n,v}}\,u_{\min}
      -128\sqrt M\,\Lambda_{n,v}\\
&\ge\frac12u_{\min}^2-128\sqrt M\,\Lambda_{n,v}\\
&\ge8B^2+524160\sqrt M\,\Lambda_{n,v}\\
&>4B^2+131072\sqrt M\,\Lambda_{n,v}.
\end{aligned}
\]
The last inequality is strict because \(\sqrt M\,\Lambda_{n,v}>0\).
% lean: CausalSmith.Stat.FinitepHomogeneityDensegamma.profile_alternative_scalar

Using the upper cutoff bound and attainment,
\[
h_{n,v}
\le4B^2+131072\sqrt M\,\Lambda_{n,v}
<W(c_{\min})
=\mathsf W_{\min}.
\]
Hence
\[
\psi_{n,v}=1
\quad\text{on }\mathcal E_{\mathrm{exp}}(P),
\qquad
0\le1-\psi_{n,v}
\le\mathbf1_{\mathcal E_{\mathrm{exp}}(P)^c}.
\]
Measurability and boundedness make the test integrable under the probability measure \(\mathbb P_{\mathrm{exp}}\). Integrating the last inequality gives
\[
\begin{aligned}
1-\mathsf R_n(P,\psi_{n,v})
&=\int(1-\psi_{n,v})\,d\mathbb P_{\mathrm{exp}}\\
&\le1-\mathbb P_{\mathrm{exp}}\!\left(\mathcal E_{\mathrm{exp}}(P)\right)\\
&=1-P^{\otimes n}\!\left(\mathcal E_{\mathrm{prof}}(P)\right)\\
&\le0.0075.
\end{aligned}
\]
% lean: CausalSmith.Stat.FinitepHomogeneityDensegamma.integral_le_bad_event
% lean: CausalSmith.Stat.FinitepHomogeneityDensegamma.whole_null_calibration_resolved

\item Finally, if \(n=2\) or \(n=3\), then \(n<4\), and the explicit small-sample branch gives
\[
\psi_{n,v}(\text{sample},\text{seed})=0
\qquad\text{for every sample and seed}.
\]
% lean: CausalSmith.Stat.FinitepHomogeneityDensegamma.ledger_small_sample
\end{enumerate}
\end{proof}

\begin{proof}[Proof of \cref{obj:thm:original-record-attainable-rate}]
\begin{enumerate}
\item Fix an admissible tuple \(v=(p,\alpha,\beta,\gamma)\). Its parameter restrictions give
\[
1<p\le2,\qquad
0<\alpha,\beta\le1,\qquad
\frac14\le\gamma\le1.
\]
Using the quantities in the statement, write
\[
D_p=1+2\alpha+\frac{\beta}{q}+\frac{S}{2\gamma}.
\]
Then
\[
0<q\le\frac12,\qquad S>0,\qquad t\ge0,\qquad D_p>1,
\qquad E_0>0,\qquad E_4>0.
\]
In particular,
\[
0<E\le\gamma\le1,\qquad
\frac{E}{S}<2,\qquad
\frac{E}{p-1}<1.
\]
Indeed, the three upper bounds follow from
\[
\frac{E}{\gamma}
\le\frac{2q}{2\gamma+q}\le1,
\qquad
\frac{E}{S}\le\frac{2}{D_p}<2,
\qquad
\frac{E}{p-1}
\le\frac{2\gamma}{p(2\gamma+q)}<1.
\]
The first uses \(q\le1/2\le2\gamma\), and the last uses \(p>1\) and \(q>0\).

For every positive argument of the geometric-series factor,
\[
0<2^{-x}<1,\qquad A(x)>0.
\]
Since \(\alpha>0\) and \(\lambda_0>0\), it follows that \(B_*\ge0\). Thus
\[
C^{\mathrm{att}}_v
=16\{16+5B_*+1024\sqrt{A_*}\}
\ge256>0.
\]
% lean: CausalSmith.Stat.FinitepHomogeneityDensegamma.CAtt_ge_256

\item The rule named \(\psi_{n,v}\) in
\cref{obj:thm:whole-null-calibration-resolved} is precisely
\(\phi^*_{n,v}\). That result supplies, for every integer \(n\ge2\),
\[
\mathsf R_n(P,\phi^*_{n,v})\le0.0075
\qquad(P\in H_0(v)).
\]
It also supplies power, for \(n\ge4\), at the calibrated separation
\[
R_{\mathrm{att}}(n,v)
=\frac{16}{3}
\left\{
16M^{-\gamma}+5B+1024M^{1/4}\sqrt{\Lambda_{n,v}}
\right\},
\qquad B=\mathfrak b_{n,v},
\]
provided \(R_{\mathrm{att}}(n,v)<D_v\).

The smallest sample sizes can be handled directly.
By \cref{obj:lem:causal-nonempty},
\[
d_0\le D_v\le\frac12.
\]
For \(n\in\{2,3\}\), the bound \(E\le1\) gives
\[
n^{-E}\ge n^{-1}\ge\frac13,
\qquad
D_v\le\frac12<\frac{256}{3}
\le C^{\mathrm{att}}_v n^{-E}.
\]
Consequently,
\[
C^{\mathrm{att}}_v n^{-E}<D_v
\quad\Longrightarrow\quad n\ge4
\qquad(n\ge2).
\]
The zero rule at \(n=2,3\) has the asserted size by the same calibration result.
We next establish the deterministic estimates needed to compare
\(R_{\mathrm{att}}(n,v)\) with the stated public level.
% lean: CausalSmith.Stat.FinitepHomogeneityDensegamma.attainment_small_sample_saturated

\item Let \(n\ge4\), and use the tuning of
\cref{obj:def:explicit-score-ledger,obj:def:blocks}. Explicitly,
\[
s=\lfloor n/2\rfloor\ge2,\qquad
a=s^{-E}\in(0,1],\qquad
M=\mathcal R(a^{-1/\gamma}),\qquad
K=\mathcal R\!\left(\max\{M,a^{-1/S}\}\right).
\]
Upper dyadic rounding obeys
\[
x_{\mathrm{dyad}}\le\mathcal R(x_{\mathrm{dyad}})
\le2x_{\mathrm{dyad}}
\qquad(x_{\mathrm{dyad}}\ge1),
\qquad
\mathcal R(2^j)=2^j\qquad(j\in\mathbb N_0),
\]
by the lower and upper ceiling bounds for \(\log_2 x_{\mathrm{dyad}}\).
It is also increasing. Hence
\[
a^{-1/\gamma}\le M\le2a^{-1/\gamma},
\qquad
K\ge a^{-1/S},
\]
and therefore
\[
M^{-\gamma}\le a,\qquad K^{-S}\le a.
\]

The bounds \(E/\gamma\le1\), \(E/S<2\), and \(E/(p-1)<1\) give
\[
a^{-1/\gamma}=s^{E/\gamma}\le s,
\qquad
a^{-1/S}=s^{E/S}\le s^2,
\qquad
1\le a^{-1/(p-1)}=s^{E/(p-1)}\le s.
\]
Thus \(M\le2s\). For the fine rank, if \(a^{-1/S}\le M\), dyadic idempotence gives \(K=M\le2s\le2s^2\); if \(M<a^{-1/S}\), rounding gives \(K\le2a^{-1/S}\le2s^2\). In either case,
\[
M\le2s,\qquad K\le2s^2.
\]

For the increment-cutoff bounds, the public quantities in either score branch are
\[
L=\log_2(K/M)\in\mathbb N_0,\qquad
R_j=2^jM\quad(0\le j\le L),
\]
\[
T_j=
\mathcal R\!\left(
\max\left\{
1,\min\left\{
a^{-1/(p-1)},
\left[
\frac{R_j^{-\alpha}}{a(R_j/K)^{\lambda_0}}
\right]^{1/(p-1)}
\right\}
\right\}
\right)
\quad(1\le j\le L).
\]
The single-cutoff score uses the main cutoff
\[
T_0=\mathcal R(a^{-1/(p-1)}).
\]
Since the argument defining each \(T_j\) lies between \(1\) and
\(a^{-1/(p-1)}\), rounding and monotonicity yield
\[
1\le T_0\le2s,\qquad
1\le T_j\le T_0\quad(1\le j\le L).
\]
These statements also cover \(L=0\), when the increment index set is empty.
% lean: CausalSmith.Stat.FinitepHomogeneityDensegamma.ledger_rate_bounds

\item We now bound the bias ledger. The main-cutoff estimate is
\[
T_0^{1-p}
\le\left(a^{-1/(p-1)}\right)^{1-p}
=a.
\]
In the single-cutoff branch, this and \(K^{-S}\le a\) immediately give
\[
B=400K^{-S}+20T_0^{1-p}\le420a\le B_*a.
\]

For the multilevel branch, \(S<\gamma\) and \(\alpha<q\). Define the positive local tail targets by
\[
z_{\mathrm{tail},j}
=\frac{R_j^{-\alpha}}{a(R_j/K)^{\lambda_0}}>0
\qquad(1\le j\le L).
\]
The cutoff formula gives
\[
T_j\ge
\min\left\{a^{-1/(p-1)},z_{\mathrm{tail},j}^{1/(p-1)}\right\}.
\]
If the minimum is the first term, its \((1-p)\)-power is \(a\); if it is the second term, that power is
\[
\left(z_{\mathrm{tail},j}^{1/(p-1)}\right)^{1-p}
=z_{\mathrm{tail},j}^{-1}
=a(R_j/K)^{\lambda_0}R_j^\alpha.
\]
Since \(1-p<0\), both cases imply
\[
T_j^{1-p}
\le a+a(R_j/K)^{\lambda_0}R_j^\alpha.
\]

The increment ranks satisfy \(M\ge1\) and \(R_L\le K\). Forward and backward geometric summation therefore give, for every \(r_{\mathrm{geom}}>0\),
\[
\begin{aligned}
\sum_{j=1}^{L}R_{j-1}^{-r_{\mathrm{geom}}}
&\le\sum_{j=0}^{L-1}2^{-jr_{\mathrm{geom}}}
\le A(r_{\mathrm{geom}}),\\
\sum_{j=1}^{L}(R_j/K)^{r_{\mathrm{geom}}}
&\le\sum_{j=0}^{L-1}2^{-jr_{\mathrm{geom}}}
\le A(r_{\mathrm{geom}}),\\
\sum_{j=1}^{L}R_j^{r_{\mathrm{geom}}}
&\le K^{r_{\mathrm{geom}}}A(r_{\mathrm{geom}}).
\end{aligned}
\]
For the second inequality, use \(R_j/K\le2^{-(L-j)}\); the third follows by factoring out \(K^{r_{\mathrm{geom}}}\). Each bound remains valid for an empty sum.

Using \(a_j=40R_{j-1}^{-\alpha}\), \(R_j=2R_{j-1}\), and \(2^\alpha\le2\), we obtain
\[
\begin{aligned}
a_jT_j^{1-p}
&\le40aR_{j-1}^{-\alpha}
 +40a\,2^\alpha(R_j/K)^{\lambda_0}\\
&\le80a\left\{R_{j-1}^{-\alpha}+(R_j/K)^{\lambda_0}\right\}.
\end{aligned}
\]
Summation consequently gives
\[
\sum_{j=1}^{L}a_jT_j^{1-p}
\le80a\{A(\alpha)+A(\lambda_0)\}.
\]
Substituting into the multilevel bias ledger proves
\[
\begin{aligned}
B
&=400K^{-S}+20T_0^{1-p}
  +10\sum_{j=1}^{L}a_jT_j^{1-p}\\
&\le\left[420+800\{A(\alpha)+A(\lambda_0)\}\right]a
=B_*a.
\end{aligned}
\]
Thus this bias estimate holds in both score branches.
% lean: CausalSmith.Stat.FinitepHomogeneityDensegamma.ledger_bias_bound

\item Consider first the single-cutoff covariance branch,
\[
S\ge\gamma\quad\text{or}\quad\alpha\ge q.
\]
If \(S\ge\gamma\), then \(E\le\gamma\le S\). If \(\alpha\ge q\), then
\[
E\le E_0\le q\le\alpha<S.
\]
Thus \(a^{-1/S}=s^{E/S}\le s\). The same two cases for fine-rank rounding used above now give
\[
K\le2s.
\]

Since \(0\le2-p\le1\) and \(T_0\le2a^{-1/(p-1)}\), the main second-moment quantity satisfies
\[
V_0=32T_0^{2-p}
\le32\,2^{2-p}a^{-t}
\le64a^{-t}.
\]
The single-cutoff formulas \(L_1=16\sqrt{V_0}\) and \(L_2=4\sqrt{KV_0}\) yield the exact identity
\[
\Lambda_{n,v}
=\frac{256V_0}{s}+\frac{32KV_0}{s(s-1)}.
\]
Because \(s\ge2\) and \(K\le2s\),
\[
32K\le64s\le128(s-1).
\]
It follows that
\[
\Lambda_{n,v}
\le\frac{384V_0}{s}
\le\frac{24576a^{-t}}{s}.
\]

The tail-channel identity and the choice \(E\le E_0\) give
\[
E_0\left(2+t+\frac{1}{2\gamma}\right)=1,
\qquad
E\left(2+t+\frac{1}{2\gamma}\right)\le1.
\]
Hence
\[
\frac{a^{-1/(2\gamma)}a^{-t}}{s}
=s^{E(1/(2\gamma)+t)-1}
\le s^{-2E}=a^2.
\]
Finally, \(M\le2a^{-1/\gamma}\) implies
\[
\begin{aligned}
\sqrt M\,\Lambda_{n,v}
&\le24576\sqrt2\,
 \frac{a^{-1/(2\gamma)}a^{-t}}{s}\\
&\le24576\sqrt2\,a^2
\le A_*a^2,
\end{aligned}
\]
where the last inequality follows from the displayed definition of \(A_*\).
% lean: CausalSmith.Stat.FinitepHomogeneityDensegamma.ledger_single_covariance_rate

\item Continue with the multilevel branch \(S<\gamma\), \(\alpha<q\).
Here \(a^{-1/\gamma}\le a^{-1/S}\). If \(a^{-1/S}\le M\), then
\[
K=M\le2a^{-1/\gamma}\le2a^{-1/S};
\]
if \(M<a^{-1/S}\), rounding gives the same final bound. Thus
\[
K\le2a^{-1/S},
\qquad
aK^\alpha
\le2^\alpha a^{1-\alpha/S}
=2^\alpha a^{\beta/S}\le2.
\]

The refinement depth satisfies
\[
2^L\le R_L\le K\le2a^{-1/S},
\qquad
L\log2\le\log2-\frac{\log a}{S}.
\]
The exponential tangent inequality, applied at \(-\log a-1\), gives
\[
-\log a
\le\exp(-\log a-1)=a^{-1}\exp(-1),
\qquad
-a\log a\le\exp(-1).
\]
Consequently,
\[
aL
\le a+\frac{-a\log a}{S\log2}
\le1+\frac{1}{\exp(1)S\log2}
=D_*.
\]

The bound \(T_j^{1-p}\le a+a(R_j/K)^{\lambda_0}R_j^\alpha\) can be rewritten using
\[
a(R_j/K)^{\lambda_0}R_j^\alpha
=aK^\alpha(R_j/K)^{\alpha+\lambda_0}.
\]
Geometric summation, \(aK^\alpha\le2\), and the fact that \(A\) decreases on positive arguments now give
\[
\begin{aligned}
\sum_{j=1}^{L}T_j^{1-p}
&\le aL+aK^\alpha A(\alpha+\lambda_0)\\
&\le D_*+2A(\alpha).
\end{aligned}
\]
The same geometric estimates supply
\[
\sum_{j=1}^{L}a_j\le80A(\alpha),
\]
and, from \(d_j=110R_{j-1}^{-s_0}+20T_j^{1-p}\),
\[
\sum_{j=1}^{L}d_j
\le110A(s_0)+20D_*+40A(\alpha).
\]
Moreover, \(T_j\le T_0\) and \(2-p\ge0\) imply
\[
\sqrt{V_j}\le\sqrt{V_0},
\qquad
\sum_{j=1}^{L}a_j\sqrt{V_j}
\le80A(\alpha)\sqrt{V_0}.
\]
Since \(T_0\ge1\), we also have \(\sqrt{V_0}\ge1\).
Inserting these bounds into the singleton ledger gives
\[
\begin{aligned}
L_1
&=16\sqrt{V_0}
 +8\sum_{j=1}^{L}(d_j+a_j\sqrt{V_j})\\
&\le
\left\{16+880A(s_0)+160D_*+960A(\alpha)\right\}
\sqrt{V_0}
=L_*\sqrt{V_0}.
\end{aligned}
\]
Using \(V_0\le64a^{-t}\), we obtain
\[
L_1^2\le64L_*^2a^{-t}.
\]
Together with \(a^{-1/(2\gamma)}a^{-t}/s\le a^2\), this proves
\[
\sqrt M\,\frac{L_1^2}{s}
\le\sqrt2\,64L_*^2
\frac{a^{-1/(2\gamma)}a^{-t}}{s}
\le\sqrt2\,64L_*^2a^2.
\]
% lean: CausalSmith.Stat.FinitepHomogeneityDensegamma.ledger_multilevel_singleton_covariance_rate

\item Still in the multilevel branch, define the backward summability exponent by
\[
\eta=1-(\alpha+\lambda_0)t.
\]
Because \(\alpha<q\) and \(t\ge0\),
\[
\eta\ge1-(q+\lambda_0)t
=\frac{(p-1)(p+6)}{4p}
=\eta_0>0.
\]
Also,
\[
0\le\alpha t=1-\eta-\lambda_0t<1.
\]

The positive-power cutoff estimate follows from
\[
T_j\le2\max\{1,z_{\mathrm{tail},j}^{1/(p-1)}\}.
\]
Indeed, using \(2^{2-p}\le2\) gives
\[
V_j=32T_j^{2-p}
\le64\max\{1,z_{\mathrm{tail},j}^{t}\},
\]
and therefore
\[
\sqrt{V_j}
\le\sqrt{64}\left(1+z_{\mathrm{tail},j}^{t/2}\right).
\]
Multiplying by \(\sqrt{R_j}\) and expanding the local target yields
\[
\sqrt{R_jV_j}
\le\sqrt{64}
\left\{
\sqrt{R_j}
+a^{-t/2}K^{\lambda_0t/2}R_j^{\eta/2}
\right\}.
\]
Backward geometric summation implies
\[
\sum_{j=1}^{L}\sqrt{R_j}
\le\sqrt K\,A(1/2),
\qquad
\sum_{j=1}^{L}R_j^{\eta/2}
\le K^{\eta/2}A(\eta/2)
\le K^{\eta/2}A(\eta_0/2).
\]
The fine-rank amplitude estimate also gives
\[
\begin{aligned}
\sqrt K
&=(aK^\alpha)^{t/2}
  a^{-t/2}K^{(1-\alpha t)/2}\\
&\le2^{t/2}a^{-t/2}K^{(1-\alpha t)/2}.
\end{aligned}
\]
Since \(\lambda_0t+\eta=1-\alpha t\), summing the guarded increment bound proves
\[
\begin{aligned}
\sum_{j=1}^{L}\sqrt{R_jV_j}
&\le\sqrt{64}
 \left\{2^{t/2}A(1/2)+A(\eta_0/2)\right\}
 a^{-t/2}K^{(1-\alpha t)/2}\\
&=G_*a^{-t/2}K^{(1-\alpha t)/2}.
\end{aligned}
\]
Now
\[
0<\frac{1-\alpha t}{2}\le\frac12,
\qquad
t+\frac{1-\alpha t}{S}=\frac{1+\beta t}{S}.
\]
Thus \(K\le2a^{-1/S}\) converts the last bound into
\[
\sum_{j=1}^{L}\sqrt{R_jV_j}
\le\sqrt2\,G_*a^{-(1+\beta t)/(2S)}.
\]

Using the definition of \(L_2\) and \((x+y)^2\le2x^2+2y^2\) for its two nonnegative summands, we get
\[
\begin{aligned}
L_2^2
&=16\left(
\sqrt{MV_0}+\sum_{j=1}^{L}\sqrt{R_jV_j}
\right)^2\\
&\le32MV_0
 +32\left(\sum_{j=1}^{L}\sqrt{R_jV_j}\right)^2\\
&\le4096s\,a^{-t}
 +64G_*^2a^{-(1+\beta t)/S}.
\end{aligned}
\]
Here the last line uses \(M\le2s\) and \(V_0\le64a^{-t}\).
Since \(s(s-1)\ge s^2/2\), it follows that
\[
\frac{2L_2^2}{s(s-1)}
\le\frac{16384a^{-t}}{s}
 +\frac{256G_*^2a^{-(1+\beta t)/S}}{s^2}.
\]

For the second noise channel, direct algebra gives
\[
2+t=\frac1q,
\qquad
2S+1+\beta t+\frac{S}{2\gamma}=D_p,
\]
and hence
\[
E_4\left(2+\frac{1+\beta t}{S}+\frac{1}{2\gamma}\right)=2.
\]
Using \(E\le E_4\), we obtain
\[
\frac{a^{-1/(2\gamma)}a^{-(1+\beta t)/S}}{s^2}
=s^{E(1/(2\gamma)+(1+\beta t)/S)-2}
\le s^{-2E}=a^2.
\]
Combining this with the tail-channel balance gives
\[
\sqrt M\,\frac{2L_2^2}{s(s-1)}
\le\sqrt2\{16384+256G_*^2\}a^2.
\]
Adding the bound \(\sqrt M\,L_1^2/s\le\sqrt2\,64L_*^2a^2\) therefore proves
\[
\begin{aligned}
\sqrt M\,\Lambda_{n,v}
&=\sqrt M\left\{\frac{L_1^2}{s}
                 +\frac{2L_2^2}{s(s-1)}\right\}\\
&\le\sqrt2\{64L_*^2+16384+256G_*^2\}a^2\\
&\le A_*a^2.
\end{aligned}
\]
This includes \(L=0\), since all the increment sums then vanish. Together with the single-cutoff bound \(\sqrt M\,\Lambda_{n,v}\le24576\sqrt2\,a^2\le A_*a^2\), the covariance estimate holds in both branches.
% lean: CausalSmith.Stat.FinitepHomogeneityDensegamma.ledger_multilevel_covariance_rate_of_singleton

\item The preceding estimates now give
\[
M^{-\gamma}\le a,\qquad
B\le B_*a,\qquad
M^{1/4}\sqrt{\Lambda_{n,v}}
=\sqrt{\sqrt M\,\Lambda_{n,v}}
\le\sqrt{A_*}\,a.
\]
Consequently,
\[
R_{\mathrm{att}}(n,v)
\le\frac{16}{3}
\{16+5B_*+1024\sqrt{A_*}\}\,a.
\]
The relation \(n\le3s\), together with \(0<E\le1\), implies
\[
a=s^{-E}\le3^E n^{-E}\le3n^{-E}.
\]
Thus
\[
R_{\mathrm{att}}(n,v)\le C^{\mathrm{att}}_v n^{-E}.
\]

If \(C^{\mathrm{att}}_v n^{-E}<D_v\), the bound \(D_v\le C^{\mathrm{att}}_v n^{-E}\) for \(n\in\{2,3\}\) ensures \(n\ge4\). The last inequality then gives
\[
R_{\mathrm{att}}(n,v)<D_v,\qquad
d(P)\ge C^{\mathrm{att}}_v n^{-E}
\ \Longrightarrow\
d(P)\ge R_{\mathrm{att}}(n,v).
\]
The power assertion of
\cref{obj:thm:whole-null-calibration-resolved} therefore yields
\[
1-\mathsf R_n(P,\phi^*_{n,v})\le0.0075
\]
for every stipulated alternative.

To obtain the capped-radius bound, observe that
\(C^{\mathrm{att}}_v n^{-E}>0\) and that both error bounds are below \(1/10\).
When \(C^{\mathrm{att}}_v n^{-E}<D_v\), the explicit test is admissible in
\cref{obj:def:testing-risk}, so
\[
B_n\!\left(v,C^{\mathrm{att}}_v n^{-E}\right)
\le0.0075\le\frac1{10}.
\]
For an empty alternative set its worst-error supremum has value zero, giving the same inequality. The defining set in
\cref{obj:def:critical-radius} therefore contains
\(C^{\mathrm{att}}_v n^{-E}\).
When \(D_v\le C^{\mathrm{att}}_v n^{-E}\), its sentinel \(D_v\) instead gives
\[
r_n^*(v)\le D_v\le C^{\mathrm{att}}_v n^{-E}.
\]
The defining set is nonempty and bounded below by zero in both cases. Hence
\[
r_n^*(v)\le C^{\mathrm{att}}_v n^{-E}
\qquad(n\ge2).
\]
% lean: CausalSmith.Stat.FinitepHomogeneityDensegamma.original_record_attainment_finite

\item Let \(n\ge N^{\mathrm{att}}_v\). By the ceiling definition,
\[
n\ge
\max\left\{4,
\left(\frac{2C^{\mathrm{att}}_v}{d_0}\right)^{1/E}\right\}.
\]
All bases are positive, and \(-E<0\). Taking this negative power therefore gives
\[
n^{-E}
\le
\left[
\left(\frac{2C^{\mathrm{att}}_v}{d_0}\right)^{1/E}
\right]^{-E}
=\left(\frac{2C^{\mathrm{att}}_v}{d_0}\right)^{-1}.
\]
Thus
\[
C^{\mathrm{att}}_v n^{-E}
\le\frac{d_0}{2}<d_0,
\]
which proves the witness-threshold assertion. The bounds \(M\le2s\), \(K\le2s^2\), \(1\le T_0\le2s\), and \(1\le T_j\le T_0\) for \(1\le j\le L\) give the tuning assertions.
% lean: CausalSmith.Stat.FinitepHomogeneityDensegamma.attainment_threshold

\item For compact uniformity, let
\[
\mathcal K\subseteq\mathcal V
\]
be compact in the product topology. At every admissible tuple,
\[
p-1>0,\quad S>0,\quad\gamma>0,\quad q>0,\quad D_p>0,
\quad s_0>0,\quad\lambda_0>0,\quad\eta_0>0.
\]
In particular, every geometric denominator appearing in the constants satisfies
\[
1-2^{-x}>0
\qquad
\left(
x\in
\{\alpha,\lambda_0,s_0,1/2,\eta_0/2\}
\right).
\]
The displayed formulas consequently show that \(E(v)\) and
\(C^{\mathrm{att}}_v\) are continuous at every admissible tuple: their operations are sums, products, quotients with nonzero denominators, minima, square roots, and real powers with positive bases.

Define the real threshold target on the admissible domain by
\[
\Xi(v)
=\max\left\{4,
\left(\frac{2C^{\mathrm{att}}_v}{d_0}\right)^{1/E(v)}
\right\}
\qquad(v\in\mathcal V).
\]
Since \(E(v)>0\), \(C^{\mathrm{att}}_v>0\), and \(d_0>0\), this function is continuous there as well.

If \(\mathcal K\) is nonempty, compactness supplies tuples
\(v_C,v_E,v_N\in\mathcal K\) satisfying
\[
C^{\mathrm{att}}_{v_C}
=\max_{v\in\mathcal K}C^{\mathrm{att}}_v,
\qquad
E(v_E)=\min_{v\in\mathcal K}E(v),
\qquad
\Xi(v_N)=\max_{v\in\mathcal K}\Xi(v).
\]
Choose
\[
C_{\mathcal K}=C^{\mathrm{att}}_{v_C},
\qquad
N_{\mathcal K}=\Xi(v_N)+1,
\qquad
e_{\mathcal K}=E(v_E)>0.
\]
The ceiling bound gives, for every \(v\in\mathcal K\),
\[
C^{\mathrm{att}}_v\le C_{\mathcal K},
\qquad
N^{\mathrm{att}}_v=\lceil\Xi(v)\rceil
<\Xi(v)+1\le N_{\mathcal K},
\qquad
e_{\mathcal K}\le E(v).
\]
If \(\mathcal K\) is empty, choose
\[
(C_{\mathcal K},N_{\mathcal K},e_{\mathcal K})=(0,0,1);
\]
the universal bounds are then vacuous.
% lean: CausalSmith.Stat.FinitepHomogeneityDensegamma.attainment_compact_uniformity

\item Finally, let \(w=(\alpha,\beta,\gamma)\in\mathcal W\), and define its matched admissible tuple by
\[
v_{\mathrm b}=(2,\alpha,\beta,\gamma)\in\mathcal V.
\]
At this tuple,
\[
q=\frac12,\qquad t=0,\qquad S=\alpha+\beta,
\]
so substitution into the two exponent formulas gives
\[
E_0(v_{\mathrm b})=\frac{2\gamma}{4\gamma+1},
\qquad
E_4(v_{\mathrm b})
=\frac{2S}{1+2S+S/(2\gamma)}.
\]
Taking their minimum proves
\[
E(v_{\mathrm b})=E_{\mathrm b}.
\]

The definitions in
\cref{obj:def:bounded-model,obj:def:bounded-null,obj:def:sharp-frontier-scales}
give
\[
\mathcal M_w^{\mathrm b}\subseteq\mathcal M_{v_{\mathrm b}},
\qquad
H_0^{\mathrm b}(w)\subseteq H_0(v_{\mathrm b}),
\qquad
\phi^{*,\mathrm b}_{n,w}=\phi^*_{n,v_{\mathrm b}},
\]
and
\[
C^{\mathrm{att}}_{(2,w)}n^{-E_{\mathrm b}}
=C^{\mathrm{att}}_{v_{\mathrm b}}n^{-E(v_{\mathrm b})}>0.
\]
Moreover, \cref{obj:prop:bounded-submodel-comparison} supplies the exact cap identity
\[
D_w^{\mathrm b}=D_{v_{\mathrm b}}\ge d_0.
\]
The original-model size bound therefore applies to every bounded null law. If
\[
C^{\mathrm{att}}_{(2,w)}n^{-E_{\mathrm b}}<D_w^{\mathrm b},
\]
the cap identity makes the original-model power condition hold at
\(v_{\mathrm b}\). Thus every bounded-model law with distance at least this level satisfies
\[
1-\mathsf R_n(P,\phi^{*,\mathrm b}_{n,w})\le0.0075.
\]

When this positive level is below \(D_w^{\mathrm b}\), the resulting bounded-model risk bound is
\[
B_n^{\mathrm b}
\!\left(w,C^{\mathrm{att}}_{(2,w)}n^{-E_{\mathrm b}}\right)
\le0.0075\le\frac1{10},
\]
so the level belongs to the defining set of
\cref{obj:def:bounded-radius}. When it is at least \(D_w^{\mathrm b}\), the sentinel gives the radius bound directly. Therefore, for every \(n\ge2\),
\[
r_n^{*,\mathrm b}(w)
\le C^{\mathrm{att}}_{(2,w)}n^{-E_{\mathrm b}}.
\]
% lean: CausalSmith.Stat.FinitepHomogeneityDensegamma.bounded_record_attainment_finite
\end{enumerate}
\end{proof}

\begin{proof}[Proof of \cref{obj:thm:tail-essential-full-record-converse}]
Fix an admissible tuple \(v=(p,\alpha,\beta,\gamma)\), and put
\(w=(\alpha,\beta,\gamma)\). Throughout, use the quantities and selected
constructions of
\cref{obj:def:sharp-frontier-scales,obj:def:converse-handle}.

\begin{enumerate}
\item
Fix \(n\ge2\), and first suppose \(F(p)<1\).
By \cref{obj:lem:rough-full-record-sharp-lower}, both selected priors are
normalized, and
\[
\begin{gathered}
\operatorname{supp}_+(\pi_{0,n,v})\subseteq H_0(v),\\
\operatorname{supp}_+(\pi_{1,n,v})
\subseteq
\{P\in\mathcal M_v:c_vn^{-E_4(p)}\le d(P)<d_0\},\\
d_0\le D_v,\qquad
\operatorname{TV}(\mathbb P_{0,n,v},\mathbb P_{1,n,v})<\frac12,\qquad
B_n(v,c_vn^{-E_4(p)})\ge\frac25.
\end{gathered}
\]
These conclusions cover every \(n\ge2\), including the paired-tent
selection below \(N_{\mathrm{low}}(v)\).

The scale identification follows by subtracting the two exponents:
\[
\begin{aligned}
E_4(p)-E_0(p)
&=
\frac{2S(2\gamma+q)-2\gamma qD_p}
     {(2\gamma+q)D_p}\\
&=
\frac{2\gamma q}{(2\gamma+q)D_p}\{F(p)-1\}.
\end{aligned}
\]
Admissibility gives \(q>0\), \(S>0\), \(\gamma>0\), and \(D_p>0\).
Thus \(F(p)<1\) implies
\[
E(p)=E_4(p),\qquad \rho_n(v)=n^{-E_4(p)}.
\]

We also verify the arm support at the selected magnitude. When
\(n\ge N_{\mathrm{low}}(v)\), the rank conditions are supplied by
\cref{obj:lem:rough-full-record-sharp-lower}, and the tuning agrees with
\cref{obj:lem:copula-frame-legality}:
\[
b=\frac{K^{-\beta}}{256},\qquad
\varepsilon_{n,v}=(16b)^{1/q}
=16^{-1/q}K^{-\beta/q},\qquad
L_{n,v}=\varepsilon_{n,v}^{-1/p}.
\]
The normalized table in \cref{obj:def:marked-table} assigns its outcome
mass to \(0,-L_{n,v},L_{n,v}\). Conditioning on either treatment arm
therefore retains all its mass on those three points. When
\(n<N_{\mathrm{low}}(v)\), the selected paired-tent construction has the
same arm-support conclusion by
\cref{obj:lem:paired-tent-full-record-lower}. Consequently, in either
selection,
\[
Q_{a,P}\bigl(\{0,-L_{n,v},L_{n,v}\}\mid x\bigr)=1
\]
for each \(a\in\{0,1\}\), for \(\lambda\)-almost every \(x\), and for every
positive-weight law in either prior.

Finally, normalization and nonnegative weights ensure that the
alternative prior has a positive-weight law. Choosing such a law \(P\)
gives
\[
c_v\rho_n(v)\le d(P)<d_0\le D_v.
\]
Moreover, \(N_{\mathrm{low}}(v)\ge2\), \(k_v>0\), and
\(c_0^{\mathrm e}>0\), so the defining minimum for \(c_v\) is positive.
Since \(n>0\), also \(\rho_n(v)>0\). Hence
\[
0<c_v\rho_n(v)<D_v,
\]
and every positive-weight alternative law has distance strictly below
\(D_v\). This proves the first group of conclusions in the rough phase.
% lean: CausalSmith.Stat.FinitepHomogeneityDensegamma.rough_phase_converse_assembly

\item
Suppose instead \(F(p)\ge1\). The exponent-difference identity from the
preceding step gives
\[
E(p)=E_0(p),\qquad
c_v=c_0^{\mathrm e},\qquad
c_v\rho_n(v)=c_0^{\mathrm e}n^{-E_0(p)}.
\]
The selector of \cref{obj:def:converse-handle} chooses the paired-tent
priors for every \(n\ge2\). Their finite sign distributions are
normalized. All null sign configurations give the same zero-effect
null law \(P_0\), so
\[
\pi_{0,n,v}=\delta_{P_0},\qquad
\mathbb P_{0,n,v}=P_0^{\otimes n}.
\]
By \cref{obj:lem:paired-tent-full-record-lower},
\[
\begin{gathered}
P_0\in H_0(v),\\
\operatorname{supp}_+(\pi_{1,n,v})
\subseteq
\{P\in\mathcal M_v:
c_0^{\mathrm e}n^{-E_0(p)}\le d(P)<d_0\},\\
d_0\le D_v,\qquad
\operatorname{TV}(\mathbb P_{0,n,v},\mathbb P_{1,n,v})<\frac12,\\
B_n(v,c_0^{\mathrm e}n^{-E_0(p)})\ge\frac25.
\end{gathered}
\]
That construction also gives the required three-point conditional arm
support at \(L_{n,v}\). Its alternative prior has a positive-weight
law; for any such law,
\[
c_v\rho_n(v)\le d(P)<d_0\le D_v.
\]
Positivity of \(c_0^{\mathrm e}\) and \(n^{-E_0(p)}\) yields
\(0<c_v\rho_n(v)<D_v\). Substituting the displayed scale identity proves
the first group of conclusions in this phase as well. The equality
case \(F(p)=1\) is included here.
% lean: CausalSmith.Stat.FinitepHomogeneityDensegamma.nonrough_phase_converse_assembly

\item
We next prove divergence of the selected magnitude, again splitting at
\(F(p)<1\). In that phase, for all
\(n\ge N_{\mathrm{low}}(v)\), the selected fine rank satisfies
\[
K\ge H_{\mathrm{fine}}n^{2/D_p}\longrightarrow\infty,
\]
and moment normalization gives
\[
L_{n,v}
=
\left(16^{-1/q}K^{-\beta/q}\right)^{-1/p}
=
16^{1/(pq)}K^{\beta/(pq)}.
\]
Both \(2/D_p\) and \(\beta/(pq)\) are positive, so
\(L_{n,v}\to\infty\). The finitely many preceding selections leave this
limit unchanged.

When \(F(p)\ge1\), the paired-tent formula applies throughout:
\[
N=2\left\lceil n^{2q/(2\gamma+q)}\right\rceil
\ge n^{2q/(2\gamma+q)}\longrightarrow\infty,
\qquad
L_{n,v}=N^{\gamma/(p-1)}.
\]
Here \(2q/(2\gamma+q)>0\) and \(\gamma/(p-1)>0\), proving the same
divergence.
% lean: CausalSmith.Stat.FinitepHomogeneityDensegamma.magnitude_tendsto_atTop

\item
Assume \(p<2\). At the matched smoothness tuple, write
\[
D_2=1+2S+\frac{S}{2\gamma}.
\]
The identity \(q^{-1}=2+(2-p)/(p-1)\) gives
\[
D_p=D_2+\frac{\beta(2-p)}{p-1}.
\]
Direct subtraction then yields the two strictly positive differences
\[
E_0(2)-E_0(p)
=
\frac{4\gamma^2(2-p)}
     {(4\gamma+1)(2\gamma p+p-1)}>0,
\]
\[
E_4(2)-E_4(p)
=
\frac{2S\beta(2-p)}
     {(p-1)D_pD_2}>0.
\]
All denominators are positive by admissibility. Since \(E(p)\) is the
minimum of its two component exponents,
\[
E(p)\le E_0(p)<E_0(2),\qquad
E(p)\le E_4(p)<E_4(2).
\]
Thus the exponent gap, defined on this face by
\[
\delta_v=E(2)-E(p),
\]
is strictly positive. For every positive integer \(n\),
\[
\frac{\rho_n(v)}{\rho_n^{\mathrm b}(w)}
=n^{\delta_v}\longrightarrow\infty.
\]
% lean: CausalSmith.Stat.FinitepHomogeneityDensegamma.matched_scale_ratio_tendsto

\item
Continue with \(p<2\). Since \(c_v>0\), the preceding limit permits an
integer \(N_{\mathrm{dom}}(v)\in\mathbb N\) such that
\[
\frac{\rho_n(v)}{\rho_n^{\mathrm b}(w)}
\ge\frac{C_w^{\mathrm b}}{c_v}
\qquad
\text{for every integer }n\ge N_{\mathrm{dom}}(v)\text{ with }n>0.
\]
Multiplying by the positive bounded scale and by \(c_v\) gives
\[
C_w^{\mathrm b}\rho_n^{\mathrm b}(w)\le c_v\rho_n(v).
\]
Define the required threshold by
\[
N(v)=\max\{2,N_{\mathrm{dom}}(v),N_w^{\mathrm b}\}.
\]
For \(n\ge N(v)\), \cref{obj:thm:original-record-attainable-rate}, applied
at \((2,w)\), and \cref{obj:prop:bounded-submodel-comparison} give
\[
C_w^{\mathrm b}\rho_n^{\mathrm b}(w)<d_0\le D_w^{\mathrm b}.
\]
The bounded total test therefore satisfies
\[
\mathsf R_n(P,\phi^{*,\mathrm b}_{n,w})\le0.0075
\quad(P\in H_0^{\mathrm b}(w)),
\]
and
\[
1-\mathsf R_n(P,\phi^{*,\mathrm b}_{n,w})\le0.0075
\quad
(P\in\mathcal M_w^{\mathrm b},\ d(P)\ge c_v\rho_n(v)).
\]
Its size is below \(1/10\), so
\[
B_n^{\mathrm b}(w,c_v\rho_n(v))\le0.0075.
\]

Let \(\pi_0,\pi_1\) be any normalized finite priors satisfying the bounded
support conditions in the statement. We establish the two-prior
inequality used to compare their mixtures. For any randomized test
\(\phi\), define its seed average on the complete record space by
\[
f_\phi:
\bigl([0,1]\times\{0,1\}\times\mathbb R\bigr)^n\longrightarrow[0,1],
\qquad
f_\phi(\mathcal D_n)=\int_{[0,1]}\phi(\mathcal D_n,U)\,\lambda(dU).
\]
Bounded product integration and finite-sum linearity give, for either
prior,
\[
\int f_\phi\,d\mathbb P_{\pi_\nu,n}
=
\int \mathsf R_n(P,\phi)\,\pi_\nu(dP),
\qquad \nu\in\{0,1\}.
\]
The bounded-function form of the total-variation inequality gives
\[
\left|
\int f_\phi\,d\mathbb P_{\pi_1,n}
-
\int f_\phi\,d\mathbb P_{\pi_0,n}
\right|
\le
\operatorname{TV}(\mathbb P_{\pi_0,n},\mathbb P_{\pi_1,n}).
\]
If \(\phi\) has size at most \(1/10\) on the bounded null, normalization
and the null support condition imply
\[
\int \mathsf R_n(P,\phi)\,\pi_0(dP)\le\frac1{10}.
\]
The alternative support condition and normalization consequently yield
\[
\begin{aligned}
\sup_{\substack{P\in\mathcal M_w^{\mathrm b}\\
                 d(P)\ge c_v\rho_n(v)}}
\{1-\mathsf R_n(P,\phi)\}
&\ge
\int \{1-\mathsf R_n(P,\phi)\}\,\pi_1(dP)\\
&\ge
\frac9{10}
-\operatorname{TV}(\mathbb P_{\pi_0,n},\mathbb P_{\pi_1,n}).
\end{aligned}
\]
Weights equal to zero contribute zero in these averages; every
positive-weight law satisfies the stipulated support condition.
Normalization ensures a positive-weight alternative law, and the zero
rejection map is an admissible test. Taking the infimum over
null-sized tests therefore gives
\[
\frac9{10}
-\operatorname{TV}(\mathbb P_{\pi_0,n},\mathbb P_{\pi_1,n})
\le B_n^{\mathrm b}(w,c_v\rho_n(v)).
\]
Combining this with the attained risk bound proves
\[
\operatorname{TV}(\mathbb P_{\pi_0,n},\mathbb P_{\pi_1,n})
\ge0.8925\ge\frac12.
\]
The threshold depends on \(v\), and the argument applies to every
normalized finite prior pair in the statement.
% lean: CausalSmith.Stat.FinitepHomogeneityDensegamma.bounded_tail_essentiality

\item
Now assume \(p=2\), and fix \(n\ge2\). Then
\(v=(2,\alpha,\beta,\gamma)\).
We first verify the original-model support conclusions for the selected
bounded priors, splitting at \(F(2)<1\).

In the rough phase, \cref{obj:lem:rough-full-record-sharp-lower} supplies
\[
\begin{gathered}
\operatorname{supp}_+(\pi^{\mathrm b}_{0,n,v})
\subseteq H_0^{\mathrm b}(w),\\
\operatorname{supp}_+(\pi^{\mathrm b}_{1,n,v})
\subseteq
\{P\in\mathcal M_w^{\mathrm b}:
c_vn^{-E_4(2)}\le d(P)<d_0\}.
\end{gathered}
\]
By the phase identity, \(\rho_n(v)=n^{-E_4(2)}\).
The model inclusion of \cref{obj:prop:bounded-submodel-comparison}
preserves the effect, and the null constancy condition is the same in
both models. Hence
\[
\begin{gathered}
\operatorname{supp}_+(\pi^{\mathrm b}_{0,n,v})\subseteq H_0(v),\\
\operatorname{supp}_+(\pi^{\mathrm b}_{1,n,v})
\subseteq
\{P\in\mathcal M_v:c_v\rho_n(v)\le d(P)\}.
\end{gathered}
\]

In the complementary phase \(F(2)\ge1\), the selected bounded priors are
the signed-binary paired-tent priors. To identify their null support,
define the balanced deterministic zero-effect law by
\[
P_{\mathrm{zero}}(dx,dA,dy)
=
\lambda(dx)\,
\frac{\delta_0+\delta_1}{2}(dA)\,
\delta_0(dy).
\]
It has propensity \(1/2\), baseline mean zero, and effect zero. Its
smoothness conditions hold because these functions are constant, and
its raw moment is zero. Its signed-binary conversion from
\cref{obj:prop:bounded-submodel-comparison} is
\[
P_{\mathrm{zero}}^{\mathrm{bin}}(dx,dA,dy)
=
\lambda(dx)\,
\frac{\delta_0+\delta_1}{2}(dA)\,
\frac{\delta_{-1}+\delta_1}{2}(dy).
\]
This conversion preserves the means and effect, and its conditional
second moment is one. Thus
\(P_{\mathrm{zero}}^{\mathrm{bin}}\in H_0(v)\).
Every selected null sign configuration gives precisely this law, so
\[
\operatorname{supp}_+(\pi^{\mathrm b}_{0,n,v})
\subseteq\{P_{\mathrm{zero}}^{\mathrm{bin}}\}\subseteq H_0(v).
\]
For the alternatives,
\cref{obj:lem:paired-tent-full-record-lower} and
\cref{obj:prop:bounded-submodel-comparison} give original-model
membership and
\[
d(P)\ge c_0^{\mathrm e}n^{-E_0(2)}
=c_v\rho_n(v)
\qquad
(P\in\operatorname{supp}_+(\pi^{\mathrm b}_{1,n,v})).
\]
This proves the original-model support conclusions in both phases.
% lean: CausalSmith.Stat.FinitepHomogeneityDensegamma.second_moment_rough_prior_support

\item
The bounded risk lower bound follows from the same phase split.
When \(F(2)<1\),
\cref{obj:lem:rough-full-record-sharp-lower} gives
\[
B_n^{\mathrm b}(w,c_vn^{-E_4(2)})\ge\frac25.
\]
When \(F(2)\ge1\),
\cref{obj:lem:paired-tent-full-record-lower} gives
\[
B_n^{\mathrm b}(w,c_0^{\mathrm e}n^{-E_0(2)})\ge\frac25.
\]
The phase identities identify each displayed separation with
\(c_v\rho_n(v)\). The risk comparison in
\cref{obj:prop:bounded-submodel-comparison} consequently yields
\[
\frac25
\le B_n^{\mathrm b}(w,c_v\rho_n(v))
\le B_n(v,c_v\rho_n(v)).
\]
% lean: CausalSmith.Stat.FinitepHomogeneityDensegamma.second_moment_testingRisk_lower

\item
We assemble the remaining second-moment conclusions. Evaluation at
\(v=(2,w)\), together with
\cref{obj:prop:bounded-submodel-comparison}, gives
\[
\begin{gathered}
\rho_n(v)=\rho_n^{\mathrm b}(w),\qquad
c_v=c_w^{\mathrm b},\qquad
C_v=C_w^{\mathrm b},\qquad
N_v=N_w^{\mathrm b},\\
D_v=D_w^{\mathrm b}.
\end{gathered}
\]
For \(F(2)<1\), the selected-prior conclusions of
\cref{obj:lem:rough-full-record-sharp-lower} apply for every \(n\ge2\),
covering the bounded rough construction above its threshold and the
signed-binary paired-tent construction below it. For \(F(2)\ge1\), use
the signed-binary conclusions of
\cref{obj:lem:paired-tent-full-record-lower}. The signed-binary model
inclusion of \cref{obj:prop:bounded-submodel-comparison} supplies bounded
membership, with null constancy retained. In either case, both selected
bounded priors are normalized and
\[
\begin{gathered}
\operatorname{supp}_+(\pi^{\mathrm b}_{0,n,v})
\subseteq H_0^{\mathrm b}(w),\\
\operatorname{supp}_+(\pi^{\mathrm b}_{1,n,v})
\subseteq
\{P\in\mathcal M_w^{\mathrm b}:
c_v\rho_n(v)\le d(P)<d_0\},\\
\operatorname{TV}(\mathbb P^{\mathrm b}_{0,n,v},
                  \mathbb P^{\mathrm b}_{1,n,v})<\frac12,\qquad
B_n^{\mathrm b}(w,c_v\rho_n(v))\ge\frac25.
\end{gathered}
\]
For any finite representation of either prior, nonnegative weights
summing to one satisfy
\[
\sum_{j:\,\omega_j>0}\omega_j=\sum_j\omega_j=1.
\]
Indeed, every omitted weight is zero. At least one weight is positive,
and combining weights of coincident laws preserves their sum.
Therefore
\[
\operatorname{supp}_+(\pi^{\mathrm b}_{\nu,n,v})\ne\varnothing,
\qquad
\sum_{P\in\operatorname{supp}_+(\pi^{\mathrm b}_{\nu,n,v})}
\pi^{\mathrm b}_{\nu,n,v}(\{P\})=1,
\qquad \nu\in\{0,1\}.
\]
Choose a positive-weight alternative law \(P\). Its placement, and the
witness lower bound from
\cref{obj:prop:bounded-submodel-comparison}, give
\[
0<c_v\rho_n(v)\le d(P)<d_0\le D_w^{\mathrm b}.
\]
% lean: CausalSmith.Stat.FinitepHomogeneityDensegamma.second_moment_bounded_receipt

\item
We obtain the lower capped-radius bound first in the bounded model.
Every successful bounded separation satisfies
\[
r_{\mathrm{succ}}\in(0,D_w^{\mathrm b}),
\qquad
B_n^{\mathrm b}(w,r_{\mathrm{succ}})\le\frac1{10}.
\]
If \(r_{\mathrm{succ}}<c_v\rho_n(v)\), containment of the separated
alternative classes gives
\[
B_n^{\mathrm b}(w,c_v\rho_n(v))
\le B_n^{\mathrm b}(w,r_{\mathrm{succ}})
\le\frac1{10},
\]
contradicting the lower bound \(2/5\). The cap
\(D_w^{\mathrm b}\) also exceeds \(c_v\rho_n(v)\). Thus every member of
the nonempty set defining the capped infimum is at least
\(c_v\rho_n(v)\), and
\[
c_v\rho_n(v)\le r_n^{*,\mathrm b}(w).
\]
Now \cref{obj:prop:bounded-submodel-comparison} gives
\[
c_v\rho_n(v)
\le r_n^{*,\mathrm b}(w)
=r_n^{*,\mathrm{bin}}(w)
\le r_n^*(v).
\]
The upper capped-radius bounds and both size bounds follow from
\cref{obj:thm:original-record-attainable-rate}:
\[
r_n^*(v)\le C_v\rho_n(v),\qquad
r_n^{*,\mathrm b}(w)\le C_w^{\mathrm b}\rho_n^{\mathrm b}(w),
\]
\[
\mathsf R_n(P,\phi^*_{n,v})\le0.0075
\quad(P\in H_0(v)),\qquad
\mathsf R_n(P,\phi^{*,\mathrm b}_{n,w})\le0.0075
\quad(P\in H_0^{\mathrm b}(w)).
\]
Substituting the matching scale and multiplier identities proves both
asserted two-sided radius bounds.
% lean: CausalSmith.Stat.FinitepHomogeneityDensegamma.second_moment_capped_lower

\item
Suppose \(C_v\rho_n(v)<D_v\). The matching identities give
\[
C_w^{\mathrm b}\rho_n^{\mathrm b}(w)<D_w^{\mathrm b}.
\]
To prove nonemptiness in the bounded model, suppose its separated class
were empty. Every bounded-model law would then satisfy
\[
d(P)<C_w^{\mathrm b}\rho_n^{\mathrm b}(w).
\]
The bounded model is nonempty, as witnessed by a positive-weight law
from the selected alternative prior. The displayed bound would imply
\[
D_w^{\mathrm b}
\le C_w^{\mathrm b}\rho_n^{\mathrm b}(w),
\]
contradicting the strict inequality. Hence there is a bounded-model law
with
\[
d(P)\ge C_w^{\mathrm b}\rho_n^{\mathrm b}(w)=C_v\rho_n(v).
\]
Its original-model membership follows from
\cref{obj:prop:bounded-submodel-comparison}, proving nonemptiness of both
separated classes.

Under this same strict-separation condition,
\cref{obj:thm:original-record-attainable-rate} supplies
\[
1-\mathsf R_n(P,\phi^*_{n,v})\le0.0075
\quad(P\in\mathcal M_v,\ d(P)\ge C_v\rho_n(v)),
\]
and
\[
1-\mathsf R_n(P,\phi^{*,\mathrm b}_{n,w})\le0.0075
\quad
(P\in\mathcal M_w^{\mathrm b},\
d(P)\ge C_w^{\mathrm b}\rho_n^{\mathrm b}(w)).
\]
% lean: CausalSmith.Stat.FinitepHomogeneityDensegamma.second_moment_face_package

\item
Finally, let \(n\ge N_v\). Positivity of the component exponents implies
\(E(2)>0\), and
\cref{obj:thm:original-record-attainable-rate} gives \(C_v>0\).
The threshold formula yields
\[
n\ge
\max\left\{4,\left(\frac{2C_v}{d_0}\right)^{1/E(2)}\right\}.
\]
Since raising positive numbers to the power \(-E(2)\) reverses their
order,
\[
\begin{aligned}
C_v\rho_n(v)
&=C_vn^{-E(2)}\\
&\le
C_v\left\{
\left(\frac{2C_v}{d_0}\right)^{1/E(2)}
\right\}^{-E(2)}
=\frac{d_0}{2}.
\end{aligned}
\]
Also, the witness-threshold conclusion of
\cref{obj:thm:original-record-attainable-rate} gives
\[
C_v\rho_n(v)<d_0\le D_v.
\]
The matching identities therefore imply
\[
C_w^{\mathrm b}\rho_n^{\mathrm b}(w)<D_w^{\mathrm b}.
\]
At \(n=2,3\), the small-sample branch of
\cref{obj:def:explicit-score-ledger} sets the total rule to zero.
Evaluating that same branch at \(v=(2,w)\) gives, for every sample \(z\),
\[
\phi^*_{n,v}(z)=0,\qquad
\phi^{*,\mathrm b}_{n,w}(z)=0.
\]
This completes all conclusions on the second-moment face.
% lean: CausalSmith.Stat.FinitepHomogeneityDensegamma.second_moment_attainment_half
\end{enumerate}
\end{proof}

\begin{proof}[Proof of \cref{obj:thm:heavy-tail-comparison}]
\begin{enumerate}
\item Fix an admissible public tuple and its matched smoothness tuple:
\[
v=(p,\alpha,\beta,\gamma)\in\mathcal V,
\qquad
w=(\alpha,\beta,\gamma),
\qquad
v_2=(2,\alpha,\beta,\gamma)\in\mathcal V.
\]
Use the quantities of \cref{obj:def:sharp-frontier-scales}, in particular
\[
q=\frac{p-1}{p},
\qquad
S=\alpha+\beta,
\qquad
E(p)=\min\{E_0(p),E_4(p)\}.
\]
Admissibility gives
\[
p-1>0,\qquad
0<q\le\frac12,\qquad
S>0,\qquad
\gamma>0,\qquad
2\gamma+q>0,\qquad
D_p>0.
\]
Indeed, \(q\le1/2\) follows by multiplying by \(p>0\) and using \(p\le2\); every term added to \(1\) in \(D_p\) is positive. Consequently,
\[
E_0(p)>0,\qquad E_4(p)>0,\qquad E(p)>0.
\]
The same conclusions apply at \(v_2\).

We first establish the phase selection needed for the lower bounds. Since
\[
pq=p-1,
\qquad
\frac1q-1=\frac1{p-1},
\]
subtracting the two channel exponents gives
\[
\begin{aligned}
E_4(p)-E_0(p)
&=\frac{2S(2\gamma+q)-2\gamma qD_p}
        {(2\gamma+q)D_p}\\
&=\frac{2\gamma q}{(2\gamma+q)D_p}
  \left(
  \frac{2S}{q}-2\alpha-\frac{\beta}{q}
  +\frac{S}{2\gamma}-1
  \right)\\
&=\frac{2\gamma q}{(2\gamma+q)D_p}\bigl(F(p)-1\bigr).
\end{aligned}
\]
The coefficient multiplying \(F(p)-1\) is strictly positive. Thus
\[
\begin{aligned}
F(p)\ge1&\ \Longrightarrow\ E_0(p)\le E_4(p)
                       \ \Longrightarrow\ E(p)=E_0(p),\\
F(p)<1&\ \Longrightarrow\ E_4(p)<E_0(p)
                       \ \Longrightarrow\ E(p)=E_4(p).
\end{aligned}
\]
These implications also hold with \(p=2\).

The definition of the lower multiplier gives
\[
N_{\mathrm{low}}(v)\ge2,
\qquad
c_0^{\mathrm e}>0,
\qquad
k_v>0,
\qquad
c_v>0,
\]
because every power occurring in either branch has a positive base. Applying this argument at \(v_2\), and applying \cref{obj:thm:original-record-attainable-rate} at both tuples, yields
\[
c_w^{\mathrm b}>0,\qquad C_v>0,\qquad C_w^{\mathrm b}>0.
\]
For every \(n\ge2\), positivity of real powers also gives
\[
\rho_n(v)>0,
\qquad
\rho_n^{\mathrm b}(w)>0.
\]
% lean: CausalSmith.Stat.FinitepHomogeneityDensegamma.phase_channel_difference
% lean: CausalSmith.Stat.FinitepHomogeneityDensegamma.phase_channel_selection
% lean: CausalSmith.Stat.FinitepHomogeneityDensegamma.cLower_pos
% lean: CausalSmith.Stat.FinitepHomogeneityDensegamma.matched_scales_pos

\item Fix an integer \(n\ge2\). We obtain the bounded lower bound by splitting according to the phase at \(v_2\).

If \(F(2)<1\), then
\[
E(2)=E_4(2),
\qquad
c_w^{\mathrm b}\rho_n^{\mathrm b}(w)
=c_w^{\mathrm b}n^{-E_4(2)}.
\]
The bounded conclusion of \cref{obj:lem:rough-full-record-sharp-lower} supplies normalized finite alternative witnesses and
\[
B_n^{\mathrm b}
 \bigl(w,c_w^{\mathrm b}\rho_n^{\mathrm b}(w)\bigr)
\ge\frac25.
\]
A normalized finite prior has a positive-weight atom: otherwise every nonnegative weight would be zero, contradicting that their sum is one. Choosing such an alternative atom gives
\[
\exists P\in\mathcal M_w^{\mathrm b}:
\quad
c_w^{\mathrm b}\rho_n^{\mathrm b}(w)\le d(P).
\]

If \(F(2)\ge1\), the public definitions and the phase selection give
\[
E(2)=E_0(2)=\frac{2\gamma}{4\gamma+1},
\qquad
c_w^{\mathrm b}=c_0^{\mathrm e}
=\frac{4^{-\gamma}}{16\sqrt3}.
\]
The signed-binary witnesses in \cref{obj:lem:paired-tent-full-record-lower} therefore give the same risk bound and a positive-weight alternative satisfying
\[
P\in\mathcal M_w^{\mathrm{bin}}
\subseteq\mathcal M_w^{\mathrm b},
\qquad
c_w^{\mathrm b}\rho_n^{\mathrm b}(w)\le d(P),
\]
where the inclusion follows from \cref{obj:prop:bounded-submodel-comparison}. Hence, in both cases,
\[
\exists P\in\mathcal M_w^{\mathrm b}:
\quad c_w^{\mathrm b}\rho_n^{\mathrm b}(w)\le d(P),
\qquad
B_n^{\mathrm b}
 \bigl(w,c_w^{\mathrm b}\rho_n^{\mathrm b}(w)\bigr)
\ge\frac25.
\]

We also verify that this bounded lower separation lies below the bounded cap. For any admissible tuple, the branch \(F(p)\ge1\) gives \(c_v=c_0^{\mathrm e}\). In the other branch, the phase identity gives \(E_0(p)-E_4(p)>0\), so
\[
N_{\mathrm{low}}(v)^{-(E_0(p)-E_4(p))}\le1,
\qquad
c_v\le
c_0^{\mathrm e}N_{\mathrm{low}}(v)^{-(E_0(p)-E_4(p))}
\le c_0^{\mathrm e}.
\]
Since \(E(p)>0\) and \(n\ge2\),
\[
0<\rho_n(v)\le1,
\qquad
0<c_v\rho_n(v)\le c_0^{\mathrm e}
\le\frac1{16\sqrt3}
<\frac1{8\sqrt2}=d_0.
\]
The strict numerical inequality follows from
\(\sqrt2<2\sqrt3\), obtained by squaring these positive quantities. Applying this bound at \(v_2\), and using
\(d_0\le D_w^{\mathrm b}\) from \cref{obj:prop:bounded-submodel-comparison}, gives
\[
0<c_w^{\mathrm b}\rho_n^{\mathrm b}(w)
<d_0\le D_w^{\mathrm b}.
\]

For the original model, \cref{obj:thm:tail-essential-full-record-converse} supplies
\[
0<c_v\rho_n(v)<D_v,
\qquad
B_n\bigl(v,c_v\rho_n(v)\bigr)\ge\frac25,
\]
and a normalized finite alternative prior. Choosing a positive-weight atom as above gives
\[
\exists P\in\mathcal M_v:
\quad c_v\rho_n(v)\le d(P)<D_v.
\]

We now pass from these failures to the capped radii. For
\(0<r_1\le r_2<D_v\), with a nonempty alternative class at \(r_2\), inclusion of the alternative classes gives, for every admissible test,
\[
\sup_{\substack{P\in\mathcal M_v\\ d(P)\ge r_2}}
 \{1-\mathsf R_n(P,\phi)\}
\le
\sup_{\substack{P\in\mathcal M_v\\ d(P)\ge r_1}}
 \{1-\mathsf R_n(P,\phi)\}.
\]
The errors lie in \([0,1]\), and the null-size constraint is unchanged. Taking infima over tests therefore yields
\[
B_n(v,r_2)\le B_n(v,r_1).
\]
The identical argument gives bounded-risk monotonicity.

If a successful separation \(\widetilde r\in(0,D_v)\) satisfied
\(\widetilde r<c_v\rho_n(v)\), monotonicity and its success condition would imply
\[
\frac25
\le B_n\bigl(v,c_v\rho_n(v)\bigr)
\le B_n(v,\widetilde r)
\le\frac1{10},
\]
a contradiction. The sentinel \(D_v\) also exceeds \(c_v\rho_n(v)\). Thus every element of the nonempty set defining the capped infimum in \cref{obj:def:critical-radius} is at least \(c_v\rho_n(v)\). The bounded argument uses its positive-weight witness and the legal separation established above. Consequently,
\[
c_v\rho_n(v)\le r_n^*(v),
\qquad
c_w^{\mathrm b}\rho_n^{\mathrm b}(w)
\le r_n^{*,\mathrm b}(w).
\]
The displayed risk bounds also imply both asserted strict inequalities against \(1/10\).
% lean: CausalSmith.Stat.FinitepHomogeneityDensegamma.bounded_matched_lower
% lean: CausalSmith.Stat.FinitepHomogeneityDensegamma.matched_lower_separation
% lean: CausalSmith.Stat.FinitepHomogeneityDensegamma.normalized_prior_positive_atom
% lean: CausalSmith.Stat.FinitepHomogeneityDensegamma.cappedRadius_lower_of_nonempty_failure

\item By \cref{obj:thm:original-record-attainable-rate}, the public tests satisfy
\[
r_n^*(v)\le C_v\rho_n(v),
\qquad
r_n^{*,\mathrm b}(w)\le C_w^{\mathrm b}\rho_n^{\mathrm b}(w).
\]
Together with \cref{obj:prop:bounded-submodel-comparison} and the preceding lower bounds, this gives
\[
\begin{aligned}
c_v\rho_n(v)
&\le r_n^*(v)\le C_v\rho_n(v),\\
c_w^{\mathrm b}\rho_n^{\mathrm b}(w)
&\le r_n^{*,\mathrm b}(w)
=r_n^{*,\mathrm{bin}}(w)
\le C_w^{\mathrm b}\rho_n^{\mathrm b}(w).
\end{aligned}
\]
The same attainment result supplies
\[
\begin{aligned}
\mathsf R_n(P,\phi^*_{n,v})&\le0.0075<\frac1{10}
&&\quad(P\in H_0(v)),\\
\mathsf R_n(P,\phi^{*,\mathrm b}_{n,w})&\le0.0075<\frac1{10}
&&\quad(P\in H_0^{\mathrm b}(w)).
\end{aligned}
\]
Under the respective conditions
\[
C_v\rho_n(v)<D_v,
\qquad
C_w^{\mathrm b}\rho_n^{\mathrm b}(w)<D_w^{\mathrm b},
\]
it supplies, respectively,
\[
\begin{aligned}
1-\mathsf R_n(P,\phi^*_{n,v})&\le0.0075<\frac1{10}
&&\quad(P\in\mathcal M_v,\ d(P)\ge C_v\rho_n(v)),\\
1-\mathsf R_n(P,\phi^{*,\mathrm b}_{n,w})&\le0.0075<\frac1{10}
&&\quad(P\in\mathcal M_w^{\mathrm b},\
d(P)\ge C_w^{\mathrm b}\rho_n^{\mathrm b}(w)).
\end{aligned}
\]

The bounded radius is positive by its lower bound. Dividing the matched bounds therefore gives
\[
\frac{c_v\rho_n(v)}
     {C_w^{\mathrm b}\rho_n^{\mathrm b}(w)}
\le\Delta_n(v)
\le
\frac{C_v\rho_n(v)}
     {c_w^{\mathrm b}\rho_n^{\mathrm b}(w)}.
\]
For \(n>0\), the law of real powers gives
\[
\frac{\rho_n(v)}{\rho_n^{\mathrm b}(w)}
=\frac{n^{-E(p)}}{n^{-E(2)}}
=n^{E(2)-E(p)}.
\]
Hence
\[
\frac{c_v}{C_w^{\mathrm b}}n^{E(2)-E(p)}
\le\Delta_n(v)
\le
\frac{C_v}{c_w^{\mathrm b}}n^{E(2)-E(p)}.
\]
Finally, \cref{obj:prop:bounded-submodel-comparison} gives
\[
0<r_n^{*,\mathrm b}(w)\le r_n^*(v),
\qquad
1\le\Delta_n(v).
\]
% lean: CausalSmith.Stat.FinitepHomogeneityDensegamma.matched_radius_ratio_bounds
% lean: CausalSmith.Stat.FinitepHomogeneityDensegamma.frontier_at

\item Since \(E(p)>0\) and \(E(2)>0\), the definitions of the scales give
\[
\rho_n(v)=n^{-E(p)}\longrightarrow0,
\qquad
\rho_n^{\mathrm b}(w)=n^{-E(2)}\longrightarrow0.
\]
The public thresholds are those of \cref{obj:def:sharp-frontier-scales}, evaluated at \(v\) and \(v_2\). Their formulas in \cref{obj:thm:original-record-attainable-rate} give
\[
N_v\ge4,\qquad N_w^{\mathrm b}\ge4.
\]
Its witness-threshold conclusion therefore applies to every natural number satisfying the corresponding threshold and yields
\[
\begin{aligned}
n\ge N_v&\ \Longrightarrow\ C_v\rho_n(v)<d_0,\\
n\ge N_w^{\mathrm b}&\ \Longrightarrow\
C_w^{\mathrm b}\rho_n^{\mathrm b}(w)<d_0.
\end{aligned}
\]
% lean: CausalSmith.Stat.FinitepHomogeneityDensegamma.rho_tendsto_zero
% lean: CausalSmith.Stat.FinitepHomogeneityDensegamma.rhoBounded_tendsto_zero
% lean: CausalSmith.Stat.FinitepHomogeneityDensegamma.frontier_at

\item We make the exponent difference explicit, always preserving \(w\) when evaluating at \(2\). Define
\[
D_2=1+2S+\frac{S}{2\gamma},
\qquad
F_2=2S+\frac{S}{2\gamma}.
\]
Substitution of \(q=(p-1)/p\) gives
\[
D_p=D_2+\frac{\beta(2-p)}{p-1},
\qquad
F(p)=F_2+\frac{(2\alpha+\beta)(2-p)}{p-1}
\ge F_2.
\]
For the supplied-propensity channel, a common denominator yields
\[
\begin{aligned}
E_0(2)-E_0(p)
&=\frac{2\gamma}{4\gamma+1}
 -\frac{2\gamma(p-1)}{2\gamma p+p-1}\\
&=\frac{4\gamma^2(2-p)}
 {(4\gamma+1)(2\gamma p+p-1)}.
\end{aligned}
\]
Both denominators are positive. For the rough channel, the identity for \(D_p\) gives
\[
\begin{aligned}
E_4(2)-E_4(p)
&=\frac{2S}{D_2}-\frac{2S}{D_p}\\
&=\frac{2S(D_p-D_2)}{D_2D_p}\\
&=\frac{2S\beta(2-p)}{(p-1)D_pD_2}.
\end{aligned}
\]
Combining these identities with the phase rules \(E(p)=E_0(p)\) for \(F(p)\ge1\) and \(E(p)=E_4(p)\) for \(F(p)<1\), applied also at \(2\), gives the exhaustive expression
\[
E(2)-E(p)=
\begin{cases}
\displaystyle
\frac{4\gamma^2(2-p)}
 {(4\gamma+1)(2\gamma p+p-1)},
&F_2\ge1,\\[6pt]
\displaystyle
E_4(2)-E_0(p),
&F_2<1\le F(p),\\[4pt]
\displaystyle
\frac{2S\beta(2-p)}{(p-1)D_pD_2},
&F(p)<1.
\end{cases}
\]
Indeed, in the first case \(F(p)\ge F_2\ge1\), so both minimum exponents select the supplied-propensity channel. In the second case the bounded exponent selects the rough channel and the finite-moment exponent selects the supplied-propensity channel. In the third case \(F_2\le F(p)<1\), so both select the rough channel.

If \(p<2\), the two channel differences displayed above are strictly positive. Therefore
\[
E(p)\le E_0(p)<E_0(2),
\qquad
E(p)\le E_4(p)<E_4(2),
\]
and hence
\[
E(p)<\min\{E_0(2),E_4(2)\}=E(2).
\]
Using the scale definitions in \cref{obj:def:sharp-frontier-scales}, we obtain
\[
\frac{\rho_n(v)}{\rho_n^{\mathrm b}(w)}
=n^{E(2)-E(p)}
\longrightarrow\infty.
\]
% lean: CausalSmith.Stat.FinitepHomogeneityDensegamma.phase_bounded_identities
% lean: CausalSmith.Stat.FinitepHomogeneityDensegamma.phase_bounded_tail_difference
% lean: CausalSmith.Stat.FinitepHomogeneityDensegamma.phase_bounded_rough_difference
% lean: CausalSmith.Stat.FinitepHomogeneityDensegamma.phase_difference_eq
% lean: CausalSmith.Stat.FinitepHomogeneityDensegamma.exponent_phase_algebra

\item Suppose \(p<2\). Apply \cref{obj:thm:tail-essential-full-record-converse} to the selected priors and magnitudes of \cref{obj:def:converse-handle}. It gives
\[
L_{n,v}\longrightarrow\infty.
\]
For every \(n\ge2\), both selected priors are normalized finite probability priors, and their positive-weight laws satisfy
\[
\begin{aligned}
\pi_{0,n,v}(\{P\})>0
&\ \Longrightarrow\ P\in H_0(v),\\
\pi_{1,n,v}(\{P\})>0
&\ \Longrightarrow\
P\in\mathcal M_v,\quad
c_v\rho_n(v)\le d(P)<D_v.
\end{aligned}
\]
For each positive-weight law in either prior, their conditional outcome kernels satisfy
\[
Q_{a,P}\bigl(\{0,-L_{n,v},L_{n,v}\}\mid x\bigr)=1
\]
for \(a\in\{0,1\}\) and \(\lambda\)-almost every \(x\). The same result supplies the complete-record bound
\[
\operatorname{TV}(\mathbb P_{0,n,v},\mathbb P_{1,n,v})
<\frac12.
\]

For any normalized finite probability prior \(\pi\), write its complete-record mixture as
\[
\mathbb P_{\pi,n}=\int P^{\otimes n}\,\pi(dP).
\]
The bounded-prior conclusion of the cited result gives a threshold \(N(v)\in\mathbb N\) such that, for every \(n\ge N(v)\) and every pair of normalized finite probability priors \(\pi_0,\pi_1\),
\[
\left.
\begin{aligned}
\pi_0(\{P\})>0&\ \Longrightarrow\ P\in H_0^{\mathrm b}(w),\\
\pi_1(\{P\})>0&\ \Longrightarrow\
P\in\mathcal M_w^{\mathrm b},\quad c_v\rho_n(v)\le d(P)
\end{aligned}
\right\}
\ \Longrightarrow\
\operatorname{TV}(\mathbb P_{\pi_0,n},\mathbb P_{\pi_1,n})
\ge\frac12.
\]
These are precisely the asserted support, magnitude, and complete-record mixture conclusions.
% lean: CausalSmith.Stat.FinitepHomogeneityDensegamma.frontier_at

\item For \(p<2\), the lower ratio bound and the strictly positive exponent difference give
\[
0<\frac{c_v}{C_w^{\mathrm b}},
\qquad
\frac{c_v}{C_w^{\mathrm b}}n^{E(2)-E(p)}
\longrightarrow\infty,
\qquad
\Delta_n(v)\longrightarrow\infty.
\]
The definition in \cref{obj:def:tail-region} includes admissibility and \(p<2\). The preceding divergence supplies its remaining condition, so
\[
\mathcal T=\{v\in\mathcal V:p<2\}.
\]

To characterize \cref{obj:def:equality-region}, suppose that an admissible tuple has a bounded radius ratio. Thus there exists \(C_\Delta\in\mathbb R\) such that
\[
\Delta_n(v)\le C_\Delta
\qquad(n\ge2).
\]
If \(p\ne2\), admissibility implies \(p<2\), and the divergence just established supplies \(N_\Delta\in\mathbb N\) with
\[
n\ge N_\Delta\ \Longrightarrow\ \Delta_n(v)>C_\Delta.
\]
At
\[
n_\Delta=\max\{2,N_\Delta\},
\]
the two bounds would give
\[
C_\Delta<\Delta_{n_\Delta}(v)\le C_\Delta,
\]
a contradiction. Thus boundedness implies \(p=2\). Conversely, when \(p=2\), evaluation at \(2\) preserves the entire tuple, and the upper ratio bound reduces to
\[
\Delta_n(v)\le\frac{C_v}{c_w^{\mathrm b}}
\qquad(n\ge2).
\]
Consequently,
\[
\mathcal E=\{v\in\mathcal V:p=2\}.
\]
% lean: CausalSmith.Stat.FinitepHomogeneityDensegamma.radiusRatio_tendsto_atTop
% lean: CausalSmith.Stat.FinitepHomogeneityDensegamma.tailRegion_characterization
% lean: CausalSmith.Stat.FinitepHomogeneityDensegamma.second_moment_radiusRatio_bounded
% lean: CausalSmith.Stat.FinitepHomogeneityDensegamma.equalityRegion_characterization

\item Define the ambient set
\[
\mathcal U
=\{(p,\alpha,\beta,\gamma)\in\mathbb R^4:p<2\}.
\]
The coordinate map to \(p\) is continuous, so \(\mathcal U\) is open in the product topology. The tuple
\[
v_{\mathrm o}=(3/2,1/2,1/2,1/2)
\]
satisfies every admissibility inequality and belongs to \(\mathcal U\). Thus
\[
v_{\mathrm o}\in\mathcal U\cap\mathcal V,
\qquad
\mathcal U\cap\mathcal V\ne\varnothing.
\]
The characterization of \(\mathcal T\) gives
\[
\mathcal U\cap\mathcal V\subseteq\mathcal T,
\]
which supplies the required open set.
% lean: CausalSmith.Stat.FinitepHomogeneityDensegamma.tailRegion_nonempty_open

\item It remains to prove compact uniformity. For each admissible tuple define the auxiliary positive base and lower envelope by
\[
\Theta(v)
=\max\left\{2,32^{D_p\gamma/(2S)}\right\}+1,
\]
\[
\underline c_{\mathrm{env}}(v)
=\min\left\{
c_0^{\mathrm e},\,
k_v,\,
c_0^{\mathrm e}\Theta(v)^{-|E_0(p)-E_4(p)|}
\right\}.
\]
Every entry in this minimum is positive, so
\[
\underline c_{\mathrm{env}}(v)>0.
\]
The ceiling bound gives
\[
2\le N_{\mathrm{low}}(v)\le\Theta(v).
\]
Since the exponent \(-|E_0(p)-E_4(p)|\) is nonpositive, and since
\(N_{\mathrm{low}}(v)\ge1\), monotonicity in the base and then in the exponent gives
\[
\begin{aligned}
\Theta(v)^{-|E_0(p)-E_4(p)|}
&\le N_{\mathrm{low}}(v)^{-|E_0(p)-E_4(p)|}\\
&\le N_{\mathrm{low}}(v)^{-(E_0(p)-E_4(p))}.
\end{aligned}
\]
On \(F(p)\ge1\), the envelope is at most
\(c_0^{\mathrm e}=c_v\). On \(F(p)<1\), it is at most each of
\[
k_v,
\qquad
c_0^{\mathrm e}
N_{\mathrm{low}}(v)^{-(E_0(p)-E_4(p))}.
\]
The branchwise definition therefore gives, throughout the admissible domain,
\[
0<\underline c_{\mathrm{env}}(v)\le c_v.
\]

The envelope is continuous at every admissible tuple. Indeed, the coordinate maps are continuous; the denominators defining \(q,D_p,E_0,E_4\) are nonzero there; and every variable power in the envelope has a positive base. Absolute value, maximum, minimum, and the remaining arithmetic operations preserve continuity.

Let \(\mathcal K\subseteq\mathcal V\) be compact and nonempty. The continuous envelope attains a minimum at some \(v_{\mathcal K}\in\mathcal K\). Define
\[
c_{\mathcal K}
=\underline c_{\mathrm{env}}(v_{\mathcal K})
=\min_{v\in\mathcal K}\underline c_{\mathrm{env}}(v)>0.
\]
Then
\[
c_{\mathcal K}
\le\underline c_{\mathrm{env}}(v)
\le c_v
\qquad(v\in\mathcal K).
\]
The compact-uniformity conclusion of \cref{obj:thm:original-record-attainable-rate} supplies real constants \(C_{\mathcal K},N_{\mathcal K}\) such that
\[
C_v\le C_{\mathcal K},
\qquad
N_v\le N_{\mathcal K}
\qquad(v\in\mathcal K).
\]
For the empty compact set, the choices
\[
c_{\mathcal K}=1,\qquad C_{\mathcal K}=0,\qquad N_{\mathcal K}=0
\]
satisfy all requirements vacuously.

Finally, let \(\mathcal K_{\mathrm b}\subseteq\mathcal W\) be compact. Its embedded parameter set is
\[
\mathcal K_2
=\{(2,\alpha,\beta,\gamma):
(\alpha,\beta,\gamma)\in\mathcal K_{\mathrm b}\}.
\]
The embedding \(w\mapsto(2,w)\) is continuous, so \(\mathcal K_2\) is compact, and all its tuples are admissible. Applying the preceding result to \(\mathcal K_2\), and using the public identities
\[
c_w^{\mathrm b}=c_{(2,w)},
\qquad
C_w^{\mathrm b}=C_{(2,w)},
\qquad
N_w^{\mathrm b}=N_{(2,w)},
\]
gives
\[
c_{\mathcal K_2}\le c_w^{\mathrm b},
\qquad
C_w^{\mathrm b}\le C_{\mathcal K_2},
\qquad
N_w^{\mathrm b}\le N_{\mathcal K_2}
\qquad(w\in\mathcal K_{\mathrm b}),
\]
with \(c_{\mathcal K_2}>0\). This establishes both compact-uniformity assertions.
% lean: CausalSmith.Stat.FinitepHomogeneityDensegamma.compactLowerEnvelope_le
% lean: CausalSmith.Stat.FinitepHomogeneityDensegamma.continuousAt_compactLowerEnvelope
% lean: CausalSmith.Stat.FinitepHomogeneityDensegamma.cLower_compact_uniformity
% lean: CausalSmith.Stat.FinitepHomogeneityDensegamma.frontier_compact_uniformity
\end{enumerate}
\end{proof}

\begin{proof}[Proof of \cref{obj:thm:oracle-frontier-broad-class}]
Fix an admissible tuple \(v\). We use the constants and exponent defined in the statement.

\begin{enumerate}
\item Comparing the restrictions in
\cref{obj:def:model,obj:def:oracle-broad-model,obj:def:null,obj:def:oracle-broad-null}
gives
\[
\mathcal M_v\subseteq\mathcal M_v^{\mathrm{e,br}},
\qquad
H_0(v)\subseteq H_0^{\mathrm{e,br}}(v).
\]
The law displayed in \cref{obj:lem:causal-nonempty} therefore supplies a member of the broader model with distance \(d_0\).

For any \(P\in\mathcal M_v^{\mathrm{e,br}}\), put
\[
\bar\tau_P=\int_{[0,1]}\tau_P(x)\,\lambda(dx).
\]
Expanding the centered square and using the effect cap gives
\[
d(P)^2
=\int_{[0,1]}\tau_P(x)^2\,\lambda(dx)-\bar\tau_P^2
\le\int_{[0,1]}\tau_P(x)^2\,\lambda(dx)
\le\frac14.
\]
Thus the set of broader-model distances is nonempty and bounded above, and
\[
d_0\le D_v^{\mathrm{e,br}}\le\frac12.
\]
Furthermore, \(\gamma\le1\) implies
\[
4^{-\gamma}\ge4^{-1}=\frac14,
\qquad
c_v^{\mathrm e}
=\frac{4^{-\gamma}}{16\sqrt3}
\ge\frac1{64\sqrt3},
\qquad
C_v^{\mathrm e}>0.
\]
% lean: CausalSmith.Stat.FinitepHomogeneityDensegamma.oracle_model_distance_bounds
% lean: CausalSmith.Stat.FinitepHomogeneityDensegamma.cOracle_uniform_lower
% lean: CausalSmith.Stat.FinitepHomogeneityDensegamma.COr_pos

\item Fix an integer \(n\ge2\), and define
\[
r_{\mathrm{low}}=c_v^{\mathrm e}\rho_n^{\mathrm e}(v).
\]
By \cref{obj:lem:paired-tent-full-record-lower}, there are a null law \(P_0\in H_0(v)\) and a normalized finite prior
\[
\pi_1=\sum_{j=1}^k\omega_j\delta_{P_j},
\qquad
\omega_j\ge0,
\qquad
\sum_{j=1}^k\omega_j=1,
\]
with at least one positive weight, such that
\[
e_{P_0}(x)=\frac12,
\qquad
\omega_j>0
\ \Longrightarrow\
P_j\in\mathcal M_v,\quad
e_{P_j}(x)=\frac12,\quad
r_{\mathrm{low}}\le d(P_j)<d_0.
\]
Define the supplied function and the complete-record mixture by
\[
f=e_{P_0},
\qquad
\overline P_1^{(n)}
=\sum_{j=1}^k\omega_jP_j^{\otimes n}.
\]
The function \(f\) is continuous and satisfies \(f(x)=1/2\) everywhere. The same cited result gives
\[
\operatorname{TV}\!\left(P_0^{\otimes n},\overline P_1^{(n)}\right)<\frac12.
\]
The inclusions already established place \(P_0\) in the broader null and every positive-weight \(P_j\) in the broader model, with \(e_{P_j}=f\).

Admissibility gives \(q>0\), \(\gamma>0\), and hence \(E_0>0\). Consequently,
\[
0<r_{\mathrm{low}}
\le d(P_j)<d_0\le D_v^{\mathrm{e,br}}
\qquad(\omega_j>0).
\]
In particular, \(r_{\mathrm{low}}\) is a legal separation with a nonempty alternative class.
% lean: CausalSmith.Stat.FinitepHomogeneityDensegamma.paired_tent_full_record_lower
% lean: CausalSmith.Stat.FinitepHomogeneityDensegamma.oracle_frontier_broad_class

\item Let \(\phi^{\mathrm e}\) be an allowed supplied-propensity test satisfying the broader-null size constraint. Freeze its function input at \(f\), and average its public randomization:
\[
\overline\phi_f(\mathcal D_n)
=\int_0^1\phi^{\mathrm e}(f,\mathcal D_n,U)\,dU,
\qquad
0\le\overline\phi_f(\mathcal D_n)\le1.
\]
Joint measurability makes this a measurable function of the sample. Since all positive-weight witnesses have supplied function \(f\), integration against the finite mixture gives
\[
\int\overline\phi_f\,dP_0^{\otimes n}
=\mathsf R_n^{\mathrm e}(P_0,\phi^{\mathrm e}),
\qquad
\int\overline\phi_f\,d\overline P_1^{(n)}
=\sum_{j=1}^k\omega_j\mathsf R_n^{\mathrm e}(P_j,\phi^{\mathrm e}).
\]
For a measurable function taking values in \([0,1]\), integrating its level-set probability differences bounds the difference of its expectations by total variation. Therefore
\[
\left|
\sum_{j=1}^k\omega_j\mathsf R_n^{\mathrm e}(P_j,\phi^{\mathrm e})
-\mathsf R_n^{\mathrm e}(P_0,\phi^{\mathrm e})
\right|
\le
\operatorname{TV}\!\left(P_0^{\otimes n},\overline P_1^{(n)}\right),
\]
and hence
\[
\sum_{j=1}^k\omega_j\mathsf R_n^{\mathrm e}(P_j,\phi^{\mathrm e})
<\frac1{10}+\frac12=\frac35.
\]
If every positive-weight atom had type-II error strictly below \(2/5\), then its rejection expectation would exceed \(3/5\). Zero-weight atoms contribute zero, so normalization would imply
\[
\sum_{j=1}^k\omega_j\mathsf R_n^{\mathrm e}(P_j,\phi^{\mathrm e})
\ge\frac35\sum_{j=1}^k\omega_j
=\frac35,
\]
a contradiction. Thus
\[
\exists j:\quad
\omega_j>0,\qquad
1-\mathsf R_n^{\mathrm e}(P_j,\phi^{\mathrm e})\ge\frac25.
\]

The feasible test class is nonempty: the jointly measurable test
\[
\phi_0^{\mathrm e}(f',\mathcal D_n,U)=0
\qquad
\bigl(f'\in C([0,1]),\ \mathcal D_n\text{ an }n\text{-record sample},\
U\in[0,1]\bigr)
\]
has rejection expectation zero under every law. For every allowed test and every law,
\[
0\le\mathsf R_n^{\mathrm e}(P,\phi^{\mathrm e})\le1.
\]
Thus the alternative errors are bounded above by one. The positive-weight atom just obtained belongs to the alternative class at \(r_{\mathrm{low}}\), so
\[
\sup_{\substack{P\in\mathcal M_v^{\mathrm{e,br}}\\
d(P)\ge r_{\mathrm{low}}}}
\left\{1-\mathsf R_n^{\mathrm e}(P,\phi^{\mathrm e})\right\}
\ge\frac25.
\]
Taking the infimum over feasible tests in
\cref{obj:def:oracle-broad-risk} proves
\[
B_n^{\mathrm{e,br}}(v,r_{\mathrm{low}})\ge\frac25.
\]
This also proves the required error witness for the particular prior constructed above.
% lean: CausalSmith.Stat.FinitepHomogeneityDensegamma.oracle_single_null_prior_error_witness
% lean: CausalSmith.Stat.FinitepHomogeneityDensegamma.oracleRejectProb_bounds
% lean: CausalSmith.Stat.FinitepHomogeneityDensegamma.oracle_frontier_broad_class

\item For every separation \(r\) with \(0<r\le r_{\mathrm{low}}\), the alternative class at \(r_{\mathrm{low}}\) is contained in the alternative class at \(r\). Both are nonempty, their error suprema are bounded, and the feasible test class is unchanged. Therefore
\[
B_n^{\mathrm{e,br}}(v,r)
\ge B_n^{\mathrm{e,br}}(v,r_{\mathrm{low}})
\ge\frac25>\frac1{10}.
\]
Every successful separation in the set defining
\cref{obj:def:oracle-broad-radius} is consequently above \(r_{\mathrm{low}}\); its sentinel \(D_v^{\mathrm{e,br}}\) is also above \(r_{\mathrm{low}}\). Taking the infimum gives
\[
r_n^{*,\mathrm{e,br}}(v)\ge r_{\mathrm{low}}.
\]
% lean: CausalSmith.Stat.FinitepHomogeneityDensegamma.oracleBroadCriticalRadius_ge_of_failed

\item We now establish calibration and power directly on the broader model. Use the construction in \cref{obj:def:oracle-handle}, with
\[
s=\lfloor n/2\rfloor,\qquad
a=s^{-E_0},\qquad
T=2^{\lceil\log_2(a^{-1/(p-1)})\rceil},\qquad
J=2^{\lceil\log_2(a^{-1/\gamma})\rceil},
\]
\[
b=20T^{1-p},
\qquad
\Lambda=\frac{80T^{2-p}}s.
\]
Here
\[
s\ge1,\qquad 0<a\le1,\qquad T\ge1,\qquad J\ge1,\qquad
b\ge0,\qquad\Lambda>0.
\]
For a broader-model law \(P\), overlap makes the clipped supplied propensity equal to \(e_P\).

The elementary clipping inequalities are
\[
|y-\ell_T(y)|\le |y|^pT^{1-p},
\qquad
\ell_T(y)^2\le |y|^pT^{2-p}.
\]
Indeed, when \(|y|\le T\), the first difference vanishes and
\[
y^2=|y|^p|y|^{2-p}\le |y|^pT^{2-p}.
\]
When \(|y|>T\), the clipping remainder is at most \(|y|\), and
\[
|y|
=|y|^p|y|^{1-p}
\le |y|^pT^{1-p},
\qquad
\ell_T(y)^2\le T^2
=T^pT^{2-p}\le |y|^pT^{2-p}.
\]
The raw moment restriction gives first-moment integrability in both arms, using \(|y|\le1+|y|^p\). Integrating the two clipping inequalities yields, for \(\lambda\)-almost every \(x\),
\[
\left|\int_{\mathbb R}(y-\ell_T(y))\,Q_{\iota_{\mathrm{arm}},P}(dy\mid x)\right|
\le10T^{1-p},
\qquad
\int_{\mathbb R}\ell_T(y)^2\,Q_{\iota_{\mathrm{arm}},P}(dy\mid x)
\le10T^{2-p}
\quad(\iota_{\mathrm{arm}}\in\{0,1\}).
\]

Define the clipped conditional score mean by
\[
g_{T,P}(x)
=
\int_{\mathbb R}\ell_T(y)\,Q_{1,P}(dy\mid x)
-
\int_{\mathbb R}\ell_T(y)\,Q_{0,P}(dy\mid x).
\]
Cancellation of the arm probabilities against the inverse-propensity denominators identifies this function as the conditional mean of \(R_i\). The conditional arm-mean identities then give
\[
|g_{T,P}(x)-\tau_P(x)|\le20T^{1-p}=b
\]
almost everywhere. The conditional second moment of the score satisfies
\[
\begin{aligned}
\mathbb E(R_i^2\mid X_i=x)
&=
\frac{\int\ell_T(y)^2\,Q_{1,P}(dy\mid x)}{e_P(x)}
+
\frac{\int\ell_T(y)^2\,Q_{0,P}(dy\mid x)}{1-e_P(x)}
\\
&\le40T^{2-p}+40T^{2-p}
=80T^{2-p}.
\end{aligned}
\]
% lean: CausalSmith.Stat.FinitepHomogeneityDensegamma.oracle_clipping_moments
% lean: CausalSmith.Stat.FinitepHomogeneityDensegamma.conditional_truncation_moments

\item Define the constant-removing matrix and the sample expectation notation by
\[
u=J^{-1/2}(1,\ldots,1)^T,\qquad
C_J=\operatorname{Id}_{\mathbb R^J}-uu^T,
\]
\[
\mathbb P_{n,P}=P^{\otimes n},
\qquad
\mathbb E_{n,P}[\cdot]=\int(\cdot)\,dP^{\otimes n}.
\]
The histogram cells in \cref{obj:def:projection} are disjoint and have Lebesgue mass \(1/J\). Consequently,
\[
\int_{[0,1]}F_J(x)F_J(x)^T\,\lambda(dx)
=\operatorname{Id}_{\mathbb R^J},
\qquad
u^TF_J(x)=1,
\qquad
C_JF_J(x)=F_J(x)-u.
\]
For every \(z_{\mathrm{dir}}\in\mathbb R^J\),
\[
\|C_Jz_{\mathrm{dir}}\|^2
=\|z_{\mathrm{dir}}\|^2-(u^Tz_{\mathrm{dir}})^2
\le\|z_{\mathrm{dir}}\|^2.
\]
Combining these identities with the conditional score second-moment bound gives
\[
\begin{aligned}
\mathbb E_{n,P}
\!\left[\bigl(z_{\mathrm{dir}}^T(F_J(X_i)-u)R_i\bigr)^2\right]
&\le80T^{2-p}
\int_{[0,1]}
\bigl(z_{\mathrm{dir}}^TC_JF_J(x)\bigr)^2\,\lambda(dx)
\\
&=80T^{2-p}\|C_Jz_{\mathrm{dir}}\|^2
\\
&\le80T^{2-p}\|z_{\mathrm{dir}}\|^2.
\end{aligned}
\]

The two disjoint evaluation blocks have equal size and consist of independent records. Their common mean and centered vectors are
\[
\mu_P
=\mathbb E_{n,P}V_1
=\mathbb E_{n,P}V_2
=C_J\int_{[0,1]}F_J(x)g_{T,P}(x)\,\lambda(dx),
\qquad
\mathsf Z_{\iota_{\mathrm{block}},P}=V_{\iota_{\mathrm{block}}}-\mu_P\quad(\iota_{\mathrm{block}}\in\{1,2\}).
\]
Variance additivity within a block and the preceding single-record bound imply
\[
\mathbb E_{n,P}
\bigl[(z_{\mathrm{dir}}^T\mathsf Z_{\iota_{\mathrm{block}},P})^2\bigr]
=
\operatorname{Var}_{n,P}(z_{\mathrm{dir}}^TV_{\iota_{\mathrm{block}}})
\le\Lambda\|z_{\mathrm{dir}}\|^2.
\]
The centered vectors are independent and have mean zero.

Define the quadratic statistic by
\[
W_P=\langle V_1,V_2\rangle.
\]
Block independence gives
\[
\mathbb E_{n,P}W_P=\|\mu_P\|^2,
\qquad
W_P-\|\mu_P\|^2
=
\langle\mu_P,\mathsf Z_{1,P}\rangle
+\langle\mu_P,\mathsf Z_{2,P}\rangle
+\langle\mathsf Z_{1,P},\mathsf Z_{2,P}\rangle.
\]
All three cross-product expectations in the square of this expansion vanish:
\[
\begin{aligned}
\mathbb E_{n,P}
\bigl[\langle\mu_P,\mathsf Z_{1,P}\rangle
      \langle\mu_P,\mathsf Z_{2,P}\rangle\bigr]&=0,\\
\mathbb E_{n,P}
\bigl[\langle\mu_P,\mathsf Z_{1,P}\rangle
      \langle\mathsf Z_{1,P},\mathsf Z_{2,P}\rangle\bigr]&=0,\\
\mathbb E_{n,P}
\bigl[\langle\mu_P,\mathsf Z_{2,P}\rangle
      \langle\mathsf Z_{1,P},\mathsf Z_{2,P}\rangle\bigr]&=0.
\end{aligned}
\]
The first equality uses independence and both zero means; the other two follow by integrating first over the block appearing linearly.

For the remaining quadratic term, independence and the directional second-moment bound give
\[
\begin{aligned}
\mathbb E_{n,P}
\bigl[\langle\mathsf Z_{1,P},\mathsf Z_{2,P}\rangle^2\bigr]
&\le\Lambda\,\mathbb E_{n,P}\|\mathsf Z_{1,P}\|^2\\
&=\Lambda\sum_{j=1}^J
\mathbb E_{n,P}(\mathsf Z_{1,P})_j^2
\le J\Lambda^2.
\end{aligned}
\]
Thus
\[
\operatorname{Var}_{n,P}(W_P)
\le2\Lambda\|\mu_P\|^2+J\Lambda^2.
\]
All quantities have finite second moments, since the clipped scores and finite histogram features are bounded.

Define the good event by
\[
\mathcal G_{n,P}
=
\left\{\mathcal D_n:
|W_P-\|\mu_P\|^2|
\le16\bigl(\sqrt\Lambda\,\|\mu_P\|+\sqrt J\,\Lambda\bigr)
\right\}.
\]
Its tolerance is strictly positive. Applying the second-moment tail bound gives
\[
\begin{aligned}
\mathbb P_{n,P}(\mathcal G_{n,P}^{\,c})
&\le
\frac{2\Lambda\|\mu_P\|^2+J\Lambda^2}
{256\bigl(\sqrt\Lambda\,\|\mu_P\|+\sqrt J\,\Lambda\bigr)^2}\\
&\le\frac1{128}<0.01.
\end{aligned}
\]
The last inequality uses the nonnegative cross term in the denominator.
% lean: CausalSmith.Stat.FinitepHomogeneityDensegamma.oracleBlockVec_variance_le
% lean: CausalSmith.Stat.FinitepHomogeneityDensegamma.oracleBlockVec_inner_integral
% lean: CausalSmith.Stat.FinitepHomogeneityDensegamma.oracleBlockVec_quadratic_variance_le
% lean: CausalSmith.Stat.FinitepHomogeneityDensegamma.oracleBlockVec_tail_le

\item Suppose \(P\in H_0^{\mathrm{e,br}}(v)\), and choose its null constant:
\[
c_{\mathrm{null},P}\in[-1/2,1/2],
\qquad
\tau_P(x)=c_{\mathrm{null},P}\quad(x\in[0,1]).
\]
The histogram coefficient of this constant is \(c_{\mathrm{null},P}u\), which \(C_J\) annihilates. Hence
\[
\mu_P
=C_J\int_{[0,1]}
F_J(x)\bigl(g_{T,P}(x)-c_{\mathrm{null},P}\bigr)\,\lambda(dx).
\]
The clipping bound and cell masses imply
\[
\left|
\left(\int F_J(x)
\bigl(g_{T,P}(x)-c_{\mathrm{null},P}\bigr)\,\lambda(dx)\right)_j
\right|
\le\sqrt J\,\frac bJ=\frac b{\sqrt J}.
\]
Summing the squared coordinate bounds and using contraction by \(C_J\) gives
\[
\|\mu_P\|^2\le
\sum_{j=1}^J\frac{b^2}{J}=b^2.
\]
On \(\mathcal G_{n,P}\),
\[
W_P
\le b^2+16\sqrt\Lambda\,b+16\sqrt J\,\Lambda
\le2b^2+80\sqrt J\,\Lambda
\le2b^2+1024\sqrt J\,\Lambda.
\]
Here
\[
16\sqrt\Lambda\,b\le b^2+64\Lambda
\]
follows from \((b-8\sqrt\Lambda)^2\ge0\), and \(\sqrt J\ge1\).

The deterministic rule in \cref{obj:def:oracle-handle} rejects precisely when \(W_P\) exceeds its displayed cutoff. Its rejection event is therefore contained in \(\mathcal G_{n,P}^{\,c}\), and averaging the independent public randomization leaves that probability unchanged. Thus
\[
\mathsf R_n^{\mathrm e}(P,\phi^{*,\mathrm e}_{n,v})
\le0.01
\qquad(P\in H_0^{\mathrm{e,br}}(v)).
\]
% lean: CausalSmith.Stat.FinitepHomogeneityDensegamma.oracle_whole_null_size_le
% lean: CausalSmith.Stat.FinitepHomogeneityDensegamma.oracle_null_population_norm_le
% lean: CausalSmith.Stat.FinitepHomogeneityDensegamma.oracle_null_cutoff_le

\item For a general broader-model law \(P\), define its histogram coefficients and their associated constant by
\[
\theta_{\mathrm{hist},P}
=\int_{[0,1]}F_J(x)\tau_P(x)\,\lambda(dx),
\qquad
c_{\mathrm{hist},P}
=\frac1{\sqrt J}\sum_{j=1}^J(\theta_{\mathrm{hist},P})_j,
\qquad
\theta_{\mathrm{cen},P}=C_J\theta_{\mathrm{hist},P}.
\]
Subtracting the constant coefficient gives
\[
\theta_{\mathrm{cen},P}
=\int F_J(x)\bigl(\tau_P(x)-c_{\mathrm{hist},P}\bigr)\,\lambda(dx)
=\int F_J(x)\bigl((\Pi_J\tau_P)(x)-c_{\mathrm{hist},P}\bigr)\,\lambda(dx).
\]
The second equality holds because histogram projection preserves every cell integral.

Effect Hölder smoothness bounds the deviation from a cell average:
\[
|\tau_P(x)-(\Pi_J\tau_P)(x)|
\le20J^{-\gamma},
\qquad
\int(\tau_P-\Pi_J\tau_P)^2\,d\lambda
\le(20J^{-\gamma})^2.
\]
Indeed, every point in the same cell lies within distance \(1/J\), so averaging the Hölder bound over that cell proves the pointwise inequality.

The histogram projection is constant on each cell, and the residual has integral zero on each cell. Its residual is therefore orthogonal to the projection and to constants. Expanding the square, and using the histogram feature isometry, yields
\[
\int(\tau_P-c_{\mathrm{hist},P})^2\,d\lambda
=
\|\theta_{\mathrm{cen},P}\|^2
+\int(\tau_P-\Pi_J\tau_P)^2\,d\lambda.
\]
Also, expansion around the mean gives
\[
\int(\tau_P-c_{\mathrm{hist},P})^2\,d\lambda
=d(P)^2+(\bar\tau_P-c_{\mathrm{hist},P})^2
\ge d(P)^2.
\]
Consequently,
\[
d(P)^2
\le\|\theta_{\mathrm{cen},P}\|^2+(20J^{-\gamma})^2
\le\bigl(\|\theta_{\mathrm{cen},P}\|+20J^{-\gamma}\bigr)^2,
\]
and hence
\[
d(P)-20J^{-\gamma}\le\|\theta_{\mathrm{cen},P}\|.
\]

Finally,
\[
\mu_P-\theta_{\mathrm{cen},P}
=C_J\int F_J(x)\bigl(g_{T,P}(x)-\tau_P(x)\bigr)\,\lambda(dx).
\]
The same cellwise coefficient bound used for the null, now applied to
\(|g_{T,P}-\tau_P|\le b\), gives
\[
\|\mu_P-\theta_{\mathrm{cen},P}\|\le b.
\]
The reverse triangle inequality therefore proves the required population signal bound:
\[
d(P)-20J^{-\gamma}-b\le\|\mu_P\|.
\]
% lean: CausalSmith.Stat.FinitepHomogeneityDensegamma.centered_effect_histogram_distance
% lean: CausalSmith.Stat.FinitepHomogeneityDensegamma.oracle_effect_histogram_coefficient_distance
% lean: CausalSmith.Stat.FinitepHomogeneityDensegamma.oracle_population_clipping_error
% lean: CausalSmith.Stat.FinitepHomogeneityDensegamma.oracle_alternative_population_norm_ge

\item If
\[
d(P)\ge20J^{-\gamma}+5b+128\sqrt{\sqrt J\,\Lambda},
\]
the preceding population bound gives
\[
\|\mu_P\|\ge4b+128\sqrt{\sqrt J\,\Lambda}.
\]
Since \(J\ge1\),
\[
\sqrt\Lambda\le\sqrt{\sqrt J\,\Lambda},
\qquad
16\sqrt\Lambda\,\|\mu_P\|\le\frac{\|\mu_P\|^2}{2}.
\]
On \(\mathcal G_{n,P}\), it follows that
\[
\begin{aligned}
W_P
&\ge\|\mu_P\|^2
-16\sqrt\Lambda\,\|\mu_P\|-16\sqrt J\,\Lambda\\
&\ge\frac{\|\mu_P\|^2}{2}-16\sqrt J\,\Lambda\\
&\ge8b^2+8176\sqrt J\,\Lambda\\
&>2b^2+1024\sqrt J\,\Lambda.
\end{aligned}
\]
The penultimate inequality expands
\((4b+128\sqrt{\sqrt J\,\Lambda})^2/2\) and discards its nonnegative cross term. The last inequality is strict because \(\Lambda>0\). Thus the rule rejects throughout \(\mathcal G_{n,P}\), proving
\[
1-\mathsf R_n^{\mathrm e}(P,\phi^{*,\mathrm e}_{n,v})\le0.01
\]
at the displayed sufficient separation.
% lean: CausalSmith.Stat.FinitepHomogeneityDensegamma.oracle_power_of_mean_bound
% lean: CausalSmith.Stat.FinitepHomogeneityDensegamma.oracle_power_cutoff_lt

\item To express this separation in the public rate, define the variance-growth exponent
\[
t=\frac{2-p}{p-1}\ge0.
\]
Direct algebra gives
\[
2+t=\frac1q,
\qquad
E_0\left(2+t+\frac1{2\gamma}\right)=1.
\]
Upward dyadic rounding costs at most a factor of two:
\[
a^{-1/(p-1)}\le T\le2a^{-1/(p-1)},
\qquad
a^{-1/\gamma}\le J\le2a^{-1/\gamma}.
\]
The lower rounding bounds and the negative bias exponents imply
\[
b=20T^{1-p}\le20a,
\qquad
J^{-\gamma}\le a.
\]
Because \(0\le2-p\le1\), the upper rounding bounds imply
\[
T^{2-p}\le2a^{-t},
\qquad
\sqrt J\le\sqrt2\,a^{-1/(2\gamma)}.
\]
The exponent balance now gives
\[
s^{-1}a^{-t-1/(2\gamma)}=a^2,
\]
and consequently
\[
\sqrt J\,\Lambda
=\frac{80\sqrt J\,T^{2-p}}s
\le160\sqrt2\,a^2.
\]
Combining these estimates yields
\[
20J^{-\gamma}+5b+128\sqrt{\sqrt J\,\Lambda}
\le
\bigl(120+128\sqrt{160\sqrt2}\bigr)a.
\]

For every \(n\ge2\), the block size satisfies \(n\le3s\). Moreover,
\[
0<E_0=\frac{2\gamma q}{2\gamma+q}\le1,
\]
so
\[
a=s^{-E_0}
\le3^{E_0}n^{-E_0}
\le3\rho_n^{\mathrm e}(v).
\]
It follows that
\[
20J^{-\gamma}+5b+128\sqrt{\sqrt J\,\Lambda}
\le C_v^{\mathrm e}\rho_n^{\mathrm e}(v).
\]
Thus every broader-model law at or above the nominal upper separation has the asserted type-II bound, in particular under the stated gate
\(C_v^{\mathrm e}\rho_n^{\mathrm e}(v)<D_v^{\mathrm{e,br}}\).
All preceding estimates apply when \(s=1\) or \(J=1\); positivity of \(s,T,J,\Lambda\) suffices throughout.
% lean: CausalSmith.Stat.FinitepHomogeneityDensegamma.oracle_exponent_balance
% lean: CausalSmith.Stat.FinitepHomogeneityDensegamma.oracle_attaining_separation_le

\item Define
\[
r_{\mathrm{up}}=C_v^{\mathrm e}\rho_n^{\mathrm e}(v)>0.
\]
If \(r_{\mathrm{up}}<D_v^{\mathrm{e,br}}\), calibration makes
\(\phi^{*,\mathrm e}_{n,v}\) feasible at level \(1/10\), and its power bound gives
\[
B_n^{\mathrm{e,br}}(v,r_{\mathrm{up}})
\le
\sup_{\substack{P\in\mathcal M_v^{\mathrm{e,br}}\\d(P)\ge r_{\mathrm{up}}}}
\left\{1-\mathsf R_n^{\mathrm e}(P,\phi^{*,\mathrm e}_{n,v})\right\}
\le0.01\le\frac1{10}.
\]
For a nonempty alternative class, this supremum is bounded between zero and one; the empty-class convention assigns it zero and yields the same upper estimate. Here nonemptiness also follows from
\(r_{\mathrm{up}}<D_v^{\mathrm{e,br}}\): otherwise every model distance would be strictly below \(r_{\mathrm{up}}\), making \(r_{\mathrm{up}}\) an upper bound for their supremum. Thus \(r_{\mathrm{up}}\) belongs to the successful-radius set.

If \(D_v^{\mathrm{e,br}}\le r_{\mathrm{up}}\), the sentinel in
\cref{obj:def:oracle-broad-radius} instead gives
\[
r_n^{*,\mathrm{e,br}}(v)
\le D_v^{\mathrm{e,br}}\le r_{\mathrm{up}}.
\]
The defining set is nonempty because it contains the sentinel, and it is bounded below by zero. In both cases,
\[
r_n^{*,\mathrm{e,br}}(v)\le r_{\mathrm{up}}.
\]
Together with the lower bound, this proves the radius inequalities. The witnesses constructed above have all the required model memberships, common function input, distance bounds, normalization, total variation comparison, and positive-weight error witness.
% lean: CausalSmith.Stat.FinitepHomogeneityDensegamma.oracleBroadCriticalRadius_le_of_test
% lean: CausalSmith.Stat.FinitepHomogeneityDensegamma.oracle_frontier_broad_class

\item For \(n\ge N_v^{\mathrm e}\), the ceiling definition implies
\[
n\ge\left(\frac{2C_v^{\mathrm e}}{d_0}\right)^{1/E_0}.
\]
The base and \(E_0\) are positive. Raising this inequality to the negative power \(-E_0\) therefore gives
\[
C_v^{\mathrm e}n^{-E_0}
\le
C_v^{\mathrm e}
\left[
\left(\frac{2C_v^{\mathrm e}}{d_0}\right)^{1/E_0}
\right]^{-E_0}
=
C_v^{\mathrm e}\left(\frac{2C_v^{\mathrm e}}{d_0}\right)^{-1}
=\frac{d_0}{2}.
\]
% lean: CausalSmith.Stat.FinitepHomogeneityDensegamma.oracle_attainment_threshold

\item For \(P\in\mathcal M_v^{\mathrm{e,br}}\), the conditional arm-mean identities hold simultaneously for \(\lambda\)-almost every \(x\):
\[
\int_{\mathbb R}y\,Q_{0,P}(dy\mid x)=m_{0,P}(x),
\qquad
\int_{\mathbb R}y\,Q_{1,P}(dy\mid x)
=m_{0,P}(x)+\tau_P(x).
\]
Overlap gives \(e_P(x)>0\) and \(1-e_P(x)>0\). Pulling the constant denominators out of the integrals and cancelling them yields
\[
\begin{aligned}
&e_P(x)\int_{\mathbb R}\frac{y}{e_P(x)}\,Q_{1,P}(dy\mid x)
+(1-e_P(x))\int_{\mathbb R}\frac{-y}{1-e_P(x)}\,Q_{0,P}(dy\mid x)\\
&\qquad=
\int_{\mathbb R}y\,Q_{1,P}(dy\mid x)
-\int_{\mathbb R}y\,Q_{0,P}(dy\mid x)
=\tau_P(x).
\end{aligned}
\]
The conditional first moments are finite almost everywhere by the raw moment bound.
% lean: CausalSmith.Stat.FinitepHomogeneityDensegamma.oracle_untruncated_mean_identity

\item Finally, the coordinate maps \(v\mapsto p\) and \(v\mapsto\gamma\) are continuous in the product topology. On the admissible domain,
\[
p>0,\qquad
q=\frac{p-1}{p}>0,\qquad
2\gamma+q>0.
\]
Hence \(q\) and \(E_0\) are continuous there, and \(E_0\) is everywhere positive.

For a compact subset \(\mathcal K\) of that domain, define
\[
E_{\min,\mathcal K}
=
\begin{cases}
\displaystyle\min_{v\in\mathcal K}E_0(v),&\mathcal K\ne\varnothing,\\[1mm]
1,&\mathcal K=\varnothing.
\end{cases}
\]
In the nonempty case, continuity and compactness supply an attained minimum, whose value is positive by admissibility. In the empty case the assertions are vacuous. Thus
\[
E_{\min,\mathcal K}>0,
\qquad
E_{\min,\mathcal K}\le E_0(v)\quad(v\in\mathcal K).
\]
Together with the global lower-multiplier bound and the fixed finite upper multiplier,
\[
c_v^{\mathrm e}\ge\frac1{64\sqrt3},
\qquad
C_v^{\mathrm e}=3\bigl(120+128\sqrt{160\sqrt2}\bigr),
\]
this proves the asserted local uniformity.
% lean: CausalSmith.Stat.FinitepHomogeneityDensegamma.oracle_compact_uniformity
\end{enumerate}
\end{proof}

\begin{proof}[Proof of \cref{obj:thm:oracle-frontier-resolved}]
Fix \(v\in\mathcal V\). Admissibility gives
\[
1<p\le2,\qquad \frac14\le\gamma\le1,\qquad
0<q=\frac{p-1}{p}\le\frac12,\qquad
2\gamma+q>0,\qquad E_0>0.
\]
We use the constants and scales in the statement.

\begin{enumerate}
\item Since \(-1\le-\gamma\) and the base \(4\) exceeds one,
\[
4^{-1}\le4^{-\gamma},
\qquad
\frac{1}{64\sqrt3}
=\frac{4^{-1}}{16\sqrt3}
\le c^{\mathrm e}_v.
\]
Also,
\[
C^{\mathrm e}_v
=3\left(120+128\sqrt{160\sqrt2}\right)>0.
\]
% lean: CausalSmith.Stat.FinitepHomogeneityDensegamma.cOracle_uniform_lower
% lean: CausalSmith.Stat.FinitepHomogeneityDensegamma.COr_pos

\item Fix an integer \(n\ge2\), and define the lower separation by
\[
\varrho_-=c^{\mathrm e}_v\rho_n^{\mathrm e}(v).
\]
By \cref{obj:lem:paired-tent-full-record-lower}, there are a null law \(P_0\in H_0(v)\) and a finite prior
\[
\pi_1=\sum_{j=1}^{k}\omega_j\delta_{P_j},
\qquad
\omega_j\ge0,\qquad
\sum_{j=1}^{k}\omega_j=1,
\qquad
\exists j:\omega_j>0.
\]
For every positive-weight support law, that result supplies
\[
P_j\in\mathcal M_v,\qquad
e_{P_j}(x)=\frac12=e_{P_0}(x)\quad(x\in[0,1]),
\qquad
\varrho_-\le d(P_j)<d_0\le D_v.
\]
It also supplies
\[
\mathbb P_{\pi_1,n}
=\sum_{j=1}^{k}\omega_jP_j^{\otimes n},
\qquad
\operatorname{TV}\!\left(P_0^{\otimes n},\mathbb P_{\pi_1,n}\right)<\frac12,
\qquad
B_n^{\mathrm e}(v,\varrho_-)\ge\frac25.
\]
% lean: CausalSmith.Stat.FinitepHomogeneityDensegamma.paired_tent_full_record_lower

Choose an index \(j_+\in\{1,\ldots,k\}\) with \(\omega_{j_+}>0\). Positivity of the lower multiplier and of \(n^{-E_0}\) yields
\[
0<\varrho_-\le d(P_{j_+})<D_v.
\]
Thus the alternatives at separation \(\varrho_-\) are nonempty. For every \(0<\varrho\le\varrho_-\), their inclusion in the alternatives at separation \(\varrho\) gives
\[
B_n^{\mathrm e}(v,\varrho_-)
\le B_n^{\mathrm e}(v,\varrho).
\]
The size constraint is the same in both infima. Consequently, every successful separation below \(D_v\) in \cref{obj:def:oracle-radius} is at least \(\varrho_-\), because
\[
B_n^{\mathrm e}(v,\varrho)\ge\frac25>\frac1{10}
\qquad(0<\varrho\le\varrho_-).
\]
The additional candidate \(D_v\) also exceeds \(\varrho_-\). Taking the infimum therefore gives
\[
c^{\mathrm e}_v\rho_n^{\mathrm e}(v)
=\varrho_-\le r_n^{*,\mathrm e}(v).
\]
% lean: CausalSmith.Stat.FinitepHomogeneityDensegamma.oracleCriticalRadius_ge_of_failed

\item We next establish the null calibration of the rule in \cref{obj:def:oracle-handle}. Use its tuning quantities
\[
s=\lfloor n/2\rfloor,\qquad a=s^{-E_0},\qquad
T=2^{\lceil\log_2(a^{-1/(p-1)})\rceil},\qquad
J=2^{\lceil\log_2(a^{-1/\gamma})\rceil},
\]
\[
b=20T^{1-p},\qquad
\Lambda=\frac{80T^{2-p}}s.
\]
Here \(s\ge1\), \(0<a\le1\), \(T\ge1\), \(J\ge1\), \(b\ge0\), and \(\Lambda>0\).
With \(F_J\) from \cref{obj:def:projection}, define
\[
u=J^{-1/2}(1,\ldots,1)^T,\qquad
C_J=I_J-uu^T.
\]
The cells have Lebesgue measure \(1/J\), so
\[
\int_0^1F_J(x)F_J(x)^T\,d\lambda(x)=I_J,
\qquad
\int_0^1F_J(x)\,d\lambda(x)=u,
\qquad
C_JF_J(x)=F_J(x)-u.
\]
Thus \(C_J\) is a symmetric orthogonal projection and
\[
\|C_J\xi_{\mathrm{dir}}\|_2\le\|\xi_{\mathrm{dir}}\|_2
\qquad(\xi_{\mathrm{dir}}\in\mathbb R^J).
\]

For \(P\in\mathcal M_v\), define the clipped conditional mean and second moment on \([0,1]\) by
\[
g_{T,P}(x)
=\int_{\mathbb R}\ell_T(y)\,Q_{1,P}(dy\mid x)
-\int_{\mathbb R}\ell_T(y)\,Q_{0,P}(dy\mid x),
\]
\[
m_{2,T,P}(x)
=\frac{\int_{\mathbb R}\ell_T(y)^2\,Q_{1,P}(dy\mid x)}{e_P(x)}
+\frac{\int_{\mathbb R}\ell_T(y)^2\,Q_{0,P}(dy\mid x)}{1-e_P(x)}.
\]
These functions are finite and measurable, since clipping is bounded and overlap makes both denominators positive.

For every \(y\in\mathbb R\),
\[
|y-\ell_T(y)|\le |y|^pT^{1-p},
\qquad
\ell_T(y)^2\le |y|^pT^{2-p}.
\]
Indeed, on \(|y|\le T\), the clipping remainder vanishes and
\(|y|^{2-p}\le T^{2-p}\); on \(|y|>T\), the remainder is at most \(|y|\) and the clipped square equals \(T^2\). The moment restriction in \cref{obj:ass:raw-moment} therefore gives, for each \(a'\in\{0,1\}\) and almost every \(x\),
\[
\left|\int_{\mathbb R}(y-\ell_T(y))\,Q_{a',P}(dy\mid x)\right|
\le10T^{1-p},
\qquad
\int_{\mathbb R}\ell_T(y)^2\,Q_{a',P}(dy\mid x)
\le10T^{2-p}.
\]
The same moment restriction ensures integrability of \(y\), using
\(|y|\le1+|y|^p\).
The arm-mean identities now yield
\[
g_{T,P}(x)-\tau_P(x)
=
\int_{\mathbb R}(y-\ell_T(y))\,Q_{0,P}(dy\mid x)
-\int_{\mathbb R}(y-\ell_T(y))\,Q_{1,P}(dy\mid x),
\]
and hence
\[
|g_{T,P}(x)-\tau_P(x)|\le b,
\qquad
m_{2,T,P}(x)
\le10T^{2-p}
\left(\frac1{e_P(x)}+\frac1{1-e_P(x)}\right)
\le80T^{2-p}
\]
almost everywhere. These are respectively the conditional mean and second moment bounds for \(R_i\), since the supplied true propensity agrees with its overlap-clipped version.
% lean: CausalSmith.Stat.FinitepHomogeneityDensegamma.oracle_clipping_moments

Define the common block mean and the two histogram coefficient vectors by
\[
\mu_P=\mathbb E_{P^{\otimes n}}[V_1]
=\mathbb E_{P^{\otimes n}}[V_2]
=C_J\theta_{T,P},
\]
\[
\theta_{T,P}=\int_0^1g_{T,P}(x)F_J(x)\,d\lambda(x),
\qquad
\theta_{\tau,P}=\int_0^1\tau_P(x)F_J(x)\,d\lambda(x).
\]
The equality of the block means follows by averaging \(s\) identically distributed records in each block.

For a null law, define its constant effect by
\[
c_P=\tau_P(0),\qquad \tau_P(x)=c_P\quad(x\in[0,1]).
\]
Because \(C_Ju=0\),
\[
\mu_P=C_J\int_0^1(g_{T,P}(x)-c_P)F_J(x)\,d\lambda(x).
\]
For every histogram coordinate \(j\),
\[
\left|
\left(\int_0^1(g_{T,P}-c_P)F_J\,d\lambda\right)_j
\right|
=
\left|\sqrt J\int_{I_{j,J}}(g_{T,P}-c_P)\,d\lambda\right|
\le\frac b{\sqrt J}.
\]
Summing the squared coordinate bounds and applying the contraction \(C_J\) gives
\[
\|\mu_P\|_2^2
\le\sum_{j=1}^{J}\frac{b^2}{J}=b^2,
\qquad
\|\mu_P\|_2\le b.
\]
% lean: CausalSmith.Stat.FinitepHomogeneityDensegamma.oracle_null_population_norm_le

For any model law and any \(\xi_{\mathrm{dir}}\in\mathbb R^J\), the conditional second moment bound and the feature isometry give
\[
\mathbb E_P\!\left[
\left\langle\xi_{\mathrm{dir}},C_JF_J(X)R\right\rangle^2
\right]
\le80T^{2-p}\int_0^1
\left\langle C_J\xi_{\mathrm{dir}},F_J(x)\right\rangle^2\,d\lambda(x)
\le80T^{2-p}\|\xi_{\mathrm{dir}}\|_2^2.
\]
Here the single-record score is defined by
\[
R=\frac{A\ell_T(Y)}{e_P(X)}
-\frac{(1-A)\ell_T(Y)}{1-e_P(X)}.
\]
Independence within a block gives, for \(b'\in\{1,2\}\),
\[
\operatorname{Var}_{P^{\otimes n}}
\!\left(\left\langle\xi_{\mathrm{dir}},V_{b'}\right\rangle\right)
=
\frac1s\operatorname{Var}_P
\!\left(\left\langle\xi_{\mathrm{dir}},C_JF_J(X)R\right\rangle\right)
\le\Lambda\|\xi_{\mathrm{dir}}\|_2^2.
\]
% lean: CausalSmith.Stat.FinitepHomogeneityDensegamma.oracleBlockVec_variance_le

Define
\[
\Delta_{b',P}=V_{b'}-\mu_P\quad(b'\in\{1,2\}),
\qquad
W_{\mathrm e}=\langle V_1,V_2\rangle.
\]
The centered vectors are independent and have mean zero. In particular,
\[
\mathbb E_{P^{\otimes n}}[W_{\mathrm e}]=\|\mu_P\|_2^2,
\]
\[
W_{\mathrm e}-\|\mu_P\|_2^2
=
\langle\mu_P,\Delta_{1,P}\rangle
+\langle\mu_P,\Delta_{2,P}\rangle
+\langle\Delta_{1,P},\Delta_{2,P}\rangle.
\]
The expectations of all three cross products in the square of this sum vanish: the product of the two linear terms factors into their zero means, and conditioning on the other block makes each linear--bilinear cross product zero. To bound the remaining bilinear square, let
\[
\mathbf e^{(j)}
=\bigl(\mathbf1_{\{j'=j\}}\bigr)_{j'=1}^{J}
\qquad(j=1,\ldots,J)
\]
be the coordinate unit vectors. Independence and the directional second moment bounds imply
\[
\mathbb E_{P^{\otimes n}}
\!\left[\langle\Delta_{1,P},\Delta_{2,P}\rangle^2\right]
\le\Lambda\,\mathbb E_{P^{\otimes n}}\|\Delta_{2,P}\|_2^2
=\Lambda\sum_{j=1}^{J}
\mathbb E_{P^{\otimes n}}
\!\left[\langle\mathbf e^{(j)},\Delta_{2,P}\rangle^2\right]
\le J\Lambda^2.
\]
Consequently,
\[
\operatorname{Var}_{P^{\otimes n}}(W_{\mathrm e})
\le2\Lambda\|\mu_P\|_2^2+J\Lambda^2.
\]
All the required moments are finite because the clipped score and the finite-rank features are bounded.
% lean: CausalSmith.Stat.FinitepHomogeneityDensegamma.oracleBlockVec_quadratic_variance_le

The squared-deviation inequality now gives
\[
\begin{aligned}
&P^{\otimes n}\!\left(
|W_{\mathrm e}-\|\mu_P\|_2^2|
>16\{\sqrt\Lambda\|\mu_P\|_2+\sqrt J\,\Lambda\}
\right)\\
&\quad\le
\frac{2\Lambda\|\mu_P\|_2^2+J\Lambda^2}
{256\{\sqrt\Lambda\|\mu_P\|_2+\sqrt J\,\Lambda\}^2}
\le\frac{2}{256}<0.01.
\end{aligned}
\]
The denominator is positive since \(J\ge1\) and \(\Lambda>0\).
Define the good event by
\[
\mathcal G_P=
\left\{
|W_{\mathrm e}-\|\mu_P\|_2^2|
\le16\{\sqrt\Lambda\|\mu_P\|_2+\sqrt J\,\Lambda\}
\right\}.
\]
Thus \(P^{\otimes n}(\mathcal G_P^c)\le0.01\).
% lean: CausalSmith.Stat.FinitepHomogeneityDensegamma.oracleBlockVec_tail_le

For a null law, on \(\mathcal G_P\) we have
\[
W_{\mathrm e}
\le b^2+16\sqrt\Lambda\,b+16\sqrt J\,\Lambda
\le2b^2+80\sqrt J\,\Lambda
\le2b^2+1024\sqrt J\,\Lambda.
\]
The middle inequality follows from
\[
16\sqrt\Lambda\,b\le b^2+64\Lambda,
\qquad
\Lambda\le\sqrt J\,\Lambda,
\]
the first being equivalent to \((b-8\sqrt\Lambda)^2\ge0\).
The rejection event is therefore contained in \(\mathcal G_P^c\). Since the rule ignores the public randomization variable,
\[
\mathsf R_n^{\mathrm e}(P,\phi^{*,\mathrm e}_{n,v})
=
P^{\otimes n}\!\left(
W_{\mathrm e}>2b^2+1024\sqrt J\,\Lambda
\right)
\le0.01
\qquad(P\in H_0(v)).
\]
These calculations also apply when \(s=1\); when \(J=1\), \(C_J=0\) and both centered block vectors vanish.
% lean: CausalSmith.Stat.FinitepHomogeneityDensegamma.oracle_whole_null_size_le

\item We establish the alternative population bound before applying the same concentration event. Define
\[
\overline\tau_P=\int_0^1\tau_P(x)\,d\lambda(x),
\qquad
c_{P,J}=\frac1{\sqrt J}\sum_{j=1}^{J}(\theta_{\tau,P})_j,
\qquad
h_{P,J}(x)=\tau_P(x)-c_{P,J}.
\]
The coefficient formula for constants and the definition of \(C_J\) give
\[
C_J\theta_{\tau,P}
=\int_0^1h_{P,J}(x)F_J(x)\,d\lambda(x).
\]
On each cell, the histogram projection is the cell average. Hence
\[
\tau_P(x)-(\Pi_J\tau_P)(x)
=J\int_{I_{j,J}}\{\tau_P(x)-\tau_P(z)\}\,d\lambda(z)
\qquad(x\in I_{j,J}),
\]
and \cref{obj:ass:effect-smoothness} gives
\[
|\tau_P(x)-(\Pi_J\tau_P)(x)|\le20J^{-\gamma},
\qquad
\int_0^1\{\tau_P-\Pi_J\tau_P\}^2\,d\lambda
\le(20J^{-\gamma})^2.
\]
Subtracting a constant commutes with \(\Pi_J\). Its cellwise residual has integral zero on every cell, and is therefore orthogonal to the cellwise constant projection. The feature isometry then yields
\[
\int_0^1h_{P,J}^2\,d\lambda
=
\|C_J\theta_{\tau,P}\|_2^2
+\int_0^1\{\tau_P-\Pi_J\tau_P\}^2\,d\lambda.
\]
Expanding the distance to the constant \(c_{P,J}\) also gives
\[
\int_0^1h_{P,J}^2\,d\lambda
=d(P)^2+(\overline\tau_P-c_{P,J})^2
\ge d(P)^2.
\]
It follows that
\[
d(P)^2
\le\|C_J\theta_{\tau,P}\|_2^2+(20J^{-\gamma})^2
\le\bigl(\|C_J\theta_{\tau,P}\|_2+20J^{-\gamma}\bigr)^2,
\]
so
\[
d(P)-20J^{-\gamma}\le\|C_J\theta_{\tau,P}\|_2.
\]
% lean: CausalSmith.Stat.FinitepHomogeneityDensegamma.centered_effect_histogram_distance

The same cellwise coefficient bound used for the null, now applied to \(g_{T,P}-\tau_P\), gives
\[
\|\mu_P-C_J\theta_{\tau,P}\|_2
=
\left\|C_J\int_0^1(g_{T,P}-\tau_P)F_J\,d\lambda\right\|_2
\le b.
\]
Combining this with the preceding lower bound and the reverse triangle inequality yields
\[
d(P)-20J^{-\gamma}-b\le\|\mu_P\|_2.
\]
% lean: CausalSmith.Stat.FinitepHomogeneityDensegamma.oracle_alternative_population_norm_ge

Define the noise scale by
\[
\zeta_{\mathrm{noise}}=\sqrt{\sqrt J\,\Lambda}>0.
\]
If
\[
d(P)\ge20J^{-\gamma}+5b+128\zeta_{\mathrm{noise}},
\]
then
\[
\|\mu_P\|_2\ge4b+128\zeta_{\mathrm{noise}}.
\]
Since \(\sqrt\Lambda\le\zeta_{\mathrm{noise}}\), this implies
\[
16\sqrt\Lambda\,\|\mu_P\|_2\le\frac{\|\mu_P\|_2^2}{2}.
\]
On \(\mathcal G_P\), therefore,
\[
\begin{aligned}
W_{\mathrm e}
&\ge\|\mu_P\|_2^2
-16\sqrt\Lambda\,\|\mu_P\|_2
-16\zeta_{\mathrm{noise}}^2\\
&\ge\frac{\|\mu_P\|_2^2}{2}-16\zeta_{\mathrm{noise}}^2\\
&\ge8b^2+8176\zeta_{\mathrm{noise}}^2
>2b^2+1024\zeta_{\mathrm{noise}}^2.
\end{aligned}
\]
Here the third inequality uses
\[
\frac{(4b+128\zeta_{\mathrm{noise}})^2}{2}
\ge8b^2+8192\zeta_{\mathrm{noise}}^2,
\]
and the last is strict because \(\zeta_{\mathrm{noise}}>0\).
Thus failure to reject is contained in \(\mathcal G_P^c\), proving
\[
1-\mathsf R_n^{\mathrm e}(P,\phi^{*,\mathrm e}_{n,v})\le0.01
\]
under the displayed separation condition.
% lean: CausalSmith.Stat.FinitepHomogeneityDensegamma.oracle_power_of_mean_bound

\item We now bound that sufficient separation by the scale in the statement. Define the variance-growth exponent by
\[
t=\frac{2-p}{p-1}\ge0.
\]
Both dyadic targets are at least one, and upward rounding gives
\[
a^{-1/(p-1)}\le T\le2a^{-1/(p-1)},
\qquad
a^{-1/\gamma}\le J\le2a^{-1/\gamma}.
\]
Since \(1-p<0\) and \(-\gamma<0\),
\[
b=20T^{1-p}\le20a,
\qquad
J^{-\gamma}\le a.
\]
Also \(0\le2-p\le1\), so
\[
T^{2-p}
\le2^{2-p}a^{-t}\le2a^{-t},
\qquad
\sqrt J\le\sqrt2\,a^{-1/(2\gamma)}.
\]
Substitution of the definitions of \(q\) and \(E_0\) gives the balance identity
\[
E_0\left(2+t+\frac1{2\gamma}\right)=1.
\]
Because \(a=s^{-E_0}\), this identity implies
\[
\frac{a^{-t-1/(2\gamma)}}s=a^2.
\]
Consequently,
\[
\sqrt J\,\Lambda
\le160\sqrt2\,\frac{a^{-t-1/(2\gamma)}}s
=160\sqrt2\,a^2,
\qquad
\zeta_{\mathrm{noise}}\le\sqrt{160\sqrt2}\,a.
\]

For every \(n\ge2\), \(n\le3s\). Moreover,
\[
2\gamma q\le2\gamma+q,\qquad 0<E_0\le1,
\qquad 3^{E_0}\le3.
\]
Monotonicity of the negative power therefore gives
\[
a=s^{-E_0}
\le3^{E_0}n^{-E_0}
\le3\rho_n^{\mathrm e}(v).
\]
Combining the approximation, clipping, and covariance bounds yields
\[
\begin{aligned}
20J^{-\gamma}+5b+128\zeta_{\mathrm{noise}}
&\le\left(120+128\sqrt{160\sqrt2}\right)a\\
&\le C^{\mathrm e}_v\rho_n^{\mathrm e}(v).
\end{aligned}
\]
Thus, whenever \(C^{\mathrm e}_v\rho_n^{\mathrm e}(v)<D_v\), every model law at distance at least \(C^{\mathrm e}_v\rho_n^{\mathrm e}(v)\) has type-II error at most \(0.01\).
% lean: CausalSmith.Stat.FinitepHomogeneityDensegamma.oracle_attaining_separation_le

\item Define the upper separation by
\[
\varrho_+=C^{\mathrm e}_v\rho_n^{\mathrm e}(v)>0.
\]
The witnesses already give \(D_v>0\). If \(\varrho_+<D_v\), the calibrated rule satisfies the size constraint \(1/10\), and its alternative errors give
\[
\sup_{P\in\mathcal M_v:\,d(P)\ge\varrho_+}
\left\{1-\mathsf R_n^{\mathrm e}(P,\phi^{*,\mathrm e}_{n,v})\right\}
\le0.01\le\frac1{10}.
\]
For an empty alternative index set, the real-valued supremum is assigned the value \(0\), so the same bound applies in that case. All worst-error suprema are nonnegative, and the admissible test family contains the calibrated rule. Taking its value in the infimum of \cref{obj:def:oracle-risk} therefore shows
\[
B_n^{\mathrm e}(v,\varrho_+)\le\frac1{10}.
\]
Hence \(\varrho_+\) is a candidate in \cref{obj:def:oracle-radius}, and
\[
r_n^{*,\mathrm e}(v)\le\varrho_+.
\]
If instead \(D_v\le\varrho_+\), the candidate \(D_v\) gives
\[
r_n^{*,\mathrm e}(v)\le D_v\le\varrho_+.
\]
Both cases establish the upper radius bound.
% lean: CausalSmith.Stat.FinitepHomogeneityDensegamma.oracleCriticalRadius_le_of_test

\item Retain the prior and null law already constructed, and let \(\phi^{\mathrm e}\) satisfy the stipulated size bound on \(H_0(v)\). Define its seed-averaged map with the common propensity supplied by
\[
\overline\phi^{\mathrm e}_{P_0}(\mathcal D_n)
=\int_0^1\phi^{\mathrm e}(e_{P_0},\mathcal D_n,u_{\mathrm{seed}})\,d\lambda(u_{\mathrm{seed}}),
\qquad
0\le\overline\phi^{\mathrm e}_{P_0}\le1.
\]
This is a measurable function on the complete-record sample space. Integrating event-probability differences over its level sets gives
\[
\left|
\int\overline\phi^{\mathrm e}_{P_0}\,d\mathbb P_{\pi_1,n}
-\int\overline\phi^{\mathrm e}_{P_0}\,dP_0^{\otimes n}
\right|
\le\operatorname{TV}\!\left(P_0^{\otimes n},\mathbb P_{\pi_1,n}\right).
\]
Indeed, each integral equals the integral over \(\vartheta_{\mathrm{lev}}\in[0,1]\) of the probability of the measurable level set
\(\{\overline\phi^{\mathrm e}_{P_0}>\vartheta_{\mathrm{lev}}\}\).

Every positive-weight support law has propensity \(e_{P_0}\). Finite linearity of integration and the null size bound thus yield
\[
\sum_{j=1}^{k}\omega_j
\mathsf R_n^{\mathrm e}(P_j,\phi^{\mathrm e})
\le
\mathsf R_n^{\mathrm e}(P_0,\phi^{\mathrm e})
+\operatorname{TV}\!\left(P_0^{\otimes n},\mathbb P_{\pi_1,n}\right)
<\frac1{10}+\frac12=\frac35.
\]
If every positive-weight index had type-II error below \(2/5\), its rejection probability would exceed \(3/5\). The zero-weight terms vanish, and normalization would then imply
\[
\sum_{j=1}^{k}\omega_j
\mathsf R_n^{\mathrm e}(P_j,\phi^{\mathrm e})
\ge\frac35\sum_{j=1}^{k}\omega_j=\frac35,
\]
a contradiction. Therefore
\[
\exists j:\quad
\omega_j>0,\qquad
1-\mathsf R_n^{\mathrm e}(P_j,\phi^{\mathrm e})\ge\frac25.
\]
This proves the asserted individual error witness for the same common-propensity prior.
% lean: CausalSmith.Stat.FinitepHomogeneityDensegamma.oracle_single_null_prior_error_witness

\item For \(n\ge N^{\mathrm e}_v\), the definition of the ceiling gives
\[
n\ge
\max\left\{2,\left(\frac{2C^{\mathrm e}_v}{d_0}\right)^{1/E_0}\right\}.
\]
Since \(E_0>0\), raising the lower bound for \(n\) to the power \(-E_0\) reverses the inequality:
\[
n^{-E_0}
\le
\left\{
\left(\frac{2C^{\mathrm e}_v}{d_0}\right)^{1/E_0}
\right\}^{-E_0}
=\left(\frac{2C^{\mathrm e}_v}{d_0}\right)^{-1}.
\]
Multiplication by \(C^{\mathrm e}_v>0\) proves
\[
C^{\mathrm e}_v\rho_n^{\mathrm e}(v)\le\frac{d_0}{2}.
\]
% lean: CausalSmith.Stat.FinitepHomogeneityDensegamma.oracle_attainment_threshold

\item Let \(P\in\mathcal M_v\). On the common almost-everywhere set where the two conditional arm-mean identities hold,
\[
\int_{\mathbb R}y\,Q_{0,P}(dy\mid x)=m_{0,P}(x),
\qquad
\int_{\mathbb R}y\,Q_{1,P}(dy\mid x)=m_{0,P}(x)+\tau_P(x).
\]
Overlap gives \(e_P(x)>0\) and \(1-e_P(x)>0\) at every \(x\). Pulling the constant denominators outside the arm integrals and cancelling them therefore yields
\[
\begin{aligned}
&e_P(x)\int_{\mathbb R}\frac{y}{e_P(x)}\,Q_{1,P}(dy\mid x)
+(1-e_P(x))
\int_{\mathbb R}\frac{-y}{1-e_P(x)}\,Q_{0,P}(dy\mid x)\\
&\qquad=
\int_{\mathbb R}y\,Q_{1,P}(dy\mid x)
-\int_{\mathbb R}y\,Q_{0,P}(dy\mid x)
=\tau_P(x).
\end{aligned}
\]
This is the claimed identity for the untruncated score.
% lean: CausalSmith.Stat.FinitepHomogeneityDensegamma.oracle_untruncated_mean_identity

\item Finally, let \(\mathcal K\subseteq\mathcal V\) be compact. The coordinate functions \(p\) and \(\gamma\) are continuous in the inherited topology. Since \(p>1\) and \(2\gamma+(p-1)/p>0\), the function
\[
v\longmapsto
E_0(v)=
\frac{2\gamma\,((p-1)/p)}
{2\gamma+(p-1)/p}
\]
is continuous and positive throughout \(\mathcal V\). If \(\mathcal K\) is nonempty, it attains a positive minimum; if \(\mathcal K\) is empty, choose the value one. Thus there is
\[
\eta_{\mathcal K}>0,
\qquad
\eta_{\mathcal K}\le E_0(v)\quad(v\in\mathcal K).
\]
The multiplier bounds already proved give
\[
c^{\mathrm e}_v\ge\frac1{64\sqrt3},
\qquad
C^{\mathrm e}_v=3\left(120+128\sqrt{160\sqrt2}\right)
\qquad(v\in\mathcal K).
\]
Using \(d_0=1/(8\sqrt2)\) from \cref{obj:lem:paired-tent-full-record-lower}, the base \(2C^{\mathrm e}_v/d_0\) exceeds one and is independent of \(v\). Therefore
\[
N^{\mathrm e}_v
\le
\left\lceil
\max\left\{
2,\left(\frac{2C^{\mathrm e}_v}{d_0}\right)^{1/\eta_{\mathcal K}}
\right\}
\right\rceil
\qquad(v\in\mathcal K).
\]
This supplies a common finite threshold bound. Positivity of the denominators and continuity of the exponent also hold at \(p=2\) and \(\gamma=1/4\), so the same compact-uniform argument includes those faces.
% lean: CausalSmith.Stat.FinitepHomogeneityDensegamma.oracle_compact_uniformity
\end{enumerate}
\end{proof}

\begin{proof}[Proof of \cref{obj:thm:exact-propensity-preservation-boundary}]
\begin{enumerate}
\item Fix \(v=(p,\alpha,\beta,\gamma)\in\mathcal V\), and use the quantities of \cref{obj:def:sharp-frontier-scales}. By \cref{obj:thm:heavy-tail-comparison,obj:thm:oracle-frontier-resolved},
\[
c_v>0,\qquad C_v>0,\qquad
c^{\mathrm e}_v\ge\frac{1}{64\sqrt3}>0,\qquad
C^{\mathrm e}_v>0.
\]
Consequently,
\[
\frac{c_v}{C^{\mathrm e}_v}>0,
\qquad
\frac{C_v}{c^{\mathrm e}_v}>0.
\]
% lean: CausalSmith.Stat.FinitepHomogeneityDensegamma.exact_propensity_preservation_boundary

\item We establish the identities for the comparison exponent. Admissibility gives
\[
p-1>0,\qquad p>0,\qquad q>0,\qquad
\gamma>0,\qquad D_p>0,
\]
and hence
\[
\frac{2\gamma q}{(2\gamma+q)D_p}>0.
\]
The phase identity and channel selection in \cref{obj:thm:heavy-tail-comparison} give
\[
E_4(p)-E_0(p)
=
\frac{2\gamma q}{(2\gamma+q)D_p}\bigl(F(p)-1\bigr),
\]
\[
F(p)\ge1\ \Longrightarrow\ E(p)=E_0(p),
\qquad
F(p)<1\ \Longrightarrow\ E(p)=E_4(p).
\]

If \(F(p)\ge1\), the positive prefactor implies
\[
E_0(p)-E_4(p)\le0.
\]
Thus, using the definition of \(G(v)\),
\[
G(v)=0=E_0(p)-E(p)
=
\frac{2\gamma q}{(2\gamma+q)D_p}\max\{0,1-F(p)\}.
\]
If \(F(p)<1\), the same identity implies
\[
E_0(p)-E_4(p)
=
\frac{2\gamma q}{(2\gamma+q)D_p}\bigl(1-F(p)\bigr)>0.
\]
Channel selection now gives
\[
G(v)=E_0(p)-E_4(p)=E_0(p)-E(p)
=
\frac{2\gamma q}{(2\gamma+q)D_p}\max\{0,1-F(p)\}.
\]
Together, these cases establish
\[
G(v)=E_0(p)-E(p),
\qquad
G(v)=
\frac{2\gamma q}{(2\gamma+q)D_p}\max\{0,1-F(p)\}.
\]
In particular,
\[
F(p)\ge1\ \Longrightarrow\ G(v)=0.
\]
% lean: CausalSmith.Stat.FinitepHomogeneityDensegamma.preservation_gap_identity

\item For every integer \(n\ge2\), the two frontier conclusions in
\cref{obj:thm:heavy-tail-comparison,obj:thm:oracle-frontier-resolved} supply
\[
c_v\rho_n(v)\le r_n^*(v)\le C_v\rho_n(v),
\qquad
c^{\mathrm e}_v\rho_n^{\mathrm e}(v)
\le r_n^{*,\mathrm e}(v)
\le C^{\mathrm e}_v\rho_n^{\mathrm e}(v).
\]
Since \(n>0\), both power scales are positive:
\[
\rho_n(v)=n^{-E(p)}>0,
\qquad
\rho_n^{\mathrm e}(v)=n^{-E_0(p)}>0.
\]
The lower bounds therefore imply
\[
r_n^*(v)>0,\qquad r_n^{*,\mathrm e}(v)>0.
\]
Dividing the original-radius lower bound by the supplied-radius upper bound, and the original-radius upper bound by the supplied-radius lower bound, yields
\[
\frac{c_v\rho_n(v)}
     {C^{\mathrm e}_v\rho_n^{\mathrm e}(v)}
\le
\frac{r_n^*(v)}{r_n^{*,\mathrm e}(v)}
\le
\frac{C_v\rho_n(v)}
     {c^{\mathrm e}_v\rho_n^{\mathrm e}(v)}.
\]
The scale quotient satisfies
\[
\frac{\rho_n(v)}{\rho_n^{\mathrm e}(v)}
=
\frac{n^{-E(p)}}{n^{-E_0(p)}}
=
n^{E_0(p)-E(p)}
=
n^{G(v)}.
\]
Consequently, for every integer \(n\ge2\),
\[
\frac{c_v}{C^{\mathrm e}_v}n^{G(v)}
\le
\frac{r_n^*(v)}{r_n^{*,\mathrm e}(v)}
\le
\frac{C_v}{c^{\mathrm e}_v}n^{G(v)}.
\]
% lean: CausalSmith.Stat.FinitepHomogeneityDensegamma.preservation_ratio_bounds

\item If \(F(p)=1\), the exponent identity gives \(G(v)=0\). Substituting \(n^0=1\) into the quotient bounds gives
\[
\frac{c_v}{C^{\mathrm e}_v}
\le
\frac{r_n^*(v)}{r_n^{*,\mathrm e}(v)}
\le
\frac{C_v}{c^{\mathrm e}_v}
\qquad(n\ge2).
\]
Both bounding constants are strictly positive by \cref{obj:thm:heavy-tail-comparison,obj:thm:oracle-frontier-resolved}.
% lean: CausalSmith.Stat.FinitepHomogeneityDensegamma.exact_propensity_preservation_boundary

\item Suppose \(F(p)<1\). The exponent identity gives
\[
G(v)=
\frac{2\gamma q}{(2\gamma+q)D_p}\bigl(1-F(p)\bigr)>0.
\]
A positive power of the natural sample size tends to infinity, and multiplication by the positive lower multiplier preserves this limit:
\[
n^{G(v)}\longrightarrow+\infty,
\qquad
\frac{c_v}{C^{\mathrm e}_v}n^{G(v)}
\longrightarrow+\infty.
\]
For all sufficiently large natural numbers \(n\), we have \(n\ge2\), so the quotient lower bound gives
\[
\frac{c_v}{C^{\mathrm e}_v}n^{G(v)}
\le
\frac{r_n^*(v)}{r_n^{*,\mathrm e}(v)}.
\]
Comparison with the diverging lower bound proves
\[
\frac{r_n^*(v)}{r_n^{*,\mathrm e}(v)}
\longrightarrow+\infty.
\]
% lean: CausalSmith.Stat.FinitepHomogeneityDensegamma.preservation_ratio_tendsto

Moreover, channel selection gives \(E(p)=E_4(p)\), and hence
\[
G(v)=E_0(p)-E(p)=E_0(p)-E_4(p).
\]
Substitution into the quotient bounds proves
\[
\frac{c_v}{C^{\mathrm e}_v}n^{E_0(p)-E_4(p)}
\le
\frac{r_n^*(v)}{r_n^{*,\mathrm e}(v)}
\le
\frac{C_v}{c^{\mathrm e}_v}n^{E_0(p)-E_4(p)}
\qquad(n\ge2).
\]
% lean: CausalSmith.Stat.FinitepHomogeneityDensegamma.exact_propensity_preservation_boundary

\item We characterize the preservation region of
\cref{obj:def:preservation-region}. Suppose first that \(v\in\mathcal R\).
There is a real upper bound \(b_{\mathrm{sup}}\) satisfying
\[
\frac{r_n^*(v)}{r_n^{*,\mathrm e}(v)}
\le b_{\mathrm{sup}}
\qquad\text{for every integer }n\ge2.
\]
If \(F(p)<1\), the divergence just proved supplies an integer \(n\ge2\) such that
\[
b_{\mathrm{sup}}
<
\frac{r_n^*(v)}{r_n^{*,\mathrm e}(v)}
\le b_{\mathrm{sup}},
\]
a contradiction. Thus \(F(p)\ge1\).

Conversely, if \(v\in\mathcal V\) and \(F(p)\ge1\), then \(G(v)=0\), so
\[
\frac{r_n^*(v)}{r_n^{*,\mathrm e}(v)}
\le\frac{C_v}{c^{\mathrm e}_v}
\qquad\text{for every integer }n\ge2.
\]
This finite upper bound places \(v\) in \(\mathcal R\). Expanding \(F(p)\) therefore gives
\[
\mathcal R
=
\left\{v\in\mathcal V:
\frac{2\alpha}{p-1}
+\frac{\beta}{q}
+\frac{\alpha+\beta}{2\gamma}\ge1
\right\}.
\]
% lean: CausalSmith.Stat.FinitepHomogeneityDensegamma.exact_propensity_preservation_boundary

\item Finally, let \(\mathcal K\subseteq\mathcal V\) be an arbitrary compact set.
The compact-uniformity conclusion of
\cref{obj:thm:heavy-tail-comparison} supplies real constants
\(\underline c_{\mathcal K},\overline C_{\mathcal K}\) satisfying
\[
\underline c_{\mathcal K}>0,
\qquad
\underline c_{\mathcal K}\le c_v,
\qquad
C_v\le\overline C_{\mathcal K}
\qquad(v\in\mathcal K).
\]
The supplied-propensity upper multiplier is independent of \(v\) and positive. Define
\[
c_{\mathrm{pres},\mathcal K}
=
\frac{\underline c_{\mathcal K}}
{3\bigl(120+128\sqrt{160\sqrt2}\bigr)},
\qquad
C_{\mathrm{pres},\mathcal K}
=
64\sqrt3\,\overline C_{\mathcal K}.
\]
Then \(c_{\mathrm{pres},\mathcal K}>0\), and for every \(v\in\mathcal K\),
\[
c_{\mathrm{pres},\mathcal K}
\le\frac{c_v}{C^{\mathrm e}_v},
\qquad
64\sqrt3\,C_v
\le C_{\mathrm{pres},\mathcal K}.
\]
These are the required compact-uniform constants.
% lean: CausalSmith.Stat.FinitepHomogeneityDensegamma.exact_propensity_preservation_boundary
\end{enumerate}
\end{proof}

\bibliographystyle{plainnat}

\bibliography{references}

\end{document}
